% =====================================================================
% AI and Cybersecurity -- Manuale del corso
% Politecnico di Torino -- A.A. 2026/2027
% =====================================================================
\documentclass[11pt,a4paper]{report}
% =====================================================================
% Preambolo -- Manuale AI and Cybersecurity
% =====================================================================
\usepackage[T1]{fontenc}
\usepackage[utf8]{inputenc}
\usepackage[italian]{babel}
\usepackage{lmodern}
\usepackage[a4paper,margin=2.4cm,headheight=14pt]{geometry}
\usepackage{microtype}
\usepackage{xcolor}
\usepackage{graphicx}
\usepackage{enumitem}
\usepackage{booktabs}
\usepackage{array}
\usepackage{tabularx}
\usepackage{longtable}
\usepackage{multirow}
\usepackage{multicol}
\usepackage{amsmath,amssymb,bm}
\usepackage{listings}
\usepackage[most]{tcolorbox}
\usepackage{tikz}
\usetikzlibrary{mindmap,trees,positioning,arrows.meta,shapes.geometric,shapes.misc,calc,fit,backgrounds}
\usepackage[edges]{forest}
\usepackage{fancyhdr}
\usepackage{hyperref}
\hypersetup{
  colorlinks=true,
  linkcolor=blue!50!black,
  urlcolor=blue!60!black,
  pdftitle={AI and Cybersecurity -- Manuale del corso}
}
\usepackage{bookmark}

\setlist{itemsep=2pt,topsep=4pt}
\renewcommand{\arraystretch}{1.25}
\setcounter{tocdepth}{1}
\setlength{\emergencystretch}{3em}

% ---------------------------------------------------------------------
% Colori e intestazioni
% ---------------------------------------------------------------------
\definecolor{polblu}{HTML}{1F3A5F}
\definecolor{defcol}{HTML}{2E6DA4}
\definecolor{esecol}{HTML}{2E8B57}
\definecolor{codebg}{HTML}{F5F7FA}
\definecolor{codekw}{HTML}{0B5394}
\definecolor{codecm}{HTML}{6A8759}
\definecolor{codestr}{HTML}{A64D79}

\pagestyle{fancy}
\fancyhf{}
\fancyhead[L]{\small\textsc{AI and Cybersecurity}}
\fancyhead[R]{\small\nouppercase{\leftmark}}
\fancyfoot[C]{\thepage}
\renewcommand{\headrulewidth}{0.4pt}

% ---------------------------------------------------------------------
% Riquadri sulle fonti (comuni a tutti i manuali della pipeline)
% ---------------------------------------------------------------------
% Punto che il docente ha indicato come rilevante per l'esame (titolo opzionale)
\newtcolorbox{esame}[1][]{breakable,colback=red!4,colframe=red!60!black,
  fonttitle=\bfseries,title={Per l'esame\ifx&#1&\else: #1\fi}}

% Nota generica (titolo opzionale)
\newtcolorbox{nota}[1][Nota]{breakable,colback=gray!6,colframe=gray!60!black,
  fonttitle=\bfseries,title={#1}}

% Informazione non confermata in aula o ricostruita da noi: da verificare
\newtcolorbox{verifica}{breakable,colback=orange!6,colframe=orange!80!black,
  fonttitle=\bfseries,title=Da verificare}

% Soluzione di quiz, esercizi e simulazioni (titolo opzionale)
\newtcolorbox{soluzione}[1][Soluzione]{breakable,colback=green!4,colframe=green!45!black,
  fonttitle=\bfseries,title={#1}}

% Integrazione dei redattori: contenuto NON detto a lezione né presente nelle slide
\newtcolorbox{integrazione}[1][]{breakable,colback=violet!5,colframe=violet!60!black,
  fonttitle=\bfseries,title={Integrazione -- non dal corso\ifx&#1&\else: #1\fi}}

% Segnaposto per parti ancora da scrivere
\newcommand{\dacompletare}[1]{%
  \begin{tcolorbox}[colback=yellow!8,colframe=yellow!60!black,
    fonttitle=\bfseries,title=Da completare]#1\end{tcolorbox}}

% ---------------------------------------------------------------------
% Riquadri didattici (dal corso)
% ---------------------------------------------------------------------
\tcbset{aicbox/.style={enhanced,breakable,boxrule=0.6pt,arc=2pt,
  left=6pt,right=6pt,top=4pt,bottom=4pt,fonttitle=\bfseries\small,
  before skip=8pt,after skip=8pt}}
\newtcolorbox{definizione}[1][]{aicbox,colback=defcol!5,colframe=defcol,
  title={Definizione\ifx&#1&\else: #1\fi}}
\newtcolorbox{esempio}[1][]{aicbox,colback=esecol!5,colframe=esecol,
  title={Esempio\ifx&#1&\else: #1\fi}}
\newtcolorbox{attenzione}[1][]{aicbox,colback=red!3,colframe=red!45!black,
  title={Attenzione\ifx&#1&\else: #1\fi}}
\newtcolorbox{sintesi}[1][]{aicbox,colback=polblu!4,colframe=polblu,
  title={In sintesi\ifx&#1&\else: #1\fi}}

% ---------------------------------------------------------------------
% Intestazione di una lezione del diario
%   \lezione{data}{titolo}{docenti}{argomenti}
% ---------------------------------------------------------------------
\newcommand{\lezione}[4]{%
  \chapter[#1 -- #2]{#2}%
  \begin{tcolorbox}[colback=blue!3,colframe=blue!50!black,
    fonttitle=\bfseries,title=Scheda lezione]
  \textbf{Data:} #1\\
  \textbf{Docenti:} #3\\
  \textbf{Argomenti:} #4
  \end{tcolorbox}}

% Rimando cliccabile a una lezione del diario: \rif{AAAA-MM-GG}{GG/MM}
\newcommand{\rif}[2]{\hyperref[lez:#1]{lez.~#2}}

% ---------------------------------------------------------------------
% Quiz: domanda numerata + opzioni (a), (b), ...
% ---------------------------------------------------------------------
\newcounter{domanda}
\newenvironment{domanda}[1][]{%
  \refstepcounter{domanda}\par\medskip\noindent
  \textbf{Domanda \thedomanda}\ifx&#1&\else\ \textit{(#1)}\fi\textbf{.}~}%
  {\par}
\newenvironment{opzioni}{\begin{enumerate}[label=(\alph*),nosep,leftmargin=2.2em]}{\end{enumerate}}

% ---------------------------------------------------------------------
% Codice Python
% ---------------------------------------------------------------------
\lstdefinestyle{py}{
  language=Python,
  basicstyle=\ttfamily\footnotesize,
  keywordstyle=\color{codekw}\bfseries,
  commentstyle=\color{codecm}\itshape,
  stringstyle=\color{codestr},
  backgroundcolor=\color{codebg},
  frame=single, rulecolor=\color{black!15},
  numbers=left, numberstyle=\tiny\color{black!45}, numbersep=6pt,
  breaklines=true, breakatwhitespace=true,
  showstringspaces=false, columns=fullflexible, keepspaces=true,
  tabsize=4, xleftmargin=14pt, framexleftmargin=12pt,
  morekeywords={self, True, False, None},
  literate={à}{{\`a}}1 {è}{{\`e}}1 {é}{{\'e}}1 {ì}{{\`i}}1 {ò}{{\`o}}1 {ù}{{\`u}}1
}
\lstset{style=py}

% ---------------------------------------------------------------------
% Macro
% ---------------------------------------------------------------------
\newcommand{\en}[1]{\textit{#1}}             % termine inglese
\newcommand{\code}[1]{\texttt{\small #1}}    % codice inline
\newcommand{\R}{\mathbb{R}}
% Fonte di un capitolo tematico: \fonti{lezioni}{slide}
\newcommand{\fonti}[2]{%
  \noindent{\small\color{black!60}\textbf{Lezioni:} #1 \quad\textbf{Slide:} #2}\par\medskip}

% ---------------------------------------------------------------------
% Mappe concettuali (forest)
% ---------------------------------------------------------------------
\forestset{
  mappa/.style={
    for tree={grow'=0, draw=black!35, rounded corners=2pt, node options={align=left},
      font=\scriptsize, inner sep=2.5pt, l sep=5mm, s sep=1.2mm,
      edge={draw=black!45, thin}, anchor=west, child anchor=west, parent anchor=east},
    forked edges,
    where level=0{fill=polblu!18, font=\small\bfseries, text width=2.4cm}{},
    where level=1{fill=defcol!10, font=\scriptsize\bfseries, text width=3.0cm}{},
    where level=2{fill=white, text width=8.6cm}{},
  },
  mappa3/.style={
    mappa,
    where level=2{fill=esecol!6, text width=3.2cm}{},
    where level=3{fill=white, text width=5.2cm}{},
  }
}


\title{AI and Cybersecurity\\[4pt]
  \large Manuale del corso -- Politecnico di Torino, A.A.~2026/2027}
\author{Marco Pedron e colleghi del corso}
\date{Ultimo aggiornamento: \today}

\begin{document}
\maketitle
\tableofcontents
% =====================================================================
% Premessa: come leggere il manuale (legenda dei riquadri e fonti)
% =====================================================================
\chapter*{Come leggere questo manuale}
\addcontentsline{toc}{chapter}{Come leggere questo manuale}

Questo manuale è scritto da studenti, per studenti: \textbf{non è materiale ufficiale del corso} e non
sostituisce le slide, i notebook e le indicazioni dei docenti. Nasce dalle registrazioni delle
lezioni (trascritte automaticamente e poi revisionate), integrate con slide e materiale dei laboratori.

\section*{Le parti}
\begin{description}[leftmargin=2em,style=nextline]
  \item[Diario delle lezioni] una voce per lezione, in ordine cronologico: sintesi, avvisi pratici,
    contenuti in breve e punti segnalati per l'esame. Serve a seguire l'avanzamento del corso.
  \item[Manuale del corso] un capitolo per argomento (dalle basi del deep learning alle architetture,
    agli embedding e all'anomaly detection), che unisce lezioni, slide e notebook in un'esposizione
    ordinata. Rimanda alle lezioni del diario.
  \item[Laboratori] come si scrive il report e, per ogni lab, contesto, consegna, metodo e insidie.
    \textbf{Non contiene soluzioni}: i lab sono valutati e vanno svolti dal proprio gruppo.
  \item[Schemi di ripasso] mappe concettuali e tabelle ``quale scelta quando''.
  \item[Quiz ed esame] formato dell'esame, quiz a scelta multipla e vero/falso, domande aperte ed
    esercizi su pipeline e codice, con soluzioni e rimando alla lezione.
  \item[Appendice] glossario e acronimi.
\end{description}

\section*{Legenda dei riquadri}

\begin{esame}
Punto che il docente ha indicato come rilevante per l'esame.
\end{esame}

\begin{integrazione}
Contenuto aggiunto dai redattori: \textbf{non detto a lezione e non presente nelle slide}
(spiegazioni, esempi, riferimenti bibliografici in più). Utile, ma da trattare con cautela all'esame.
\end{integrazione}

\begin{verifica}
Informazione non confermata: passaggio poco chiaro nella registrazione, dato ricostruito o calcolato.
\end{verifica}

\begin{nota}
Nota generica (fonti usate, avvertenze, collegamenti).
\end{nota}

\begin{soluzione}
Soluzione di quiz, simulazioni e laboratori.
\end{soluzione}

Riquadri didattici, con contenuti presi dal corso:

\begin{definizione}[termine]
Definizione di un concetto.
\end{definizione}

\begin{esempio}
Esempio fatto a lezione o nelle slide.
\end{esempio}

\begin{attenzione}
Errore comune o insidia segnalata dal docente.
\end{attenzione}

\begin{sintesi}
Riepilogo dei punti chiave.
\end{sintesi}

\dacompletare{Parte ancora da scrivere (per esempio un argomento non ancora trattato a lezione).}

I termini tecnici restano in inglese, come nelle slide e all'esame, con la traduzione alla prima
occorrenza; il glossario in appendice raccoglie le corrispondenze.

\section*{Fonti}
Registrazioni delle lezioni (Prof.~Danilo Giordano; Prof.~Matteo Boffa per NLP e ultima parte del
corso), slide, notebook e testi dei laboratori pubblicati sul portale della didattica. Slide, notebook,
registrazioni e dataset non sono inclusi nel repository.


% =====================================================================
\part{Diario delle lezioni}
% Una lezione per file, in ordine cronologico (nome file: AAAA-MM-GG.tex).
% Per aggiungere una lezione: copia lezioni/_template.tex e aggiungi qui l'\input.
% =====================================================================
% =====================================================================
\lezione{21/09/2026}{Presentazione del corso e perché l'AI nella cybersecurity}
  {Prof.~Danilo Giordano}
  {organizzazione del corso, esame e laboratori; minacce e threat intelligence; perché AI/ML;
   AI, ML, reti neurali e deep learning; breve storia; shallow vs deep learning}
\label{lez:2026-09-21}
% =====================================================================

\begin{nota}[Fonte]
Trascrizione automatica della registrazione, revisionata. Numeri e dati dell'esame sono ripresi dalle
slide \emph{0-Intro-Course} e \emph{1-intro-ai-for-cybersecurity} (pp.~1--25), confrontati con quanto
detto in aula.
\end{nota}

\section{Sintesi}
Prima lezione: il docente presenta obiettivi, organizzazione, laboratori ed esame, poi apre la parte
tecnica. L'obiettivo dichiarato non è imparare dettagli fini, ma capire \textbf{come usare AI e ML per
la cybersecurity} e \textbf{cosa c'è sotto il cofano} dei modelli. La parte tecnica motiva l'uso
dell'AI (attacchi più veloci, complessi e costosi), distingue AI, machine learning, reti neurali e deep
learning e confronta la pipeline \emph{shallow} (con feature engineering manuale) con quella
\emph{deep} (con feature extraction automatica). Approfondimento nel capitolo~\ref{cap:ai-cyber}.

\section{Avvisi pratici}
\begin{itemize}
  \item Docenti: Prof.~Danilo Giordano (gran parte del corso) e Prof.~Matteo Boffa (NLP e ultima
    parte). 6 CFU, circa 60 ore tra lezioni e laboratori.
  \item Orari: lunedì 13:00--14:30 (aula 2C) e mercoledì 11:30--14:30 (aula 2S, usato anche per i
    lab). Il martedì è uno slot di recupero, usato \textbf{solo se annunciato} sul portale.
  \item Slide, notebook e testi dei lab escono sul portale man mano. Le slide sono volutamente
    verbose e, secondo il docente, \textbf{bastano anche per il voto massimo}; i libri sono facoltativi.
  \item Ambiente consigliato: \textbf{Google Colab}, per avere una GPU senza averla sul portatile.
  \item Gruppi di lab \textbf{auto-organizzati} da formare prima del primo lab (iscrizione su Moodle).
  \item L'uso di strumenti come ChatGPT è ammesso come aiuto, ma non deve sostituire il lavoro:
    allo scritto non ci sono.
\end{itemize}

\section{Contenuti}

\subsection{Esame e laboratori}
L'esame ha due parti indipendenti, che si possono superare in appelli diversi; si registra quando
entrambe sono sufficienti. Il voto finale è $0{,}6 \times$ scritto $+\ 0{,}4 \times$ media dei report,
arrotondato solo alla fine. Dettagli nel capitolo~\ref{cap:formato}.
\begin{itemize}
  \item \textbf{Scritto} su Moodle con Lockdown Browser e portatile personale, 100 minuti: due
    domande a scelta multipla senza penalità, quattro gruppi di vero/falso con penalità, due domande
    aperte di teoria, tre esercizi di difficoltà crescente su pipeline di ML e codice (capire se un
    codice funziona e perché, non scriverlo da zero).
  \item \textbf{Laboratori}: quattro (introduzione ai framework, 3~h; supervised, anomaly detection,
    NLP, 6~h ciascuno). Per ognuno un report di \textbf{massimo 10 pagine} più il codice; si valuta il
    report, il codice si guarda solo in caso di dubbi. Ogni report vale da 0 a 32 e deve essere
    almeno 18; sopra il 18 non si può riconsegnare per migliorare. I voti valgono due anni
    accademici.
  \item \textbf{Lab 1} solo formativo ma obbligatorio: se consegnato entro il \textbf{1° dicembre}
    con una soluzione ragionevole vale \textbf{+1 punto} sul voto finale. Gli altri report vanno
    consegnati \textbf{10 giorni prima} dell'appello.
\end{itemize}

\begin{esame}
Allo scritto c'è anche del codice: bisogna saper leggere una pipeline in Python/PyTorch e trovarne
gli errori. Nei report conta la \textbf{metodologia} (cosa, come, perché, e se il risultato è buono),
non il codice in sé.
\end{esame}

\subsection{Perché l'AI nella cybersecurity}
Più dispositivi connessi, spesso economici e poco protetti, e attacchi più frequenti: servono difese
più veloci e automatiche. Il docente distingue \textbf{threat} (pericolo potenziale) e
\textbf{attack} (minaccia sfruttata) e tre livelli di conoscenza: ciò che sappiamo (patch, approccio
reattivo), ciò che sappiamo di non sapere (threat intelligence, approccio proattivo) e ciò che non
sappiamo di non sapere (zero-day). I numeri delle slide mostrano attacchi più rapidi (esfiltrazione
mediana da 9 a 2 giorni tra 2022 e 2024), più complessi e più costosi. L'AI aiuta con
\textbf{automazione}, \textbf{intelligenza} (correlazioni tra molte variabili) e \textbf{robustezza},
ma aiuta anche gli attaccanti: è un gioco del gatto e del topo.

\subsection{AI, machine learning, reti neurali, deep learning}
Insiemi annidati: l'AI emula il ragionamento umano anche con regole scritte a mano; il ML
\textbf{impara dai dati} e quindi conosce solo ciò che ha visto; le reti neurali sono un algoritmo di
ML a layer; il deep learning usa reti con più hidden layer e architetture diverse per tipo di dato
(FFNN per tabelle, RNN per sequenze, CNN per immagini, GNN per grafi). Segue una breve storia, dal
perceptron agli LLM, e le ragioni del successo recente: tanti dati, GPU, modelli pre-addestrati.

\subsection{Shallow learning e deep learning}
Nella pipeline KDD classica l'analista costruisce le feature a mano (esempio delle frodi bancarie);
nel deep learning la rete estrae le feature da sola e lo sforzo si sposta sul \emph{model
engineering}. Il preprocessing resta necessario in entrambi i casi. Contro del deep learning:
scarsa trasparenza e interpretabilità.

\begin{esame}
Non usare il deep learning quando basta una soglia: ``non sparare a una mosca con un cannone''.
Saper dire \textbf{quando} ha senso il deep learning e quando no.
\end{esame}

\section{Domande aperte e punti da verificare}
\begin{verifica}
\begin{itemize}
  \item Punteggio massimo dello scritto: in aula il docente dice \textbf{33 punti}, la slide
    \emph{0-Intro-Course} riporta 32. Da confermare sul portale.
  \item Dimensione dei gruppi: in aula \textbf{4 persone} (``novità di quest'anno''); una slide dice 3.
\end{itemize}
\end{verifica}

% =====================================================================
\lezione{23/09/2026}{Paradigmi di apprendimento e modi di usare i modelli}
  {Prof.~Danilo Giordano}
  {training, validation e test; overfitting e underfitting; supervised, unsupervised,
   semi-supervised, self-supervised e reinforcement learning; ensemble e stacking; federated,
   transfer learning e modelli generativi; AI per difesa e attacco; adversarial machine learning}
\label{lez:2026-09-23}
% =====================================================================

\begin{nota}[Fonte]
Trascrizione automatica della registrazione, revisionata. I contenuti seguono le slide
\emph{1-intro-ai-for-cybersecurity} (pp.~26--54); i dati numerici sono ripresi dalle slide.
\end{nota}

\section{Sintesi}
Dopo un breve riepilogo su \emph{shallow} e \emph{deep learning} (più efficace e flessibile, ma
poco trasparente e poco interpretabile: per questo esiste l'\en{explainable AI}, fuori
programma), il docente introduce il lessico di base (training, validation, test) e passa in
rassegna i paradigmi di apprendimento, ciascuno con esempi cyber, benefici e problemi. Chiude con i
modi di usare i modelli (ensemble, federated, transfer, generativi) e con le tre facce del rapporto
tra AI e sicurezza: AI per difendere, AI per attaccare, attacchi ai modelli stessi. L'intera
lezione completa il primo deck di slide. Approfondimento nel capitolo~\ref{cap:ai-cyber}.

\section{Contenuti}

\subsection{Training, validation e test}
Il \textbf{training} (o \en{fitting}: nelle librerie si chiama \code{model.fit}) costruisce il modello
sul training set. Il \textbf{validation set}, mai usato per addestrare, serve a scegliere gli
\textbf{iperparametri} (per esempio l'architettura della rete). Il \textbf{test set}, mai visto né in
training né in validation, serve solo a controllare che le prestazioni restino simili a quelle di
validation, cioè che il modello sia \textbf{generale}. L'analogia dell'esame: le lezioni sono il
training, un tema mai visto è la validazione; chi impara a memoria dieci temi passati senza capire
la materia va in \textbf{overfitting}, chi non impara nulla nemmeno sul training è in
\textbf{underfitting}. Nella realtà spesso un test set adeguato non c'è nemmeno.

\begin{esame}
Il docente chiede esplicitamente di ricordare la distinzione tra i tre insiemi e i loro ruoli: è
alla base di qualunque modello affidabile e torna in tutto il corso (learning curve, scelta degli
iperparametri, capitolo~\ref{cap:addestramento}).
\end{esame}

\subsection{Supervised learning}
Con un dataset \textbf{etichettato} il modello impara a mappare le feature sulla label (classificazione
se categorica, regressione se numerica). La label è la \textbf{ground truth}: costosa, data da un umano,
e deve essere \textbf{molto affidabile}. Esempi: filtro antispam, classificazione di malware,
previsione del volume di traffico, in versione shallow (feature costruite) o deep (dato grezzo). Il
modello conosce solo ciò che ha visto: l'attaccante si adatta (dal ``principe nigeriano'' al finto
amico che chiede soldi), i dati cambiano nel tempo (\textbf{drift}) e le classi malevole sono rare
(\textbf{sbilanciamento}). Rimedi: drift detection con ri-addestramento, pesi delle classi nella loss,
metriche adatte, ricampionamento.

\subsection{Unsupervised learning}
Nessuna label: si cercano strutture nei dati, con dati economici da raccogliere (tutto il traffico
della rete del Politecnico). Metodo \textbf{descrittivo}: \textbf{clustering} (gruppi di IP che si
comportano in modo simile, per esempio una botnet, poi da studiare) e \textbf{anomaly detection}
(impari il normale, segnali ciò che se ne discosta; capitolo~\ref{cap:anomaly}). Problemi: anche qui
c'è drift (la botnet cambia macchine, l'esfiltrazione viene ``spalmata'' nel tempo) e soprattutto
\textbf{l'interpretazione}: senza ground truth le metriche misurano differenze, non correttezza, e i
falsi allarmi crescono. Lo sforzo risparmiato nell'etichettare si paga dopo, nell'analisi; aiutano
l'analisi manuale e la XAI. Il docente cita il ``mantra'' di Jain e Dubes sulla validazione del
clustering come ``arte oscura''.

\subsection{Semi-supervised learning}
Un piccolo insieme etichettato addestra un primo modello, che etichetta il grande insieme non
etichettato (\emph{pseudo-label}); il modello finale si addestra su entrambi. Esempi: IDS con pochi
attacchi etichettati e molto traffico normale; rilevamento di dispositivi infetti da botnet. Unisce la
ground truth del supervised al basso costo dell'unsupervised, ma il campione etichettato può non
essere rappresentativo (pochissimi malevoli) e i dati non etichettati evolvono.

\subsection{Self-supervised learning}
Nessuno in aula lo conosceva: il docente sottolinea che è il paradigma alla base dei modelli moderni,
LLM inclusi, e che il corso ne parla molto. Sono i dati stessi a fornire la supervisione. In un
\emph{pretext task} un \textbf{encoder} comprime l'input in un \textbf{bottleneck} a bassa dimensione
(la \en{feature extraction} la fa il modello); nel \emph{downstream task} un \textbf{decoder} ricostruisce
l'input, oppure un classificatore usa la rappresentazione. Se l'output è molto diverso dall'input, la
rappresentazione non è ancora buona. Esempi: denoising delle foto dello smartphone (idea
dell'\en{autoencoder}), analisi dei log (una mossa diversa da quella attesa è sospetta), malware,
immagini false con GAN (generatore contro discriminatore). Benefici: nessun costo di etichettatura,
feature automatiche, robustezza alle piccole variazioni. Problemi: costo computazionale enorme,
servono dataset enormi \textbf{e di qualità} (cento libri pieni di errori non insegnano l'italiano),
bias appresi, valutazione difficile.

\subsection{Reinforcement learning}
Un \textbf{agente} agisce in un \textbf{ambiente} (stato, azioni) per massimizzare una
\textbf{reward} a lungo termine, per tentativi ed errori; i dati li genera l'agente stesso. Esempi:
regole di firewall o soglie di IDS adattive, simulazione di attacchi, controllo di una lampadina smart
che accetta comandi in chiaro. Problemi: costo, progettazione della reward (se è sbagliata l'agente
smette di esplorare), apprendimento lento e instabile, scarsa interpretabilità.

\subsection{Come si usano i modelli}
\begin{itemize}
  \item \textbf{Ensemble}: più modelli addestrati su dati, sottoinsiemi o feature diversi, con voto di
    maggioranza o media (la random forest). Nel \textbf{model stacking} un \emph{meta-learner} impara
    di quale modello base fidarsi e quando. Più robusto, ma costoso e poco utile se i modelli base sono
    correlati.
  \item \textbf{Federated learning}: ogni dispositivo addestra sui propri dati e invia al server
    \textbf{il modello, non i dati}; il modello aggregato torna ai dispositivi (esempio dello spam
    filter addestrato sui telefoni degli studenti). Privacy e GDPR, scalabilità; ma costo sull'edge,
    dati eterogenei e \textbf{model poisoning}.
  \item \textbf{Transfer learning}: si riusa un modello pre-addestrato e se ne ri-addestrano di solito
    solo gli ultimi layer (\en{fine-tuning}) con pochi dati. Esempio: STAMINA di Intel, che trasforma
    i binari di malware in immagini in scala di grigi e fa fine-tuning di ResNet-50
    (capitolo~\ref{cap:cnn-gnn}). Rischi: \emph{domain mismatch} e bias ereditati.
  \item \textbf{Modelli generativi}: campionano nuovi esempi dalla distribuzione appresa (dati
    sintetici, varianti di malware, agenti che cercano vulnerabilità). Problemi: costo, valutazione
    difficile, bias, uso improprio (deepfake, phishing), privacy dei dati inviati ai chatbot, questioni
    etiche, legali e sociali.
\end{itemize}

\begin{esame}
Commentando l'uso di ChatGPT e del \en{vibe coding}, il docente ricorda che allo scritto c'è una
domanda di codice da risolvere senza strumenti: bisogna saper capire da soli se un codice è giusto o
sbagliato, e perché.
\end{esame}

\subsection{AI e cybersecurity}
Tre famiglie. \textbf{AI per la difesa}: IDS e analisi di malware senza firme statiche, novelty
detection per gli zero-day, SIEM e SOAR potenziati, predictive security e analisi comportamentale
(secondo le slide oltre il 90\% dei SOC dipende ancora da processi manuali, e chi usa AI e
automazione risparmia in media 2,22 milioni di dollari per breach). \textbf{AI per l'attacco}:
social engineering con deepfake e phishing generato e tradotto, scoperta automatica di vulnerabilità,
evasione; l'AI generativa abbassa la barriera d'ingresso. \textbf{Adversarial machine learning}:
\emph{data poisoning}, \emph{adversarial examples} (rumore impercettibile sull'input, non sul modello:
la demo del cane classificato come ``hot dog'' o ``chiesa''), aggiramento dei \emph{guardrail} degli
LLM (l'``exploit della nonna'', jailbreak a pochi giorni da ogni rilascio; gli agenti con tool
allargano la superficie d'attacco) e \emph{model inversion}, che ricostruisce i dati di training dal
modello, anche nel federated learning. Approfondimento nel capitolo~\ref{cap:ai-cyber}.

\section{Domande aperte e punti da verificare}
\begin{verifica}
\begin{itemize}
  \item Parlando dei report scritti con ChatGPT il docente accenna a una possibile chiamata
    ``all'orale''. Il formato d'esame presentato il 21/09 prevede solo scritto e report: chiarire se
    esiste un orale di verifica sui report sospetti.
\end{itemize}
\end{verifica}

% =====================================================================
\lezione{28/09/2026}{Il perceptron e la struttura di una rete neurale}
  {Prof.~Danilo Giordano}
  {input, hidden e output layer; perceptron, pesi e bias; funzioni di attivazione;
   layer fully connected; necessità della non linearità; teorema di approssimazione universale;
   profondità contro larghezza; logit, sigmoid e softmax}
\label{lez:2026-09-28}
% =====================================================================

\begin{nota}[Fonte]
Trascrizione automatica della registrazione, revisionata. I contenuti seguono le slide
\emph{2-foundament-deep-learning} (pp.~1--18); numeri e formule sono ripresi dalle slide.
\end{nota}

\section{Sintesi}
Inizia il secondo deck: che cosa c'è ``sotto il cofano'' di una rete neurale e che cosa sono
davvero i trilioni di parametri. Il docente parte dal singolo neurone (il perceptron: somma pesata più
funzione di attivazione), mostra perché servono il bias e le attivazioni non lineari, impila i
neuroni in layer e layer in reti, e arriva al teorema di approssimazione universale e alla scelta di
reti profonde invece che larghe. Chiude con le attivazioni dell'ultimo layer (sigmoid, softmax,
identità). Approfondimento nel capitolo~\ref{cap:reti-neurali}.

\section{Avvisi pratici}
\begin{itemize}
  \item La settimana successiva sono previste lezioni anche il \textbf{martedì} (oltre a lunedì),
    dedicate alle reti feed-forward anche dal lato del codice; il \textbf{mercoledì} il primo
    laboratorio. Annuncio ufficiale sul portale.
  \item Su Moodle verrà aperta la sezione per la creazione dei gruppi: dalla settimana successiva si
    lavora in gruppo.
\end{itemize}

\section{Contenuti}

\subsection{La struttura di una rete}
Input layer (un neurone per feature, nel caso tabellare), hidden layer (dove la rete ``ragiona''; con
tre o più si parla di deep learning) e output layer (la decisione finale). Le reti accettano solo
numeri: testo, grafi e immagini vanno convertiti (un pixel sono tre valori RGB). La conversione è
in genere standard; le parti difficili sono il preprocessing e la scelta dell'architettura, da trovare
sperimentalmente (\emph{model engineering}).

\subsection{Il perceptron}
Il neurone pesa le feature, ne fa la somma e applica un'attivazione: $y = f(\mathbf{w}^T\mathbf{x})$, con
$x_0 = 1$. I \textbf{pesi} sono i parametri della rete: indicano l'importanza di ogni feature,
partono casuali e si imparano dai dati; le $x_i$ invece non si toccano. Il \textbf{bias} $w_0$ non
dipende da alcuna feature e serve a traslare la funzione: senza bias, con attivazione lineare, tutte
le rette di separazione passano per l'origine. Nel caso 2D i pesi inclinano e spostano il confine,
ma l'output grezzo non è limitato: per ottenere una probabilità serve un'attivazione limitata.

\subsection{Funzioni di attivazione}
Servono a imporre proprietà all'output, introdurre non linearità, decidere quando un neurone è
attivo e migliorare l'ottimizzazione. \textbf{Sigmoid}: output in $[0,1]$, adatta alla classificazione
binaria, ma negli hidden layer causa \en{vanishing gradient}. \textbf{Tanh}: limitata in $[-1,1]$, la
si ritrova nelle RNN. \textbf{ReLU}: zero per input negativi (neurone inattivo, quindi sparsità),
illimitata sopra; è la più importante negli hidden layer perché mitiga il vanishing gradient, ma
moltiplicando molti valori grandi il gradiente può \textbf{esplodere}. \textbf{Leaky ReLU}: piccola
pendenza anche per input negativi. Il docente collega la sparsità agli LLM, che non usano tutti i
parametri per ogni richiesta. Tabella completa nel capitolo~\ref{cap:reti-neurali}.

\begin{esame}
Il docente avverte che le formule di questo deck \textbf{non verranno chieste} all'esame: non
importa ricordare la forma esatta della sigmoid, importa capire \textbf{il concetto} che c'è dietro
(a cosa serve ogni attivazione, perché la non linearità, perché profondità).
\end{esame}

\subsection{Margine, layer e non linearità}
Tra due confini lineari che separano bene il training è preferibile quello con \textbf{margine}
maggiore, idea che ritorna nelle SVM. Due perceptron sullo stesso input formano un \textbf{layer
fully connected} ($\mathbf{y} = f(W\mathbf{x})$, 6 parametri nell'esempio a due input e due neuroni);
un secondo layer combina gli output del primo. Una rete è quindi una sequenza di prodotti tra matrici
alternati ad attivazioni. Ogni layer produce una rappresentazione più astratta del dato. Senza
attivazioni non lineari, però, mille layer collassano in \textbf{un solo layer lineare}, e il mondo non è
lineare (basta guardare la distribuzione dei voti di un esame). L'attivazione non è la stessa ovunque:
ReLU negli hidden layer, sigmoid o softmax in uscita, perché i due punti della rete risolvono problemi
diversi.

\subsection{Approssimazione universale: profonde, non larghe}
Il teorema di approssimazione universale (Cybenko e Hornik, 1989) garantisce che una rete con un solo
hidden layer e un numero finito di neuroni approssima qualunque funzione continua su un compatto. Con
la ReLU ogni neurone aggiunge un punto di rottura: nell'esempio delle slide la curva è approssimata
sempre meglio con 4, 9 e 19 neuroni. Il teorema però non dice \emph{quanti} neuroni servono. Il numero
di regioni lineari cresce in fretta con neuroni e feature: con 500 neuroni e 100 feature si superano
le $10^{107}$ regioni con appena 51.001 parametri.

\begin{esame}
Messaggio su cui il docente insiste: preferire reti \textbf{più profonde e meno larghe}. I layer,
separati da non linearità, combinano la loro complessità e costruiscono rappresentazioni sempre più
astratte, il che migliora la generalizzazione. Attenzione però: una rete troppo complessa va in
overfitting (capitolo~\ref{cap:regolarizzazione}).
\end{esame}

\subsection{Logit e attivazioni in uscita}
L'output prima dell'ultima attivazione è il \textbf{logit}, non normalizzato. In classificazione
binaria basta un logit e la \textbf{sigmoid}; nel multi-classe si usa un neurone per classe (ne
basterebbero $C-1$) e la \textbf{softmax}, che produce probabilità con somma 1 e amplifica le piccole
differenze tra logit (esempio: CNN per la rappresentazione dell'immagine, FFNN per classificare
cane, gatto, uccello o tartaruga). In regressione l'uscita è l'identità. In teoria in classificazione
l'attivazione finale c'è sempre; nel software può mancare perché è inclusa altrove, come si vedrà con
la loss di PyTorch (capitolo~\ref{cap:ffnn}).

\section{Domande aperte e punti da verificare}
\begin{verifica}
\begin{itemize}
  \item Soglia del deep learning: la slide parla di ``3 o più'' hidden layer, mentre in aula la frase
    (``more [than] three'', con un esempio a quattro layer) è ambigua. Vale la definizione della slide.
\end{itemize}
\end{verifica}

% =====================================================================
\lezione{30/09/2026}{Addestrare una rete: ottimizzazione, loss e valutazione}
  {Prof.~Danilo Giordano}
  {forward pass, loss ed epoch; mini-batch; gradient descent e SGD; learning rate e scheduler;
   momentum, normalizzazione e Adam; backpropagation e regola della catena; BCE e cross-entropy;
   vanishing gradient; procedura di training; metriche; underfitting, overfitting e regolarizzazione}
\label{lez:2026-09-30}
% =====================================================================

\begin{nota}[Fonte]
Trascrizione automatica della registrazione, revisionata. I contenuti seguono le slide
\emph{2-foundament-deep-learning} (pp.~19--51); formule e numeri sono ripresi dalle slide.
\end{nota}

\section{Sintesi}
Lezione densa che completa il deck sui fondamenti: come si passa da pesi casuali a pesi utili. Il
docente costruisce il ciclo di addestramento (forward pass, loss, backpropagation, aggiornamento),
confronta gli ottimizzatori dal gradient descent ad Adam, spiega la backpropagation come regola della
catena e mostra con numeri il vanishing gradient. Poi passa alla pratica: procedura di training,
dimensione del batch, scelta della loss, metriche di valutazione con dati sbilanciati e diagnosi di
underfitting e overfitting, con una prima panoramica della regolarizzazione. Approfondimenti nei
capitoli~\ref{cap:addestramento} e~\ref{cap:regolarizzazione}.

\section{Avvisi pratici}
\begin{itemize}
  \item Entro un giorno viene pubblicato su Moodle il modulo per la creazione dei gruppi. Il primo
    laboratorio è il mercoledì successivo: conviene accordarsi almeno informalmente prima. Termine per
    registrare ufficialmente i gruppi: \textbf{fine ottobre}.
\end{itemize}

\section{Contenuti}

\subsection{Il ciclo di addestramento}
Si indica con $\phi$ l'insieme di pesi e bias. A ogni passo: forward pass, calcolo della
\textbf{loss} tra predizione e label, backpropagation dell'errore, aggiornamento dei pesi secondo un
\textbf{ottimizzatore}. Un'\textbf{epoch} è un passaggio dell'intero training set. Si aggiorna dopo ogni
\textbf{mini-batch}, campionato a caso \textbf{senza reinserimento} (ogni punto conta una volta per
epoch): aggiornamenti più frequenti, meno memoria, migliore generalizzazione. La loss (per esempio
MSE) dipende solo da $\phi$, perché dati e label sono fissi; una loss pari a zero deve insospettire.
Con un trilione di parametri non c'è soluzione in forma chiusa: si parte a caso e ci si muove secondo
una policy.

\subsection{Ottimizzatori}
\begin{itemize}
  \item \textbf{Gradient descent}: passo nella direzione opposta al gradiente, con visione solo
    locale. Funziona con loss convesse, ma si blocca nei minimi locali, è sensibile
    all'inizializzazione e usa tutti i dati. Il docente racconta un suo progetto in cui una piccola
    rete cambiava drasticamente prestazioni al solo variare dell'inizializzazione: si scelse una SVM,
    un po' peggiore ma stabile.
  \item \textbf{SGD}: gradiente su un mini-batch. Traiettoria rumorosa (``cammina come un
    ubriaco''), che può anche risalire e proprio per questo uscire dai minimi locali; più veloce ed
    economica, generalizza bene, ma non converge in senso classico.
  \item \textbf{Learning rate} $\alpha$: il gradiente dà la direzione, $\alpha$ la lunghezza del passo
    (non la dimensione del batch). Con uno \emph{scheduler} si parte alto per esplorare e si riduce per
    rifinire.
  \item \textbf{Momentum}: media pesata con i passi precedenti ($\beta$), traiettoria più liscia;
    $\beta$ troppo basso o troppo alto rallenta comunque la convergenza.
  \item \textbf{Normalizzazione del gradiente}: stesso passo in tutte le direzioni, al prezzo di
    oscillazioni vicino al minimo.
  \item \textbf{Adam}: momentum sul gradiente e sul suo quadrato, con \emph{bias correction} nei primi
    passi. È l'ottimizzatore che si userà di più; $\alpha$, $\beta$, $\gamma$ restano iperparametri.
\end{itemize}
Ottimizzatore, batch size, learning rate e momentum sono \textbf{iperparametri}: si scelgono
addestrando molti modelli (\en{grid search}), ``più arte che scienza''. In aula si sono viste due demo
interattive sull'effetto degli iperparametri.

\subsection{Backpropagation e vanishing gradient}
La backpropagation è la \textbf{regola della catena} applicata all'indietro sul grafo
computazionale: il gradiente di un layer riusa quello del layer successivo moltiplicato per il
gradiente locale. Il docente la svolge passo per passo sull'esempio $\hat y = \phi_1\phi_2 x$ con loss
quadratica. Per la classificazione binaria la loss è la \textbf{binary cross-entropy}, in cui la label
fa da selettore tra $-\log\hat y$ e $-\log(1-\hat y)$; nel multi-classe si usa la cross-entropy.
Poiché i gradienti si moltiplicano layer dopo layer, attivazioni limitate come la sigmoid li fanno
\textbf{svanire}: nell'esempio a quattro layer il gradiente cala di circa tre ordini di grandezza fino
al primo hidden layer, mentre con la ReLU resta stabile. I primi layer restano così quasi uguali
all'inizializzazione.

\begin{esame}
Anche qui il docente precisa che all'esame \textbf{non si chiede di calcolare gradienti} né di
svolgere la matematica della backpropagation: si chiede di capire il concetto (che cosa fa, perché
serve, che cosa causa il vanishing gradient).
\end{esame}

\subsection{L'addestramento in pratica}
I passi sono: definire problema e dati (obiettivo e quindi loss, preprocessing, split, batch size),
scegliere l'architettura, scegliere ottimizzatore e loss, addestrare e validare. Batch di dimensione
potenza di 2: 1 è \emph{streaming} e molto rumoroso, 16--64 va bene per dataset piccoli, 128--512 è il
compromesso più comune, da 1024 in su solo con dataset molto grandi. Loss tipiche: MSE e MAE (meno
sensibile agli outlier) per la regressione, BCE per il binario, cross-entropy per il multi-classe.
Dati di qualità restano fondamentali: ``garbage in, garbage out''.

\subsection{Metriche}
La loss guida l'addestramento, ma il modello migliore si sceglie con le \textbf{metriche}: accuracy,
e per classe precision, recall e F1. Esempio del classificatore che risponde sempre ``goodware'' su
9.900 goodware e 100 malware: accuracy del 99\% e recall sui malware pari a zero. Se invece risponde
sempre ``malware'', la recall sale a 1 ma la precision crolla all'1\% e nessuno si fida più degli
allarmi. Dettagli nel capitolo~\ref{cap:addestramento}.

\begin{esame}
Con dati sbilanciati, il caso tipico in cybersecurity, l'accuracy inganna. Si sceglie il compromesso
tra precision e recall in base al costo degli errori: per la classe malware recall alta con una
precision almeno decente; in caso di dubbio, la F1.
\end{esame}

\subsection{Underfitting, overfitting e regolarizzazione}
Si leggono le \textbf{learning curve} di training e validation; sul validation set la loss si calcola
senza mai aggiornare i pesi. \textbf{Underfitting}: loss alta anche in training (attenzione alla
scala); si rimedia con preprocessing migliore o modello più capace. \textbf{Overfitting}: la loss di
training continua a scendere mentre quella di validation risale. Rimedi, con un termine di penalità
nella loss o agendo sul processo: \textbf{L1} (pesi a zero, sparsità) e \textbf{L2} (pesi piccoli,
\code{weight\_decay}) sui pesi; \textbf{dropout} sull'architettura; \textbf{early stopping} sul
processo, tenendo il modello all'epoch $t$ prima che la validation peggiori; \textbf{batch
normalization} sui dati (lo stesso esempio appare diverso in batch diversi, di solito prima
dell'attivazione); \emph{noise injection}, che nel corso non si userà. Approfondimento nel
capitolo~\ref{cap:regolarizzazione}.

\section{Domande aperte e punti da verificare}
\begin{verifica}
\begin{itemize}
  \item Scadenza dei gruppi: il 21/09 erano da formare prima del primo lab, oggi il termine indicato
    è la fine di ottobre (``circa un mese''). Controllare la data esatta sul modulo Moodle.
  \item In aula L1 e L2 sono descritte come ``aggiungere un po' di rumore'' all'aggiornamento dei
    pesi: tecnicamente sono penalità deterministiche nella loss; la vera aggiunta di rumore è la
    \emph{noise injection}.
\end{itemize}
\end{verifica}

% =====================================================================
\lezione{05/10/2026}{Data handling: preprocessing, trasformazione e split}
  {Prof.~Danilo Giordano}
  {data handling nella pipeline deep; rumore, outlier, ridondanza, valori mancanti; one-hot, label
   encoding, embedding; normalizzazione e standardizzazione; hold-out, k-fold, leave-one-out;
   sbilanciamento, costo degli errori, data leakage; progettazione di una FFNN; tensori e DataLoader}
\label{lez:2026-10-05}
% =====================================================================

\begin{nota}[Fonte]
Trascrizione automatica della registrazione, revisionata. Numeri, tabelle e best practice sono
ripresi dalle slide \emph{3-deep-learning-methodologies} (pp.~1--22), confrontate con quanto detto
in aula.
\end{nota}

\section{Sintesi}
Si apre il blocco sulle metodologie di deep learning. Il messaggio di fondo è che il deep learning
\textbf{non fa magie sui dati}: anche in una pipeline deep il \en{data handling} (split,
preprocessing, trasformazione) prende il 50--60\% del tempo. La lezione passa in rassegna la pulizia
dei dati, la codifica delle feature categoriche e numeriche, i protocolli di split e i quattro
problemi da gestire nello split (distribuzione delle classi, costo degli errori, dimensione del
dataset, data leakage). Chiude con una panoramica delle architetture, la checklist per progettare
una FFNN e i primi strumenti PyTorch per caricare i dati. Approfondimento nel
capitolo~\ref{cap:ffnn}; per gli embedding il capitolo~\ref{cap:embedding}.

\section{Avvisi pratici}
\begin{itemize}
  \item La lezione di \textbf{martedì 6 ottobre} (slot di recupero) è confermata, come annunciato
    sul portale.
  \item Il \textbf{Lab~1} (PyTorch e FFNN) ha 3~ore in aula; serve a prendere confidenza con
    ottimizzatori, scheduler, momentum ecc. Il report consegnato entro il \textbf{1° dicembre}, se
    ragionevole, vale \textbf{+1 punto}; un report di due pagine vuote va riconsegnato
    (capitolo~\ref{cap:lab1}).
  \item Per mercoledì: dare un'occhiata alla documentazione di \textbf{Pandas} e PyTorch.
  \item Da questa parte in poi le lezioni alternano circa un'ora e mezza di teoria a esempi di codice.
\end{itemize}

\section{Contenuti}

\subsection{Il deep learning non basta: data handling}
La sequenza corretta è \textbf{split} $\to$ \textbf{preprocessing} $\to$ \textbf{trasformazione}
$\to$ \en{model engineering}. Il docente avverte che in azienda spesso non si potranno caricare i
dati su ChatGPT o strumenti simili (dati sensibili): bisogna saperli trattare da soli.

\subsection{Preprocessing}
\begin{itemize}
  \item \textbf{Rumore}: alcune architetture lo assorbono; troppo rumore porta a underfitting.
  \item \textbf{Outlier}: esempio del reddito normalizzato tra 0 e 1 con Elon Musk come massimo: tutti
    gli altri finiscono schiacciati vicino a 0, indistinguibili da un bambino senza reddito.
  \item \textbf{Ridondanza}: feature \en{proxy} l'una dell'altra e duplicati (questi ultimi danno
    peso artificiale a certi punti).
  \item \textbf{Valori mancanti}: una rete non accetta un input incompleto. Quattro opzioni: eliminare
    i punti (attenzione a \emph{quali}: potrebbero essere proprio la classe malevola), stimare il
    valore, eliminare la feature, sostituire con tutti i valori possibili pesati per probabilità.
\end{itemize}

\subsection{Trasformazione delle feature}
\begin{itemize}
  \item \textbf{One-hot}: un neurone per valore; niente ordini artificiali, ma esplode ad alta
    cardinalità (oltre 171.000 parole inglesi).
  \item \textbf{Label encoding}: un neurone, un intero per categoria; adatto alle feature ordinali e
    alla \textbf{label} (la cross-entropy non calcola distanze tra classi). Esempio delle
    \textbf{porte}: 80, 81, 8080 sono numeri ma vanno trattate come \textbf{categoriche}.
  \item \textbf{Embedding}: vettore denso in uno spazio latente di dimensione fissa (es.\ 32, 64,
    128); concetti simili finiscono vicini (\en{man}/\en{woman} come \en{king}/\en{queen}); richiede
    abbastanza dati.
  \item \textbf{Numeriche}: normalizzazione min-max $z=(x-x_{min})/(x_{max}-x_{min})$, con intervallo
    fisso ma sensibile agli outlier; standardizzazione $z=(x-\mu)/\sigma$, robusta agli outlier ma
    senza limiti precisi.
\end{itemize}

\subsection{Split e protocolli di valutazione}
Training per i pesi, validation per gli iperparametri, test per verificare la generalità (se il
test va male si riparte da capo, non si ritoccano gli iperparametri). Campionamento
\textbf{stratificato} (stesse proporzioni di classe in ogni partizione) o \textbf{bootstrap} (con
reinserimento, per dataset piccoli). Protocolli:
\begin{itemize}
  \item \textbf{hold-out} (es.\ 70/20/10), per dataset grandi;
  \item \textbf{k-fold cross-validation}: $k$ addestramenti con gli stessi iperparametri; se le $k$
    prestazioni sono simili, gli iperparametri vanno bene. Il modello finale si \textbf{ri-addestra su
    tutti i fold}, non si sceglie il fold ``migliore'' (sarebbe fortuna);
  \item \textbf{leave-one-out} ($k=n$), solo per dataset molto piccoli.
\end{itemize}

\subsection{I quattro problemi dello split}
\begin{enumerate}
  \item \textbf{Classi sbilanciate}: oversampling (SMOTE, ADASYN), undersampling (anche
    \en{cluster-based}), class weights, focal loss e soprattutto la metrica giusta: media
    \textbf{macro} invece di \en{micro}.
  \item \textbf{Costo degli errori}: un falso negativo su un malware costa più di un falso positivo;
    loss pesata e metriche sulla classe rilevante.
  \item \textbf{Dataset piccolo}: cross-validation, data augmentation, modelli pre-addestrati con
    fine-tuning, approccio semi-supervised.
  \item \textbf{Data leakage}: correlazione temporale (interpolare $t$ conoscendo $t-1$ e $t+1$ non è
    prevedere), feature che fanno da scorciatoia (tutti i DoS sulla porta 80), identificatori (la
    matricola per predire il voto). Con un andamento temporale si usa uno split basato sul tempo
    (\en{sliding} o \en{expanding window}).
\end{enumerate}

\begin{esame}
Domanda classica: \textbf{si può normalizzare (o codificare) usando tutto il dataset? No}, sarebbe
data leakage. Prima si divide, poi si fa \code{fit} delle trasformazioni \textbf{solo sul training}
e solo \code{transform} su validation e test (un 150 in validation con massimo di training 100
diventa 1,5, ed è giusto così). Il docente insiste anche su: usare la metrica giusta con classi
sbilanciate e non usare uno split casuale quando l'ordine temporale conta.
\end{esame}

\subsection{Architetture e progettazione di una FFNN}
Panoramica: FFNN per dati tabellari (e come classificatore finale di altre reti), RNN per sequenze,
CNN per immagini, GNN per grafi. Checklist per una FFNN: tutto in forma numerica (attenzione alle
feature numeriche ma categoriche); input layer dimensionato dal preprocessing (2 numeriche + una
categorica con 10 valori in one-hot = 12 neuroni); pochi hidden layer, prima in profondità che in
larghezza, neuroni in potenze di 2; output in base al compito (in PyTorch con
\code{CrossEntropyLoss} niente softmax nel modello); inizializzazione automatica con seed fissati;
poi mini-batch, loss, ottimizzatore, learning rate (partire da 0,001 con scheduler) e
regolarizzazione (dropout 0,2--0,5, L2 0,01) in caso di overfitting. Tabella completa delle best
practice nel capitolo~\ref{cap:ffnn}.

\subsection{Caricare i dati in PyTorch}
Il \textbf{tensore} è una matrice multidimensionale di un solo tipo (\code{dtype}); gli zeri
torneranno utili per le maschere in NLP. \code{TensorDataset(X, y)} unisce feature e label;
\code{DataLoader(..., batch\_size, shuffle)} crea i mini-batch. Con \code{shuffle=True} l'ordine
cambia a ogni epoch (migliore generalizzazione), ma ogni punto resta visto una sola volta per epoch.
Con dati sequenziali non si mescolano i punti della stessa serie.

\begin{nota}[Domanda lasciata in sospeso]
Serve mescolare anche validation e test? Il docente rimanda la risposta alla lezione successiva
(\rif{2026-10-06}{06/10}): no, perché lì i pesi non si aggiornano.
\end{nota}

\section{Domande aperte e punti da verificare}
\begin{verifica}
\begin{itemize}
  \item Nella slide (e a voce) $\sigma$ della standardizzazione è chiamata ``varianza'': è la
    \textbf{deviazione standard}.
  \item Una versione della slide sullo sbilanciamento dice ``\en{micro avg instead of macro}'': la
    versione completa corregge in ``\en{macro avg instead of micro}'', coerente con quanto detto in
    aula.
\end{itemize}
\end{verifica}

% =====================================================================
\lezione{06/10/2026}{Feed-forward neural network in PyTorch}
  {Prof.~Danilo Giordano}
  {shuffle in validation e test; classe \code{nn.Module}; logit e \code{CrossEntropyLoss}; training
   loop con e senza mini-batch; \code{model.eval()}; early stopping; dropout, batch norm e weight
   decay nel codice; inizializzazione dei pesi; GPU e riproducibilità; TensorFlow Playground;
   notebook d'esempio}
\label{lez:2026-10-06}
% =====================================================================

\begin{nota}[Fonte]
Trascrizione automatica della registrazione (lezione di recupero del martedì), revisionata. Codice e
parametri sono ripresi dalle slide \emph{3-deep-learning-methodologies} (pp.~23--33) e dal notebook
FFNN mostrato in aula.
\end{nota}

\section{Sintesi}
Lezione pratica che prepara il Lab~1. Dopo un breve ripasso della checklist per progettare una FFNN,
il docente risponde alla domanda lasciata in sospeso (lo shuffle in validation e test) e mostra, pezzo
per pezzo, come si scrive una FFNN in PyTorch: definizione della rete come classe, loss e
ottimizzatore, training loop prima senza e poi con mini-batch, early stopping, regolarizzazione,
inizializzazione, GPU e seed. Seguono una prova su TensorFlow Playground e il notebook d'esempio con
una pipeline completa su un dataset sintetico. Approfondimento nei capitoli~\ref{cap:ffnn},
\ref{cap:addestramento} e \ref{cap:regolarizzazione}.

\section{Avvisi pratici}
\begin{itemize}
  \item \textbf{Lab~1 mercoledì 7 ottobre}: dataset tabellare piccolo e semplice, solo FFNN.
  \item Ambiente: PyTorch su \textbf{Google Colab} (consigliato, non obbligatorio: CPU e GPU, notebook
    condivisibili nel gruppo); vanno bene anche un Jupyter locale o Anaconda. La GPU servirà
    soprattutto dal Lab~2.
  \item Lab~2 userà lo stesso dato come tabella, serie temporale e grafo; seguiranno anomaly
    detection e NLP.
  \item Il notebook d'esempio mostra \textbf{come usare la libreria}, non come risolvere un problema
    reale: non va copiato così com'è nei lab.
\end{itemize}

\section{Contenuti}

\subsection{Shuffle in validation e test}
No: la stocasticità serve all'SGD, cioè solo al training. In validation e test i pesi non si
aggiornano, quindi l'ordine non cambia le predizioni.

\subsection{Rete, loss e ottimizzatore}
Una rete è una classe che eredita da \code{nn.Module}: \code{\_\_init\_\_} descrive i layer,
\code{forward} come li attraversano i dati. Una rete fatta solo di \code{nn.Linear} \textbf{non ha
non linearità} (equivale a un solo layer lineare): serve un'attivazione come \code{nn.ReLU}.
L'output sono i \textbf{logit}, non normalizzati; la loss (\code{criterion}, es.\
\code{nn.CrossEntropyLoss}, eventualmente con \code{weight}) e l'ottimizzatore (es.\ \code{AdamW} con
learning rate e \code{weight\_decay}) completano il quadro. Appena creato, il modello ha pesi
casuali: non è ancora un modello.

\begin{esame}
\begin{itemize}
  \item Nell'ultimo layer in teoria serve un'attivazione (softmax nel multi-classe), ma
    \code{nn.CrossEntropyLoss} si aspetta i \textbf{logit} e applica la softmax al suo interno:
    nel modello non va messa. Prima di usare una loss si legge la documentazione.
  \item \code{model.eval()} \textbf{deve esserci} in validation e test (il docente parla di ``errore
    enorme''): spegne il dropout e fa usare alla batch norm le statistiche memorizzate.
    \code{torch.no\_grad()} invece disattiva solo il calcolo dei gradienti. Non sono la stessa cosa.
\end{itemize}
\end{esame}

\subsection{Il training loop}
Per ogni epoch: \code{model.train()}, \code{zero\_grad}, forward, loss, \code{backward},
\code{step}, salvataggio della loss; poi \code{model.eval()} e validation sotto
\code{torch.no\_grad()}. Con i mini-batch \code{zero\_grad/backward/step} stanno
\textbf{dentro} il ciclo sui batch (i pesi si aggiornano a ogni batch) e la loss di epoch è la
\textbf{media pesata} per la dimensione dei batch (con 90 punti e batch da 64 l'ultimo batch ha 26
punti). Nella prima learning curve d'esempio la loss scende ancora senza plateau: il training è
stato fermato troppo presto. La loss dice se il modello converge; per giudicarlo servono le metriche
(capitolo~\ref{cap:addestramento}).

\subsection{Early stopping e regolarizzazione}
Early stopping con \code{min\_delta} (miglioramento minimo significativo) e \code{patience} (epoch
consecutive senza miglioramento); si salva lo stato del \textbf{modello migliore} e lo si ricarica alla
fine, invece di restituire quello dell'ultima epoch. Il dropout si definisce come un layer: a ogni
passo spegne neuroni diversi scelti a caso (non li ``congela'' una volta per tutte) e in validation
sono tutti attivi; valori ragionevoli 0,2--0,5. La batch norm è un layer (\code{BatchNorm1d}); la
regolarizzazione L2 è il \code{weight\_decay} dell'ottimizzatore. Dettagli nel
capitolo~\ref{cap:regolarizzazione}.

\subsection{Inizializzazione, GPU e riproducibilità}
L'inizializzazione automatica sceglie il metodo in base all'attivazione; nell'esempio manuale con
pesi dell'ordine di $10^{-5}$ si va incontro al vanishing gradient: si può fare, ma porta più
problemi che benefici. Modello e dati vanno sullo stesso \code{device} (\code{cuda} se disponibile) e
riportati su CPU quando servono fuori da PyTorch; la GPU va usata solo se serve. Per la
riproducibilità si fissano i seed di tutto ciò che è casuale; su hardware diverso i risultati possono
comunque cambiare un po', e se dipendono molto dal seed il modello è instabile.

\subsection{TensorFlow Playground}
Con buone feature basta anche un solo layer; un learning rate di 10 rende il training rumorosissimo,
uno minuscolo lentissimo; con attivazione lineare il problema non lineare non si risolve; una rete
profonda con sigmoid resta bloccata. Morale: basta una scelta sbagliata per non imparare nulla.

\subsection{Il notebook d'esempio}
Dataset sintetico con 2 feature e 3 classi, sporcato con outlier, valori mancanti e duplicati. Passi:
rimozione di duplicati e mancanti tenendo traccia di quante righe si perdono; \code{LabelEncoder}
per la label; visualizzazione; split 60/20/20 con \code{train\_test\_split}; \code{StandardScaler}
fittato solo sul training; modello lineare (loss quasi ferma, come previsto), poi non lineare (3
hidden layer da 128 con ReLU), mini-batch, early stopping, dropout, batch norm e weight decay. Con
dropout la validation loss può stare sotto quella di training. Sulla curva molto rumorosa dell'ultima
variante il docente avverte: una validation loss così in un report fa pensare che la soluzione non sia
stata valutata.

\subsection{Quali iperparametri toccare}
Non si sa a priori: dipende da dataset, split e architettura. Si lavora soprattutto sul
\textbf{learning rate} (partendo dall'ordine di $10^{-3}$, eventualmente con scheduler); momentum e
\code{betas} si possono lasciare di default; il weight decay si tocca in caso di overfitting. Nel
Lab~1 si chiede proprio di sperimentare per farsi un'idea.

\section{Domande aperte e punti da verificare}
\begin{verifica}
\begin{itemize}
  \item Il docente dice che senza \code{model.eval()} si commette un errore grave; non dice
    esattamente cosa succede. Integrazione dei redattori: i pesi non cambiano senza
    \code{optimizer.step()}, ma il dropout resta attivo e la batch norm aggiorna le sue statistiche
    con i dati di validation (capitolo~\ref{cap:regolarizzazione}).
\end{itemize}
\end{verifica}

% =====================================================================
\lezione{12/10/2026}{Recurrent neural network e LSTM}
  {Prof.~Danilo Giordano}
  {perché servono le RNN; hidden state; varianti input/output; RNN mono e bidirezionale; cella RNN e
   classification head; relazione di ricorrenza; BPTT; vanishing ed exploding gradient; dipendenze a
   lungo termine; LSTM; dati in forma di tensore 3D; progettazione di una RNN; da tabella a sequenze}
\label{lez:2026-10-12}
% =====================================================================

\begin{nota}[Fonte]
Trascrizione automatica della registrazione, revisionata. Formule, tabelle e best practice sono
riprese dalle slide \emph{3-deep-learning-methodologies} (pp.~34--71) e dal notebook
\emph{RNN-Pollution-Multivariate}.
\end{nota}

\section{Sintesi}
Dopo una battuta sul Lab~1 (una guida passo passo per imparare a usare una rete, più che una vera
pipeline di data science), si entra nelle architetture deep vere e proprie con le
\textbf{recurrent neural network}. Una FFNN tratta ogni punto in modo indipendente; una RNN mantiene
uno \textbf{stato} che evolve nel tempo e riassume la serie. La lezione costruisce la RNN a partire
dalla FFNN, introduce la \en{backpropagation through time} e i suoi problemi di gradiente, motiva e
descrive l'LSTM, poi passa alla preparazione dei dati (tensore 3D, lunghezze diverse) e alla
progettazione. Approfondimento nel capitolo~\ref{cap:rnn}.

\section{Avvisi pratici}
\begin{itemize}
  \item Il Lab~2 si concentra sul caso \textbf{many-to-one} (una serie in input, una label in output).
  \item L'esempio di codice sulle sequenze prosegue mercoledì (\rif{2026-10-14}{14/10}).
\end{itemize}

\section{Contenuti}

\subsection{Perché le RNN}
Da un solo punto non si può predire la traiettoria: manca la storia. Le RNN gestiscono input
sequenziali di lunghezza variabile e mantengono uno \textbf{hidden state} che rappresenta la serie
fino a $t-1$ (``\en{my brother plays}\ldots'': il contesto suggerisce cosa viene dopo; in modo
semplificatissimo è l'idea della predizione della parola successiva). Applicazioni: traduzione,
previsione di serie (scostamenti forti dalla previsione possono indicare un'anomalia), classificazione
di serie, riconoscimento vocale, \en{part-of-speech tagging}.

\subsection{Varianti e architetture}
Quattro configurazioni: \textbf{one-to-one}, \textbf{many-to-one} (sentiment, sequenza di log per
rilevare un attacco), \textbf{one-to-many} (image captioning), \textbf{many-to-many} (traduzione).
La RNN \textbf{mono-direzionale} va bene per il tempo reale ma fatica con le dipendenze lunghe; la
\textbf{bidirezionale} usa due RNN con pesi diversi in direzioni opposte, cattura passato e futuro ma
costa di più e richiede la sequenza completa (es.\ un'email da classificare come phishing).

\subsection{Dentro la cella}
\[
h_t = F_1\big(W_{hh}h_{t-1} + W_{xh}x_t + b_{hh}\big), \qquad
\hat y_t = F_2\big(W_{yh}h_t + b_{yh}\big)
\]
con $F_1$ tipicamente \textbf{tanh}. Il passaggio dalla FFNN alla RNN: la stessa rete applicata a ogni
istante, poi collegata propagando lo stato. Relazione di ricorrenza $h_t = f_W(x_t, h_{t-1})$:
\textbf{stessa funzione e stessi pesi} a ogni passo. Nel many-to-one solo l'ultima cella ha la
classification head.

\begin{esame}
La parte ricorrente fa \textbf{feature extraction}: lo stato finale comprime la serie in pochi numeri
(un \textbf{embedding} della serie), e la head (da un layer lineare a una FFNN) lo usa per
classificare. Le due parti si addestrano \textbf{insieme}. Il docente insiste anche sul capire
\textbf{perché} la BPTT è più complessa della backpropagation classica.
\end{esame}

\subsection{BPTT e problemi di gradiente}
La \en{backpropagation through time} somma le loss dei passi, propaga l'errore attraverso output e
stati fino all'inizio della sequenza e accumula i gradienti sui pesi condivisi. Il gradiente rispetto
agli stati iniziali contiene molti fattori di $W_{hh}$: valori $>1$ portano all'\textbf{exploding
gradient} (rimedio: \en{gradient clipping}), valori $<1$ al \textbf{vanishing gradient} (rimedi:
attivazione, inizializzazione, architettura). Esempio delle dipendenze lunghe: ``\en{the clouds are in
the}\ldots'' si risolve con il passato recente; una lunga storia che parte dalla Francia e finisce con
``parlo fluentemente\ldots'' no.

\subsection{LSTM}
Proposta nel 1997, mantiene un \textbf{cell state} $C_t$ (memoria a lungo termine) e un
\textbf{hidden state} $h_t$ (breve termine), regolati da gate con sigmoid: \textbf{forget gate}
(quanto passato tenere), \textbf{input gate} e memoria candidata con tanh (quanto presente
aggiungere: un articolo come ``\en{the}'' conta poco), \textbf{output gate} (quanto del cell state
diventa $h_t$). Gestisce dipendenze lunghe al prezzo di molti più pesi. Quanto passato ricordi non si
sa a priori: dipende da dati, training e dimensioni degli stati. I transformer saranno trattati nella
parte di NLP.

\subsection{Dati e progettazione}
L'input di una RNN è un tensore 3D \code{(batch, seq\_len, features)}. Una serie mono-dimensionale si
scrive \code{[[1],[2],[3]]}: \code{[1,2,3]} sarebbe un solo istante con 3 feature. Il numero di
feature per istante è fisso; il numero di passi può variare tra serie, ma dentro un mini-batch serve
la stessa lunghezza (padding o troncamento, solo dentro il batch). Checklist: tipo di RNN, lunghezza
delle sequenze, memoria lunga o no, mono o bidirezionale; input = numero di \textbf{feature}; hidden
size in potenze di 2; 1--2 layer ricorrenti; head lineare o FFNN; ottimizzazione come per le FFNN
(Adam, learning rate 0,001 con scheduler, shuffle solo tra serie indipendenti).

\subsection{Da tabella a sequenze}
Nel dataset sull'inquinamento ogni riga è un'ora (35.040 righe, 7 feature): non ci sono ancora serie.
Una funzione \code{create\_sequences} prende finestre di 7 righe consecutive e come label il valore
di inquinamento successivo, producendo il tensore 3D per il training. L'esempio completo si vede nella
lezione del \rif{2026-10-14}{14/10}.

\section{Domande aperte e punti da verificare}
\begin{verifica}
\begin{itemize}
  \item Label della finestra: a voce il docente dice che è quella dell'``ultimo elemento, il
    settimo''; nel notebook è \code{labels[i + seq\_length]}, cioè il valore \textbf{successivo} alla
    finestra (si predice l'ora dopo). Vale il notebook.
\end{itemize}
\end{verifica}

% =====================================================================
\lezione{14/10/2026}{RNN in PyTorch, CNN e graph neural network}
  {Prof.~Danilo Giordano}
  {lunghezze diverse e padding dinamico; \code{collate\_fn} e packed sequence; \code{out} e
   \code{hidden}; notebook pollution e FordA; CNN (convoluzione, pooling, reti famose, STAMINA); GNN
   (permutation invariance ed equivariance, message passing, GCN, GraphSAGE, GAT, over-smoothing,
   task, split transductive e inductive, limiti)}
\label{lez:2026-10-14}
% =====================================================================

\begin{nota}[Fonte]
Trascrizione automatica della registrazione, revisionata. Codice, formule e numeri sono ripresi dalle
slide \emph{3-deep-learning-methodologies} (pp.~72--123) e dai notebook
\emph{RNN-Pollution-Multivariate} e \emph{RNN-LSTM-FordA}.
\end{nota}

\section{Sintesi}
Lezione lunga che chiude la parte sulle architetture per la classificazione. Prima si completa la
pratica delle RNN: come gestire serie di lunghezza diversa con il padding dinamico nel mini-batch e le
\en{packed sequence}, e due notebook (regressione sull'inquinamento, classificazione FordA). Poi una
breve introduzione alle \textbf{CNN}, utile anche per introdurre il concetto di convoluzione, e infine
la teoria delle \textbf{graph neural network}. Approfondimento nei capitoli~\ref{cap:rnn} e
\ref{cap:cnn-gnn}.

\section{Avvisi pratici}
\begin{itemize}
  \item Il \textbf{Lab~2} si terrà nei due mercoledì successivi (21 e 28 ottobre). Per liberare il
    mercoledì, la settimana prossima ci sarà lezione \textbf{lunedì e martedì} (19 e 20 ottobre);
    l'avviso sarà pubblicato sul portale.
  \item Nel Lab~2 le sequenze di API call descrivono un goodware o un malware: servirà la
    rappresentazione dell'\textbf{intera serie}; la parte GNN dei lab fa solo graph classification.
  \item Le GNN in PyTorch si vedranno lunedì (\rif{2026-10-19}{19/10}).
\end{itemize}

\section{Contenuti}

\subsection{RNN in pratica: lunghezze diverse}
Una rete è moltiplicazione di matrici: una serie più corta lascia un ``buco''. Tagliare tutto alla
lunghezza minima butta informazione; paddare con valori inventati (es.\ la media) fa credere alla rete
che la serie finisca così. Soluzione: un \code{Dataset} personalizzato e una \code{collate\_fn},
passata al \code{DataLoader}, che padda le serie \textbf{solo fino alla lunghezza massima del
batch} e restituisce anche le lunghezze originali. Basta paddare nel batch perché tutti i passi
condividono gli stessi pesi. Nel modello, \code{pack\_padded\_sequence} dice alla RNN quali istanti
sono padding, così vengono saltati.

\begin{esame}
Il padding serve solo alla matematica (matrice completa): \textbf{non deve entrare} nello stato finale
né nella backpropagation, altrimenti la rete si addestra su dati inventati. Il padding vale 0, che
potrebbe essere un valore reale: per questo la RNN deve sapere dove inizia.
\end{esame}

\subsection{\code{out} e \code{hidden}}
\code{hidden} contiene lo stato finale di ogni layer (\code{hidden[-1]} è quello dell'ultimo layer, la
rappresentazione più astratta) e si usa per i task a livello di \textbf{sequenza}, come la
classificazione; \code{out} contiene gli stati dell'ultimo layer a ogni passo e serve per i task a
livello di \textbf{passo} (previsione del valore successivo, many-to-many). Esempio della scalata:
\code{hidden} riassume tutto il viaggio, \code{out} lo descrive passo per passo. Varianti:
bidirezionale (head di dimensione \code{2*hidden\_size}, concatenando le due direzioni) e LSTM
(restituisce hidden e cell state; per classificare si usa l'hidden).

\subsection{I due notebook}
\begin{itemize}
  \item \textbf{Pollution} (regressione): guardare i dati prima di usarli (direzione del vento come
    stringa, valori sospetti come vento a 600~km/h); niente split casuale che rompe le serie; scaling
    fittato sul training; RNN + layer lineare senza attivazione finale, loss MSE. Il training loop è
    identico a quello della FFNN; le loss delle RNN sono più rumorose; manca il plateau; l'MSE va letto
    rispetto alla scala dei valori.
  \item \textbf{FordA} (classificazione di serie univariate): per mostrare il padding si accorciano a
    480 le prime 100 serie (le altre sono lunghe 500). Risultato: accuracy intorno al 51\% su un dataset
    bilanciato, predicendo sempre la stessa classe, con RNN, LSTM e Bi-RNN. Non perché le RNN non
    funzionino: head troppo semplice, hidden size piccola, training breve.
\end{itemize}

\subsection{Convolutional neural network}
Usi in cybersecurity: malware come immagine, siti clone da screenshot, intrusion detection su serie
(anche CNN-LSTM), deepfake. Struttura: layer convoluzionali che estraggono feature in modo gerarchico
(bordi, forme, oggetti) e una FFNN finale. Il layer è \textbf{convoluzione} (filtri scorrevoli con
pesi addestrabili; $N$ filtri producono $N$ strati di profondità; padding ai bordi), \textbf{ReLU} e
\textbf{pooling} (es.\ max pooling $2\times2$: una matrice $4\times4$ diventa $2\times2$). Reti
famose: LeNet-5 (1998), AlexNet (2012, ReLU e dropout), ResNet-50 (2015, skip connection).
\textbf{STAMINA} (Intel): byte del binario come pixel in scala di grigi, immagine $224\times224$,
ResNet-50 pre-addestrata con fine-tuning, accuracy oltre il 95\%.

\subsection{Graph neural network}
Molti dati sono grafi (botnet, frodi, grafi di API call). Una FFNN su matrice di adiacenza e feature
dipende dall'ordine dei nodi, non scala e non gestisce grafi di dimensione diversa; CNN e RNN
presuppongono una struttura fissa. Si generalizza la convoluzione: \textbf{un nodo è descritto dai
suoi vicini}. Due proprietà: \textbf{permutation invariance} a livello di grafo,
$f(A_1,X_1)=f(A_2,X_2)$, e \textbf{permutation equivariance} a livello di nodo,
$Pf(A,X)=f(PAP^T,PX)$. Con il \textbf{message passing} ogni nodo aggrega i messaggi dei vicini
(che a loro volta hanno aggregato i propri) con pesi \textbf{condivisi tra tutti i nodi}; un layer è
messaggio + aggregazione indipendente dall'ordine (somma, media, max) + il messaggio del nodo stesso +
non linearità.

\begin{esame}
Saper spiegare \textbf{perché} una GNN è invariante/equivariante (pesi condivisi e aggregazione che non
dipende dall'ordine) e perché è \textbf{induttiva} (gli stessi parametri descrivono nodi mai visti).
Il docente insiste anche sull'\textbf{attenzione} di GAT, che tornerà nella parte di NLP.
\end{esame}

Tre architetture: \textbf{GCN} (tutti i vicini, normalizzati per il grado), \textbf{GraphSAGE}
(campiona un numero fisso di vicini, aggregazione in due stadi: adatta a grafi grandi), \textbf{GAT}
(pesa ogni vicino con un coefficiente di attenzione appreso). Con troppi layer i \en{receptive field}
si sovrappongono e gli embedding diventano tutti simili (\textbf{over-smoothing}): scegliere il
numero di layer in base al receptive field necessario, aumentare l'espressività dentro il layer,
aggiungere MLP prima e dopo, skip connection. Task: node classification, link prediction, graph
classification. Split \textbf{transductive} (un solo grafo, si dividono le label) o
\textbf{inductive} (grafi diversi negli split, il setting dei lab). Limiti: over-smoothing,
scalabilità, grafi eterogenei e dinamici.

\section{Domande aperte e punti da verificare}
\begin{verifica}
\begin{itemize}
  \item Date del Lab~2 e lezione del martedì: ricavate da ``\en{next week}'' e ``\en{the next one and
    the one after}'' detti in aula; confermare con l'avviso sul portale.
  \item Nel notebook pollution si normalizza anche la label; il docente osserva che probabilmente non
    andrebbe fatto.
\end{itemize}
\end{verifica}

% =====================================================================
\lezione{19/10/2026}{GNN in pratica e feature categoriche}
  {Prof.~Danilo Giordano}
  {GNN con PyTorch Geometric (dati, GCN, GraphSAGE, GAT, training loop); scelte progettuali e iperparametri
   delle GNN; feature categoriche: ordinali e nominali; ID sequenziali, feature engineering, one-hot, embedding;
   token \code{<UNK>}}
\label{lez:2026-10-19}
% =====================================================================

\begin{nota}[Fonte]
Trascrizione automatica della registrazione, revisionata, confrontata con le slide
\emph{3-deep-learning-methodologies} (pp.~124--136), \emph{4-embedding-introduction} (pp.~1--22) e con il
notebook sulle GNN mostrato in aula.
\end{nota}

\section{Sintesi}
La lezione chiude la parte sulle architetture con il notebook delle \textbf{GNN} e apre il tema delle
\textbf{feature categoriche}, necessario per il Lab~2. Il filo conduttore: RNN e GNN servono a
\textbf{estrarre una rappresentazione} quando costruirla a mano (serie temporali, grafi) è difficile. Nella
seconda parte il docente costruisce una ``roadmap'' per trasformare le categorie in numeri sensati per il
modello, fino agli \textbf{embedding appresi}, che verranno approfonditi il giorno dopo
(\rif{2026-10-20}{20/10}). Approfondimenti nei capitoli~\ref{cap:cnn-gnn} e~\ref{cap:embedding}.

\section{Avvisi pratici}
\begin{itemize}
  \item Il \textbf{Lab~2} (supervisionato, dataset con feature categoriche) si svolge \textbf{mercoledì 21/10 e
    mercoledì 28/10}: servono le GNN viste oggi e la codifica delle feature categoriche
    (capitolo~\ref{cap:lab2}).
\end{itemize}

\section{Contenuti}

\subsection{GNN in PyTorch Geometric}
Il preprocessing è lo stesso delle altre reti (one-hot per le feature categoriche dei nodi, normalizzazione
o standardizzazione delle numeriche). Un grafo è un oggetto \code{Data} di \code{torch\_geometric} con
feature dei nodi \code{x}, archi \code{edge\_index} di forma \code{[2, n\_archi]} e label \code{y}: gli archi
sono \textbf{diretti}, quindi un collegamento percorribile nei due sensi diventa due archi. Il
\code{DataLoader} di \code{torch\_geometric.loader} unisce grafi di dimensioni diverse in un unico grafo
disconnesso, \textbf{senza padding}: in un batch di due grafi \code{x} ha come prima dimensione la somma dei
nodi (nel notebook 288). Per visualizzare un grafo piccolo si usa NetworkX (\code{to\_networkx},
\code{spring\_layout}).

Il modello: due layer \code{GCNConv} (due hop), \code{global\_mean\_pool} per passare dagli embedding dei
nodi a \textbf{un embedding per grafo}, poi una head lineare. GraphSAGE e GAT si ottengono cambiando il
layer (\code{SAGEConv}, \code{GATConv} con più \en{heads}: con \code{concat=True} l'output ha dimensione
\code{hidden\_dim} $\times$ \code{heads}). Il training loop è quasi identico ai precedenti, con due
differenze: la label è un attributo del batch (\code{batch.y}) e la loss va pesata con
\code{batch.num\_graphs}, non con \code{x.size(0)}, che conta i nodi. Sul dataset del notebook le tre
architetture superano il 90\%; la GAT converge meglio ma è molto più lenta. Con il \code{NeighborLoader}
GraphSAGE campiona i vicini in training (es.\ 10 al primo hop, 5 al secondo), mentre in validation e test
usa tutti i vicini.

\begin{attenzione}[da ricordare per il Lab~2]
Nel training loop di una GNN per la classificazione di grafi la loss si accumula moltiplicando per il
\textbf{numero di grafi} del batch, non per il numero di nodi (288 invece di 2, nell'esempio). Il docente
ricorda che questo codice servirà nel lab del mercoledì.
\end{attenzione}

\subsection{Scelte progettuali per una GNN}
Le slide elencano le decisioni da prendere: tipo di GNN (GCN per grafi non troppo grandi, GraphSAGE per grafi
grandi e nodi ad alto grado, GAT per rappresentazioni migliori con l'attenzione), task (nodo, arco o grafo),
scala (full-batch o mini-batch con campionamento), aggregazione, dimensione dell'embedding, \textbf{numero di
layer} (2--3, legato al \en{receptive field}; troppi portano all'\textbf{over-smoothing}), head di predizione
con pooling per i grafi. Le best practice sugli iperparametri (potenze di 2 per l'hidden size, Adam,
learning rate 0,01 o 0,001, dropout 0,2--0,5) sono nel capitolo~\ref{cap:cnn-gnn}.

\subsection{Feature categoriche: dalle categorie ai numeri}
Le reti richiedono numeri e assumono che le relazioni tra i numeri abbiano senso: assegnare numeri a caso
alle categorie non va bene. Per le \textbf{label} basta quasi sempre l'ID sequenziale. Per le feature:
\begin{itemize}
  \item \textbf{ordinali} (gravità, livello di accesso): si ordinano con la \textbf{conoscenza del dominio},
    non in ordine alfabetico, e si mappano su 1, 2, 3;
  \item \textbf{nominali} (es.\ AES, DES, RC4, RSA): un ID sequenziale introduce un ordine falso e mette il
    modello su una pista sbagliata; il \textbf{feature engineering} (sicuro/insicuro) dà una chiave di
    lettura ma può essere sbagliata o cambiare nel tempo; il \textbf{one-hot} non introduce bias (vettori
    ortogonali ed equidistanti) ma soffre la \textbf{maledizione della dimensionalità} e va rifatto se il
    vocabolario si estende; gli \textbf{embedding} mappano ogni categoria su un vettore \textbf{denso} di
    dimensione $H \ll V$, appreso con la backpropagation risolvendo un task (quindi solo con reti neurali).
\end{itemize}
Il docente mostra ``dietro le quinte'' come, epoch dopo epoch, la rappresentazione 2D degli algoritmi si
separa tra sicuri e insicuri. Le categorie mai viste in training finiscono nel token \textbf{\code{<UNK>}}.

\begin{esame}[la roadmap delle feature categoriche]
Saper scegliere e giustificare la codifica: ordine intrinseco $\Rightarrow$ ID sequenziali; poche
categorie $\Rightarrow$ one-hot; si può definire un task per una rete $\Rightarrow$ embedding; altrimenti
feature engineering. Un ID sequenziale su una feature nominale ``non si dovrebbe, ma si può'' se si sa
cosa si fa: il rischio è introdurre un bias e pagare scelte sub-ottimali con più dati ed epoch.
\end{esame}

\section{Domande aperte e punti da verificare}
\begin{verifica}
\begin{itemize}
  \item La slide indica come distanza tra due vettori one-hot distinti il valore 2: è la distanza euclidea
    \emph{al quadrato}; la distanza euclidea è $\sqrt{2}$.
  \item I tempi di training delle tre GNN (GCN, GraphSAGE, GAT) detti in aula dipendono dall'esecuzione del
    notebook: vanno presi come ordini di grandezza.
\end{itemize}
\end{verifica}

% =====================================================================
\lezione{20/10/2026}{Imparare gli embedding; dropout e batch norm nel dettaglio}
  {Prof.~Danilo Giordano}
  {embedding supervisionati con \code{nn.Embedding}: dati tabellari, sequenze, architetture composte;
   shortcut learning; embedding self-supervised (autoencoder, contrastive learning); dropout e batch
   normalization in training e in valutazione; linee guida per i report dei lab}
\label{lez:2026-10-20}
% =====================================================================

\begin{nota}[Fonte]
Trascrizione automatica della registrazione, revisionata, confrontata con le slide
\emph{4-embedding-introduction} (pp.~23--51), \emph{2a-dropout\_batchnorm} e con il documento di linee guida
per i report presentato a fine lezione.
\end{nota}

\section{Sintesi}
La lezione completa il discorso sulle feature categoriche iniziato il giorno prima (\rif{2026-10-19}{19/10}):
come si \textbf{impara} un embedding. Prima nel caso supervisionato, con tre esempi di complessità
crescente (tabella, sequenza, grafi nel tempo), poi nel caso \textbf{self-supervised}, dove il segnale non
viene dalle label ma dalle proprietà interne dei dati. Seguono un breve deck che chiarisce come
\textbf{dropout} e \textbf{batch normalization} si comportano in training e in valutazione, e la
presentazione delle \textbf{linee guida per i report}. Approfondimenti nei capitoli~\ref{cap:embedding},
\ref{cap:regolarizzazione} e~\ref{cap:lab-report}.

\section{Avvisi pratici}
\begin{itemize}
  \item Il Lab~2 inizia \textbf{il giorno dopo} (mercoledì 21/10): gli embedding visti oggi servono lì.
  \item Sul portale è caricato il documento con le \textbf{raccomandazioni per i report} (dettagli nel
    capitolo~\ref{cap:lab-report}). Punti richiamati in aula:
    \begin{itemize}
      \item chi non è ancora in un gruppo deve unirsi a uno \textbf{entro fine mese}; ci sono gruppi da due o
        tre persone da contattare;
      \item consegna \textbf{10 giorni prima} dello scritto; PDF di \textbf{massimo 10 pagine} con lab ID,
        group ID, nomi ed email nell'intestazione; LaTeX o Word indifferentemente, con un font leggibile;
      \item il notebook deve \textbf{girare ed essere già eseguito} (con gli output visibili), ma il report
        deve essere autosufficiente: non si va a cercare i risultati nel notebook;
      \item novità di quest'anno: un \textbf{logbook degli esperimenti} sintetico, che \textbf{non conta}
        nelle 10 pagine;
      \item ogni report vale da 0 a 32 e serve almeno 18; si valuta il ragionamento, non lo 0,81 contro lo
        0,82; figure leggibili e tabelle al posto di elenchi di numeri nel testo.
    \end{itemize}
\end{itemize}

\section{Contenuti}

\subsection{Embedding supervisionati}
Ipotesi: label disponibili, feature numeriche e categoriche, e una \textbf{rete neurale} (un random forest
classifica ma non fa feature extraction). \code{nn.Embedding(V, H)} è una lookup table addestrabile con $V$
righe (le categorie) e $H$ colonne (dimensione dell'embedding, un iperparametro); la backpropagation ne
aggiorna i valori come per gli altri pesi.
\begin{itemize}
  \item \textbf{Tabella} (metadati di email, task phishing/normale): il dominio del mittente ha cardinalità
    enorme e nessun ordine $\Rightarrow$ embedding. Si fa lo split, si costruisce il \textbf{vocabolario
    solo sul training} (validation e test simulano il futuro) con \code{<UNK>} per i valori mai visti, si
    mappano gli ID, si separano categoriche e numeriche, e nel modello l'embedding si \textbf{concatena}
    alle numeriche prima della FFNN (input del primo layer: \code{embed\_dim + num\_numeric}).
  \item \textbf{Sequenza} (IP che contattano una darknet, task: attacco di botnet e quale): embedding su
    ogni elemento, poi LSTM con input di dimensione \code{embed\_dim}; con \code{padding\_idx} l'embedding
    del padding non viene aggiornato. È lo schema del lab.
  \item \textbf{Architettura composta} (movimento laterale): embedding dei servizi $\to$ GNN per ogni
    snapshot $\to$ RNN sulla serie di snapshot $\to$ predizione, tutto addestrato con lo stesso segnale
    supervisionato.
\end{itemize}

\begin{esame}[un embedding supervisionato vale per il suo task]
Il docente insiste: l'embedding appreso con le label è \textbf{ottimo per quel task e sub-ottimo per
altri}. Numericamente lo si può riusare, ma è stato costruito per un altro obiettivo. E può imparare
\textbf{scorciatoie}: nell'esempio husky contro lupo il classificatore impara lo sfondo (erba o neve),
non l'animale.
\end{esame}

\subsection{Embedding self-supervised}
Senza label il segnale viene dalle \textbf{proprietà interne} dei dati. Due strade: un
\textbf{autoencoder}, in cui l'encoder comprime e il decoder ricostruisce (un problema di regressione senza
label; al test si congela l'encoder e si addestra una piccola head con poche label); il
\textbf{contrastive learning}, in cui le versioni aumentate di un'immagine (positivi) devono avere embedding
vicini all'originale (anchor) e un'immagine diversa (negativo) lontano, con la \textbf{contrastive loss}.
Pre-addestramento su grandi corpora (\en{pretext}) e fine-tuning con poche label (\en{downstream}) sono la
base dei successi recenti, e sono preziosi in cybersecurity dove le label scarseggiano. Le augmentation
però iniettano conoscenza, quindi bias (la domanda della slide sullo scoiattolo verde).

\subsection{Dropout e batch normalization nel dettaglio}
\begin{itemize}
  \item \textbf{Dropout}: in training ogni neurone è spento con probabilità $p$, con una maschera nuova per
    ogni campione e ogni mini-batch, e i sopravvissuti sono scalati di $1/(1-p)$; in valutazione è
    l'identità. Ogni neurone impara solo dai campioni in cui era acceso.
  \item \textbf{Batch norm}: in training standardizza ogni feature con media e varianza \textbf{del
    mini-batch}, poi applica $\gamma$ e $\beta$ \textbf{appresi}, e aggiorna media e varianza
    \textbf{mobili} (calcolate, non apprese); in valutazione usa le statistiche mobili congelate, così il
    layer successivo riceve dati con la stessa distribuzione vista in training.
\end{itemize}
Tabelle ed esempio numerico nel capitolo~\ref{cap:regolarizzazione}.

\begin{integrazione}[\code{model.eval()} e \code{torch.no\_grad()}]
È \code{model.eval()} che cambia il comportamento di dropout e batch norm; \code{torch.no\_grad()}
disattiva soltanto il calcolo dei gradienti. In valutazione servono entrambi.
\end{integrazione}

\section{Domande aperte e punti da verificare}
\begin{verifica}
\begin{itemize}
  \item Domanda aperta della slide: un modello addestrato con augmentation che tolgono il colore
    riconoscerebbe come anomala la foto di uno scoiattolo verde? Le slide la lasciano come spunto di
    ragionamento, senza risposta ufficiale.
  \item Formato e lunghezza del logbook degli esperimenti: da controllare nel documento sul portale.
\end{itemize}
\end{verifica}

% =====================================================================
\lezione{02/11/2026}{Anomaly detection: concetti, one-class SVM e k-means}
  {Prof.~Danilo Giordano}
  {anomaly, outlier e novelty detection; perché in cybersecurity; complessità e tipi di anomalie (point,
   contextual, collective); sfide per AI/ML; supervised vs unsupervised; score e label; shallow learning;
   SVM e kernel trick; one-class SVM e parametro $\nu$; clustering e k-means; SSE e metodo del gomito}
\label{lez:2026-11-02}
% =====================================================================

\begin{nota}[Fonte]
Trascrizione automatica della registrazione, revisionata, confrontata con le slide
\emph{5-Supervised-vs-Unsupervised-shallow} (pp.~1--33).
\end{nota}

\begin{verifica}
Tra la lezione del 20/10 e questa ci sono state attività non registrate (le due sessioni del Lab~2 e
le eventuali lezioni di quei giorni). Il diario non ne riporta il contenuto.
\end{verifica}

\section{Sintesi}
Prima lezione dell'ultimo blocco del Prof.~Giordano, dedicato all'\textbf{anomaly detection}; dopo
inizierà l'NLP con il Prof.~Boffa. Il punto di partenza: i classificatori falliscono quando le label
mancano, sono molto sbilanciate o inaffidabili, o quando si vogliono scoprire \textbf{novità}. La lezione
definisce il problema, le sue complessità e le sfide per il ML, confronta l'approccio supervisionato e
quello non supervisionato e inizia gli algoritmi \emph{shallow}: SVM, \textbf{one-class SVM} e
\textbf{k-means}. Si ferma alla scelta di $K$; l'uso di k-means per l'anomaly detection è rimandato alla
lezione successiva. Approfondimento nel capitolo~\ref{cap:anomaly}.

\section{Avvisi pratici}
\begin{itemize}
  \item Il docente chiede un riscontro sul Lab~2 (``difficile?'') e lo difende come utile, anche se
    impegnativo.
  \item La parte su k-means per l'anomaly detection riprende ``domani o mercoledì'': quella lezione non è
    nelle registrazioni (vedi \rif{2026-11-09}{09/11}).
\end{itemize}

\section{Contenuti}

\subsection{Che cos'è e perché serve}
L'anomaly detection cerca \textbf{pattern che deviano dal comportamento atteso}. Va distinta la
\textbf{novelty detection}: un pattern nuovo può essere la \textbf{nuova normalità} (l'uso della VPN
durante il COVID). In cybersecurity le anomalie segnalano attività malevole (zero-day, intrusioni,
esfiltrazione dopo un movimento laterale), guasti e configurazioni errate, frodi bancarie. Serve perché
per uno \textbf{zero-day} non può esistere un classificatore, perché gli attaccanti si adattano alle
regole (gatto e topo) e per avere un \textbf{early warning} prima del danno.

\subsection{Complessità, tipi di anomalie e sfide}
Le complessità sono quattro: \textbf{unknowingness} (le anomalie sono ignote finché non accadono),
\textbf{rarità e sbilanciamento} (e sbagliare su un'anomalia costa di più), \textbf{classi eterogenee}
(lo stesso tipo di attacco si presenta in forme diverse, es.\ DNS tunnelling e HTTP/2 push) e
\textbf{pattern che evolvono}. I tipi di anomalia sono tre:
\begin{itemize}
  \item \textbf{point}: un'istanza anomala da sola (login da un paese mai visto);
  \item \textbf{contextual}: anomala solo dato il contesto (un backup di due ore alle 3 di domenica);
  \item \textbf{collective}: un insieme di istanze normali che insieme formano l'anomalia (tentativi di
    login ognuno sotto la soglia, ma in sequenza).
\end{itemize}
Le sfide per AI/ML: recall bassa e molti falsi positivi (se la maggior parte degli allarmi è falsa,
nessuno si fida più del modello), dati ad alta dimensionalità e dipendenti, poche label, rumore, anomalie
complesse e bisogno di \textbf{spiegazioni}. Approccio \textbf{supervisionato}: accurato con buone label
ma non generalizza a nuove anomalie; \textbf{non supervisionato}: assume che il normale domini, trova
minacce nuove ma ha più falsi positivi ed è meno interpretabile. L'output di un algoritmo può essere uno
\textbf{score}, che permette di ordinare le istanze e di lavorare a budget (esempio Intesa Sanpaolo:
le prime 10.000 del mese), oppure una \textbf{label}, il cui numero si controlla con gli iperparametri.

\begin{esame}[un'anomalia non è per forza un attacco]
Il docente torna più volte sul contesto: un rilevatore che blocca subito tutte le VPN perché il loro uso
è ``anomalo'' fermerebbe tutti i lavoratori da remoto. Saper distinguere anomalia, novità e attacco, e
riconoscere point, contextual e collective anomaly negli esempi.
\end{esame}

\subsection{Shallow learning: SVM e one-class SVM}
Gli algoritmi shallow sono interpretabili, vanno bene sui dati tabellari e sulle anomalie puntuali e
danno una \textbf{baseline}; sono deboli sulle anomalie contestuali e collettive, che richiedono feature
costruite a mano (es.\ tentativi falliti aggregati su mezz'ora). La \textbf{SVM} (1995) cerca l'iperpiano
che separa le classi \textbf{massimizzando il margine}; il \textbf{kernel trick} porta i dati in uno spazio
in cui diventano separabili (lineare, polinomiale, \textbf{RBF}, il più usato). La \textbf{one-class SVM}
non riceve label: costruisce un confine attorno ai dati normali, e ciò che sta fuori è anomalo. Il
parametro $\bm{\nu}$ è la \textbf{frazione attesa di outlier}. Si può usare in due modi: addestrata
\textbf{solo su dati normali} (``supervisionata'' tra virgolette: le label servono solo a scegliere i
dati, $\nu$ piccolo) o su \textbf{tutti i dati} senza label (serve una stima della frazione di anomalie).

\subsection{Clustering e k-means}
Nel clustering le anomalie sono i \textbf{cluster piccoli e lontani}, i punti \textbf{lontani dal
centroide} del proprio cluster o i punti \textbf{senza cluster}. K-means è partizionale e basato su
centroidi: ogni punto appartiene a un cluster, funziona con cluster sferici di dimensioni simili.
Dipende dall'\textbf{inizializzazione} casuale, quindi si esegue più volte con lo stesso $K$ e si tiene il
risultato migliore. $K$ si sceglie con l'\textbf{SSE} e il \textbf{metodo del gomito}, restando nella
zona del gomito: più cluster abbassano l'SSE ma vanno studiati uno per uno, senza label.

\begin{esame}[SSE che risale]
Il docente chiede di ricordarlo: se nel grafico SSE contro $K$ compare un picco verso l'alto,
probabilmente un'inizializzazione sfortunata ha lasciato un \textbf{cluster vuoto}. Si riesegue k-means.
\end{esame}

\section{Domande aperte e punti da verificare}
\begin{verifica}
\begin{itemize}
  \item La slide sulla one-class SVM dice che con $\nu$ più piccolo si hanno meno outlier e confini
    ``più stretti'': ``meno outlier'' è corretto, ma lasciando fuori meno punti la regione normale si
    allarga. Da chiarire con il docente.
  \item Contenuto delle lezioni e dei lab non registrati tra il 20/10 e il 02/11: non ricostruito.
\end{itemize}
\end{verifica}

% =====================================================================
\lezione{09/11/2026}{Anomaly detection in pratica: dal clustering agli autoencoder}
  {Prof.~Danilo Giordano}
  {notebook \code{anomaly\_detection} su KDD Cup 1999: studio del dataset e preprocessing; PCA (con e senza
   scaling) e t-SNE; contaminazione; one-class SVM, LOF, k-means, DBSCAN, Isolation Forest; autoencoder con
   loss di ricostruzione mista; encoder come estrattore di feature; confronto finale e validazione}
\label{lez:2026-11-09}
% =====================================================================

\begin{nota}[Fonte]
Trascrizione automatica della registrazione, revisionata, confrontata con il notebook
\code{anomaly\_detection} (versione pubblicata sul portale) e con le slide
\emph{5-Supervised-vs-Unsupervised-shallow} (pp.~33--57) e \emph{6-Supervised-vs-Unsupervised-deep}.
\end{nota}

\begin{verifica}
La lezione dà per vista la teoria di k-means per l'anomaly detection, DBSCAN, validazione del clustering,
LOF, PCA, Isolation Forest (slide~5, pp.~33--57) e dell'anomaly detection con il deep learning (slide~6).
La lezione del 02/11 si ferma alla scelta di $K$ e annuncia il seguito ``domani o mercoledì'': quella
parte \textbf{non è nelle registrazioni} disponibili. Il diario non ricostruisce cosa sia stato detto: la
teoria, presa dalle slide, è nel capitolo~\ref{cap:anomaly}.
\end{verifica}

\section{Sintesi}
Ultima lezione del Prof.~Giordano sull'anomaly detection, tutta sul notebook che prepara il Lab~3. Il
docente non commenta ogni riga (le funzioni di servizio sono molte e il codice è semplice), ma mostra
\textbf{cosa fare e come leggere i risultati}. Si applicano allo stesso dataset tutti i metodi del blocco,
shallow e deep, con lo stesso schema: fit, score, soglia, valutazione con le label. Il messaggio
principale: ciò che il modello vede in training conta più del modello stesso; i metodi addestrati su
\textbf{dati normali puliti} (one-class SVM, autoencoder) vanno molto meglio di quelli addestrati su
tutto. Approfondimenti nei capitoli~\ref{cap:anomaly} e~\ref{cap:lab3}.

\section{Avvisi pratici}
\begin{itemize}
  \item Prima parte del \textbf{Lab~3} (anomaly detection) \textbf{mercoledì 11/11}.
  \item Il notebook è solo dimostrativo: iperparametri fissati senza tuning, nessuna fase di validazione.
\end{itemize}

\begin{esame}[nel lab va fatto ciò che il notebook non fa]
Il docente lo chiede esplicitamente: nel lab \textbf{training}, \textbf{tuning sul validation set} e una
sola \textbf{valutazione finale sul test}. Nel notebook i metodi si confrontano sul test solo perché
nessuno è stato ottimizzato; gli autoencoder usano il validation per monitorare il training, quindi un
confronto equo richiede un set che nessun metodo ha visto.
\end{esame}

\section{Contenuti}

\subsection{Il dataset e il preprocessing}
KDD Cup 1999: una riga per connessione, 41 feature generiche (non pensate per l'anomaly detection) e una
label con il tipo di attacco. Le label \textbf{non si usano in training}: servono solo a verificare se si
ritrovano le anomalie note. Prima si studia il dataset: rimuovendo i \textbf{duplicati} resta circa un
terzo delle righe; le categorie di \code{service} e \code{flag} sono in gran parte rarissime; \code{normal}
domina, \code{neptune} (SYN flood) è molto frequente e molti attacchi sono rari. Scelte del notebook: si
toglie \code{neptune} (crea regioni dense, quindi non è un'anomalia per densità), si crea una label
binaria, si fa uno split \textbf{stratificato} (circa 6\% di anomalie) \textbf{prima} di encoding e
scaling, e le categorie rare confluiscono in \code{other}.

\begin{esame}[studiare il dataset e non buttare la label]
``Prima di qualsiasi pipeline, studiate il problema e il dataset.'' La colonna con la label originale
\textbf{non va eliminata}: serve dopo per interpretare cluster ed errori. I risultati vanno rappresentati
in modo che si veda ciò che si vuole mostrare (asse logaritmico o una tabella per conteggi molto
sbilanciati).
\end{esame}

\subsection{Visualizzare: PCA e t-SNE}
Senza scaling delle numeriche, una sola componente della PCA spiega il 99,9\% della varianza: non è una
scoperta ma un \textbf{errore}, perché una feature in byte domina tutte le altre. Con
\code{StandardScaler} (fit solo sul training) la prima componente spiega circa il 10\% e se ne scelgono una
ventina al ginocchio della curva. Con t-SNE le coordinate non hanno significato: si può dire se un gruppo
di attacchi è separato, non quanto è lontano. Nel lab la \textbf{perplexity} va esplorata su un intervallo
ampio (fino a 50 e oltre), non 1, 2, 3.

\subsection{I metodi shallow}
La \textbf{contaminazione} (frazione attesa di anomalie, qui presa dalle label, nella realtà da stimare)
trasforma lo score in una decisione.
\begin{itemize}
  \item \textbf{One-class SVM} solo sui normali (kernel RBF, $\nu = 0{,}001$): buona precision sulle
    anomalie, recall migliorabile. Addestrata su tutti i dati con $\nu$ uguale alla contaminazione:
    precision e recall delle anomalie crollano. Il non supervisionato è più difficile.
  \item \textbf{LOF}: risultato pessimo (in aula 36 anomalie trovate su più di 600). Molti attacchi sono
    raffiche di connessioni quasi identiche, cioè gruppi piccoli ma densi: il dataset sembra adatto ad
    anomalie globali più che locali. Prima di scartare il metodo, provare altri \code{n\_neighbors}.
  \item \textbf{K-means}: $K$ con l'SSE (nel grafico in aula l'SSE risale: servivano più run per ogni $K$),
    silhouette media bassa, tabella cluster $\times$ label per interpretare (un cluster contiene tutti e
    solo i \code{teardrop}); score = distanza dal centroide. Gli attacchi che formano un cluster proprio
    sfuggono.
  \item \textbf{DBSCAN}: $\varepsilon$ dal \en{k-distance plot}, che però non mostra un ginocchio netto;
    molti normali finiscono nel rumore e la precision è molto bassa. È la maledizione della
    dimensionalità in pratica.
  \item \textbf{Isolation Forest}: non usa distanze, è velocissimo ed è il migliore tra gli shallow
    interamente non supervisionati.
\end{itemize}

\begin{attenzione}[la ROC-AUC può ingannare]
Su dati molto sbilanciati una ROC-AUC alta non basta: il docente si affida più a precision, recall e F1
della classe anomala. La matrice di confusione si usa solo per leggere i risultati, non perché il
problema sia una classificazione.
\end{attenzione}

\subsection{Autoencoder}
Input misto (numeriche scalate e one-hot), quindi \textbf{loss di ricostruzione mista}: errore quadratico
sulle numeriche più cross-entropy su ogni blocco one-hot; da adattare alle feature del proprio problema.
Addestrato \textbf{solo sui normali}, l'errore delle anomalie sul validation cresce con le epoch e
l'errore di ricostruzione diventa un ottimo score: è il metodo migliore del notebook. Addestrato su
\textbf{normali e anomalie}, impara a ricostruire anche parte di queste e le prestazioni peggiorano. Una
versione più profonda con batch norm e dropout non migliora: un modello che generalizza ``troppo''
ricostruisce bene anche le anomalie. Usare l'\textbf{encoder come estrattore di feature} e applicare la
one-class SVM nello spazio latente qui non aiuta: il bottleneck scarta proprio l'informazione che rende
anomalo un punto.

\begin{sintesi}[da portare nel lab]
Dati normali puliti in training; scaling prima di PCA e metodi a distanza; intervalli ampi per
perplexity, $\varepsilon$, \code{n\_neighbors}; più run di k-means per ogni $K$; tuning sul validation,
una sola valutazione sul test.
\end{sintesi}

\section{Domande aperte e punti da verificare}
\begin{verifica}
\begin{itemize}
  \item I valori numerici mostrati in aula (ROC-AUC, F1, numero di anomalie trovate) differiscono in
    qualche punto dalla versione del notebook pubblicata sul portale: per i numeri fare riferimento a
    quest'ultima (tabella nel capitolo~\ref{cap:lab3} o~\ref{cap:anomaly}).
  \item La teoria di DBSCAN, LOF, Isolation Forest, validazione del clustering e deep anomaly detection
    (slide~5, pp.~33--57, e slide~6) non è nelle registrazioni: eventuali indicazioni d'esame date in aula
    su quegli argomenti non sono note.
\end{itemize}
\end{verifica}


% =====================================================================
\part{Manuale del corso}
% Capitoli per argomento: slide + diario + integrazioni segnalate.
% =====================================================================
\chapter{AI per la cybersecurity: concetti e paradigmi}
\label{cap:ai-cyber}
\fonti{\rif{2026-09-21}{21/09}, \rif{2026-09-23}{23/09}}{1-intro-ai-for-cybersecurity, pp.\ 1--54}

Obiettivo del corso: non imparare dettagli fini, ma capire \textbf{come usare AI e ML per la
cybersecurity} e \textbf{cosa c'è sotto il cofano} dei modelli che si usano, per sapere perché
funzionano (e quando no).

\section{Il problema: superficie d'attacco e costi in crescita}
\begin{itemize}
  \item Ogni anno si connettono più dispositivi ICT a Internet: circa \textbf{50 miliardi entro il 2029}.
  \item Molti sono \textbf{low-cost} (lampadine e prese smart, IoT) con poco budget per la sicurezza.
  \item Attacchi e data breach sono più frequenti, i dispositivi più complessi, i costi di difesa
        dell'ordine dei \textbf{trilioni di dollari}.
\end{itemize}
$\Rightarrow$ Serve una difesa \textbf{più veloce, più efficiente e automatizzata}.

\section{Cos'è la cybersecurity}
\begin{definizione}[threat e attack]
\begin{itemize}
  \item \textbf{Threat (minaccia)}: qualsiasi pericolo \emph{potenziale} che può sfruttare una vulnerabilità di un sistema ICT. Non è ancora successo.
  \item \textbf{Attack (attacco)}: quando un utente malevolo \emph{sfrutta effettivamente} una vulnerabilità.
\end{itemize}
\end{definizione}

La cybersecurity non è solo software (antivirus, firewall): comprende anche \textbf{attività e
policy} (es.\ formare il personale) e \textbf{strumenti, pratiche e linee guida} per proteggere
reti, software e dati da attacchi, danni e accessi non autorizzati. In sintesi: identificare
minacce e attacchi e le relative \textbf{tattiche di difesa}.

\subsection{Cosa sappiamo e cosa non sappiamo}
\begin{itemize}
  \item \textbf{What we know} (gestione del rischio tradizionale): minacce note che sappiamo come
        contrastare, es.\ applicare una patch. Approccio \textbf{reattivo}, molto lavoro umano.
  \item \textbf{What we know we don't know}: sappiamo di poter essere vulnerabili e cerchiamo
        attivamente il problema (es.\ scan delle porte aperte). È il campo della
        \textbf{Cyber Threat Intelligence (CTI)}, approccio \textbf{proattivo}.
  \item \textbf{What we don't know we don't know}: rischi che non conosciamo finché non si
        manifestano, cioè gli \textbf{zero-day}. È il problema vero.
\end{itemize}

\textbf{Threat intelligence}: raccogliere e analizzare dati sul proprio sistema per scoprire
vulnerabilità prima dell'attaccante e stare un passo avanti. Con la sola analisi umana
\textbf{non si regge} la velocità di diffusione delle minacce né la complessità dei sistemi
(es.\ centinaia di dispositivi connessi al Wi-Fi in una sola aula).

\section{Perché AI/ML: qualche numero}
\begin{center}
\small
\begin{tabularx}{\textwidth}{l X}
\toprule
\textbf{Più veloci} & Tempo mediano tra compromissione ed esfiltrazione: \textbf{9 giorni (2022) $\to$ 2 giorni (2024)} [Palo Alto Networks].\\
& \en{Breakout time} (accesso $\to$ movimento laterale): \textbf{84 min $\to$ 62 min} [CrowdStrike].\\
& Modelli recenti (es.\ \en{Mythos Preview}, Anthropic, 2026) individuano e sfruttano zero-day nei principali sistemi operativi e browser.\\
\textbf{Più complessi} & Tattiche MITRE: \textbf{9 (2015) $\to$ 14 (2024)}; attacchi \en{hands-on-keyboard} (interattivi): +60\% l'anno.\\
\textbf{Più costosi} & Costo globale del cybercrime: \textbf{9,5 trilioni USD} nel 2024 (+15\% l'anno); costo medio di un data breach: \textbf{4,88 milioni USD} (+10\%) [IBM].\\
\bottomrule
\end{tabularx}
\end{center}

\begin{attenzione}[il gioco del gatto e del topo]
L'AI aiuta chi difende (identificazione, prevenzione e risposta più rapide) ma \textbf{aiuta anche
chi attacca} (analisi dei sistemi, nuovi malware, varianti che eludono l'antivirus). Difensore e
attaccante diventano entrambi più ``intelligenti''.
\end{attenzione}

\subsection{Cosa porta l'AI}
L'AI impara dai \textbf{dati storici}: se modella il comportamento tipico della rete, quando il
comportamento cambia può sollevare un allarme. Questo rende la difesa \textbf{adattiva e proattiva}.
\begin{itemize}
  \item \textbf{Automazione}: compiti ripetitivi (leggere log, cercare pattern) eseguiti sempre con
        la stessa qualità, senza stanchezza; tempi di risposta più rapidi, meno carico per gli analisti.
  \item \textbf{Intelligenza}: coglie correlazioni tra decine o centinaia di variabili che un umano non
        vede; individua minacce nuove tramite pattern e anomalie (\en{anomaly detection}); migliora nel tempo
        se ha \textbf{dati di qualità} (l'obiettivo ideale sono gli zero-day).
  \item \textbf{Robustezza}: rileva anche piccole deviazioni, resiste a rumore e dati mancanti, deve resistere
        a tentativi di manipolazione. Esempio: Kevin Mitnick negli anni '80 otteneva dati sensibili per telefono
        in 20 minuti ingannando i dipendenti; un sistema automatico sarebbe stato meno manipolabile.
\end{itemize}

\section{AI, Machine Learning, Neural Network, Deep Learning}
\begin{center}
\begin{tikzpicture}[every node/.style={align=center, font=\small}]
  \draw[fill=defcol!8, draw=defcol] (0,0) ellipse (6.2 and 2.6);
  \draw[fill=defcol!16, draw=defcol] (0.8,-0.25) ellipse (4.8 and 2.0);
  \draw[fill=defcol!26, draw=defcol] (1.6,-0.45) ellipse (3.4 and 1.4);
  \draw[fill=defcol!40, draw=defcol] (2.4,-0.6) ellipse (1.9 and 0.8);
  \node at (-4.4,0.9) {\textbf{AI}};
  \node at (-2.7,0.4) {\textbf{ML}};
  \node at (-0.6,0.1) {\textbf{Neural}\\\textbf{Network}};
  \node at (2.4,-0.6) {\textbf{Deep}\\\textbf{Learning}};
\end{tikzpicture}
\end{center}

\begin{definizione}[Artificial Intelligence]
Capacità generale dei computer di \textbf{emulare il pensiero umano} (ragionamento, problem solving,
decisioni autonome). \textbf{Non richiede per definizione di imparare dai dati}: un sistema a regole
(``se è lungo e sottile è una penna, se è più largo è uno smartphone'') è AI. La regola la dà l'umano.
\end{definizione}

\begin{definizione}[Machine Learning]
Sottoinsieme dell'AI che \textbf{impara dai dati} (o dall'esperienza) creando un
\textbf{modello data-driven} tramite il \textbf{training}: non diamo la regola, diamo 100 penne e 100
smartphone e l'algoritmo ne ricava le caratteristiche.
\end{definizione}

Conseguenze importanti:
\begin{itemize}
  \item \textbf{Il ML conosce solo ciò che ha visto in passato.} Se gli mostri un oggetto mai visto
        (es.\ delle cuffie) non lo riconosce. Gli eventi rari sono difficili da modellare.
  \item In cybersecurity è critico: i malware evolvono (v1, v2, v3\ldots), quindi un modello generale
        è difficile o impossibile, e servono aggiornamenti periodici.
  \item \textbf{Più dati è meglio solo se sono di qualità}: \en{garbage in, garbage out}. Dieci temi d'esame
        di cui nove di altri corsi non aiutano a preparare questo esame.
\end{itemize}

\begin{definizione}[Neural Network]
Algoritmo di ML ispirato al cervello: \textbf{neuroni} (nodi) che elaborano l'input, organizzati in
\textbf{layer} (input, hidden, output) e collegati da \textbf{sinapsi} (link) che passano
l'informazione tra layer. Gli \textbf{hidden layer} sono la parte che ``ragiona''; l'output layer dà
la risposta finale (es.\ vettore di probabilità \{cane: 0,99, gatto: \ldots\}).
\end{definizione}
Sono molto versatili (denoising delle foto, LLM, guida autonoma, riconoscimento immagini) e spesso
superano gli algoritmi \en{shallow} (decision tree, SVM, random forest), specialmente su compiti come
le immagini.

\begin{definizione}[Deep Learning]
Reti neurali con \textbf{3 o più hidden layer} che elaborano grandi quantità di dati. La vera differenza
rispetto alle reti \en{shallow} però non è solo il numero di layer: esistono \textbf{architetture diverse}
in base al tipo di dato. Il lavoro si sposta sul \textbf{model engineering}.
\end{definizione}

\begin{center}
\small
\begin{tabular}{lll}
\toprule
\textbf{Tipo di dato} & \textbf{Architettura} & \textbf{Esempio}\\
\midrule
Vettori numerici / tabelle & Feed-Forward NN (FFNN) & classificare malware da feature\\
Serie temporali / sequenze & Recurrent NN (RNN) & frase, traffico nel tempo\\
Grafi & Graph NN (GNN) & Internet, social network, botnet\\
Immagini (e non solo) & Convolutional NN (CNN) & malware convertito in immagine\\
\bottomrule
\end{tabular}
\end{center}

\section{Breve storia}
\begin{center}
\small
\begin{tabularx}{\textwidth}{l X}
\toprule
1940 & ENIAC, primo computer general purpose (valvole).\\
1950 & Alan Turing propone il \textbf{Turing test}.\\
1955 & Arthur Samuel: ``dare alle macchine la capacità di imparare senza programmarle esplicitamente''; primo giocatore di dama.\\
1958 & Il \textbf{perceptron} (proposto nel 1943) usato per la prima volta in una vera rete neurale.\\
1960 & Metodi \textbf{bayesiani}: diventano lo standard rispetto alle reti neurali.\\
1970--80 & \textbf{AI winter}: poca ricerca, tecnologia non pronta.\\
1980 & \textbf{Backpropagation}: l'algoritmo chiave per addestrare le reti.\\
1985 & Crescita delle reti neurali con più hidden layer.\\
1995 & \textbf{SVM}: rigorose matematicamente, efficienti, superano le reti neurali dell'epoca.\\
1997 & Deep Blue batte Kasparov a scacchi (AI a regole, non ML).\\
2015 & ImageNet: le reti superano l'accuratezza umana; AlphaGo batte Fan Hui a Go.\\
2017 & Google introduce i \textbf{transformer}.\\
2018 & Era degli LLM: GPT (OpenAI), BERT (Google).\\
2021--24 & DALL-E, GPT-4, DeepSeek\ldots\\
\bottomrule
\end{tabularx}
\end{center}

\subsection{Quanto sono complessi i modelli di oggi}
\begin{itemize}
  \item Oltre \textbf{1 trilione di parametri} (``un'equazione con un trilione di coefficienti da tarare'').
  \item Hanno bisogno di molti dati \textbf{buoni}: c'è il rischio concreto di esaurire i dati disponibili.
  \item GPT-4: oltre \textbf{60 milioni di dollari}, più di \textbf{20.000 GPU}, $\sim$100 miliardi di petaFLOP.
        Solo grandi aziende possono permetterselo: non si addestra un LLM sul proprio portatile.
\end{itemize}

\subsection{Perché il deep learning esplode proprio ora}
\begin{enumerate}
  \item \textbf{Tanti dati} (foto, traffico di rete\ldots), ma \textbf{pochi etichettati}: etichettare
        richiede lavoro umano lento e costoso (i CAPTCHA servivano anche a etichettare immagini).
  \item \textbf{GPU}: il cuore delle reti è la moltiplicazione di matrici, che le GPU parallelizzano.
  \item \textbf{Prestazioni} più alte e \textbf{flessibilità} delle architetture.
  \item \textbf{Modelli pre-addestrati} da altri, riutilizzabili e ri-addestrabili (es.\ BERT nei lab).
\end{enumerate}

\section{Shallow learning vs deep learning}
\subsection{La pipeline KDD (shallow)}
Il \textbf{Knowledge Discovery from Data} prevede:
\begin{center}
\begin{tikzpicture}[node distance=0.35cm, every node/.style={font=\scriptsize, align=center},
  blk/.style={draw=polblu, fill=polblu!6, rounded corners=2pt, minimum height=0.9cm, text width=1.55cm}]
  \node[blk] (a) {Data};
  \node[blk, right=of a] (b) {Selection};
  \node[blk, right=of b] (c) {Pre-\\processing};
  \node[blk, right=of c] (d) {Transfor-\\mation};
  \node[blk, right=of d, fill=red!60!black!10, draw=red!60!black] (e) {Feature\\engineering};
  \node[blk, right=of e] (f) {Machine\\learning};
  \node[blk, right=of f] (g) {Interpre-\\tation};
  \foreach \x/\y in {a/b,b/c,c/d,d/e,e/f,f/g} \draw[-{Stealth}] (\x) -- (\y);
\end{tikzpicture}
\end{center}
\begin{itemize}
  \item \textbf{Selection}: si tiene solo il sottoinsieme di dati rilevante (per prevedere il voto
        d'esame serve la carriera, non la residenza).
  \item \textbf{Preprocessing}: \en{data integration} (riconciliare fonti, metadati, conflitti, ridondanza)
        e \en{data cleaning} (rumore, outlier, inconsistenze, valori mancanti, es.\ chi non dichiara il peso
        in un sondaggio). Include la normalizzazione delle scale: età 0--100 e reddito 0--10 milioni non sono
        confrontabili senza standardizzare, altrimenti nella distanza conta solo il reddito.
  \item \textbf{Feature engineering}: l'analista costruisce le feature a mano.
\end{itemize}

\begin{esempio}[frodi bancarie, Intesa Sanpaolo]
Ogni conto era descritto da centinaia di transazioni al mese. Il team ha costruito feature come:
entrate totali del mese, spesa abituale, giorni tra l'accredito e l'esaurimento del denaro, spesa
media settimanale. Poi un modello ML classificava il comportamento come benigno o fraudolento.
Con il deep learning si potrebbe dare in input direttamente la \textbf{serie temporale} delle transazioni.
\end{esempio}

\subsection{La pipeline deep}
Stessa pipeline, \textbf{senza feature engineering}: le feature sono \textbf{estratte dalla rete}
(\en{feature extraction}). Il deep learning è soprattutto \textbf{rappresentazione dei dati}: la rete
impara come rappresentare l'input nel modo migliore.

\begin{center}
\small
\begin{tabularx}{\textwidth}{X X}
\toprule
\textbf{Shallow learning} & \textbf{Deep learning}\\
\midrule
Preprocessing & Preprocessing (\textbf{non si evita!})\\
Feature engineering manuale: \textbf{forte sforzo umano}, molti tentativi & Feature extraction automatica: la rete crea feature spesso migliori di quelle umane\\
Modelli più semplici e \textbf{interpretabili} (``è un'auto perché ha 4 ruote'') & \textbf{Model engineering}: scegliere l'architettura, molti iperparametri\\
Human-driven, labor-intensive & Semi-automatico, computationally intensive\\
\bottomrule
\end{tabularx}
\end{center}

\textbf{Pro del deep learning}
\begin{itemize}
  \item \textbf{Modelli complessi}: pattern complessi su dataset enormi (intrusion detection, malware, anomalie).
  \item \textbf{Efficacia}: superano il plateau del shallow e generalizzano meglio.
  \item \textbf{Casi d'uso}: tabelle, testo, immagini, serie temporali, grafi; ad esempio per un'immagine è
        difficilissimo definire a mano feature che distinguano cane e gatto.
\end{itemize}
\textbf{Contro}
\begin{itemize}
  \item \textbf{Mancanza di trasparenza}: come lavora internamente il modello? (\en{black box})
  \item \textbf{Mancanza di interpretabilità}: perché ha preso quella decisione? Su quali parti
        dell'input si è basato? Se un IDS basato su ML segnala un'attività, come la traduco in una
        \textbf{regola di firewall}? Esistono tecniche di \en{explainable AI} (XAI), fuori dal programma.
\end{itemize}

\begin{esame}[non sparare a una mosca con un cannone]
Il deep learning non è sempre la risposta. Se in una serie temporale c'è un picco evidente, basta una
\textbf{soglia}: evita l'\en{over-engineering}. A volte le soluzioni semplici sono efficaci quanto quelle
complesse, o di più.
\end{esame}

\section{Training, validation e test}
\begin{definizione}[training, validation, testing]
\begin{itemize}
  \item \textbf{Training} (o \en{fitting}, \en{building}): creare il modello facendolo imparare dai dati.
        Nelle librerie è la funzione \code{model.fit(...)}. Usa il \textbf{training set}.
  \item \textbf{Validation}: scegliere i migliori \textbf{iperparametri} (es.\ architettura della rete)
        e valutare l'overfitting, su un \textbf{validation set} \textbf{mai usato per il training}.
  \item \textbf{Testing}: valutare le prestazioni finali su dati \textbf{mai visti} né in training né in
        validation. Serve solo a una cosa: verificare che le prestazioni siano \textbf{simili a quelle in
        validation}, cioè che il modello sia \textbf{generale}.
\end{itemize}
\end{definizione}

\begin{esempio}[l'analogia dell'esame]
Le lezioni sono il training set. Un tema d'esame che non hai mai visto è il validation set: se vai bene,
hai imparato abbastanza. Se impari a memoria le risposte di 10 temi passati senza capire la materia,
all'esame vero (con domande nuove) andrai male: questo è \textbf{overfitting}. L'opposto,
non aver imparato nulla neanche sui dati di training, è \textbf{underfitting}.
\end{esempio}

Il sogno sarebbe un test set enorme; in realtà spesso non si ha nemmeno un test set adeguato.

\section{Panoramica dei paradigmi}
\begin{center}
\small
\renewcommand{\arraystretch}{1.2}
\begin{tabularx}{\textwidth}{>{\bfseries}p{2.6cm} p{2.1cm} X X}
\toprule
Paradigma & Etichette? & Obiettivo & Esempi cyber\\
\midrule
Supervised & Sì & imparare una mappa input $\to$ output dai dati etichettati & spam, classificazione malware, previsione del traffico\\
Unsupervised & No & trovare pattern o struttura nei dati & cluster di URL, anomalie nel traffico\\
Semi-supervised & In parte & sfruttare poche etichette e tanti dati non etichettati & IDS, rilevamento botnet\\
Self-supervised & No & imparare dalla struttura dei dati stessi & analisi di log, malware, generazione di immagini false\\
Reinforcement & No, usa \en{reward} & imparare per tentativi interagendo con un ambiente & tuning di firewall/IDS, attacco e difesa, controllo di un dispositivo\\
\bottomrule
\end{tabularx}
\end{center}

\section{Supervised learning}
\begin{definizione}
Si usa un \textbf{dataset etichettato} per addestrare un modello che impara i pattern delle feature
di ogni etichetta, e poi predice l'etichetta di dati nuovi. È un metodo \textbf{predittivo}:
la \textbf{label} è la \textbf{ground truth}, data e validata da un umano, quindi \textbf{costosa} e
deve essere \textbf{affidabile}.
\end{definizione}
\begin{itemize}
  \item Label \textbf{categorica} $\Rightarrow$ \textbf{classificazione}; label \textbf{numerica} $\Rightarrow$ \textbf{regressione}.
        Molti algoritmi funzionano per entrambe con poche modifiche.
  \item \textbf{Il classificatore impara solo ciò che c'è nei dati passati}: addestrato sul meteo estivo,
        non sa cosa sia $-10$°C.
\end{itemize}

\begin{center}
\small
\begin{tabularx}{\textwidth}{l X X}
\toprule
\textbf{Esempio} & \textbf{Shallow} & \textbf{Deep}\\
\midrule
Spam filter & feature da header, oggetto, testo, metadati & testo grezzo di email e header\\
Malware & DLL importate, chiamate di funzione, sequenze di API call & eseguibile visto come immagine o in nuove rappresentazioni\\
Volume di traffico (regressione) & feature dalla serie (media/max/min ultima ora) & serie temporale direttamente; un picco inatteso può indicare esfiltrazione\\
\bottomrule
\end{tabularx}
\end{center}

\subsection{Problemi e soluzioni}
\begin{itemize}
  \item L'attaccante si adatta: dal ``principe nigeriano'' che chiede soldi al falso ``migliore amico'' che chiede 500€:
        il filtro aveva imparato che la feature importante era ``principe nigeriano''. Malware leggermente modificato
        (API diverse, ordine diverso) non viene più riconosciuto. Traffico che imita quello normale.
  \item \textbf{Drift nei dati}: i dati cambiano nel tempo (le foto di oggi sono 10 volte più pesanti di 3 anni fa)
        o perché l'attaccante li fa cambiare. Ciò che è nuovo viene assegnato a una classe nota (il ``panda'' diventa ``gatto'').
  \item \textbf{Sotto-rappresentazione / class imbalance}: alcune classi hanno pochissimi esempi (una sola foto di panda
        tra migliaia di cani, gatti e polli). In cybersecurity la maggior parte dei dati è benigna, pochi sono malevoli.
\end{itemize}
\textbf{Soluzioni}
\begin{itemize}
  \item \textbf{Drift detection}: monitorare le metriche nel tempo (se ho ancora le label, vedo le prestazioni peggiorare
        $\Rightarrow$ ri-addestrare); se non ho label, confrontare con metodi unsupervised i dati passati e presenti
        $\Rightarrow$ se sono diversi, ri-etichettare e ri-addestrare.
  \item \textbf{Gestione dello sbilanciamento}: \en{class weights} nella loss (sbagliare un malware costa più di un falso allarme),
        metriche adatte, ricampionamento, dati sintetici, tecniche di validazione diverse (vedi cap.~\ref{cap:ffnn}).
\end{itemize}

\section{Unsupervised learning}
\begin{definizione}
Crea modelli che scoprono \textbf{pattern nascosti senza label} e senza intervento umano diretto.
È un metodo \textbf{descrittivo}: descrive la struttura dei dati. Raccogliere dati non etichettati è
\textbf{economico} (es.\ catturare tutti i pacchetti della rete del Politecnico).
\end{definizione}
\begin{itemize}
  \item \textbf{Clustering}: raggruppa punti simili, es.\ trovare gruppi di computer coinvolti in un attacco coordinato
        (\textbf{botnet}) senza conoscerne gli IP; poi si studia il cluster per capire perché è diverso.
  \item \textbf{Anomaly detection}: impari cos'è normale; ciò che si discosta è anomalo (99 foto di cani, gatti e polli
        e un panda: il panda è l'anomalia).
\end{itemize}
\textbf{Problemi}
\begin{itemize}
  \item \textbf{Drift}: le botnet cambiano (nuove macchine infettate, vecchie abbandonate); l'attaccante \textbf{spalma}
        l'esfiltrazione su 5 periodi invece di fare un picco, imitando il traffico normale.
  \item \textbf{Analisi dei pattern}: senza ground truth i risultati sono \textbf{difficili da interpretare}; le metriche
        valutano differenze tra i dati, non la correttezza; più \textbf{falsi allarmi}. Lo sforzo che non fai prima
        (etichettare) lo fai dopo (interpretare). Con il deep learning le dimensioni latenti sono ancora meno interpretabili.
\end{itemize}
\textbf{Soluzioni}: drift detection (ricreare periodicamente il modello e confrontare), analisi manuale, metodi di XAI.

\begin{attenzione}[il ``mantra'' dell'unsupervised]
``La validazione delle strutture di clustering è la parte più difficile e frustrante della cluster analysis.
Senza un forte sforzo in questa direzione, la cluster analysis resterà un'arte oscura accessibile solo ai veri
credenti con esperienza e grande coraggio.'' (Jain \& Dubes, 1988)
\end{attenzione}

\section{Semi-supervised learning}
Usa dati etichettati \textbf{e} non etichettati per un problema \textbf{supervisionato}:
\begin{enumerate}
  \item un \textbf{piccolo dataset etichettato} guida l'estrazione di feature e addestra un primo modello;
  \item il modello etichetta il \textbf{grande dataset non etichettato} $\Rightarrow$ dataset \textbf{pseudo-etichettato};
  \item il modello finale si addestra su dati etichettati + pseudo-etichettati.
\end{enumerate}
Si può usare anche con classificatori shallow. Esempi: IDS addestrato con pochi attacchi etichettati e tanto traffico
normale (se si addestra solo su traffico normale, il resto sarà ``normale'' o ``sconosciuto/anomalo''); rilevamento di
dispositivi infetti da botnet con dati benigni, botnet note e dati non etichettati dell'organizzazione.

\begin{center}
\small
\begin{tabularx}{\textwidth}{X X}
\toprule
\textbf{Benefici} & \textbf{Problemi}\\
\midrule
Unisce la ground truth del supervised e il basso costo dei dati dell'unsupervised & \textbf{Dimensioni}: 1000 etichettati su 1 miliardo; il campione può non essere rappresentativo\\
Migliore generalizzazione: le piccole variazioni non coperte dalle label vengono compensate dai dati non etichettati & Sbilanciamento: se i malevoli sono 1000 su un miliardo, nel campione da 1000 ce ne saranno forse 2\\
 & \textbf{Evoluzione}: i dati non etichettati cambiano; servono approcci ibridi con anomaly detection\\
\bottomrule
\end{tabularx}
\end{center}

\section{Self-supervised learning (SSL)}
\begin{definizione}
Genera un \textbf{segnale supervisionato a partire da dati non etichettati}: sono i dati stessi a fare da
supervisione. È un \textbf{sottoinsieme dell'unsupervised}, ma mentre l'unsupervised \textbf{descrive} i dati,
il self-supervised impara a \textbf{rappresentarli}. È alla base degli \textbf{LLM} moderni.
\end{definizione}

\textbf{Come funziona}
\begin{enumerate}
  \item \textbf{Pretext task}: il modello impara una rappresentazione significativa. Un \textbf{encoder} comprime
        l'input ad alta dimensione (una foto da 30 megapixel) in uno spazio a bassa dimensione, il
        \textbf{bottleneck}. Per dire ``è un cane'' non serve ogni pixel, servono naso, bocca, pelo, quattro zampe:
        è la fase di \en{feature extraction}, fatta dal modello.
  \item \textbf{Downstream task}: si usano le feature del bottleneck per il compito reale:
        \begin{itemize}
          \item un \textbf{decoder} ricostruisce l'input e si confronta con l'originale;
          \item un \textbf{classificatore} predice qualcosa a partire dalla rappresentazione.
        \end{itemize}
\end{enumerate}
\textbf{Intuizione del segnale}: se l'output è molto diverso dall'input, la rappresentazione non è buona
$\Rightarrow$ continua ad addestrare.

\begin{esempio}[denoising delle foto dello smartphone]
La foto mossa o notturna (rumorosa) passa nell'encoder, che impara a tenere l'informazione importante (forma,
dimensioni, colori) e scartare il rumore; il decoder ricrea la foto pulita. Non serve sapere se è un cane o una chiesa:
\textbf{la supervisione viene dall'immagine stessa}. Se la ricostruzione è pessima, forse il bottleneck è troppo piccolo.
È l'idea dell'\en{autoencoder}.
\end{esempio}

\textbf{Applicazioni}
\begin{itemize}
  \item \textbf{Analisi dei log}: la sequenza tipica è apri file $\to$ modifica $\to$ salva $\to$ chiudi. Se dopo la chiusura
        compare ``invia via email'' e il modello prediceva altro, l'evento reale può essere un'anomalia.
  \item \textbf{Malware}: rappresentazioni di eseguibili per migliorare la classificazione con poche etichette.
  \item \textbf{Immagini false} (GAN): un \textbf{generatore} crea immagini sintetiche, un \textbf{discriminatore} cerca di
        capire se sono vere o false; il segnale del discriminatore è la supervisione. Se il discriminatore non distingue più,
        le immagini sembrano vere. (Uso malevolo: falsi annunci Airbnb con foto di case inesistenti.)
\end{itemize}

\begin{center}
\small
\begin{tabularx}{\textwidth}{X X}
\toprule
\textbf{Benefici} & \textbf{Problemi}\\
\midrule
Nessun costo di etichettatura & \textbf{Costo computazionale}: migliaia di GPU per mesi\\
Feature extraction automatica per molti tipi di dato (rete, codice, testo, immagini) & Serve un \textbf{dataset enorme e di qualità} (100 libri pieni di errori non ti insegnano l'italiano)\\
Funziona meglio con dataset enormi & \textbf{Bias} nei dati appreso dal modello (es.\ studi sul sistema giudiziario USA)\\
Più robusto all'evoluzione dei dati: le piccole variazioni si riconoscono & Senza label, \textbf{valutare le prestazioni} è difficile\\
\bottomrule
\end{tabularx}
\end{center}

\section{Reinforcement learning (RL)}
Un \textbf{agente} esegue \textbf{azioni} in un \textbf{ambiente} per massimizzare una \textbf{ricompensa a lungo termine}.
\begin{itemize}
  \item \textbf{Stato}: situazione corrente dell'ambiente; \textbf{azioni}: operazioni possibili;
        \textbf{reward}: aumenta se l'azione è corretta, diminuisce altrimenti (la progettiamo noi).
  \item Si impara per \textbf{tentativi ed errori}: i dati li crea l'agente stesso.
  \item Esempi: agente che adatta dinamicamente regole di firewall o soglie dell'IDS; simulare attacchi per rafforzare le
        difese; imparare a controllare una \textbf{lampadina smart} che accetta comandi in chiaro (idea di come uno stalker
        potrebbe abusarne).
\end{itemize}
\textbf{Benefici}: impara dall'esperienza senza label, risolve problemi complessi e dinamici, migliora continuamente,
alta autonomia. \textbf{Problemi}: costo computazionale; \textbf{progettare la reward} è difficile (se sovrastimi un
miglioramento, l'agente smette di esplorare); apprendimento lento e \textbf{instabile} (la lampadina che si accende e
spegne); bassa interpretabilità.

\section{Come si usano i modelli}
\subsection{Modello singolo vs ensemble}
\begin{itemize}
  \item \textbf{Ensemble}: più modelli addestrati su dati, sottoinsiemi o feature diversi; decisione finale per
        \textbf{voto di maggioranza} (classificazione) o \textbf{media} (regressione). Es.: \textbf{random forest} = insieme di alberi.
  \item \textbf{Model stacking}: le predizioni dei modelli base diventano l'input di un \textbf{meta-learner} che impara
        quale modello è più affidabile in quale situazione. Es.: rilevamento attacchi in sistemi industriali con modelli
        deep di base e un meta-learner shallow.
  \item I modelli base possono avere architetture diverse (FFNN su tabelle, RNN su serie, GNN su grafi).
\end{itemize}
\begin{center}
\small
\begin{tabularx}{\textwidth}{X X}
\toprule
\textbf{Modello singolo} & \textbf{Ensemble}\\
\midrule
+ semplice da addestrare e distribuire & + prestazioni migliori (più ``opinioni'')\\
+ più interpretabile (es.\ decision tree) & + più robusto al rumore, generalizza meglio\\
$-$ overfitting se dataset piccolo o modello troppo complesso & $-$ se i modelli base sono \textbf{correlati} (dati simili) il guadagno è limitato e aumenta il bias\\
$-$ generalizzazione limitata, plateau di prestazioni & $-$ costoso da creare e ottimizzare; poco adatto a dataset piccoli\\
\bottomrule
\end{tabularx}
\end{center}

\subsection{Federated learning}
Un \textbf{modello globale} si costruisce combinando modelli \textbf{addestrati localmente} sui dispositivi
(\en{edge}: smartphone, IoT, server locali) con i soli dati locali. Al server centrale si invia \textbf{il modello,
non i dati}; il modello aggregato viene rimandato ai dispositivi.
\begin{esempio}[spam filter degli studenti]
``Mandatemi tutte le vostre email'' $\Rightarrow$ nessuno accetta. ``Fate girare il modello sul vostro telefono e
mandatemi solo il modello'' $\Rightarrow$ accettabile. Combino 40.000 modelli e ve li rimando: ognuno riceve la conoscenza
dello spam visto dagli altri.
\end{esempio}
\textbf{Benefici}: \textbf{privacy} (utile anche con vincoli GDPR), \textbf{scalabile e distribuito}, \textbf{visibilità}
più ampia (punti di vista eterogenei). \textbf{Problemi}: costo del training sull'edge (batteria), \textbf{dati eterogenei}
(il modello globale può diventare una ``media'' che perde conoscenze specifiche, sbilanciamenti), \textbf{model poisoning}:
un partecipante malevolo invia un modello addestrato a considerare legittime le sue email di phishing.

\subsection{Transfer learning}
Si riusa un modello addestrato per un compito e lo si \textbf{ri-addestra} (\en{fine-tuning}) per un compito diverso ma
\emph{correlato}, con \textbf{molti meno dati}. Di solito si ri-addestrano \textbf{solo gli ultimi layer}.
(Analogia: voi ``modelli'' già formati in ingegneria vi specializzate in questo corso in 60 ore; uno studente del primo
anno ci metterebbe molto di più.)
\begin{esempio}[STAMINA, Intel]
Il binario del malware $\to$ sequenza di byte $\to$ pixel in scala di grigi $\to$ immagine; si riusa
\textbf{ResNet-50} (pre-addestrata su immagini) e la si fa fine-tuning su malware etichettati. Funziona, anche se il
dominio è lontanissimo. Dettagli nel cap.~\ref{cap:cnn-gnn}.
\end{esempio}
\textbf{Benefici}: training più rapido, meno calcolo ed energia, meno dati, prestazioni migliori su dataset piccoli,
funziona tra domini. \textbf{Problemi}: \textbf{domain mismatch} (un modello su cani/gatti/auto può andare male su
raggi X o log), il fine-tuning può comunque costare, il modello eredita i \textbf{bias} di quello pre-addestrato.

\subsection{Modelli generativi}
Generano \textbf{nuovi campioni non presenti nel training set} campionando dalla distribuzione appresa (ChatGPT,
DALL-E\ldots). Spesso addestrati in self-supervised, poi rifiniti con poca supervisione e/o RL.
Esempi cyber: \textbf{varianti sintetiche di malware} per anticipare attacchi; \textbf{agenti} che cercano e sfruttano
zero-day. \textbf{Benefici}: dati sintetici (anche per la privacy), creazione di contenuti, generalità.
\textbf{Problemi}: costo; \textbf{valutazione difficile} (le metriche classiche non si applicano; il modello non è
progettato per un compito, quindi va scoperto in quali sbaglia); bias; \textbf{uso improprio} (deepfake, phishing, malware);
\textbf{privacy} (i dati inviati possono essere usati per l'addestramento: un brevetto non ancora depositato dato a un chatbot
``per correggere l'inglese'' finisce nei dati); questioni etiche e legali (diritti sui contenuti generati); questioni
sociali (perdita di competenze, \en{vibe coding}, copiare all'esame).

\section{AI e cybersecurity: difesa, attacco, adversarial ML}
Tre famiglie:
\begin{enumerate}
  \item \textbf{AI per la difesa}: migliorare le difese e adattarsi più in fretta;
  \item \textbf{AI per l'attacco}: attacchi più sofisticati, realistici, veloci;
  \item \textbf{Adversarial machine learning}: attaccare \textbf{i modelli stessi}.
\end{enumerate}

\subsection{Difesa}
\begin{itemize}
  \item \textbf{Threat detection \& prevention}: IDS basati su AI (fingerprint o pattern tipici), analisi automatica di malware
        senza firme statiche, \textbf{novelty detection} per zero-day, federated learning tra organizzazioni.
  \item \textbf{Threat hunting \& incident response}: \textbf{SIEM} potenziati (correlano log, trovano anomalie in tempo reale);
        \textbf{SOAR} (orchestrano e automatizzano la risposta agli incidenti).
  \item \textbf{Predictive security}: scansione di sistemi e configurazioni per prevedere vulnerabilità (es.\ un server HTTP
        aperto verso l'esterno che non dovrebbe esserci); \textbf{analisi comportamentale} per insider threat.
\end{itemize}
Dati citati: oltre il 90\% dei SOC dipende ancora da processi manuali (``fatica di Sisifo''); chi usa AI e automazione
risparmia in media \textbf{2,22 milioni USD} per breach (IBM).

\subsection{Attacco}
\begin{itemize}
  \item \textbf{Social engineering avanzato}: deepfake audio/video; email di phishing generate da LLM, migliorate per
        tentativi e \textbf{tradotte} in lingue che l'attaccante non conosce (allarga la superficie di vittime).
  \item \textbf{Scoperta automatica di vulnerabilità} e loro sfruttamento; ciò che all'attaccante richiedeva giorni,
        l'AI lo fa in fretta.
  \item \textbf{Tecniche di evasione}: modificare leggermente le feature dell'attacco (es.\ spalmare l'esfiltrazione) per
        sembrare traffico normale.
  \item L'AI generativa \textbf{abbassa la barriera d'ingresso}: non serve saper scrivere malware, basta saper usare il modello.
\end{itemize}

\subsection{Adversarial machine learning}
\begin{itemize}
  \item \textbf{Data poisoning} (in training): iniettare dati malevoli perché il comportamento malevolo diventi ``normale''.
  \item \textbf{Adversarial examples}: modificare leggermente l'\textbf{input} (non il modello) aggiungendo rumore
        impercettibile, così che il modello sbagli (il cane che diventa ``hot dog'' o ``chiesa''). Basta conoscere input
        e output del modello e un po' del suo funzionamento.
  \item \textbf{Model exploitation}: aggirare i \textbf{guardrail} degli LLM (jailbreak, ``exploit della nonna'',
        \en{prompt injection}). Ogni nuovo modello viene jailbreakato in pochi giorni, anche se diventa sempre più difficile.
        Gli \textbf{agenti} con skill e tool (es.\ usare Word) aumentano la \textbf{superficie d'attacco}.
  \item \textbf{Privacy}: \textbf{model inversion}, ricostruire dati di training dagli output del modello. Anche nel
        federated learning, un modello abbastanza grande può far ricostruire le email di partenza.
\end{itemize}

\begin{sintesi}
\begin{itemize}
  \item Attacchi più veloci, complessi e costosi rendono necessaria una difesa \textbf{automatizzata, adattiva e proattiva}: l'AI
        impara il comportamento tipico dai dati storici e segnala le deviazioni (fino, idealmente, agli zero-day).
  \item AI $\supset$ ML $\supset$ reti neurali $\supset$ deep learning. Il ML impara dai dati e \textbf{conosce solo ciò che ha visto};
        la qualità dei dati conta più della quantità (\en{garbage in, garbage out}).
  \item Shallow: feature engineering manuale, modelli interpretabili. Deep: feature extraction automatica, model engineering,
        black box. Il deep learning non è sempre la risposta: evita l'over-engineering.
  \item Training (fit dei pesi) / validation (iperparametri, overfitting) / test (verifica della generalità su dati mai visti).
  \item Paradigmi: supervised (label costose, drift, sbilanciamento), unsupervised (dati economici, risultati difficili da
        interpretare), semi-supervised (pseudo-label), self-supervised (pretext/downstream, base degli LLM), reinforcement (reward).
  \item Uso dei modelli: ensemble e stacking, federated learning (si condivide il modello, non i dati), transfer learning
        (fine-tuning con pochi dati), modelli generativi.
  \item AI per la difesa, AI per l'attacco, e attacchi ai modelli stessi (poisoning, adversarial examples, jailbreak, model inversion).
\end{itemize}
\end{sintesi}

\chapter{Dal perceptron alla rete neurale}
\label{cap:reti-neurali}
\fonti{\rif{2026-09-28}{28/09}}{2-foundament-deep-learning, pp.\ 1--18}

\section{Il quadro generale}
\begin{itemize}
  \item \textbf{Input layer}: riceve i dati. Con dati tabellari, ogni riga è un punto e ogni colonna una feature
        ($x_1, x_2, \ldots$): un neurone di input per feature.
  \item \textbf{Hidden layer}: la parte che ``ragiona'' e modella il problema. Con 3 o più hidden layer si parla di
        \textbf{deep learning}.
  \item \textbf{Output layer}: prende la decisione finale (es.\ cane, gatto o delfino).
\end{itemize}

\begin{attenzione}[le reti mangiano solo numeri]
Testo, grafi e immagini vanno convertiti in rappresentazioni numeriche. Un'immagine è una griglia di pixel e
ogni pixel è descritto da 3 numeri (RGB). Questa conversione è di solito standard; le parti difficili sono:
\begin{itemize}
  \item il \textbf{preprocessing} che rende i dati puliti e utili;
  \item l'\textbf{architettura}: lo stesso problema può essere visto come tabella (FFNN), serie temporale (RNN) o grafo (GNN),
        e va scoperto sperimentalmente quale funziona meglio, quanti layer e quanti neuroni usare
        (\textbf{model engineering}).
\end{itemize}
\end{attenzione}

\section{Il perceptron}
\begin{definizione}[perceptron]
L'unità più semplice di una rete neurale. Prende un input con più feature e:
\begin{enumerate}
  \item pesa ogni feature con un \textbf{peso} $w_i$;
  \item calcola la \textbf{somma pesata};
  \item applica una \textbf{funzione di attivazione} $f(\cdot)$ e produce l'output:
\end{enumerate}
\[
y = f\Big(\sum_{i=0}^{n} w_i x_i\Big) = f(\mathbf{w}^T \mathbf{x}), \qquad x_0 = 1 \text{ per definizione}.
\]
\end{definizione}

\begin{itemize}
  \item I \textbf{pesi} sono i \textbf{parametri} della rete (i ``trilioni di parametri'' degli LLM): indicano
        l'\textbf{importanza} di ogni feature per il compito. All'inizio sono \textbf{casuali} e vengono
        \textbf{appresi dai dati}.
  \item Le $x_i$ non si toccano: sono i valori delle feature, dati dal mondo esterno. Si impara solo $\mathbf{w}$.
  \item Il \textbf{bias} $w_0$ è un peso come gli altri ma \textbf{non dipende da una feature} (moltiplica $x_0=1$):
        serve a \textbf{traslare} la funzione.
  \item Nei layer successivi l'input non sono le feature ma gli \textbf{output del layer precedente}; la struttura del neurone
        è identica.
\end{itemize}

\begin{esempio}[perché serve il bias]
Con attivazione lineare ($f(x)=x$) e un solo input, $y = w_1 x_1$: posso cambiare solo la \textbf{pendenza}, e tutte le
rette passano per l'origine. Se il confine tra croci e punti non passa per l'origine, non posso modellarlo. Aggiungendo
$w_0$ ottengo $y = w_1x_1 + w_0$ e posso spostare la retta.
\end{esempio}

\begin{integrazione}[legame con modelli classici]
Senza $f(\cdot)$ il perceptron è una \textbf{regressione lineare}; con $f = \sigma$ (sigmoid) è una
\textbf{regressione logistica}.
\end{integrazione}

\subsection{Caso 2D con attivazione lineare}
$y = w_1x_1 + w_2x_2 + w_0$ descrive una famiglia di funzioni (un piano/una retta di separazione).
\begin{itemize}
  \item Se tutti i pesi sono 0, l'output è sempre 0: la rete non predice nulla.
  \item Tarando $w_1$ separo lo spazio lungo $x_1$; con $w_2$ inclino il confine; con $w_0$ lo sposto.
  \item L'output grezzo non è limitato: con $x_2=10$ può valere molto più di 10. Per avere una
        \textbf{probabilità} serve una funzione limitata in $[0,1]$.
\end{itemize}

\section{Funzioni di attivazione}
Servono a:
\begin{enumerate}
  \item \textbf{imporre proprietà} all'output (es.\ sigmoid $\Rightarrow$ output in $[0,1]$);
  \item introdurre \textbf{non linearità} nel modello;
  \item decidere \textbf{quando il neurone è attivo} (ReLU: se l'input è $\leq 0$ il neurone è spento);
  \item altri motivi di ottimizzazione: convergenza più rapida, \en{sparsity}.
\end{enumerate}

\begin{center}
\small
\renewcommand{\arraystretch}{1.35}
\begin{tabularx}{\textwidth}{l l l X}
\toprule
\textbf{Funzione} & \textbf{Formula} & \textbf{Codominio} & \textbf{Uso tipico e note}\\
\midrule
Sigmoid & $\sigma(y)=\dfrac{1}{1+e^{-y}}$ & $(0,1)$ & output di classificazione \textbf{binaria}; limitata $\Rightarrow$ \en{vanishing gradient} se usata negli hidden layer\\
Tanh & $\tanh(y)$ & $(-1,1)$ & limitata; usata nelle \textbf{RNN}\\
ReLU & $\max(0,y)$ & $[0,\infty)$ & la più usata negli \textbf{hidden layer}; non limitata superiormente; neuroni inattivi $\Rightarrow$ sparsità\\
Leaky ReLU & $y$ se $y>0$, $\alpha y$ altrimenti & $\R$ & piccola pendenza per $y<0$: anche i neuroni ``negativi'' portano informazione\\
Softmax & $\dfrac{e^{z_i}}{\sum_j e^{z_j}}$ & $(0,1)$, somma 1 & output \textbf{multi-classe}\\
Lineare & $f(y)=y$ & $\R$ & output di \textbf{regressione}\\
\bottomrule
\end{tabularx}
\end{center}

\subsection{Sigmoid}
Mappa qualsiasi valore reale in $[0,1]$: anche se la somma pesata vale $-10.000$, l'output è limitato.
$P(\text{pos}\mid x) = \sigma(y)$ e $P(\text{neg}\mid x) = 1 - P(\text{pos}\mid x)$, quindi la somma è 1 per costruzione.
Soglia tipica: $0{,}5$. Con tutti i pesi a 0 l'output è $\sigma(0)=0{,}5$: il modello non sa decidere.

\textbf{Limiti della sigmoid}
\begin{enumerate}
  \item \textbf{Solo binaria}: per cane/gatto/delfino non basta.
  \item \textbf{Vanishing gradient}: in addestramento si manda indietro (\en{backpropagation}) un segnale che dipende dalle
        derivate delle attivazioni. Moltiplicando valori $<1$ per molti layer, il segnale \textbf{svanisce} e i primi layer non
        imparano (dettagli nel cap.~\ref{cap:addestramento}).
\end{enumerate}

\subsection{ReLU}
Se l'input è $\leq 0$ l'output è 0 (neurone \textbf{inattivo} per quella combinazione di feature); altrimenti l'output è
proporzionale all'input. Non essendo limitata, il segnale backpropagato resta abbastanza forte anche dopo molti layer.
Benefici: mitiga il vanishing gradient $\Rightarrow$ \textbf{convergenza più rapida}; \textbf{sparsità}: non tutti i neuroni
servono per tutti gli input.

\begin{esempio}[sparsità e LLM]
Un LLM da trilioni di parametri non li usa tutti a ogni richiesta: alcuni neuroni sono specializzati (es.\ ``inglese'' più
che ``francese''). È collegato all'idea di \en{mixture of experts} e di \en{routing} dell'informazione nella rete.
\end{esempio}

\begin{attenzione}[ReLU non risolve tutto]
Moltiplicando molti valori $>1$ il gradiente può \textbf{esplodere} (\en{exploding gradient}): il problema opposto.
Succede anche con reti non troppo profonde. Reti diverse usano attivazioni diverse o combinazioni di attivazioni.
\end{attenzione}

\section{Attivazioni in uscita e logit}
\begin{definizione}[logit]
L'output della rete \textbf{prima} dell'attivazione finale: valori non normalizzati in $(-\infty, +\infty)$.
Vanno trasformati in probabilità valide (in $[0,1]$ e con somma 1).
\end{definizione}

\subsection{Classificazione binaria}
La rete produce \textbf{un logit} $z$ (per convenzione della classe positiva). Si applica la sigmoid:
$P(\text{pos}\mid x) = \sigma(z) = \frac{1}{1+e^{-z}}$ e $P(\text{neg}\mid x) = 1 - P(\text{pos}\mid x)$.
Con logit $-10.000$ la probabilità di positivo è $\approx 0$; con $+10.000$ è $\approx 1$.

\subsection{Classificazione multi-classe}
Servono almeno $C-1$ output, ma in pratica si usa \textbf{un neurone per classe} ($C$ logit). Si applica la
\textbf{softmax}:
\[
P(c_i \mid x) = \frac{e^{z_i}}{\sum_{j} e^{z_j}}, \qquad P(c_i\mid x) \in [0,1], \qquad \sum_i P(c_i \mid x) = 1.
\]
\begin{itemize}
  \item È una generalizzazione della sigmoid e usa \textbf{tutti gli output insieme}.
  \item Converte i logit in probabilità e \textbf{amplifica} le piccole differenze tra logit.
  \item Esempio: immagine $\to$ CNN (rappresentazione) $\to$ FFNN (classificazione) $\to$ 4 logit $\to$ softmax $\to$
        \{cane, gatto, uccello, tartaruga\}; si sceglie la probabilità più alta.
\end{itemize}

\subsection{Regressione}
Si predicono numeri reali: di solito non serve imporre proprietà e l'attivazione di uscita è l'\textbf{identità}
$f(x)=x$. È praticamente l'unico caso in cui ha senso un'attivazione lineare.

\begin{attenzione}[teoria vs software]
In teoria in classificazione l'ultima attivazione c'è sempre. In \textbf{PyTorch} però, con \code{CrossEntropyLoss},
\textbf{non} si mette la softmax nel modello perché è già dentro la loss (cap.~\ref{cap:ffnn}). È una scelta della libreria:
leggi sempre la documentazione.
\end{attenzione}
\section{Il margine: quale confine è migliore?}
Tra due classificatori lineari che separano correttamente i punti di training, è preferibile quello con
\textbf{margine maggiore} dai punti: un nuovo punto vicino al confine viene classificato correttamente con più probabilità.
Un confine troppo ``attaccato'' a una classe sbaglierà facilmente nei dintorni. È l'idea delle \textbf{SVM}
(cap.~\ref{cap:anomaly}).

\section{Verso reti più grandi}
\subsection{Impilare perceptron: il layer fully connected}
Due perceptron con lo stesso input $(x_1,x_2)$ e pesi diversi ($\mathbf{w}$ e $\mathbf{q}$) danno due output indipendenti:
\[
\mathbf{y} = \begin{pmatrix} y_1\\ y_2\end{pmatrix}
= f\left(\begin{pmatrix} w_0 & w_1 & w_2\\ q_0 & q_1 & q_2\end{pmatrix}
\begin{pmatrix} 1\\ x_1\\ x_2\end{pmatrix}\right) = f(W\mathbf{x})
\]
Questo è un \textbf{layer lineare (fully connected)} con \textbf{6 parametri}. Per confronto, i modelli ``piccoli''
di oggi ne hanno 7 miliardi (es.\ Mistral 7B).

\begin{definizione}[numero di parametri]
È il numero totale di \textbf{pesi} (bias inclusi). Un layer lineare da $n_{in}$ a $n_{out}$ neuroni ha
$n_{in} \cdot n_{out} + n_{out}$ parametri.
\end{definizione}

\subsection{Aggiungere layer}
Un secondo layer combina $y_1, y_2$ con pesi $\mathbf{s}$ e produce $z$:
\[
z = f\big(\mathbf{s}^T f(W^T\mathbf{x})\big)
\]
\textbf{Una rete neurale è moltiplicazione di matrici} alternata a funzioni di attivazione: nessuna magia, ma con
trilioni di pesi.

\textbf{Perché tanti layer?} Ogni layer crea una rappresentazione \textbf{più astratta} dei dati: i primi layer
lavorano sul dato grezzo, i successivi catturano i concetti principali (sarà evidentissimo con le CNN).

\textbf{L'attivazione è la stessa in tutti i layer?} No. Negli hidden layer serve evitare il vanishing gradient
$\Rightarrow$ ReLU (non limitata); nell'ultimo layer serve una probabilità $\Rightarrow$ sigmoid o softmax.

\section{Perché la non linearità è indispensabile}
Se $f(x)=x$, allora $z = \mathbf{s}^T W^T \mathbf{x} = W'\mathbf{x}$:
\begin{enumerate}
  \item potevo usare direttamente $W' = W\mathbf{s}$ e ottenere lo stesso output;
  \item il secondo layer non serviva;
  \item il modello resta \textbf{lineare}.
\end{enumerate}
Mille layer lineari equivalgono a \textbf{un solo layer lineare} (\en{collapse}). Il mondo non è lineare (pensa alla
distribuzione dei voti di un esame), quindi servono le attivazioni non lineari tra un layer e l'altro.

Con ReLU nella rete del docente (hidden ReLU + output sigmoid) il confine di decisione non è più una retta: si formano
\textbf{più regioni}. Il numero di regioni dipende da quanti neuroni ci sono e dalla presenza della non linearità.

\section{Teorema di approssimazione universale}
\begin{definizione}[Cybenko 1989; Hornik et al.\ 1989]
Per ogni funzione continua $g$ definita su un sottoinsieme compatto di $\R^n$ e per ogni $\epsilon > 0$, esiste una
rete feed-forward con \textbf{un solo hidden layer} e un numero finito di neuroni che approssima $g$ con errore inferiore a $\epsilon$.
\end{definizione}
\begin{itemize}
  \item Ogni neurone ReLU introduce un \textbf{punto di rottura} (\en{breakpoint}): la funzione approssimata è lineare a tratti.
        Con un input e $D$ neuroni si ottengono fino a $D+1$ \textbf{regioni lineari}. Nell'esempio delle slide, la curva viene
        approssimata sempre meglio con 4, 9 e infine \textbf{19 neuroni} ($\Rightarrow$ 20 regioni lineari): una rete
        \en{shallow} basta.
  \item Il numero di neuroni di una rete shallow misura la sua \textbf{capacità}.
  \item Il problema: il teorema \textbf{non dice quanti neuroni servono} (4? 19? 20.000?) né quali pesi usare.
\end{itemize}

\begin{esame}[meglio profonde che larghe]
Si usano reti \textbf{più profonde e più strette} invece di un solo layer larghissimo: i layer, separati da non linearità,
\textbf{combinano la loro complessità} e creano rappresentazioni sempre più astratte $\Rightarrow$ migliore
\textbf{generalizzazione}. ``Il deep learning è rappresentazione dei dati.''\\
Attenzione però: una rete troppo complessa può andare in \textbf{overfitting} (cap.~\ref{cap:addestramento}).
\end{esame}

\subsection{Complessità di una rete}
Il numero di regioni lineari cresce rapidamente con il numero di neuroni e di \textbf{feature di input}.
Con 500 neuroni e 100 feature si superano le $10^{107}$ regioni, con \textbf{51.001 parametri}: tantissime regioni, eppure
la rete è minuscola per gli standard attuali (Mistral 7B ha 7 miliardi di parametri).


\begin{sintesi}
\begin{itemize}
  \item Il perceptron calcola $y=f(\mathbf{w}^T\mathbf{x})$: pesi appresi dai dati, bias $w_0$ per traslare il confine,
        attivazione $f$ per imporre proprietà e non linearità.
  \item Attivazioni: ReLU negli hidden layer (mitiga il vanishing gradient, sparsità); sigmoid (binaria), softmax
        (multi-classe) o identità (regressione) in uscita. I \textbf{logit} sono gli output prima dell'attivazione finale.
  \item Un layer fully connected è $f(W\mathbf{x})$, con $n_{in}n_{out}+n_{out}$ parametri; una rete è moltiplicazione di
        matrici alternata a non linearità. Senza non linearità, mille layer collassano in uno solo.
  \item Teorema di approssimazione universale: basta un hidden layer, ma non dice quanti neuroni né quali pesi.
        In pratica si preferiscono reti \textbf{profonde e strette}, che generalizzano meglio.
\end{itemize}
\end{sintesi}

\chapter{Addestramento: ottimizzazione, loss e valutazione}
\label{cap:addestramento}
\fonti{\rif{2026-09-30}{30/09}}{2-foundament-deep-learning, pp.\ 19--49}
\section{Il ciclo di addestramento}
Chiamiamo $\phi$ l'insieme di \textbf{tutti i pesi e bias} della rete. All'inizio sono casuali e vanno tarati sul
\textbf{training set} (il validation set serve solo a verificare che i pesi funzionino anche su altri dati).
\begin{enumerate}
  \item \textbf{Forward pass}: i dati attraversano la rete dall'input all'output e si ottengono le predizioni $\hat{y}$.
  \item Si calcola l'\textbf{errore} tra predizione $\hat{y}$ e label vera $y$: è la \textbf{loss}.
  \item \textbf{Backpropagation}: il segnale d'errore viene mandato indietro, dall'output verso l'input, per capire come
        cambiare i pesi layer per layer.
  \item \textbf{Aggiornamento dei pesi} secondo un \textbf{ottimizzatore}.
\end{enumerate}
\begin{definizione}[epoch]
Un passaggio completo dell'\textbf{intero training set} attraverso la rete. Si addestra per molte epoch
(10, 100, 1000\ldots a seconda di dati e compito).
\end{definizione}

\section{La funzione di loss}
\begin{definizione}[loss function $\mathcal{L}(\phi)$]
La funzione obiettivo che misura l'errore e che si vuole \textbf{minimizzare}. Dipende dal compito
(classificazione, regressione\ldots). Esempio per la regressione, il \textbf{Mean Squared Error}:
\[
\mathcal{L}(\phi) = \mathrm{MSE}(\phi) = \frac{1}{N}\sum_{i=1}^{N}\big(y_i - \phi^T\mathbf{x}_i\big)^2
\]
\end{definizione}
\begin{itemize}
  \item Formalmente è $\mathcal{L}(\phi, X, y)$, ma $X$ e $y$ sono \textbf{fissi}: non si possono modificare i dati per far
        quadrare le predizioni (non trasformi la foto del gatto in un cane perché il modello ha detto ``cane''). Si agisce solo su $\phi$.
  \item Loss $=0$ vorrebbe dire predizioni perfette: non succede mai, e se succede bisogna \textbf{insospettirsi} (overfitting).
\end{itemize}

\subsection{Soluzione in forma chiusa vs iterativa}
Per modelli semplici (regressione lineare) la loss è una parabola convessa: si risolve analiticamente con
$\partial\mathcal{L}/\partial\phi = 0$, oppure provando valori diversi di $\phi$ e prendendo il migliore. Con un trilione di
parametri non si può fare né l'una né l'altra cosa: si parte da $\phi$ \textbf{casuale} e ci si \textbf{muove} secondo
una \textbf{policy} (un algoritmo di ottimizzazione).

\subsection{Loss per la classificazione binaria: Binary Cross-Entropy}
\[
\mathcal{L} = -\,y\log\hat y - (1-y)\log(1-\hat y), \qquad y\in\{0,1\},\ \hat y = \text{model}(x)\in[0,1]
\]
La label $y$ fa da \textbf{selettore}:
\begin{itemize}
  \item $y=1$: $\mathcal{L} = -\log\hat y$ $\Rightarrow$ loss bassa se $\hat y\approx 1$, molto alta se $\hat y\approx 0$;
  \item $y=0$: $\mathcal{L} = -\log(1-\hat y)$ $\Rightarrow$ loss bassa se $\hat y\approx 0$, molto alta se $\hat y\approx 1$.
\end{itemize}
Per il multi-classe si generalizza nella \textbf{cross-entropy}.

\subsection{Quale loss per quale problema}\label{sec:quale-loss}
\begin{center}
\small
\renewcommand{\arraystretch}{1.4}
\begin{tabularx}{\textwidth}{l l X}
\toprule
\textbf{Problema} & \textbf{Loss} & \textbf{Note}\\
\midrule
Regressione & $\mathrm{MSE} = \frac1N\sum (y_i-\hat y_i)^2$ & semplice, \textbf{differenziabile} ovunque; il quadrato amplifica gli outlier\\
Regressione & $\mathrm{MAE} = \frac1N\sum |y_i-\hat y_i|$ & \textbf{meno sensibile agli outlier}; derivata problematica in 0\\
Binaria & BCE $= -\frac1N\sum [y_i\log\hat y_i + (1-y_i)\log(1-\hat y_i)]$ & differenziabile, molto usata\\
Multi-classe & CE $= -\sum_{i=1}^{C} y_i\log\hat y_i$ & $y$ one-hot della label vera, $\hat y$ output della softmax\\
\bottomrule
\end{tabularx}
\end{center}

\section{Gradient descent (GD)}
\[
\phi_{t+1} = \phi_t - \alpha \nabla_\phi \mathcal{L}(\phi_t)
\]
\begin{itemize}
  \item Il \textbf{gradiente} indica la direzione in cui la loss \textbf{sale}; si va nella direzione \textbf{opposta}.
  \item La visione è solo \textbf{locale}: non si conosce la forma globale della loss, solo come cambia intorno al punto attuale.
  \item Funziona benissimo con funzioni \textbf{convesse} (un solo minimo); con funzioni \textbf{non convesse} può fermarsi in un
        \textbf{minimo locale}.
\end{itemize}

\subsection{Il learning rate}
\begin{definizione}[learning rate $\alpha$]
Controlla \textbf{quanto} si aggiornano i pesi a ogni passo: è la dimensione del passo nella direzione indicata dal
gradiente. Il gradiente dice \emph{la direzione}, $\alpha$ dice \emph{quanto muoversi}.
\end{definizione}
\begin{itemize}
  \item $\alpha$ piccolo (es.\ $0{,}0001$): passi minuscoli, ottimizzazione fine ma lenta.
  \item $\alpha$ grande: passi lunghi, si esplora in fretta ma si rischia di \textbf{scavalcare} il minimo e oscillare.
\end{itemize}


\begin{esempio}[regressione lineare $y = x\phi_1 + \phi_0$]
Con 12 punti di training, partendo da un punto casuale e seguendo la direzione di massima discesa, in circa 4 iterazioni si
arriva vicino al minimo: la retta (il modello) migliora a ogni passo.
\end{esempio}

\subsection{Quando fermarsi}
\begin{itemize}
  \item il \textbf{miglioramento è trascurabile} (10\%, 3\%, 1\%, 0,1\%, 0,0000001\%: l'ultimo non vale il tempo);
  \item \textbf{budget}: numero massimo di epoch o di tempo;
  \item si è raggiunto un minimo (miglioramento nullo);
  \item la loss di \textbf{validation} inizia a peggiorare (\en{early stopping}, cap.~\ref{cap:regolarizzazione}).
\end{itemize}

\subsection{Limiti del GD}
\begin{enumerate}
  \item \textbf{Sensibile all'inizializzazione}: inizializzazioni diverse $\Rightarrow$ soluzioni diverse.
  \item Non funziona bene con loss \textbf{non convesse} (sempre, nella pratica): si blocca in minimi locali.
  \item Usa \textbf{tutti i dati} per decidere la direzione.
\end{enumerate}
\begin{esempio}[esperienza del docente]
Previsione con una rete piccola (2 layer, $\sim$64 neuroni) contro una SVM: la rete era leggermente migliore, ma
cambiando \textbf{solo l'inizializzazione casuale} le prestazioni cambiavano drasticamente. La casualità non si controlla,
quindi la rete era inaffidabile: si è scelta la SVM, un po' peggiore ma \textbf{stabile}.
\end{esempio}

\section{Mini-batch e stochastic gradient descent}
Usare tutti i dati insieme prima di ogni aggiornamento ha due problemi:
\begin{enumerate}
  \item \textbf{tempo}: con 1 TB di dati devo far passare tutto il dataset prima di correggere i pesi una sola volta;
  \item \textbf{memoria}: il dataset non entra nella RAM della GPU, quindi andrebbe caricato a pezzi e si aspetterebbe
        la fine di tutti i pezzi per aggiornare.
\end{enumerate}
\begin{definizione}[mini-batch]
Un sottoinsieme di punti \textbf{campionati a caso senza reinserimento} dal training set. I pesi si aggiornano
\textbf{dopo ogni mini-batch}; l'epoch resta il tempo per processare \textbf{tutto} il dataset.
\end{definizione}
\begin{itemize}
  \item \textbf{Senza reinserimento}: ogni punto compare in \textbf{un solo} mini-batch per epoch, così tutti i punti pesano
        uguale (un punto ripetuto in 100 mini-batch diventerebbe artificialmente importante).
  \item Vantaggi: aggiornamenti \textbf{più frequenti}, \textbf{meno memoria}, migliore \textbf{generalizzazione} e
        \textbf{scalabilità}. È la strategia che userete sempre.
\end{itemize}

\subsection{Stochastic gradient descent (SGD)}
Il gradiente si calcola su un \textbf{mini-batch casuale} $B_t$ alla volta:
\[
\phi_{t+1} = \phi_t - \alpha \sum_{i\in B_t}\frac{\partial \mathcal{L}_i(\phi_t)}{\partial\phi}
\]
\begin{itemize}
  \item Mini-batch diversi indicano direzioni diverse: \textbf{in media} si scende, ma il singolo passo può non essere quello
        di massima discesa, e \textbf{può anche salire} rispetto alla loss globale.
  \item Proprio questo permette di \textbf{uscire dai minimi locali} e saltare da una ``valle'' all'altra.
  \item La traiettoria è rumorosa (``cammina come un ubriaco''), ma nell'esempio delle slide arriva al minimo globale
        anche partendo dall'inizializzazione in cui GD si bloccava.
  \item È \textbf{più veloce} del GD: con 3 mini-batch decide la direzione in un terzo del tempo; sbaglia a volte, ma
        il più delle volte va bene.
\end{itemize}

\textbf{Proprietà dell'SGD}
\begin{enumerate}
  \item Aggiunge rumore ma la loss migliora sull'intera epoch.
  \item Mini-batch senza reinserimento: tutti gli esempi contribuiscono in egual modo.
  \item Più veloce del GD: aggiorna senza aspettare la fine del dataset.
  \item Meno costoso: il gradiente su un sottoinsieme; il mini-batch entra in memoria.
  \item Può (in principio) \textbf{uscire dai minimi locali} (non è garantito).
  \item Riduce il rischio di bloccarsi vicino a \en{saddle point}: qualche batch avrà sempre un gradiente significativo.
  \item In pratica trova parametri che \textbf{generalizzano bene}.
\end{enumerate}
\begin{attenzione}
L'SGD non ``converge'' nel senso classico (può risalire). L'idea: vicino al minimo globale tutti i punti sono ben
descritti, quindi il gradiente è piccolo qualunque batch si scelga.
\end{attenzione}

\subsection{Learning rate scheduler}
Spesso l'SGD si usa con $\alpha$ \textbf{alto all'inizio}, ridotto di un fattore costante ogni $N$ epoch:
\begin{itemize}
  \item all'inizio si \textbf{esplora} lo spazio dei parametri e si salta tra valli;
  \item alla fine si è vicini al minimo e si fa \textbf{fine-tuning} con passi piccoli.
\end{itemize}
Un $\alpha$ \textbf{estremamente} grande resta comunque dannoso: si salta da una parte all'altra senza imparare.
La dimensione del passo dipende da $\alpha$, non dalla dimensione del batch.

\subsection{Dimensione del mini-batch}\label{sec:batch-size}
Si usano di solito \textbf{potenze di 2} per efficienza computazionale.
\begin{center}
\small
\begin{tabularx}{\textwidth}{l X X}
\toprule
\textbf{Batch size} & \textbf{Pro} & \textbf{Contro / quando}\\
\midrule
1 (streaming) & aggiornamenti rapidi, pochissima memoria & molto rumoroso, convergenza lenta (troppi aggiornamenti, spesso sbagliati)\\
16--64 (piccolo) & buona generalizzazione con dataset piccoli & aggiornamenti lenti; per dataset piccoli e hardware limitato\\
128--512 (medio) & buon compromesso prestazioni/tempo & il più comune\\
1024+ (grande) & training veloce, gradiente più liscio & generalizzazione peggiore; solo con dataset molto grandi\\
\bottomrule
\end{tabularx}
\end{center}

\section{Momentum}
Invece di decidere ogni mini-batch da zero, si \textbf{ricorda la direzione dei passi precedenti}:
\[
m_{t+1} = \beta\, m_t + (1-\beta)\sum_{i\in B_t}\frac{\partial\mathcal{L}_i(\phi_t)}{\partial\phi},
\qquad \phi_{t+1} = \phi_t - \alpha\, m_{t+1}
\]
\begin{itemize}
  \item $\beta$ = quanto pesa il \textbf{passato} rispetto al presente; $1-\beta$ = quanto pesa il presente.
  \item Effetto: \textbf{traiettoria più liscia}, meno oscillazioni nelle valli.
  \item $\beta$ troppo basso: quasi nessuna memoria $\Rightarrow$ rumore (passeggiata casuale).
        $\beta$ troppo alto: si segue quasi solo il passato $\Rightarrow$ aggiornamenti lenti, si continua in una direzione
        che ora non è più la migliore. In entrambi i casi la convergenza rallenta.
\end{itemize}

\section{Learning rate fisso e normalizzazione del gradiente}
Con $\alpha$ fisso si fanno passi \textbf{grandi} sui parametri con gradiente grande (già molto esplorati) e passi
\textbf{piccoli} su quelli con gradiente piccolo (che andrebbero esplorati di più).
\begin{itemize}
  \item $\alpha$ piccolo: si converge in fretta sul parametro che influenza di più la loss, l'altro si muove pochissimo.
  \item $\alpha$ grande: entrambi si muovono, ma in modo molto instabile.
\end{itemize}
\textbf{Normalizzazione del gradiente} (usando tutti i dati):
\[
m_{t+1} = \frac{\partial\mathcal{L}(\phi_t)}{\partial\phi}, \qquad
v_{t+1} = \Big(\frac{\partial\mathcal{L}(\phi_t)}{\partial\phi}\Big)^2, \qquad
\phi_{t+1} = \phi_t - \alpha\frac{m_{t+1}}{\sqrt{v_{t+1}}+\epsilon}
\]
\begin{itemize}
  \item $\epsilon$ piccolo evita la divisione per zero.
  \item Il risultato dà in pratica il \textbf{segno} del gradiente per ogni coordinata: ci si muove di una distanza fissa
        $\alpha$, dando \textbf{la stessa importanza a tutte le direzioni}.
  \item Prezzo: si può \textbf{oscillare} attorno al minimo.
\end{itemize}

\section{Adam (Adaptive Moment Estimation)}
Combina momentum e normalizzazione: momentum sia sul gradiente sia sul gradiente al quadrato (di solito su mini-batch).
\begin{align*}
m_{t+1} &= \beta\, m_t + (1-\beta)\,\frac{\partial\mathcal{L}(\phi_t)}{\partial\phi} && \text{(primo momento)}\\
v_{t+1} &= \gamma\, v_t + (1-\gamma)\Big(\frac{\partial\mathcal{L}(\phi_t)}{\partial\phi}\Big)^2 && \text{(secondo momento)}\\
\hat m_{t+1} &= \frac{m_{t+1}}{1-\beta^{t+1}}, \qquad \hat v_{t+1} = \frac{v_{t+1}}{1-\gamma^{t+1}} && \text{(bias correction)}\\
\phi_{t+1} &= \phi_t - \alpha\,\frac{\hat m_{t+1}}{\sqrt{\hat v_{t+1}}+\epsilon}
\end{align*}
\begin{itemize}
  \item $\beta, \gamma \in [0,1)$ sono i coefficienti di momentum.
  \item \textbf{Bias correction}: $m$ e $v$ partono da 0, quindi all'inizio sono stime irrealisticamente piccole. Il fattore
        $1-\beta^{t+1}$ compensa: è significativo nei primi passi e trascurabile dopo (l'esponente fa decadere la correzione).
  \item Risultato: traiettoria liscia, convergenza più rapida e stabile, buona anche con gradienti sparsi.
  \item Costo: $\alpha$, $\beta$, $\gamma$ hanno valori di default ma sono \textbf{iperparametri} da tarare.
  \item \textbf{Adam/AdamW} è l'ottimizzatore che userete di più.
\end{itemize}
\begin{integrazione}[AdamW]
AdamW è Adam con \en{weight decay} (regolarizzazione L2) \textbf{disaccoppiato} dall'aggiornamento adattivo; in PyTorch si
attiva con il parametro \code{weight\_decay}.
\end{integrazione}

\begin{sintesi}[ottimizzatori]
\begin{center}
\small
\begin{tabular}{lll}
\toprule
\textbf{Ottimizzatore} & \textbf{Risolve} & \textbf{Introduce}\\
\midrule
GD & --- & minimi locali, lento, tutti i dati\\
SGD & velocità, memoria, minimi locali & traiettoria rumorosa\\
SGD + momentum & rumore della traiettoria & iperparametro $\beta$\\
Normalizzazione & passi sbilanciati tra parametri & oscillazioni vicino al minimo\\
Adam & tutto quanto sopra & iperparametri $\alpha,\beta,\gamma$\\
\bottomrule
\end{tabular}
\end{center}
Ottimizzatore, batch size, learning rate e coefficienti di momentum sono \textbf{iperparametri}: influenzano le prestazioni
ma sono distinti dai parametri e dall'architettura. Sceglierli è ``più arte che scienza'': si addestrano molti modelli e si
sceglie il migliore (\en{hyperparameter search}, es.\ \en{grid search} su una griglia di valori).
\end{sintesi}

\section{Backpropagation}
Ogni combinazione di loss e modello richiederebbe un calcolo del gradiente diverso, e derivare in un colpo solo rispetto a un
trilione di parametri è impraticabile. La \textbf{backpropagation} è ``solo un nome elegante'' per la \textbf{regola della
catena}:
\[
\frac{\partial f}{\partial x} = \frac{\partial f}{\partial g}\cdot\frac{\partial g}{\partial x}
\qquad\text{ovvero}\qquad \big(f(g(x))\big)' = f'(g(x))\,g'(x)
\]
applicata dalla fine del \textbf{grafo computazionale} all'indietro. Il gradiente di un layer si ottiene
\textbf{riusando quello del layer successivo} moltiplicato per il gradiente locale: si lavora un layer alla volta.

\begin{definizione}[grafo computazionale]
Grafo orientato in cui ogni nodo è un'operazione e ogni arco il flusso di dati. Esempio $y = wx + q$:
prima $a = wx$, poi $y = a + q$. PyTorch lo costruisce e lo deriva in automatico.
\end{definizione}

\begin{esempio}[backprop passo per passo]
Un punto $(x,y)$, modello (strano) $\hat y = \phi_1\phi_2 x$, loss $\mathcal{L} = (\phi_1\phi_2x - y)^2$.
Variabili intermedie: $a=\phi_1\phi_2$, \ $b = ax$, \ $c = b - y$, \ $\mathcal{L} = c^2$.

\textbf{Forward}: si calcola $\mathcal{L}$ dagli input $\phi_1,\phi_2,x,y$.

\textbf{Backward}:
\begin{align*}
\frac{\partial\mathcal{L}}{\partial c} &= 2c\\
\frac{\partial\mathcal{L}}{\partial b} &= \frac{\partial\mathcal{L}}{\partial c}\frac{\partial c}{\partial b} = 2c\cdot 1 = 2c\\
\frac{\partial\mathcal{L}}{\partial a} &= \frac{\partial\mathcal{L}}{\partial b}\frac{\partial b}{\partial a} = 2c\,x\\
\frac{\partial\mathcal{L}}{\partial \phi_1} &= \frac{\partial\mathcal{L}}{\partial a}\frac{\partial a}{\partial \phi_1} = 2c\,x\,\phi_2 = 2(\phi_1\phi_2x-y)\,x\,\phi_2\\
\frac{\partial\mathcal{L}}{\partial \phi_2} &= 2c\,x\,\phi_1 = 2(\phi_1\phi_2x-y)\,x\,\phi_1
\end{align*}
\end{esempio}

\section{L'importanza dell'attivazione: vanishing gradient}
Il gradiente di ogni layer è il prodotto dei gradienti dei layer successivi: se un layer dà $0{,}1$ e il successivo $0{,}5$,
si ottiene $0{,}05$. Le attivazioni \textbf{limitate} (sigmoid, tanh) producono derivate piccole $\Rightarrow$ il prodotto
di tanti numeri piccoli tende a zero.
\begin{definizione}[vanishing gradient]
Durante la backpropagation i gradienti diventano troppo piccoli e ``svaniscono'': i pesi dei layer vicini all'input ricevono
aggiornamenti minimi e restano quasi uguali all'inizializzazione casuale. Allora a cosa servono quei layer?
\end{definizione}
\begin{esempio}[sigmoid vs ReLU, 4 layer]
Dall'output (L4) al primo hidden layer (L1): con la \textbf{sigmoid} il gradiente cala di circa \textbf{3 ordini di
grandezza}; con la \textbf{ReLU} resta molto più stabile.
\end{esempio}
\begin{integrazione}
La derivata della sigmoid vale al massimo $0{,}25$ (in $y=0$): per questo, anche nel caso migliore, ogni layer sigmoid
riduce il gradiente di almeno un fattore 4.
\end{integrazione}
La famiglia ReLU permette un flusso di gradiente migliore. Nelle RNN però vedremo che la ReLU risolve un problema ma ne crea
un altro (cap.~\ref{cap:rnn}).

\section{Il flusso di lavoro in pratica}
\begin{enumerate}
  \item \textbf{Definire problema e dati}: obiettivo (classificazione, regressione, clustering, anomaly detection,
        generazione\ldots) $\Rightarrow$ loss; preprocessing (pulizia, augmentation, bilanciamento, trasformazione); split
        train/validation/test; batch size.
  \item \textbf{Model engineering}: scegliere l'architettura (tipo di dato, quanti layer, quanti neuroni).
  \item Scegliere \textbf{ottimizzatore} e \textbf{loss}.
  \item \textbf{Addestrare} e (se possibile) \textbf{validare}, facendo \en{hyperparameter tuning}: valutare le prestazioni e
        studiare underfitting e overfitting.
\end{enumerate}
Batch size e loss si scelgono secondo le tabelle delle sezioni~\ref{sec:batch-size} e~\ref{sec:quale-loss}.

\section{Metriche di valutazione}
\begin{attenzione}[loss $\neq$ modello migliore]
La loss serve alla macchina per addestrarsi, ma tra due modelli \textbf{quello con loss più bassa non è per forza il
migliore}: posso avere loss bassa sbagliando poco sulla classe facile (normale) e molto su quella che conta (malware).
La scelta finale si fa con le \textbf{metriche}.
\end{attenzione}

\begin{definizione}[metriche di classificazione]
\[
\text{Accuracy} = \frac{\text{predizioni corrette}}{\text{tutte le predizioni}}
\]
Per ogni classe $C_i$:
\[
\text{Precision}(C_i) = \frac{TP_i}{TP_i + FP_i}, \qquad
\text{Recall}(C_i) = \frac{TP_i}{TP_i + FN_i}, \qquad
F1(C_i) = \frac{2\cdot P\cdot R}{P + R}
\]
\begin{itemize}
  \item \textbf{Precision}: degli elementi \emph{predetti} come $C_i$, quanti lo sono davvero? (punti di altre classi predetti come $C_i$ = $FP$)
  \item \textbf{Recall}: degli elementi che \emph{sono} $C_i$, quanti ne ho trovati? (punti di $C_i$ predetti come altro = $FN$)
  \item \textbf{F1}: media armonica, \textbf{schiacciata verso il valore più basso} dei due.
\end{itemize}
\end{definizione}

\begin{esempio}[il classificatore pigro]
9.900 goodware e 100 malware. Il classificatore risponde sempre ``goodware'':
\begin{itemize}
  \item accuracy $= 9900/10000 = 99\%$ (sembra ottimo);
  \item malware: recall $= 0/100 = 0$; precision $= 0/0$, non definita. Il classificatore è inutile.
\end{itemize}
Se risponde sempre ``malware'': recall(malware) $= 1$, ma precision(malware) $= 100/10000 = 0{,}01$: quasi tutti falsi
allarmi, e nessuno si fiderà più del modello.
\end{esempio}

\begin{esame}[quale metrica?]
Con dati \textbf{sbilanciati} l'accuracy non va usata. Si sceglie un compromesso tra precision e recall in base al costo
degli errori: per la classe malware si vuole \textbf{recall alta} con una precision \textbf{decente} (altrimenti nessuno
crede agli allarmi). In caso di dubbio, l'\textbf{F1}.
\end{esame}

\section{Learning curve: underfitting e overfitting}
Si osservano le \textbf{learning curve}: loss di training e di validation in funzione delle epoch.
\begin{itemize}
  \item Il training set serve a \textbf{cambiare} i pesi; sul validation set si calcola la loss \textbf{senza mai} modificare i pesi.
  \item Situazione buona: entrambe le loss scendono e si stabilizzano; quella di training è un po' più bassa, quella di
        validation simile.
\end{itemize}

\begin{definizione}[underfitting]
Il modello è \textbf{troppo semplice}: la loss resta \textbf{alta anche sul training set}, scende poco, e la validation è
similmente alta. Le metriche sono pessime. \emph{Attenzione alla scala}: scendere da 1 a 0,9 è poco, da 0,15 a 0,05 è molto.
\end{definizione}

\begin{definizione}[overfitting]
Il modello è perfetto (o quasi) \textbf{solo sui dati di training}: la loss di training continua a scendere (anche verso 0)
mentre quella di \textbf{validation, dopo essere scesa, risale}. Il modello ha memorizzato il dataset invece di imparare.
\end{definizione}

\subsection{Rimedi all'underfitting}
\begin{itemize}
  \item \textbf{Preprocessing migliore}: togliere rumore, gestire outlier, cambiare l'encoding delle feature categoriche.
  \item \textbf{Model engineering}: più complessità (più neuroni), architettura diversa, più epoch, batch size diverso.
\end{itemize}

\subsection{Rimedi all'overfitting}
Si agisce con la \textbf{regolarizzazione} (penalità sui pesi, dropout, early stopping, batch normalization, noise injection),
trattata nel cap.~\ref{cap:regolarizzazione}. Se la rete è enorme e memorizza i dati, la cura è ridurne la complessità.

\begin{sintesi}
\begin{itemize}
  \item Ciclo: forward pass $\to$ loss $\to$ backpropagation (regola della catena sul grafo computazionale) $\to$ aggiornamento
        dei pesi con un ottimizzatore; un'epoch è un passaggio su tutto il training set.
  \item Loss: MSE/MAE per la regressione, BCE per la binaria, cross-entropy per il multi-classe. Si agisce solo sui pesi $\phi$.
  \item GD $\to$ SGD su mini-batch senza reinserimento (più veloce, può uscire dai minimi locali) $\to$ momentum $\to$
        normalizzazione del gradiente $\to$ \textbf{Adam/AdamW} (l'ottimizzatore di riferimento). Learning rate, batch size e
        coefficienti di momentum sono iperparametri.
  \item Attivazioni limitate $\Rightarrow$ vanishing gradient: i primi layer non imparano; la famiglia ReLU lo mitiga.
  \item Si sceglie il modello con le \textbf{metriche}, non con la loss: con classi sbilanciate niente accuracy, ma
        precision/recall/F1 per la classe che conta.
  \item Learning curve: underfitting = loss alta anche in training; overfitting = validation loss che risale mentre la training
        scende (rimedi nel cap.~\ref{cap:regolarizzazione}).
\end{itemize}
\end{sintesi}

\chapter{Regolarizzazione: dropout, batch normalization e oltre}
\label{cap:regolarizzazione}
\fonti{\rif{2026-09-30}{30/09}, \rif{2026-10-06}{06/10}, \rif{2026-10-20}{20/10}}{2-foundament-deep-learning, pp.\ 50--51;
2a-dropout\_batchnorm, pp.\ 1--7; notebook FFNN}

\section{Overfitting e regolarizzazione}
Una rete va in \textbf{overfitting} quando è (quasi) perfetta solo sui dati di training: la loss di training continua a scendere,
anche verso 0, mentre quella di \textbf{validation}, dopo essere scesa, \textbf{risale} (learning curve e diagnosi nel
cap.~\ref{cap:addestramento}). Il modello ha memorizzato il dataset invece di imparare.

\begin{definizione}[regolarizzazione]
Insieme di tecniche che aggiungono \textbf{vincoli o penalità} al processo di apprendimento per ridurre l'overfitting e migliorare la
\textbf{generalizzazione}. È cruciale nel deep learning perché le reti neurali hanno una grande capacità di \textbf{memorizzare} i
dati di training. Le penalità si aggiungono alla loss:
\[
\mathcal{L} = \mathcal{L}_{std} + \text{penalty term}.
\]
\end{definizione}

Le slide raggruppano le tecniche secondo la ``leva'' su cui agiscono:
\begin{center}
\small
\begin{tabularx}{\textwidth}{l l X}
\toprule
\textbf{Si agisce su\ldots} & \textbf{Tecnica} & \textbf{Idea}\\
\midrule
Pesi & L1 (Lasso), L2 (Ridge) / weight decay & penalità nella loss sul valore dei pesi\\
Architettura & Dropout & spegnere a caso neuroni durante il training\\
Processo di training & Early stopping & fermarsi quando la validation peggiora\\
Dati & Batch normalization & normalizzare gli input dei layer con le statistiche del mini-batch\\
Dati & Noise injection & rumore casuale su input, pesi o attivazioni\\
\bottomrule
\end{tabularx}
\end{center}

\section{Regolarizzazione sui pesi: L1 e L2}
\begin{itemize}
  \item \textbf{L1 (Lasso)}: penalità proporzionale al \textbf{valore assoluto} dei pesi, $\lambda\sum_i|\phi_i|$. Spinge alcuni pesi
        \textbf{a zero} $\Rightarrow$ sparsità, una sorta di selezione delle feature: parti diverse della rete si specializzano.
  \item \textbf{L2 (Ridge)}: penalità proporzionale al \textbf{quadrato} dei pesi, $\lambda\sum_i\phi_i^2$. Scoraggia i
        \textbf{pesi grandi} e distribuisce l'importanza su più feature.
  \item $\lambda$ è la \textbf{forza della regolarizzazione} (un iperparametro).
  \item In PyTorch la L2 si ottiene con il parametro \code{weight\_decay} dell'ottimizzatore (es.\ \code{AdamW}, vedi
        cap.~\ref{cap:addestramento}); un valore tipico di partenza indicato nelle linee guida di progettazione è $0{,}01$
        (cap.~\ref{cap:ffnn}).
\end{itemize}
\begin{integrazione}[sul ``rumore'']
A lezione L1/L2 sono state descritte come ``aggiungere un po' di rumore'' all'aggiornamento dei pesi. È un'intuizione:
tecnicamente sono \textbf{penalità} deterministiche nella loss che impediscono ai pesi di adattarsi perfettamente al training
set. L'aggiunta di rumore vera e propria è la \en{noise injection} (sezione~\ref{sec:noise-injection}).
\end{integrazione}

\section{Dropout}
\begin{definizione}[dropout]
Durante ogni iterazione di training si ``spegne'' (si pone a zero l'output di) un sottoinsieme \textbf{casuale} di neuroni.
I neuroni non possono diventare troppo dipendenti gli uni dagli altri $\Rightarrow$ la rete impara \textbf{feature più robuste}.
In predizione si usano \textbf{tutti} i neuroni.
\end{definizione}

\subsection{Cosa succede in training}
\begin{itemize}
  \item A ogni forward pass di training, l'output di ogni neurone è posto a 0 con probabilità $p$.
  \item La scelta è \textbf{indipendente per ogni campione} del mini-batch, e nuova a ogni mini-batch: ogni campione è elaborato da una
        \textbf{sotto-rete diversa}. Un neurone spento ora può essere acceso dopo (non si ``congela'' una volta per tutte).
  \item Gli output sopravvissuti sono \textbf{scalati di} $1/(1-p)$, così il valore atteso dell'output resta uguale con e senza dropout.
  \item Obiettivo: i neuroni non possono fare affidamento su altri neuroni specifici $\Rightarrow$ feature più robuste, meno overfitting.
  \item Il dropout \textbf{non rimuove} neuroni né pesi: azzera temporaneamente degli output.
\end{itemize}

\subsection{Dropout e backpropagation}
\begin{itemize}
  \item Forward, backward e aggiornamento avvengono \textbf{una volta per mini-batch}; il gradiente usato è la \textbf{media} sul mini-batch.
  \item Tutti i pesi vengono aggiornati, ma \textbf{ogni neurone impara solo dai campioni in cui era acceso}: si riduce la capacità di
        apprendimento e ogni neurone vede campioni diversi.
\end{itemize}

\subsection{Training vs validation/test}
\begin{center}
\small
\begin{tabularx}{\textwidth}{X X}
\toprule
\textbf{Training} (\code{model.train()}) & \textbf{Validation/Test} (\code{model.eval()})\\
\midrule
neuroni casuali spenti, nuova maschera a ogni passaggio & tutti i neuroni usati, nessuna maschera\\
sopravvissuti scalati di $1/(1-p)$ & nessuno scaling: il dropout diventa l'identità\\
output casuale: stesso input, output diversi & output deterministico, usa i pesi appresi\\
output 2, $p=0{,}5$ $\to$ 0 oppure 4 (in media 2) & output 2 $\to$ 2 (stesso valore atteso)\\
\bottomrule
\end{tabularx}
\end{center}
\begin{attenzione}[\code{no\_grad} non basta]
\code{torch.no\_grad()} disattiva solo i gradienti: è \code{model.eval()} che spegne il dropout. Se si valuta senza
\code{model.eval()}, le predizioni restano casuali.
\end{attenzione}

\subsection{Quanto dropout?}
Con $p=0{,}5$ metà dei neuroni del layer ha output 0 durante un passo di training. Secondo il docente $0{,}5$ è probabilmente troppo:
conviene provare valori tra $0{,}2$ e $0{,}5$ (stessa indicazione nelle tabelle di progettazione di FFNN, RNN e GNN:
cap.~\ref{cap:ffnn}, \ref{cap:rnn}, \ref{cap:cnn-gnn}).

\begin{integrazione}[validation loss più bassa della training con dropout]
Nel notebook d'esempio, con dropout $0{,}5$ la validation loss risulta \textbf{più bassa} di quella di training. È normale: in
training il dropout spegne neuroni e peggiora la loss, in validation la rete lavora al completo.
\end{integrazione}

\section{Batch normalization}
Un layer tra due layer che \textbf{ri-normalizza} i dati mentre attraversano la rete: normalizza gli input di uno o più layer per
\textbf{stabilizzare e velocizzare} il training. Riduce l'overfitting rendendo il training \textbf{meno sensibile all'inizializzazione}
dei pesi e si può applicare anche \textbf{prima} della funzione di attivazione.

\subsection{Cosa succede in training}
Per ogni \textbf{feature}, sull'output del layer precedente (\code{N} campioni $\times$ \code{F} feature):
\begin{enumerate}
  \item media $\mu_B$ e varianza $\sigma^2_B$ \textbf{del mini-batch corrente} (statistiche calcolate lungo il batch);
  \item standardizzazione: $\hat x = \dfrac{x-\mu_B}{\sqrt{\sigma^2_B+\epsilon}}$ $\Rightarrow$ media 0, varianza 1;
  \item scala e traslazione: $y = \gamma\hat x + \beta$, con $\gamma, \beta$ \textbf{appresi} (parametri della rete, aggiornati dalla
        backpropagation): la rete impara quanto conta la normalizzazione.
\end{enumerate}
\begin{itemize}
  \item Batch diversi $\Rightarrow$ $\mu, \sigma$ diversi: l'output di un campione \textbf{dipende dagli altri campioni del batch}
        (da qui il nome).
  \item A ogni mini-batch si aggiornano anche \textbf{media e varianza mobili} (\en{running}), \textbf{non apprese} ma calcolate:
        \[ \text{running} = 0{,}9\cdot\text{running} + 0{,}1\cdot\text{statistica del batch} \]
\end{itemize}

\subsection{Training vs validation/test}
\begin{itemize}
  \item \textbf{Training}: usa $\mu, \sigma$ del mini-batch corrente e aggiorna le running statistics.
  \item \textbf{Validation/Test}: usa le \textbf{running mean/variance memorizzate} (congelate). La BN \textbf{non sparisce}: diventa una
        scala e traslazione fissa. Senza, il layer successivo riceverebbe in validation dati con una distribuzione diversa da quella
        vista in training.
\end{itemize}
\begin{esempio}[stesso valore, output diversi]
Valore $x=6$ di una feature, $\gamma=1$, $\beta=0$:
\begin{center}
\small
\begin{tabular}{llcc}
\toprule
\textbf{Fase} & \textbf{Batch} & \textbf{Statistiche} & \textbf{Output}\\
\midrule
Training & A $=[2,4,6,8]$ & $\mu=5$, $\sigma\approx2{,}24$ & $+0{,}45$\\
Training & B $=[6,10,12,14]$ & $\mu=10{,}5$, $\sigma\approx2{,}96$ & $-1{,}52$\\
Validation/Test & qualsiasi & running $\mu=5$, $\sigma\approx2{,}24$ & $+0{,}45$\\
\bottomrule
\end{tabular}
\end{center}
Stesso input, output diversi in training: questo ``rumore di batch'' dà un lieve effetto di \textbf{regolarizzazione}.
\end{esempio}

\begin{nota}[Dropout e batch norm in PyTorch]
\begin{center}
\small
\begin{tabularx}{\linewidth}{l X X l}
\toprule
& \code{model.train()} & \code{model.eval()} & \textbf{Appreso in backprop}\\
\midrule
Dropout & maschera casuale per campione, scala $1/(1-p)$ & identità: tutti i neuroni, nessuno scaling & niente ($p$ fisso)\\
BatchNorm & $\mu,\sigma$ del batch + aggiorna le running stats & running $\mu,\sigma$ (congelate) & $\gamma, \beta$\\
\bottomrule
\end{tabularx}
\end{center}
\end{nota}

\section{Early stopping}
Si monitora la loss (o la performance) sul \textbf{validation set}: quando inizia a peggiorare mentre quella di training continua a
migliorare, si \textbf{interrompe} il training, prima che il modello si adatti troppo ai dati di training. Si usa il modello
\textbf{prima} dell'overfitting (all'epoch $t$, non $t+10$).
\begin{itemize}
  \item \code{min\_delta}: miglioramento minimo della validation loss per considerarlo significativo.
  \item \code{patience}: numero massimo di epoch consecutive senza miglioramento significativo prima di fermarsi (5, 10, 20\ldots).
  \item Più lunga la patience, più a lungo si addestra; più grande il min\_delta, più forte deve essere il miglioramento.
        Una richiesta troppo severa (calo di 0,5 in 2 epoch) ferma tutto subito.
  \item \code{best\_model\_state}: si \textbf{salvano i pesi del modello migliore} e li si ricarica alla fine. Senza, si
        restituirebbe il modello dell'ultima epoch (già in overfitting), oppure bisognerebbe ri-addestrare da zero fino all'epoch
        migliore: ore perse con dataset grandi.
\end{itemize}
\begin{lstlisting}
import copy
best_val, wait = float('inf'), 0
patience, min_delta = 20, 1e-4
best_state = None
for epoch in range(num_epochs):
    ...                                   # training + validation, calcola val_loss
    if val_loss < best_val - min_delta:   # miglioramento significativo
        best_val, wait = val_loss, 0
        best_state = copy.deepcopy(model.state_dict())
    else:
        wait += 1
        if wait >= patience:
            break                         # stop
model.load_state_dict(best_state)         # ripristina il modello migliore
\end{lstlisting}

\section{Noise injection}\label{sec:noise-injection}
Si aggiunge \textbf{rumore casuale} a input, pesi o attivazioni durante il training per rendere il modello più robusto.
È citata nelle slide ma non la useremo nel corso.

\section{Regolarizzazione nel codice}
In PyTorch dropout e batch norm sono \textbf{layer} come gli altri, definiti in \code{\_\_init\_\_}; la L2 è un parametro
dell'ottimizzatore.
\begin{lstlisting}
class RegularizedNN(nn.Module):
    def __init__(self, input_size, num_classes, p=0.5):
        super().__init__()
        self.fc1 = nn.Linear(input_size, 128)
        self.fc2 = nn.Linear(128, 128)
        self.bn2 = nn.BatchNorm1d(128)        # batch normalization
        self.drop = nn.Dropout(p)             # dropout: si definisce come un layer
        self.fc3 = nn.Linear(128, num_classes)
        self.relu = nn.ReLU()
    def forward(self, x):
        x = self.relu(self.fc1(x))
        x = self.relu(self.bn2(self.fc2(x)))  # BN prima della prossima trasformazione
        x = self.drop(x)
        return self.fc3(x)

optimizer = optim.AdamW(model.parameters(), lr=1e-3, weight_decay=1e-2)  # L2
\end{lstlisting}

\begin{esempio}[cosa mostra il notebook d'esempio]
Sul dataset giocattolo della lezione del \rif{2026-10-06}{06/10} (pipeline completa nel cap.~\ref{cap:ffnn}):
\begin{itemize}
  \item \textbf{early stopping} (patience 20, min\_delta piccolo): il training si ferma molto prima;
  \item \textbf{dropout 0,5}: curva più stabile; la validation loss è più bassa di quella di training;
  \item \textbf{batch norm + weight decay}: curva molto più rumorosa. Le prestazioni restano buone solo perché l'esempio è facile:
        ``se vedo una validation loss così nel vostro report, penso che non abbiate valutato la soluzione''.
\end{itemize}
\end{esempio}

\section{Quando la regolarizzazione non basta}
\begin{attenzione}
Se l'overfitting è forte perché la rete è enorme e memorizza i dati, i rimedi ``leggeri'' non bastano: riduci layer e
neuroni, oppure ripensa preprocessing e architettura.
\end{attenzione}
Tra gli iperparametri, il \code{weight\_decay} è quello da toccare \textbf{se c'è overfitting}; il learning rate resta il primo su cui
lavorare (cap.~\ref{cap:ffnn}).

\begin{sintesi}
\begin{itemize}
  \item La regolarizzazione aggiunge vincoli o penalità all'apprendimento ($\mathcal{L} = \mathcal{L}_{std} + \text{penalty}$) per
        combattere l'overfitting.
  \item \textbf{Pesi}: L1 (sparsità, pesi a zero) e L2 (pesi piccoli; in PyTorch \code{weight\_decay}).
  \item \textbf{Dropout}: in training azzera ogni output con probabilità $p$ (maschera per campione, scala $1/(1-p)$); in
        \code{eval()} è l'identità. Valori tipici $0{,}2$--$0{,}5$.
  \item \textbf{Batch norm}: standardizza ogni feature con $\mu_B,\sigma_B$ del mini-batch, poi $\gamma\hat x+\beta$ appresi; in
        \code{eval()} usa le running statistics congelate. Stabilizza il training, lieve effetto regolarizzante.
  \item \textbf{Early stopping}: \code{patience} e \code{min\_delta} sulla validation loss, e si ripristinano i pesi migliori.
  \item Ricorda \code{model.train()}/\code{model.eval()}: è ciò che accende e spegne dropout e statistiche di batch.
\end{itemize}
\end{sintesi}

\chapter{Data handling e feed-forward neural network}
\label{cap:ffnn}
\fonti{\rif{2026-10-05}{05/10}, \rif{2026-10-06}{06/10}}{3-deep-learning-methodologies, pp.~1--33; notebook FFNN (Lab~1)}

Questo capitolo segue la pipeline di un progetto deep dall'inizio alla fine: prima i dati (pre-processing, trasformazione,
split), poi la progettazione di una \en{feed-forward neural network} (FFNN) e infine la sua implementazione in PyTorch.
La teoria di ottimizzazione, loss e metriche è nel cap.~\ref{cap:addestramento}; dropout, batch norm, weight decay ed early
stopping nel cap.~\ref{cap:regolarizzazione}.

% =====================================================================
\section{Il deep learning non fa magie sui dati}
\label{sec:ffnn-pipeline}
La pipeline deep semplifica la feature extraction, ma \textbf{non è completamente automatica}: i dati vanno conosciuti e
lavorati prima di darli al modello. Il \textbf{data handling} (split, preprocessing, trasformazione) occupa il
\textbf{50--60\% del tempo} anche in un progetto deep (\rif{2026-10-05}{05/10}).

\begin{attenzione}[dati sensibili]
In azienda spesso \textbf{non potrete} caricare i dati su ChatGPT o simili: fareste uscire informazioni sensibili
(``e vi cercherete un altro lavoro''). Dovete saperli trattare voi.
\end{attenzione}

\begin{center}
\begin{tikzpicture}[node distance=0.5cm, every node/.style={font=\small, align=center},
  blk/.style={draw=polblu, fill=polblu!6, rounded corners=2pt, minimum height=0.9cm, minimum width=2.4cm}]
  \node[blk] (s) {Split};
  \node[blk, right=of s] (p) {Pre-processing};
  \node[blk, right=of p] (t) {Transformation};
  \node[blk, right=of t, fill=esecol!10, draw=esecol] (m) {Model\\engineering};
  \draw[-{Stealth}] (s) -- (p); \draw[-{Stealth}] (p) -- (t); \draw[-{Stealth}] (t) -- (m);
  \node[fit=(s)(p)(t), draw=polblu, dashed, inner sep=6pt, label={[font=\scriptsize]above:Data handling}] {};
\end{tikzpicture}
\end{center}
L'ordine corretto è \textbf{prima lo split}, poi preprocessing e trasformazione appresi sul training set (vedi
sez.~\ref{sec:leak-norm}).

% =====================================================================
\section{Pre-processing (con Pandas)}
\begin{itemize}
  \item \textbf{Rumore}: alcune architetture lo puliscono, altrimenti influenza il modello; troppo rumore $\Rightarrow$ modello troppo
        generico $\Rightarrow$ underfitting. In genere il rumore pesa meno degli outlier.
  \item \textbf{Outlier}: valori estremi non rappresentativi.
\end{itemize}
\begin{esempio}[Elon Musk]
Normalizzando il reddito tra 0 e 1: 0 è il bambino di 3 anni senza reddito, 1 è Elon Musk. Tutti noi finiamo
schiacciati vicino a 0, indistinguibili dal bambino: un solo outlier rovina la feature. Meglio trattarlo come outlier;
il nuovo massimo (es.\ 1 milione) copre tutto l'intervallo.
\end{esempio}
\begin{itemize}
  \item \textbf{Ridondanza}: feature che sono \textbf{proxy} l'una dell'altra (prezzo e prezzo dopo le tasse) complicano solo la rete;
        i \textbf{duplicati} aumentano artificialmente l'importanza di certi punti e riducono la generalizzazione.
  \item \textbf{Valori mancanti}: una rete non può ricevere un punto con un valore mancante (il neurone di input riceve per forza un valore). Opzioni:
  \begin{enumerate}
    \item \textbf{eliminare i punti}: rapido, ma attenzione a \emph{quali} punti togli: se l'1\% con valori mancanti è proprio
          la classe malevola, la elimini;
    \item \textbf{stimare il valore}: dalla distribuzione della feature, o meglio dalle altre feature del punto;
    \item \textbf{ignorare la feature}: se manca nel 50\% dei punti, si toglie la colonna e si tengono tutti i punti;
    \item \textbf{sostituire con tutti i valori possibili} pesati per probabilità, creando modelli diversi: più complesso e costoso.
  \end{enumerate}
\end{itemize}

% =====================================================================
\section{Trasformazione dei dati}
\begin{center}
\small
\begin{tabularx}{\textwidth}{l X}
\toprule
\textbf{Tipo} & \textbf{Esempi}\\
\midrule
Categorico \textbf{nominale} & colori, paese, algoritmo di cifratura, \textbf{porta di rete}, parole\\
Categorico \textbf{ordinale} & titolo di studio, taglia, livello di accesso, livello di minaccia\\
Numerico & temperatura, altezza, byte inviati, numero di pacchetti, tentativi di password\\
\bottomrule
\end{tabularx}
\end{center}

\subsection{Feature categoriche}
Qui il quadro introdotto il 05/10; la trattazione completa (ordinali vs nominali, embedding, token \code{<UNK>}) è nel
cap.~\ref{cap:embedding}.
\begin{definizione}[one-hot encoding]
Ogni valore diventa un vettore binario con tanti elementi quanti sono i valori possibili: tutti 0 tranne un 1.
\textbf{Un neurone di input per ogni valore possibile.}
\end{definizione}
\begin{center}
\small
\begin{tabular}{cc|ccc}
\toprule
ID & Algoritmo & EA\_AES & EA\_RSA & EA\_DES\\
\midrule
1 & AES & 1 & 0 & 0\\
2 & RSA & 0 & 1 & 0\\
3 & DES & 0 & 0 & 1\\
4 & AES & 1 & 0 & 0\\
\bottomrule
\end{tabular}
\end{center}
\textbf{Pro}: semplice ed efficace con poche categorie, non impone ordini (il colore degli occhi non ha un ordine utile).
\textbf{Contro}: esplode con alta cardinalità: il vocabolario inglese ha oltre 171.000 parole $\Rightarrow$ 171.000 neuroni
di input solo per le parole. Non converge.

\begin{definizione}[label encoding]
Assegna un intero a ogni categoria (Low = 1, Medium = 2, High = 3). \textbf{Un solo neurone di input.}
\end{definizione}
\textbf{Pro}: semplice, efficiente in memoria, e se la feature è ordinale \textbf{insegna} alla rete che Low è vicino a
Medium e lontano da High. \textbf{Contro}: \textbf{implica un ordine} che può ingannare il modello.
\begin{esempio}[le porte non sono numeri]
Le porte 80, 81 e 8080 sono numeri, ma trattarle come tali insegna che 80 e 81 sono ``vicine'' e 8080 ``lontanissima''. In realtà
8080 è un'alternativa comune a 80 per i siti web, mentre 81 non c'entra. Le porte sono \textbf{categoriche}.
\end{esempio}
\textbf{Uso tipico}: feature ordinali e \textbf{la label} della classificazione: l'ordine non conta perché la
cross-entropy non calcola distanze tra classi, e i logit danno la probabilità di ogni classe.

\begin{definizione}[embedding]
Si crea uno \textbf{spazio latente} in cui ogni valore della feature è mappato su un \textbf{vettore denso} di dimensione
fissa (es.\ 32, 64, 128). \textbf{Un neurone di input per ogni coordinata} dello spazio latente. Si può fare solo con
architetture deep.
\end{definizione}
\begin{itemize}
  \item \textbf{Pro}: rappresenta feature ad \textbf{alta cardinalità} con un vettore piccolo (171.000 parole in 128 dimensioni);
        \textbf{concetti simili finiscono vicini}: la distanza tra ``man'' e ``woman'' è simile a quella tra ``king'' e ``queen''.
  \item \textbf{Contro}: può essere costoso e richiede \textbf{abbastanza dati} per l'addestramento.
  \item Differenza con il one-hot: il one-hot è \textbf{sparso} (un solo 1); l'embedding è \textbf{denso}, ogni coordinata è
        un numero diverso da 0 con un significato appreso. Con il label encoding alfabetico ``King'' (K) e ``Queen'' (Q)
        sarebbero lontani senza motivo.
\end{itemize}
Approfondimento completo nel cap.~\ref{cap:embedding}.

\subsection{Feature numeriche (con scikit-learn)}
Un neurone di input per feature, ma i numeri grezzi creano problemi: meglio trasformarli.
\begin{center}
\renewcommand{\arraystretch}{1.6}
\begin{tabularx}{\textwidth}{l c X}
\toprule
\textbf{Tecnica} & \textbf{Formula} & \textbf{Pro / Contro}\\
\midrule
Normalizzazione (min-max) & $z = \dfrac{x - x_{min}}{x_{max}-x_{min}}$ & + intervallo fisso $[0,1]$\newline $-$ \textbf{sensibile agli outlier}\\
Standardizzazione (z-score) & $z = \dfrac{x-\mu}{\sigma}$ & + funziona anche con outlier\newline $-$ nessun limite preciso\\
\bottomrule
\end{tabularx}
\end{center}
\begin{attenzione}[errore nella slide]
Nella slide (e a voce) $\sigma$ è chiamata ``varianza'': è la \textbf{deviazione standard} (radice della varianza).
\end{attenzione}
Perché gli outlier contano con il min-max: se un valore normale diventa $0{,}00001$, il neurone fa $x\cdot w$ e la feature
diventa trascurabile a meno di pesi enormi.

% =====================================================================
\section{Splitting dei dati}
\begin{itemize}
  \item \textbf{Training}: impara i pesi.
  \item \textbf{Validation}: sceglie gli iperparametri in base alle metriche.
  \item \textbf{Test} (se possibile): verifica che il modello sia generale, cioè con prestazioni simili alla validation.
        Se non lo è, si riparte da zero; il test \textbf{non} si usa per scegliere iperparametri.
\end{itemize}
Problemi da gestire: \textbf{distribuzione delle classi}, \textbf{costo degli errori}, \textbf{dimensione del dataset},
\textbf{data leakage}.

\subsection{Tecniche di campionamento}
\begin{itemize}
  \item \textbf{Stratified sampling} senza reinserimento: ogni partizione mantiene le \textbf{stesse proporzioni di classe}
        del dataset. Con 40 rossi e 60 verdi divisi a metà: 20+30 e 20+30. Così il validation set contiene esempi di tutte le
        classi e posso valutare il modello su ognuna.
  \item \textbf{Bootstrap}: campionamento \textbf{con reinserimento}, utile con dataset o classi molto piccoli (un punto può
        comparire in più split).
\end{itemize}

\subsection{Protocolli di valutazione}
\begin{definizione}[hold-out]
Partizioni fisse, es.\ \textbf{70\% train, 20\% validation, 10\% test}. Il training è il più grande perché serve per tarare i
pesi; la validation deve bastare a giudicare gli iperparametri; il test serve solo a verificare. Adatto a \textbf{dataset grandi};
si può ripetere qualche volta con split diversi, non molte (costa).
\end{definizione}
\begin{center}
\small
\begin{tabular}{lccc}
\toprule
& \textbf{Train (70\%)} & \textbf{Val (20\%)} & \textbf{Test (10\%)}\\
\midrule
Rossi (40) & 28 & 8 & 4\\
Verdi (60) & 42 & 12 & 6\\
\bottomrule
\end{tabular}
\end{center}

\begin{definizione}[k-fold cross-validation]
Si divide il dataset in $k$ \textbf{fold} stratificati disgiunti; si addestra su $k-1$ fold e si valida sul restante;
si ripete $k$ volte cambiando il fold di validation, \textbf{con gli stessi iperparametri}.
\end{definizione}
\begin{itemize}
  \item Si ottengono $k$ prestazioni: se sono \textbf{simili}, gli iperparametri sono buoni e il modello generalizza; se variano
        molto, quegli iperparametri non vanno bene.
  \item Modello finale: non si sceglie ``il fold migliore'' (sarebbe solo fortuna sul fold di validation): si
        \textbf{ri-addestra su tutti i fold} con gli iperparametri scelti, e poi lo si valuta sul test.
  \item Stima affidabile, ma \textbf{non adatta a dataset molto grandi} (10 training da 80 milioni di punti) né a dataset
        molto piccoli.
\end{itemize}

\begin{definizione}[leave-one-out]
Cross-validation con $k = n$ (numero di punti): si addestra su tutti tranne uno, che si usa per la validazione; si ripete
$n$ volte. Solo per dataset \textbf{molto piccoli} (es.\ 100 pazienti in ambito medico).
\end{definizione}

\subsection{Problema 1: distribuzione delle classi}
Classi molto sbilanciate (99\% vs 1\%) creano bias: il modello descrive bene solo la classe maggioritaria.
\begin{itemize}
  \item \textbf{A livello di dati (resampling)}:
  \begin{itemize}
    \item \textbf{Oversampling} della minoritaria: duplicare o generare campioni sintetici.
          \textbf{SMOTE} interpola tra campioni minoritari esistenti; \textbf{ADASYN} è una variante che genera più campioni
          per le istanze difficili.
    \item \textbf{Undersampling} della maggioritaria: rimozione casuale, oppure \textbf{cluster-based}: si fa clustering e si
          campiona da ogni cluster, così si mantiene tutto lo spettro della classe.
  \end{itemize}
  \item \textbf{A livello di algoritmo}:
  \begin{itemize}
    \item \textbf{Class weights}: pesi più alti nella loss per la classe minoritaria (es.\ 0,8 vs 0,2): un errore su di essa costa di più.
    \item \textbf{Focal loss}: loss modificata che si concentra sui campioni più difficili, con pesi di classe e fattore di focus.
  \end{itemize}
  \item \textbf{Metrica giusta}: con il 99\% di maggioritaria, predire sempre quella dà 99\% di accuracy. Usare metriche per classe o
        la media \textbf{macro} (stesso peso a ogni classe) invece della \en{micro} (metriche: cap.~\ref{cap:addestramento}).
\end{itemize}
\begin{attenzione}
Una versione della slide dice ``micro avg instead of macro'': è un refuso, corretto nella versione completa
(``macro avg instead of micro''). Con classi sbilanciate si usa la \textbf{macro}.
\end{attenzione}

\subsection{Problema 2: costo degli errori}
Segnalare un software come possibile malware (falso positivo) è meno grave di non vedere un malware (falso negativo).
Rimedi: metrica adatta alla classe più rilevante, class weights, loss pesata o focal loss.

\subsection{Problema 3: dimensione del dataset}
Dataset troppo piccolo per l'hold-out: cross-validation o leave-one-out; \textbf{data augmentation} (SMOTE, ADASYN);
oppure non affrontare il problema end-to-end: \textbf{modello pre-addestrato + fine-tuning}, o approccio
\textbf{semi-supervised} se ci sono dati non etichettati.

\subsection{Problema 4: data leakage}
\begin{definizione}[data leakage]
Informazioni del training che influenzano la validazione, rendendo il modello artificialmente bravo e non generale.
\end{definizione}
\begin{itemize}
  \item \textbf{Correlazione temporale}: se in training ho i valori a $t-1$ (0) e $t+1$ (10), predire $t$ (5) in validation è
        facile, ma nella realtà quando devo predire $t$ non ho $t+1$. Imparo a interpolare, non a prevedere il futuro.
  \item \textbf{Feature che creano bias}: se tutti i DoS del dataset sono sulla porta TCP 80, il modello impara ``porta 80 = DoS'',
        una scorciatoia (\en{shortcut}) che non vale nella realtà. Verifica togliendo la feature: se crolla, la usava come scorciatoia.
  \item \textbf{Non solo il tempo}: un'immagine divisa in 3 pezzi con $c_0$ e $c_2$ in training e $c_1$ in validation è troppo facile.
  \item \textbf{Identificatori}: ID o pseudo-ID (es.\ la matricola per predire il voto) non vanno mai usati in training.
\end{itemize}
\textbf{Split basato sul tempo}: il passato in training, il futuro in validation.
\begin{itemize}
  \item \textbf{Sliding window}: finestra di training di lunghezza fissa che scorre.
  \item \textbf{Expanding window}: la finestra di training si allarga includendo tutta la storia disponibile.
\end{itemize}

\subsection{Normalizzazione ed encoding: su quali dati?}
\label{sec:leak-norm}
\begin{esame}[domanda classica]
\textbf{Si può normalizzare usando tutto il dataset? No.} Il dataset contiene anche validation e test: usarli per calcolare
min/max o media/deviazione standard fa entrare nel modello informazioni sulla loro distribuzione $\Rightarrow$
\textbf{data leakage}. Nella realtà non si conosce la distribuzione dei dati futuri.

Procedura corretta: \textbf{split} $\to$ \textbf{fit} delle trasformazioni \textbf{solo sul training} $\to$ \textbf{transform}
(senza fit) di validation e test.
\end{esame}
\begin{esempio}
Min-max imparato sul training (valori 0--100): 0 $\to$ 0, 100 $\to$ 1. Nel validation c'è 150 $\Rightarrow$ diventa
\textbf{1,5}, ed è corretto così. Ricalcolare il massimo come 150 introdurrebbe leakage.
\end{esempio}
\begin{lstlisting}
from sklearn.preprocessing import StandardScaler
scaler = StandardScaler()
X_train = scaler.fit_transform(X_train)   # fit SOLO sul training
X_val   = scaler.transform(X_val)         # solo transform
X_test  = scaler.transform(X_test)        # solo transform
\end{lstlisting}

% =====================================================================
\section{Panoramica delle architetture}
\begin{itemize}
  \item \textbf{FFNN}: nate per dati numerici (righe e colonne), es.\ malware da feature predefinite. Si usano anche dopo altre
        reti come classificatore finale (es.\ CNN codifica l'immagine, FFNN classifica).
  \item \textbf{RNN}: dati sequenziali; una catena di FFNN che ricorda la storia (prossimo valore del traffico, prossima parola).
        Cap.~\ref{cap:rnn}.
  \item \textbf{CNN}: immagini; filtri convoluzionali che estraggono feature (malware come immagine + FFNN). Cap.~\ref{cap:cnn-gnn}.
  \item \textbf{GNN}: grafi; aggregano l'informazione dei vicini (malware, social network, frodi). Cap.~\ref{cap:cnn-gnn}.
\end{itemize}
In tutte, quando si parla di rappresentazione, si parla di \textbf{embedding}, spazio latente, bottleneck.

% =====================================================================
\section{Progettare una feed-forward neural network}
\label{sec:ffnn-progetto}
\subsection{Preprocessing e trasformazione}
\begin{itemize}
  \item Tutto deve essere numerico: vanno codificate le feature categoriche, la label, date e orari.
  \item \textbf{Attenzione}: alcune feature sono numeri ma \textbf{categoriche} (le porte).
  \item Feature numeriche: normalizzazione/standardizzazione migliorano le prestazioni.
\end{itemize}

\subsection{Architettura}
\begin{enumerate}
  \item \textbf{Input layer}: numero di nodi in base a preprocessing e trasformazione.
        \emph{Esempio}: 2 feature numeriche + 1 categorica con 10 valori in one-hot $\Rightarrow$ $2 + 10 = 12$ neuroni.
  \item \textbf{Hidden layer}: partire con pochi e aumentare se serve; \textbf{prima in profondità che in larghezza}
        (ma non 100 layer da 2 neuroni). Neuroni per layer: \textbf{potenze di 2} per efficienza.
        Attivazioni: ReLU, Leaky ReLU\ldots\ (sigmoid/tanh dipende dalla profondità).
  \item \textbf{Output layer}: in base al compito: regressione (1 neurone), binaria, multi-classe (\textbf{un neurone per classe}).
        Attivazione in base a problema e loss (sigmoid, softmax). In PyTorch con \code{CrossEntropyLoss} non si mette:
        è dentro la loss (sez.~\ref{sec:ffnn-classe}).
  \item \textbf{Inizializzazione dei pesi}: automatica; impostare i \textbf{seed} per la riproducibilità. Manuale solo per
        esperimenti mirati: rischio di pesi enormi in una zona e minuscoli in un'altra, usando di fatto mezza rete.
\end{enumerate}

\subsection{Ottimizzazione}
\begin{enumerate}[resume]
  \item Come dare i dati: tutti insieme o \textbf{mini-batch} (GD, SGD\ldots)?
  \item Quale \textbf{loss}? Se sbilanciato, con pesi? (Con gli autoencoder creeremo loss personalizzate, cap.~\ref{cap:anomaly}.)
  \item Quale \textbf{ottimizzatore}, quante epoch, quale learning rate, quale momentum? (Cap.~\ref{cap:addestramento}.)
  \item In caso di \textbf{overfitting}: class weights, dropout, layer di normalizzazione (cap.~\ref{cap:regolarizzazione}).
        Se è forte (rete enorme che memorizza), ridurre layer/neuroni o ripartire dal preprocessing.
\end{enumerate}

\begin{center}
\small
\renewcommand{\arraystretch}{1.2}
\begin{tabularx}{\textwidth}{>{\bfseries}l X}
\toprule
Iperparametro & Best practice (dalle slide)\\
\midrule
\# layer & partire con 3--5 hidden layer, aumentare se serve\\
\# neuroni per layer & potenze di 2 (64, 128, 256), evitare di esagerare\\
Attivazione & ReLU, LeakyReLU\ldots\ negli hidden; softmax/sigmoid in output\\
Inizializzazione & automatica in PyTorch (dipende dall'attivazione); manuale possibile\\
Batch size & 32--256, da tarare sperimentalmente\\
Loss & in base al compito\\
Ottimizzatore & Adam (default); SGD + momentum per la generalizzazione\\
Learning rate & partire da 0,001; usare un LR scheduler\\
Epoch \& early stopping & monitorare la validation loss; il numero di epoch cambia molto con ottimizzatore e LR\\
Regolarizzazione & dropout (0,2--0,5) + L2 (0,01)\\
\bottomrule
\end{tabularx}
\end{center}

% =====================================================================
\section{Caricare i dati in PyTorch}
\begin{definizione}[tensore]
Matrice multidimensionale di numeri di \textbf{un solo tipo}. Può essere un numero (rango 0), un vettore, una matrice o una
matrice multidimensionale. Si può specificare il \code{dtype} (es.\ le label in \code{torch.long}, le feature in
\code{torch.float32}).
\end{definizione}
\begin{center}
\small
\begin{tabularx}{\textwidth}{l X}
\toprule
\textbf{Metodo} & \textbf{Uso}\\
\midrule
\code{torch.tensor()} & conversione generale da liste, tuple, matrici\\
\code{torch.zeros()}, \code{torch.ones()}, \code{torch.arange()} & tensori specifici; gli zeri servono per le \textbf{maschere} (NLP)\\
\code{torch.rand()}, \code{torch.randn()}, \code{torch.randint()} & tensori casuali\\
\bottomrule
\end{tabularx}
\end{center}

\begin{itemize}
  \item \code{TensorDataset(X, y)}: unisce feature e label in un dataset. Usato da solo (senza mini-batch) rischia di non stare
        in memoria.
  \item \code{DataLoader(dataset, batch\_size=..., shuffle=...)}: crea i \textbf{mini-batch}.
  \item \code{shuffle=True}: mescola il dataset \textbf{all'inizio di ogni epoch}; i punti arrivano al modello in ordine diverso
        $\Rightarrow$ migliore generalizzazione, meno overfitting. Ogni punto resta visto una sola volta per epoch.
\end{itemize}
\begin{lstlisting}
from torch.utils.data import TensorDataset, DataLoader
X_train_t = torch.tensor(X_train, dtype=torch.float32)
y_train_t = torch.tensor(y_train, dtype=torch.long)     # indici di classe
train_ds = TensorDataset(X_train_t, y_train_t)
train_loader = DataLoader(train_ds, batch_size=64, shuffle=True)
val_loader   = DataLoader(val_ds,   batch_size=64, shuffle=False)
\end{lstlisting}

\begin{attenzione}[shuffle e sequenze]
Se i punti di una sequenza sono su righe diverse e vanno valutati come sequenza, \textbf{non} si può mescolare: si rompe la
serie. Si possono invece mescolare \textbf{sequenze diverse} tra loro (cap.~\ref{cap:rnn}).
\end{attenzione}
\begin{esame}[shuffle su validation e test?]
\textbf{No.} La stocasticità serve all'\textbf{SGD}, cioè solo al training. In validation e test i pesi \textbf{non vanno
aggiornati} (sarebbe un errore grave): il modello è fisso e l'ordine dei dati non cambia le predizioni. Il mescolamento
renderebbe solo più difficile leggere e confrontare le loss. (Domanda posta il 05/10, risposta data il \rif{2026-10-06}{06/10}.)
\end{esame}

% =====================================================================
\section{Definire la rete: una classe}
\label{sec:ffnn-classe}
In PyTorch una rete è una classe che eredita da \code{nn.Module} e ha (almeno) due metodi:
\begin{itemize}
  \item \code{\_\_init\_\_}: descrive l'\textbf{architettura} (layer, dropout, normalizzazione, attivazioni). Riceve parametri come
        il numero di neuroni di input e di output. \code{super().\_\_init\_\_()} dice che è una rete neurale.
  \item \code{forward}: descrive \textbf{come i dati attraversano} la rete.
\end{itemize}

\begin{lstlisting}
import torch
import torch.nn as nn
import torch.optim as optim

class SimpleNN(nn.Module):
    def __init__(self, input_size, num_classes):
        super().__init__()
        self.fc1 = nn.Linear(input_size, 32)    # hidden layer: 32 neuroni
        self.fc2 = nn.Linear(32, num_classes)   # input = output del layer precedente
    def forward(self, x):
        x = self.fc1(x)
        x = self.fc2(x)
        return x                                # logit
\end{lstlisting}
\begin{esame}[c'è non linearità?]
\textbf{No}: sono solo layer lineari, quindi la rete equivale a un solo layer lineare (la composizione di funzioni lineari
è lineare, cap.~\ref{cap:reti-neurali}). Serve un'attivazione tra i layer.
\end{esame}

\begin{lstlisting}
class NonLinearNN(nn.Module):
    def __init__(self, input_size, num_classes):
        super().__init__()
        self.fc1 = nn.Linear(input_size, 32)
        self.relu = nn.ReLU()
        self.fc2 = nn.Linear(32, num_classes)
    def forward(self, x):
        x = self.relu(self.fc1(x))   # prima il layer, poi l'attivazione
        return self.fc2(x)           # logit: niente softmax qui!
\end{lstlisting}
Layer diversi possono avere attivazioni diverse; come applicarle si decide nel \code{forward}.

\begin{definizione}[logit]
Output grezzi e non normalizzati dell'ultimo layer, in $[-\infty, +\infty]$. Con 2 punti e 3 classi l'output ha forma
$2 \times 3$. Per vedere le probabilità: \code{torch.softmax(logits, dim=1)}.
\end{definizione}

\begin{esame}[nell'ultimo layer serve sempre un'attivazione?]
In teoria \textbf{sì}, per passare dai logit alle probabilità (es.\ softmax nel multi-classe). In PyTorch però
\code{nn.CrossEntropyLoss} si aspetta \textbf{i logit non normalizzati} e applica internamente la softmax: \textbf{non}
va messa nel modello. Prima di usare una loss, \textbf{leggi la documentazione}.
\end{esame}

% =====================================================================
\section{Loss e ottimizzatore}
\begin{lstlisting}
criterion = nn.CrossEntropyLoss()     # multi-classe (vale anche per il binario)
                                      # weight=torch.tensor([...]) per classi sbilanciate
model = NonLinearNN(input_size=2, num_classes=3)   # pesi casuali: non e' ancora un modello!
optimizer = optim.AdamW(model.parameters(), lr=0.01)
                                      # altri parametri: betas, weight_decay (L2), ...
\end{lstlisting}
\begin{itemize}
  \item Il \textbf{criterion} è la loss. La label deve essere \code{torch.long} con gli indici di classe.
  \item All'ottimizzatore si passano i \textbf{parametri del modello} da ottimizzare e il learning rate.
  \item Le loss disponibili sono molte; più avanti nel corso si creeranno loss personalizzate. Teoria di loss e ottimizzatori
        nel cap.~\ref{cap:addestramento}.
\end{itemize}

% =====================================================================
\section{Il training loop (senza mini-batch, solo didattico)}
\begin{lstlisting}
num_epochs = 200
train_losses, val_losses = [], []
for epoch in range(num_epochs):
    # ---- TRAINING ----
    model.train()                       # 1. modalita' training
    optimizer.zero_grad()               # 2. azzera i gradienti precedenti
    outputs = model(X_train)            # 3. forward
    loss = criterion(outputs, y_train)  # 4. loss (softmax + cross-entropy interne)
    loss.backward()                     # 5. backpropagation: calcola i gradienti
    optimizer.step()                    # 6. aggiorna i pesi
    train_losses.append(loss.item())    # 7. salva la loss per la learning curve
    # ---- VALIDATION ----
    model.eval()                        # modalita' valutazione
    with torch.no_grad():               # niente calcolo dei gradienti (piu' veloce)
        val_loss = criterion(model(X_val), y_val)
    val_losses.append(val_loss.item())
\end{lstlisting}

\begin{attenzione}[\code{model.eval()} e \code{torch.no\_grad()} non sono la stessa cosa]
Il docente insiste: \textbf{\code{model.eval()} deve esserci} in validation e test.
\begin{itemize}
  \item \code{model.eval()} cambia il comportamento dei layer che si comportano diversamente in training: il
        \textbf{dropout} si spegne e la \textbf{batch norm} usa le statistiche memorizzate (cap.~\ref{cap:regolarizzazione}).
  \item \code{torch.no\_grad()} disattiva soltanto il calcolo dei gradienti (risparmio di tempo e memoria).
\end{itemize}
\end{attenzione}
\begin{integrazione}[cosa succede davvero se dimentichi \code{model.eval()}]
I pesi \textbf{non} vengono aggiornati finché non chiami \code{optimizer.step()}; però il dropout resta attivo (predizioni
casuali) e la batch norm normalizza con le statistiche del batch di validation, \textbf{aggiornando} le sue
\en{running statistics} con dati di validation: una forma di leakage. In ogni caso è un errore.
\end{integrazione}

\textbf{Leggere la learning curve}: da 1,1 a 0,4 in 50 epoch, con training loss un po' più alta della validation (possibile se
la validation è casualmente più facile). Il problema è che \textbf{non c'è plateau}: il modello stava ancora imparando, il
training è stato fermato troppo presto. Non fermarti a 10 epoch perché ``funziona''. E la loss dice solo se il modello
converge: per giudicarlo servono le \textbf{metriche} (accuracy, precision, recall\ldots, cap.~\ref{cap:addestramento}).

% =====================================================================
\section{Il training loop con mini-batch}
\begin{lstlisting}
for epoch in range(num_epochs):
    model.train()
    running = 0.0
    for xb, yb in train_loader:              # itera sui mini-batch
        xb, yb = xb.to(device), yb.to(device)
        optimizer.zero_grad()                # zero_grad, backward e step
        loss = criterion(model(xb), yb)      # sono DENTRO il ciclo sui batch:
        loss.backward()                      # i pesi si aggiornano a ogni
        optimizer.step()                     # mini-batch
        running += loss.item() * xb.size(0)  # pesa la loss per la dimensione del batch
    train_losses.append(running / len(train_loader.dataset))

    model.eval()                             # una volta, prima del ciclo
    running = 0.0
    with torch.no_grad():                    # una volta
        for xb, yb in val_loader:            # mini-batch anche in validation (memoria)
            xb, yb = xb.to(device), yb.to(device)
            running += criterion(model(xb), yb).item() * xb.size(0)
    val_losses.append(running / len(val_loader.dataset))
\end{lstlisting}
\begin{itemize}
  \item Se \code{zero\_grad/backward/step} fossero fuori dal ciclo sui batch, non si sfrutterebbero gli aggiornamenti frequenti, che
        sono il vero vantaggio dei mini-batch (oltre al risparmio di memoria).
  \item \textbf{Perché moltiplicare per la dimensione del batch?} La loss di un batch è una media sul batch. Con 90 punti e batch da 64
        si hanno batch da 64 e da 26: senza pesatura conterebbero uguale. Si calcola la \textbf{media pesata} e si divide per la
        dimensione del dataset.
\end{itemize}

% =====================================================================
\section{Early stopping e regolarizzazione nel codice}
\label{sec:ffnn-regol-codice}
Nel notebook del \rif{2026-10-06}{06/10} la stessa FFNN viene poi arricchita con \textbf{early stopping} (\code{patience},
\code{min\_delta}, salvataggio di \code{best\_state}), \textbf{dropout} e \textbf{batch norm} definiti come layer in
\code{\_\_init\_\_}, e \textbf{weight decay} come parametro dell'ottimizzatore (\code{AdamW(..., weight\_decay=1e-2)}).
Teoria, codice completo (\code{RegularizedNN}, ciclo di early stopping) e lettura delle curve sono nel
cap.~\ref{cap:regolarizzazione}.

% =====================================================================
\section{Inizializzazione dei pesi}
Si può definire una funzione di inizializzazione: \textbf{automatica} (PyTorch sceglie il metodo in base all'attivazione) o
\textbf{manuale}.
\begin{lstlisting}
def init_weights(m):
    if isinstance(m, nn.Linear):
        nn.init.normal_(m.weight, mean=0.0, std=1e-5)   # pesi minuscoli: esperimento!
        nn.init.zeros_(m.bias)
model.apply(init_weights)
\end{lstlisting}
Con l'inizializzazione automatica i pesi sono dell'ordine di $\pm 0{,}3$--$1$; con quella manuale dell'esempio sono
$10^{-4}$--$10^{-6}$: vanno fatti crescere moltissimo per servire $\Rightarrow$ \textbf{vanishing gradient}.
\textbf{Morale}: si può fare, ma è più probabile avere problemi che benefici. Usala solo per esperimenti mirati.

% =====================================================================
\section{GPU e riproducibilità}
\begin{lstlisting}
device = torch.device('cuda' if torch.cuda.is_available() else 'cpu')
model = model.to(device)                 # sposta il modello
xb, yb = xb.to(device), yb.to(device)    # sposta i dati (batch per batch)
preds = model(xb).argmax(dim=1).cpu()    # riporta su CPU per numpy/sklearn

import random, numpy as np
random.seed(42); np.random.seed(42)
torch.manual_seed(42); torch.cuda.manual_seed_all(42)
\end{lstlisting}
\begin{itemize}
  \item \code{cuda} = GPU disponibile. Le GPU accelerano molto, ma spesso non ci sono o sono a tempo limitato: usale solo se servono
        (un dataset piccolo non ne ha bisogno; nel Lab 2 potrebbe servire).
  \item Modello \textbf{e} dati devono stare sullo stesso device. Riporta i tensori su CPU (\code{.cpu()}) quando servono per elaborazioni
        esterne. PyTorch libera la memoria dei tensori non più referenziati; evita trasferimenti CPU/GPU inutili (costano).
  \item \textbf{Riproducibilità}: imposta i seed di tutto ciò che è casuale (inizializzazione, numeri casuali, split).
        Su hardware diverso (anche server diversi su Colab) i risultati possono comunque differire leggermente.
        Se le prestazioni dipendono molto dal seed, il modello è \textbf{instabile}: va migliorata l'architettura.
\end{itemize}

% =====================================================================
\section{TensorFlow Playground: cosa si impara}
\url{https://playground.tensorflow.org/} permette di giocare con una piccola rete su dataset 2D.
\begin{itemize}
  \item Con buone \textbf{feature ingegnerizzate} basta anche un solo layer (con abbastanza neuroni per creare regioni non lineari).
  \item Learning rate \textbf{10}: training rumorosissimo (salti enormi). Learning rate \textbf{minuscolo}: forse prima o poi converge.
  \item Attivazione \textbf{lineare}: il problema non lineare non si risolve mai.
  \item Rete profonda con \textbf{sigmoid}: la loss resta \textbf{congelata}.
\end{itemize}
\textbf{Morale}: farla funzionare anche in un caso semplice non è banale: basta una scelta sbagliata (attivazione, learning rate)
per non imparare nulla.

% =====================================================================
\section{Il notebook d'esempio: una pipeline completa}
Dataset sintetico con 2 feature ($x, y$), 3 classi (A, B, C), sporcato apposta con outlier, valori mancanti e righe duplicate
(\rif{2026-10-06}{06/10}).
\begin{enumerate}
  \item \textbf{Import}: \code{torch}, \code{torch.nn}, \code{torch.optim}; \code{matplotlib}; \code{numpy}, \code{pandas};
        \code{sklearn} per scaling, split e metriche.
  \item \textbf{Duplicati}: si rimuovono di solito (creano risultati artificiali attorno a un punto). Ma se sono della classe
        minoritaria e reali, potrebbero essere utili: valuta caso per caso.
  \item \textbf{Valori mancanti}: prima controlla quante righe perdi. Tieni traccia delle statistiche (righe prima e dopo ogni passo):
        perdere l'1\% può essere irrilevante o gravissimo se è la classe minoritaria.
  \item \textbf{Label}: \code{LabelEncoder} (A, B, C $\to$ 0, 1, 2). L'ordine artificiale non conta perché si usa la cross-entropy, non
        una distanza come l'MSE.
  \item \textbf{Visualizzazione} (2 dimensioni, si può): rosso attorno al giallo $\Rightarrow$ problema non lineare.
        Per la sola visualizzazione si può usare tutto il dataset; per imparare si usa \textbf{solo il training}.
  \item \textbf{Split}: \code{train\_test\_split(..., test\_size=0.4, random\_state=...)}, poi di nuovo a metà sul 40\% $\Rightarrow$
        60\% train, 20\% validation, 20\% test.
  \item \textbf{Scaling}: \code{StandardScaler} (ci sono outlier che non si vogliono rimuovere). \code{fit\_transform} sul training,
        solo \code{transform} su validation e test. Fare \code{fit} su validation/test introdurrebbe leakage e darebbe al modello
        distribuzioni diverse tra training e validation.
  \item \textbf{Modello lineare} (un solo \code{nn.Linear}): pesi $3\times 2$ + 3 bias. La loss scende solo da 1,45 a 1,25:
        prestazioni pessime, come previsto.
  \item \textbf{Modello non lineare}: 3 hidden layer da 128 neuroni + ReLU: numeri ragionevoli.
  \item \textbf{Mini-batch} (\code{shuffle=True} solo per il training): loss da 0,55 a 0,25, più rumorosa per la stocasticità.
  \item \textbf{Early stopping, dropout, batch norm + weight decay}: effetti sulle curve discussi nel cap.~\ref{cap:regolarizzazione}
        (in sintesi: early stopping si ferma molto prima; la curva con BN + weight decay è molto rumorosa, ``se vedo una validation
        loss così nel vostro report, penso che non abbiate valutato la soluzione'').
\end{enumerate}
\begin{attenzione}[non copiare il notebook]
Il codice d'esempio mostra \emph{come usare la libreria}, non come risolvere un problema reale. Ogni dataset è diverso e nella
realtà non avrete due variabili con classi così ben separate.
\end{attenzione}

\subsection{Quali iperparametri ottimizzare? (domanda in aula)}
Non si può sapere in anticipo quali valori funzionano: dipendono da dataset, split e architettura (4 o 5 layer, numero di neuroni).
Nel Lab 1 vi si chiede di giocare con diversi parametri (anche il \code{weight\_decay}) per sviluppare sensibilità
(cap.~\ref{cap:lab1}). In generale:
\begin{itemize}
  \item si lavora soprattutto sul \textbf{learning rate}, eventualmente con uno \textbf{scheduler};
  \item \code{betas}/momentum si possono lasciare di default;
  \item il weight decay si tocca se c'è overfitting.
\end{itemize}
Lab: PyTorch su Colab (CPU e GPU, collaborazione); va bene anche un server Jupyter proprio o Anaconda. Il Lab 2 userà lo stesso
dataset come tabella, serie temporale e grafo (cap.~\ref{cap:lab2}); poi anomaly detection e NLP.

% =====================================================================
\section{Errori tipici nel codice}
\begin{attenzione}[errori di codice da riconoscere -- raccolta dei redattori]
Lo scritto include domande sul codice (\rif{2026-09-21}{21/09}). Raccolta dei redattori, a partire da lezioni, slide e notebook, degli errori da saper riconoscere nella pipeline FFNN:
\begin{enumerate}
  \item \code{scaler.fit} (o \code{fit\_transform}) su tutto il dataset o su validation/test $\Rightarrow$ data leakage
        (sez.~\ref{sec:leak-norm});
  \item \code{shuffle=True} sul \code{DataLoader} di validation/test (inutile), o su righe che formano una sequenza (rompe la serie);
  \item rete con soli \code{nn.Linear} senza attivazioni $\Rightarrow$ equivale a un modello lineare;
  \item softmax nel \code{forward} prima di \code{nn.CrossEntropyLoss} (la loss si aspetta i logit);
  \item label non in \code{torch.long} con la cross-entropy;
  \item \code{model.eval()} mancante in validation/test (dropout e batch norm restano in modalità training); \code{torch.no\_grad()}
        non lo sostituisce;
  \item \code{optimizer.zero\_grad()}, \code{backward()} e \code{step()} fuori dal ciclo sui mini-batch;
  \item \code{optimizer.step()} in validation (i pesi non vanno mai aggiornati fuori dal training);
  \item loss media per epoch non pesata con \code{xb.size(0)} (l'ultimo batch, più piccolo, conta come gli altri);
  \item early stopping senza salvare e ricaricare \code{best\_state} (si restituisce il modello dell'ultima epoch);
  \item modello e dati su device diversi.
\end{enumerate}
\end{attenzione}

\begin{sintesi}
\begin{itemize}
  \item Il data handling (split, pre-processing, trasformazione) è la parte più lunga: \textbf{prima lo split}, poi fit delle
        trasformazioni \textbf{solo sul training}.
  \item Categoriche: one-hot (poche categorie), label encoding (ordinali e label), embedding (alta cardinalità, cap.~\ref{cap:embedding}).
        Numeriche: min-max (sensibile agli outlier) o z-score.
  \item Split: hold-out (dataset grandi), k-fold (stima affidabile, poi ri-addestramento su tutti i fold), leave-one-out (dataset minuscoli);
        attenzione a sbilanciamento (resampling, class weights, macro avg), costo degli errori, dimensione, leakage (tempo, scorciatoie, ID).
  \item FFNN: input = n.\ di feature dopo l'encoding; 3--5 hidden layer, potenze di 2, ReLU; output in base al task, \textbf{senza softmax}
        con \code{CrossEntropyLoss}; Adam, LR 0,001.
  \item PyTorch: \code{nn.Module} con \code{\_\_init\_\_} e \code{forward}; \code{TensorDataset} + \code{DataLoader} (shuffle solo in training);
        loop \code{train()} $\to$ \code{zero\_grad} $\to$ forward $\to$ loss $\to$ \code{backward} $\to$ \code{step}; validation con
        \code{eval()} \textbf{e} \code{no\_grad()}.
  \item Leggere la learning curve: senza plateau il training è stato fermato troppo presto; la loss non basta, servono le metriche.
\end{itemize}
\end{sintesi}

\chapter{Recurrent neural network e LSTM}
\label{cap:rnn}
\fonti{\rif{2026-10-12}{12/10}, \rif{2026-10-14}{14/10}}{3-deep-learning-methodologies, pp.~34--76; notebook RNN (1-RNN-Pollution-Multivariate, 2-RNN-LSTM-FordA)}

Le RNN estendono la FFNN (cap.~\ref{cap:ffnn}) ai dati \textbf{sequenziali}. Il capitolo va dalla teoria (cella, BPTT,
gradiente, LSTM) alla progettazione e all'implementazione in PyTorch, con i due notebook mostrati a lezione.

% =====================================================================
\section{Perché servono le RNN}
Dato un solo punto (posizione, raggio, colore\ldots) non si può predire la sua \textbf{traiettoria}: manca la storia.
Una FFNN riceve un \textbf{vettore di dimensione fissa} e restituisce un vettore di dimensione fissa; i punti sono trattati
in modo indipendente e i pesi si aggiornano indipendentemente dall'ordine. Se invece do la \textbf{sequenza} delle posizioni
precedenti a una rete che gestisce sequenze, posso indovinare la direzione.

\begin{definizione}[recurrent neural network]
Classe di reti che gestisce input di \textbf{lunghezza variabile} e dati sequenziali $x(t)$. A differenza di una FFNN
mantiene uno \textbf{stato} che \textbf{evolve nel tempo} e descrive l'andamento della serie.
\end{definizione}
Esempio: ``my brother plays \ldots'': sapendo cosa c'è prima, mi aspetto un gioco o uno strumento. In modo
semplificatissimo, è l'idea di predire la parola successiva (alla base di GPT).

\textbf{Applicazioni}: traduzione automatica (una parola dipende dalle precedenti e non sempre corrisponde a una sola parola
nell'altra lingua), previsione di serie temporali (se il valore reale si discosta molto dalla previsione: anomalia, forse un
attacco), classificazione di serie, riconoscimento vocale, \en{part-of-speech tagging} (capire che ``volleyball'' è un nome).

% =====================================================================
\section{FFNN vs RNN in breve}
\begin{itemize}
  \item Una FFNN riceve $x$ e restituisce $y$; il punto successivo viene elaborato indipendentemente.
  \item Una RNN aggiunge un \textbf{self-loop}: riceve $x(t)$ \textbf{e} lo stato nascosto precedente $h(t-1)$, e restituisce
        l'output $y(t)$ e il nuovo stato $h(t)$.
  \item Servono \textbf{nuovi pesi}: oltre a $w$ (importanza delle feature) c'è $w'$ (importanza del passato).
        La complessità aumenta.
  \item Anche le RNN si organizzano in più layer (più layer = rappresentazione più astratta); in pratica 1--2 layer.
\end{itemize}
\begin{definizione}[hidden state]
Rappresentazione interna della serie fino a $t-1$. L'output dipende da $x_t$ \textbf{e} da $h_{t-1}$; il nuovo $h_t$
è aggiornato con l'informazione presente.
\end{definizione}

\begin{center}
\begin{tikzpicture}[every node/.style={font=\small}, cell/.style={draw=polblu, fill=polblu!8, rounded corners=2pt,
  minimum width=1.1cm, minimum height=0.8cm}]
  % forma compatta
  \node[cell] (c) at (0,0) {RNN};
  \node (x) at (0,-1.3) {$x_t$}; \node (y) at (0,1.3) {$\hat y_t$};
  \draw[-{Stealth}] (x) -- (c); \draw[-{Stealth}] (c) -- (y);
  \draw[-{Stealth}] (c.east) to[out=0,in=-90,looseness=1] ++(0.7,0.6) to[out=90,in=0] ++(-0.7,0.5)
        node[above right, font=\scriptsize] {$h_{t-1}$} to[out=180,in=90] ++(-0.7,-0.5) to[out=-90,in=180] (c.west);
  \node at (1.8,0) {$=$};
  % srotolata
  \foreach \i/\t in {0/0,1/1,2/2,3/t} {
    \node[cell] (u\i) at (3+1.8*\i,0) {RNN};
    \node (ux\i) at (3+1.8*\i,-1.3) {$x_{\t}$};
    \node (uy\i) at (3+1.8*\i,1.3) {$\hat y_{\t}$};
    \draw[-{Stealth}] (ux\i) -- (u\i); \draw[-{Stealth}] (u\i) -- (uy\i);
  }
  \draw[-{Stealth}] (u0) -- node[above, font=\scriptsize] {$h_0$} (u1);
  \draw[-{Stealth}] (u1) -- node[above, font=\scriptsize] {$h_1$} (u2);
  \draw[-{Stealth}, dashed] (u2) -- node[above, font=\scriptsize] {$\cdots$} (u3);
\end{tikzpicture}
\end{center}
A sinistra la forma compatta con il self-loop, a destra la versione \textbf{srotolata} (\en{unrolled}): stessa cella, stessi
pesi, ripetuta nel tempo.

% =====================================================================
\section{Varianti input/output}
\begin{center}
\small
\renewcommand{\arraystretch}{1.25}
\begin{tabularx}{\textwidth}{>{\bfseries}l X X}
\toprule
Tipo & Descrizione & Esempi\\
\midrule
One-to-one & un timestamp alla volta, senza passato: è una rete standard & classificazione o regressione su un singolo campione\\
Many-to-one & una sequenza in input, \textbf{un} output; lo stato finale codifica l'intera sequenza & sentiment analysis (``terrifically'' bad/good: serve tutta la frase), sequenza di log per rilevare un attacco in corso\\
One-to-many & un input espanso su più passi di output & image captioning (``un uomo che gioca a baseball''), piano di risposta a partire da un alert\\
Many-to-many & sequenza $\to$ sequenza, lunghezze anche diverse & traduzione, classificazione di ogni frame del traffico in tempo reale\\
\bottomrule
\end{tabularx}
\end{center}
Nel lab ci si concentra sul \textbf{many-to-one}: una serie in input, una label (es.\ malware o no) in output.

% =====================================================================
\section{Architetture: mono e bidirezionale}
\subsection{RNN mono-direzionale}
Elabora i dati in sequenza, dal passato al futuro, aggiornando lo stato con l'input corrente e lo stato precedente.
\begin{itemize}
  \item \textbf{Pro}: adatta alla \textbf{predizione in tempo reale}; relativamente semplice da addestrare (un solo layer, pochi pesi).
  \item \textbf{Contro}: fatica con le \textbf{dipendenze a lungo termine} (dimentica l'inizio di una serie lunga) per
        vanishing/exploding gradient.
  \item \textbf{Training}: \en{Backpropagation Through Time} (BPTT), variante della backprop che tara sia i pesi dell'input sia quelli
        dello stato.
  \item \textbf{Applicazioni}: classificazione supervisionata del traffico (benigno/malevolo), anomaly detection non supervisionata.
\end{itemize}

\subsection{RNN bidirezionale (Bi-RNN)}
Due RNN \textbf{indipendenti} (con \textbf{pesi diversi}, appresi insieme) elaborano la stessa sequenza in direzioni opposte
(avanti e indietro), mantenendo due stati nascosti; gli output vengono poi combinati.
\begin{itemize}
  \item \textbf{Pro}: cattura contesto passato \textbf{e futuro} $\Rightarrow$ rappresentazione migliore (traduzione di una frase completa).
  \item \textbf{Contro}: costo computazionale più alto; \textbf{non adatta al tempo reale} (il futuro non c'è).
  \item Si usa quando si ha \textbf{la sequenza completa}: un'email, quando arriva, è completa $\Rightarrow$ classificazione di
        phishing; riconoscimento di pattern nei log per insider threat.
  \item BPTT applicata separatamente alle due direzioni.
\end{itemize}

% =====================================================================
\section{Struttura interna di una cella RNN}
\[
h_t = F_1\big(W_{hh}h_{t-1} + W_{xh}x_t + b_{hh}\big)
\]
\begin{itemize}
  \item $h_{t-1}$: stato precedente (la ``traiettoria'' fino a $t-1$); $x_t$: input corrente;
  \item $W_{xh}$, $W_{hh}$: matrici dei pesi di input e di stato (ognuna con i propri pesi); $b_{hh}$: bias;
  \item $F_1$: attivazione, tipicamente \textbf{tanh} (non ReLU) per limitare i problemi di gradiente;
  \item per ora l'output della cella $y_t$ coincide con $h_t$.
\end{itemize}

\subsection{Dalla cella alla predizione: la classification head}
$y_t$ (o $h_t$) ha tanti valori quanti la \textbf{hidden size}. Per la predizione finale si aggiunge almeno un layer lineare
(eventualmente dropout/normalizzazione e un'attivazione):
\[
\hat y_t = F_2\big(W_{yh}h_t + b_{yh}\big)
\]
\begin{itemize}
  \item \textbf{Many-to-many}: una head per ogni $x_t$.
  \item \textbf{Many-to-one}: solo l'\textbf{ultima} cella ha la head, perché il suo output ``conosce'' l'intera serie.
\end{itemize}
\begin{esame}[la RNN fa rappresentazione]
La parte ricorrente fa \textbf{feature extraction}: lo stato nascosto alla fine della serie comprime l'andamento della serie in,
ad esempio, 8 numeri: è un \textbf{embedding} della serie. La head (da un layer lineare a una FFNN complessa) usa quella
rappresentazione per classificare. Le due parti si addestrano \textbf{insieme} (ottimizzazione congiunta), non una dopo l'altra.
\end{esame}

\subsection{Dalla FFNN alla RNN, passo per passo}
\begin{enumerate}
  \item Una FFNN con un layer e più output, applicata a un singolo istante $x_t$.
  \item La stessa rete applicata a ogni istante: ogni $\hat y_t$ dipende solo dal suo $x_t$. L'errore è calcolato separatamente per ogni
        rete; $x_0$ e $x_1$ non influenzano $\hat y_2$. \textbf{Nessun concetto di sequenza.}
  \item Si collegano le reti con un link: lo stato $h_t$ viene propagato al passo successivo. Ora
        \[\hat y_t = f(x_t, h_{t-1})\]
        dipende da input presente e memoria passata.
  \item Disegno compatto: la cella con l'arco che torna su se stessa; la catena è la versione \textbf{srotolata} (\en{unrolled}).
\end{enumerate}
\begin{definizione}[relazione di ricorrenza]
\[ h_t = f_W(x_t, h_{t-1}) \]
\textbf{La stessa funzione e lo stesso insieme di pesi} si usano a ogni passo: non sono reti diverse una dopo l'altra, è la stessa
rete ripetuta nel tempo.
\end{definizione}

\subsection{Grafo computazionale nel tempo}
\begin{itemize}
  \item $W_{xh}$ mappa l'input $x_t$ nello stato $h_t$;
  \item $W_{hh}$ definisce come $h_{t-1}$ influenza $h_t$;
  \item $W_{hy}$ mappa lo stato $h_t$ nell'output $y_t$.
\end{itemize}
Tutte queste matrici sono \textbf{condivise tra i passi temporali} e vengono aggiornate con gradienti \textbf{accumulati lungo la
sequenza}.

% =====================================================================
\section{Backpropagation Through Time (BPTT)}
Nella backprop classica (cap.~\ref{cap:reti-neurali}): forward, loss, derivata rispetto a ogni parametro, aggiornamento. Con il tempo:
\begin{enumerate}
  \item si \textbf{sommano le loss} di tutti i passi per ottenere la loss totale della sequenza (in many-to-one c'è una sola predizione,
        ma i pesi vanno comunque ottimizzati lungo tutto il tempo);
  \item l'errore si backpropaga \textbf{sia attraverso gli output sia attraverso gli stati}, dalla fine all'inizio della sequenza;
  \item si \textbf{accumulano i gradienti} per aggiornare i pesi condivisi.
\end{enumerate}
\begin{attenzione}
Per tarare i pesi non basta considerare l'ultimo istante: serve la moltiplicazione con tutti gli stati e gli input precedenti,
indietro nel tempo, anche molto indietro. Per questo è molto più complesso.
\end{attenzione}

\subsection{Il flusso del gradiente}
Il gradiente rispetto a $h_0$ contiene \textbf{molti fattori di $W_{hh}$} e calcoli di gradiente ripetuti. I passi temporali non
sono layer indipendenti (è la stessa rete), ma per il gradiente si comportano come una rete profondissima.
\begin{itemize}
  \item Molti valori $>1$: \textbf{exploding gradient} $\Rightarrow$ i pesi cambiano drasticamente, come con un learning rate enorme;
        apprendimento inaffidabile. Rimedio: \textbf{gradient clipping} (non lasciare crescere il gradiente oltre una soglia).
  \item Molti valori $<1$: \textbf{vanishing gradient} $\Rightarrow$ il passato lontano non conta, la rete non lo ricorda.
        Rimedi: (1) funzione di attivazione, (2) inizializzazione dei pesi, (3) architettura (troppi layer ricorrenti peggiorano).
  \item Nelle RNN si trova spesso \textbf{tanh} invece di ReLU: essendo limitata evita l'esplosione, ma restando $<1$ può far svanire il gradiente.
\end{itemize}
\begin{esame}[exploding vs vanishing nelle RNN]
Il prodotto ripetuto per $W_{hh}$ lungo la sequenza è la causa di entrambi. Exploding $\to$ gradient clipping (e tanh);
vanishing $\to$ attivazione, inizializzazione, architettura, e soprattutto celle con gate (LSTM).
\end{esame}

% =====================================================================
\section{Il problema delle dipendenze a lungo termine}
\begin{esempio}
``The clouds are in the \ldots'': la parola mancante (\emph{sky}) dipende da poche parole vicine. Serve ricorrenza, ma il passato
utile è recente.

``Sono cresciuto in Francia, poi mi sono trasferito in Italia, ho fatto le elementari, giocato a basket, mi sono trasferito a Roma,
ho passato un anno in Australia\ldots'' e dopo un'ora e mezza: ``e parlo fluentemente \ldots''. L'informazione utile (Francia) è
troppo lontana e viene dimenticata. Lo stesso vale per un libro di 1000 pagine in cui conta la prima.
\end{esempio}
\textbf{Mitigazioni}: un'attivazione adatta (ReLU? ovunque? no, altrimenti esplode), una buona inizializzazione, e la soluzione
robusta: una cella più ``intelligente'', la \textbf{gated cell} (LSTM, GRU\ldots; il corso si concentra sull'LSTM). I
\textbf{transformer} (la T di GPT) saranno trattati nella parte NLP.

% =====================================================================
\section{LSTM (Long Short-Term Memory)}
Proposta nel 1997 (Hochreiter e Schmidhuber) per mitigare il vanishing (e l'exploding) gradient.
\begin{itemize}
  \item Mantiene \textbf{due stati}: il \textbf{cell state} $C_t$ (memoria \textbf{a lungo termine}) e l'\textbf{hidden state} $h_t$
        (memoria \textbf{a breve termine}).
  \item Usa dei \textbf{gate} per controllare il flusso: cosa tenere del presente, come aggiornare selettivamente il cell state,
        cosa mandare in output.
  \item Si addestra con BPTT, ma con un flusso di gradiente \textbf{parzialmente ininterrotto} attraverso il cell state.
  \item \textbf{Pro}: gestisce dipendenze a lungo raggio. \textbf{Contro}: molto più costosa di una RNN semplice (più componenti, più
        pesi, più non linearità).
\end{itemize}

\subsection{I quattro passaggi}
\begin{enumerate}
  \item \textbf{Forget gate}: quanto del passato tenere. Prende $h_{t-1}$ e $x_t$, calcola $f_t\in[0,1]$ con una \textbf{sigmoid} e lo
        moltiplica per $C_{t-1}$. $f_t\approx1$: ricorda; $f_t\approx 0$: dimentica (moltiplica per 0 la memoria lunga).
  \item \textbf{Input gate}: quanta nuova informazione aggiungere. $i_t\in[0,1]$ con una sigmoid. $i_t\approx 0$: il passo corrente non
        conta (es.\ un articolo, ``the'', non dice nulla della storia).
  \item \textbf{Update (candidate)}: crea una memoria candidata $\tilde C_t$ con \textbf{tanh}; $i_t$ decide quanta finisce nel cell state.
  \item \textbf{Output gate}: $o_t$ (sigmoid) decide quanto del cell state (passato per tanh) diventa $h_t$, usato per predire $y_t$ o
        passato al passo successivo.
\end{enumerate}
\begin{integrazione}[equazioni standard dell'LSTM]
\begin{align*}
f_t &= \sigma(W_f[h_{t-1}, x_t] + b_f) & i_t &= \sigma(W_i[h_{t-1}, x_t] + b_i)\\
\tilde C_t &= \tanh(W_C[h_{t-1}, x_t] + b_C) & C_t &= f_t \odot C_{t-1} + i_t \odot \tilde C_t\\
o_t &= \sigma(W_o[h_{t-1}, x_t] + b_o) & h_t &= o_t \odot \tanh(C_t)
\end{align*}
$\odot$ = prodotto elemento per elemento. Tre delle attivazioni sono sigmoid: qui servono funzioni \textbf{limitate} perché fanno da
``rubinetti'' tra 0 e 1.
\end{integrazione}
\textbf{Quanto passato ricorda un'LSTM?} Non si sa in anticipo: dipende dal training, dai dati (se il lungo termine serve, la rete lo
terrà) e dalla dimensione di hidden e cell state.

% =====================================================================
\section{Preparare i dati per una RNN}
\label{sec:rnn-dati}
I dati devono essere numeri \textbf{e} sequenze, rappresentati come \textbf{tensore 3D}
\code{(batch, seq\_len, features)}.
\begin{esempio}
\code{[[[1,2,3],[4,5,6]], [[7,8,9],[10,11,12]]]}
\begin{itemize}
  \item ogni vettore (\code{[1,2,3]}) è un istante $x(t)$ e ogni elemento una feature;
  \item ogni matrice (\code{[[1,2,3],[4,5,6]]}) è una serie: 2 istanti descritti da 3 feature;
  \item l'insieme delle matrici è il batch di serie.
\end{itemize}
\end{esempio}
\begin{attenzione}[serie mono-dimensionali]
\code{[1,2,3]} per PyTorch è \textbf{un istante con 3 feature}. Una serie di 3 istanti con una sola feature si scrive
\code{[[1],[2],[3]]}: ogni elemento in un vettore (\code{x.unsqueeze(-1)}, sez.~\ref{sec:rnn-varianti}).
\end{attenzione}

\subsection{Da dati tabellari a sequenze}
Dataset di inquinamento con una riga per ora (35.040 righe, 7 feature + label): ogni riga è un istante, non c'è ancora nessuna serie.
Una funzione \code{create\_sequences} costruisce le serie:
\begin{itemize}
  \item si sceglie la lunghezza (\code{seq\_length = 7}) e si prendono 7 righe consecutive;
  \item la label è quella dell'elemento \textbf{successivo} alla finestra (si predice l'inquinamento all'ora dopo);
  \item si ottengono 35.033 matrici, ognuna con 7 istanti descritti da 7 feature: il tensore 3D per il training.
\end{itemize}
\begin{lstlisting}
def create_sequences(features, labels, seq_length=7):
    X, y = [], []
    for i in range(len(features) - seq_length):
        X.append(features[i:i + seq_length])   # finestra di 7 istanti
        y.append(labels[i + seq_length])       # valore successivo
    return np.array(X), np.array(y)            # X: (N, 7, n_features)
\end{lstlisting}

\subsection{Serie di lunghezza diversa}
In una FFNN un punto con 4 feature \textbf{non} può entrare in un input layer da 5: il quinto neurone non riceverebbe nulla.
Nelle RNN:
\begin{itemize}
  \item il \textbf{numero di feature} per istante deve essere \textbf{sempre lo stesso};
  \item il \textbf{numero di passi temporali può variare} tra serie: un grande vantaggio.
\end{itemize}
Dentro lo stesso \textbf{mini-batch} però le serie devono avere la stessa lunghezza: si usa il \textbf{padding} delle più corte
oppure si \textbf{tronca} tutto a una lunghezza fissa. In generale non serve paddare tutto il dataset: \textbf{solo dentro il batch}
(sez.~\ref{sec:rnn-padding}).

% =====================================================================
\section{Progettare una RNN}
\begin{enumerate}
  \item \textbf{Tipo}: many-to-one, one-to-many, many-to-many? Quanto sono lunghe le sequenze (5 o 5000)? Serve memoria lunga (LSTM)?
        Ha senso leggere anche all'indietro (bidirezionale) o no (mono)?
  \item \textbf{Input}: numero di nodi = numero di \textbf{feature} (non di passi temporali).
  \item \textbf{Hidden size}: quanti neuroni per rappresentare hidden state (e cell state). Più neuroni = più informazione memorizzabile,
        ma più dati per addestrare e più complessità.
  \item \textbf{Hidden layer}: quanti? Più layer = rappresentazione più astratta ma training più difficile. Attivazione tipica: tanh.
  \item \textbf{Padding}: serve? In generale no, dentro il batch sì. Il padding non va considerato come parte reale della serie.
  \item \textbf{Cosa mettere prima dell'output}: un solo layer lineare o una FFNN?
  \item \textbf{Output}: numero di classi, attivazione in base a problema e loss.
  \item \textbf{Ottimizzazione}: mini-batch (shuffle? solo se le serie sono indipendenti), loss (pesi se sbilanciato), ottimizzatore,
        epoch, learning rate; in caso di overfitting, come per le FFNN (cap.~\ref{cap:regolarizzazione}); se forte, ripartire dai dati.
\end{enumerate}

\begin{center}
\small
\renewcommand{\arraystretch}{1.2}
\begin{tabularx}{\textwidth}{>{\bfseries}l X}
\toprule
Iperparametro & Best practice (dalle slide)\\
\midrule
Architettura & mono o bidirezionale, RNN o LSTM in base a cosa si impara dall'ordine all'indietro e a quanto deve essere lunga la memoria\\
\# layer & 1--2 layer ricorrenti, aumentare se serve\\
Hidden size & potenze di 2 (64, 128, 256), non esagerare\\
Attivazione & tanh negli hidden; softmax/sigmoid/lineare in output\\
Inizializzazione & automatica; manuale possibile anche per hidden/cell state\\
Lunghezze diverse & padding o troncamento\\
Batch size & 32--256\\
Shuffle & se le sequenze non sono correlate tra loro\\
Ottimizzatore, LR & Adam (SGD + momentum per la generalizzazione); partire da 0,001 con scheduler\\
Epoch \& early stopping & monitorare la validation loss; il numero di epoch cambia molto con ottimizzatore e LR\\
Regolarizzazione & dropout (0,2--0,5), come layer separato o dentro la RNN; L2 (0,01)\\
Dopo la RNN & valutare se aggiungere layer lineari/non lineari o altre architetture\\
\bottomrule
\end{tabularx}
\end{center}

% =====================================================================
\section{RNN in PyTorch: lunghezze diverse e padding}
\label{sec:rnn-padding}
Una rete è moltiplicazione di matrici. Con tre serie di 2 feature e lunghezze 3, 2 e 3, al terzo passo una matrice ha un
``buco'': non si può moltiplicare per la matrice dei pesi (\rif{2026-10-14}{14/10}).
\begin{itemize}
  \item \textbf{Tagliare} tutto alla lunghezza minima: si butta informazione.
  \item \textbf{Paddare} con valori inventati (es.\ la media): la rete crederebbe che la serie finisca davvero con quei valori.
\end{itemize}
Serve qualcosa di più furbo: padding \textbf{dinamico}, \textbf{solo dentro il mini-batch}, e la RNN deve \textbf{sapere} quali
istanti sono padding.

\subsection{Dataset personalizzato e \code{collate\_fn}}
\begin{lstlisting}
from torch.utils.data import Dataset, DataLoader
from torch.nn.utils.rnn import pad_sequence, pack_padded_sequence

class TimeSeriesDataset(Dataset):
    def __init__(self, series, labels):      # series: lista di tensori (len_i, n_feat)
        self.series, self.labels = series, labels
    def __len__(self):
        return len(self.series)
    def __getitem__(self, idx):               # una serie e la sua label
        return self.series[idx], self.labels[idx]

def collate_fn(batch):
    seqs, labels = zip(*batch)
    lengths = torch.tensor([len(s) for s in seqs])        # lunghezze originali
    padded = pad_sequence(seqs, batch_first=True, padding_value=0.0)  # pad alla max del batch
    return padded, lengths, torch.tensor(labels)

loader = DataLoader(ds, batch_size=64, shuffle=True, collate_fn=collate_fn)
\end{lstlisting}
\begin{itemize}
  \item \code{collate\_fn} riceve i dati del mini-batch, trova la lunghezza massima \textbf{del batch} e padda le serie più corte.
        Restituisce anche le \textbf{lunghezze originali}.
  \item Si passa al \code{DataLoader}, che la usa \textbf{solo se serve}.
  \item Perché basta paddare nel batch? Tutti i passi temporali condividono gli stessi pesi: serve solo che la matrice del batch sia
        completa per moltiplicarla. Batch diversi possono avere lunghezze diverse.
  \item Il padding vale 0, che potrebbe essere un valore reale: per questo la RNN deve sapere dove inizia il padding.
\end{itemize}

\subsection{Il modello}
\begin{lstlisting}
class RNNClassifier(nn.Module):
    def __init__(self, input_size, hidden_size, num_layers, num_classes):
        super().__init__()
        self.rnn = nn.RNN(input_size,             # n. di feature per istante
                          hidden_size,            # dimensione dello stato (embedding)
                          num_layers=num_layers,  # layer ricorrenti impilati
                          batch_first=True)       # input (batch, seq_len, features)
        self.fc = nn.Linear(hidden_size, num_classes)   # classification head
    def forward(self, x, lengths):
        packed = pack_padded_sequence(x, lengths.cpu(), batch_first=True,
                                      enforce_sorted=False)
        out, hidden = self.rnn(packed)    # hidden: (num_layers, batch, hidden_size)
        return self.fc(hidden[-1])        # stato finale dell'ULTIMO layer
\end{lstlisting}
\begin{itemize}
  \item \code{pack\_padded\_sequence}: dice alla RNN quali istanti sono reali e quali padding, così li \textbf{salta}. Il padding serve
        solo alla matematica (matrice completa) e \textbf{non entra} nello stato finale né nella backpropagation: fondamentale, altrimenti
        si addestrerebbe anche sul padding.
  \item \code{hidden[-1]}: lo stato finale dell'\textbf{ultimo} layer, la rappresentazione più astratta della serie.
  \item \code{batch\_first=True}: senza, PyTorch si aspetta \code{(seq\_len, batch, features)}.
\end{itemize}

\subsection{\code{out} vs \code{hidden}}
\begin{center}
\small
\begin{tabularx}{\textwidth}{l X X}
\toprule
& \code{hidden} (\code{h\_n}) & \code{out}\\
\midrule
Cosa contiene & lo stato finale di \textbf{ogni layer} dopo tutta la serie & gli stati dell'\textbf{ultimo layer} a \textbf{ogni passo}\\
Forma & \code{(num\_layers, batch, hidden\_size)} & \code{(batch, seq\_len, hidden\_size)}\\
Uso & task a livello di \textbf{sequenza}: classificazione con head lineare/FFNN (\code{hidden[-1]}) & task a livello di \textbf{passo}: previsione del valore successivo, many-to-many\\
Con padding & con le packed sequence cattura l'ultimo passo \textbf{valido} & l'ultimo passo valido cambia da serie a serie: va selezionato, non preso alla cieca\\
\bottomrule
\end{tabularx}
\end{center}
\begin{esempio}[la scalata]
Sali in montagna registrando battito, pressione e temperatura. \code{hidden} è la rappresentazione dell'intero viaggio (in cima);
\code{out} ti dà la rappresentazione passo per passo, utile se vuoi prevedere il prossimo valore. Nel Lab 2 le sequenze di API call
descrivono un goodware o un malware: serve la rappresentazione \textbf{dell'intera serie} (cap.~\ref{cap:lab2}).
\end{esempio}

\subsection{Varianti: bidirezionale e LSTM}
\label{sec:rnn-varianti}
\begin{lstlisting}
# Bidirezionale
self.rnn = nn.RNN(input_size, hidden_size, num_layers, batch_first=True,
                  bidirectional=True)
self.fc = nn.Linear(hidden_size * 2, num_classes)     # 2 direzioni -> 2x
...
h = torch.cat((hidden[-2], hidden[-1]), dim=1)        # avanti + indietro
logits = self.fc(h)

# LSTM
self.lstm = nn.LSTM(input_size, hidden_size, num_layers, batch_first=True)
...
out, (h_n, c_n) = self.lstm(packed)     # restituisce hidden E cell state
logits = self.fc(h_n[-1])               # per classificare si usa l'hidden state

# Serie univariata: (batch, time) -> (batch, time, 1)
x = x.unsqueeze(-1)                     # [1,2,3] -> [[1],[2],[3]]
\end{lstlisting}

\begin{attenzione}[errori tipici nel codice RNN -- raccolta dei redattori]
\begin{itemize}
  \item Bidirezionale con \code{nn.Linear(hidden\_size, ...)}: l'input della head è \code{hidden\_size * 2}.
  \item \code{out, h = self.lstm(x)}: l'LSTM restituisce \code{(h\_n, c\_n)}; per classificare si usa \code{h\_n[-1]}, non il cell state.
  \item Serie univariata senza \code{unsqueeze(-1)}: \code{[1,2,3]} viene letto come un istante con 3 feature.
  \item Padding senza \code{pack\_padded\_sequence}: la rete si addestra anche sugli zeri di padding.
  \item \code{out[:, -1, :]} con serie paddate: l'ultimo passo può essere padding; con le packed sequence usare \code{hidden[-1]}.
  \item \code{input\_size} pari alla lunghezza della serie invece che al numero di feature per istante.
  \item \code{shuffle=True} su righe che compongono una serie (si mescolano solo serie intere e indipendenti).
\end{itemize}
\end{attenzione}

% =====================================================================
\section{Notebook 1: inquinamento (regressione multivariata)}
\begin{itemize}
  \item Label: livello di inquinamento (numero) $\Rightarrow$ \textbf{regressione}. 7 feature, una riga per ora, train e test già separati.
  \item \textbf{Guardare i dati prima di usarli}: nessun nullo; la direzione del vento è una stringa (SE, NE, cv\ldots) codificata in
        ordine di bussola con un label encoding (semplice, ma forse il one-hot sarebbe meglio).
  \item Plot di serie e distribuzioni: inquinamento fino a 1000 (outlier?) e vento a 600 km/h (\textbf{errore del sensore}?): va deciso
        cosa tenere.
  \item Validation: \textbf{non} con split casuale (romperebbe le serie: il punto 2 in validation tra 1 e 3 in training). Forse meglio
        creare prima le serie e poi dividerle.
  \item Scaling fittato sul training. Nell'esempio si normalizza anche la label: probabilmente non andrebbe fatto.
  \item Serie di lunghezza 7 tutte uguali $\Rightarrow$ basta \code{TensorDataset}. RNN + layer lineare, \textbf{nessuna attivazione finale}
        (regressione), loss \textbf{MSE}. Il training loop è identico a quello della FFNN (cap.~\ref{cap:ffnn}).
  \item Iperparametri del notebook: hidden size 64, output 1, batch 64, Adam con LR 0,001, 200 epoch.
  \item Le loss delle RNN sono \textbf{più rumorose} di quelle delle FFNN; qui non c'è plateau (andava addestrata di più). L'MSE va
        interpretato rispetto alla scala dei valori (nel notebook, su dati scalati: training da circa 0,0056 a 0,0048, validation
        intorno a 0,0055).
\end{itemize}
\begin{verifica}
Nel notebook la validation è ottenuta con \code{df.iloc[train\_size:]}, cioè con uno split \textbf{temporale} (le ultime righe del
training) fatto \textbf{prima} di creare le sequenze, non con uno split casuale. Le osservazioni a lezione sul rischio dello split
casuale restano valide come avvertimento generale.
\end{verifica}

% =====================================================================
\section{Notebook 2: FordA (classificazione di serie)}
\begin{itemize}
  \item Serie \textbf{univariate} con label che indica se il motore funziona; validation = 20\% del training.
  \item Per mostrare il padding si tolgono 20 elementi alle prime 100 serie (480 invece di 500).
  \item Con \code{batch\_size=100}: il primo batch ha tutte serie da 480 (nessun padding), il secondo da 500.
        Con 128: il primo batch contiene serie da 480 e da 500 $\Rightarrow$ padding a 500; i 20 istanti di padding non contano
        per i pesi.
  \item Modelli confrontati: RNN mono, LSTM e Bi-RNN, tutti con hidden size 64, 2 layer, Adam LR 0,001, \code{CrossEntropyLoss}.
  \item Risultato: \textbf{accuracy 51\%} su un dataset bilanciato, con predizione sempre della classe 0 (``tanto vale lanciare una moneta''),
        con RNN, LSTM e Bi-RNN. Non perché le RNN non funzionino: la head è troppo semplice, forse la hidden size è piccola, il
        training è breve. Su CPU è anche molto lento.
\end{itemize}
\begin{integrazione}[leggere la loss bloccata a 0,693]
Nell'output del notebook train e validation loss restano per tutte le 200 epoch intorno a $0{,}693 = \ln 2$: è la cross-entropy di un
classificatore binario che assegna probabilità $0{,}5$ a entrambe le classi. Una loss ferma esattamente lì segnala che il modello
\textbf{non sta imparando nulla}, prima ancora di guardare l'accuracy.
\end{integrazione}
\begin{verifica}
A lezione si è attribuito il risultato anche a un ``training breve'', ma nel notebook le epoch sono 200 (circa 260--390 s su GPU per
modello). Il problema sembra piuttosto di ottimizzazione/architettura (serie lunghe 500 passi, vanishing gradient, head minima).
\end{verifica}

\begin{sintesi}
\begin{itemize}
  \item RNN: stato nascosto $h_t = F_1(W_{hh}h_{t-1} + W_{xh}x_t + b)$ con \textbf{pesi condivisi} nel tempo; tanh tipica; varianti
        one-to-one, many-to-one (focus del corso), one-to-many, many-to-many; mono (tempo reale) o bidirezionale (sequenza completa).
  \item La parte ricorrente produce un \textbf{embedding} della serie; la head classifica; si addestrano insieme.
  \item BPTT: loss sommate nel tempo, gradiente attraverso output e stati; prodotti ripetuti di $W_{hh}$ $\Rightarrow$ exploding
        (gradient clipping) o vanishing gradient.
  \item LSTM: cell state (lungo termine) + hidden state (breve termine); forget, input, candidate (tanh), output gate (sigmoid).
  \item Dati: tensore \code{(batch, seq\_len, features)}; input size = n.\ di feature; lunghezze diverse $\Rightarrow$ padding nel batch
        con \code{collate\_fn} + \code{pack\_padded\_sequence}; classificazione con \code{hidden[-1]}, previsione passo-passo con \code{out}.
  \item Notebook: pollution (regressione, MSE, nessuna attivazione finale); FordA (51\%: loss ferma a $\ln 2$, il modello non impara).
\end{itemize}
\end{sintesi}

\chapter{Convolutional e graph neural network}
\label{cap:cnn-gnn}
\fonti{\rif{2026-10-14}{14/10}, \rif{2026-10-19}{19/10}}{3-deep-learning-methodologies, pp.~77--136; notebook GNN (PyTorch Geometric)}

Dopo FFNN (dati tabellari, cap.~\ref{cap:ffnn}) e RNN (sequenze, cap.~\ref{cap:rnn}), le ultime due architetture del corso:
le CNN per le \textbf{immagini} e le GNN per i \textbf{grafi}. Tutte servono a \textbf{estrarre una rappresentazione} dei dati
quando costruirla a mano è difficile; le GNN creano embedding di nodi, archi o grafi, e i pesi condivisi garantiscono che
l'embedding non dipenda dall'ordine di lettura.

% =====================================================================
\section{Convolutional neural network (CNN)}
Nate per le \textbf{immagini} (grazie anche alle label raccolte con i CAPTCHA), sono usate anche in cybersecurity:
\begin{enumerate}
  \item \textbf{malware} trasformato in immagine e classificato;
  \item \textbf{siti malevoli} da screenshot: un sito clone si riconosce anche visivamente (logo diverso, grafica imperfetta),
        utile contro tecniche come redirection, iframe nascosti, content hiding;
  \item \textbf{intrusion detection}: CNN su serie temporali, anche combinate (CNN-LSTM, CNN-GRU);
  \item verifica di \textbf{deepfake}.
\end{enumerate}

\subsection{Come funziona}
\begin{itemize}
  \item \textbf{Input}: immagine in forma numerica (ogni pixel un numero, o tre in RGB).
  \item \textbf{Feature extraction} con layer convoluzionali, in modo \textbf{gerarchico}:
  \begin{itemize}
    \item layer iniziali (\en{low-level}): bordi, angoli, texture (linee);
    \item layer intermedi (\en{mid-level}): forme, parti di oggetti;
    \item layer finali (\en{high-level}): rappresentazioni astratte, parti di oggetti o oggetti interi (naso, occhi, orecchie).
  \end{itemize}
  \item \textbf{Fully connected (FFNN)}: usa le feature estratte per classificare, con softmax (``cane: 0,9'').
\end{itemize}

\subsection{Il layer convoluzionale}
\begin{center}
\begin{tikzpicture}[node distance=0.6cm, every node/.style={font=\small},
  blk/.style={draw=polblu, fill=polblu!6, rounded corners=2pt, minimum height=0.8cm, minimum width=2.4cm, align=center}]
  \node[blk] (c) {Convolution\\{\scriptsize filtri scorrevoli}};
  \node[blk, right=of c] (a) {Activation\\{\scriptsize ReLU elemento per elemento}};
  \node[blk, right=of a] (p) {Pooling\\{\scriptsize downsampling}};
  \draw[-{Stealth}] (c) -- (a); \draw[-{Stealth}] (a) -- (p);
\end{tikzpicture}
\end{center}

\begin{definizione}[tensore: rango e forma]
\textbf{Rango} = numero di dimensioni (scalare 0, vettore 1, matrice 2\ldots). \textbf{Forma} (\en{shape}) = numero di elementi per
dimensione (vettore lungo 5: \code{[5]}; matrice $3\times5$: \code{[3,5]}). Un'immagine è un tensore di rango 3 con forma
\code{[d, h, w]}: altezza, larghezza e \textbf{profondità} $d$ (1 per scala di grigi, 3 per RGB).
\end{definizione}

\textbf{Convoluzione}
\begin{itemize}
  \item Un \textbf{filtro} (\en{kernel}), es.\ $3\times3$, scorre sull'immagine. Per ogni posizione si moltiplicano elemento per elemento
        i valori dei pixel con i valori del filtro e si sommano: il risultato è il valore del pixel in output.
  \item Così un pixel è descritto \textbf{insieme al suo intorno} (per capire se un punto fa parte di una linea servono i vicini).
  \item I valori dei filtri sono i \textbf{pesi addestrabili}. Ogni layer ha \textbf{centinaia o migliaia di filtri}.
  \item Ai \textbf{bordi} manca l'intorno: si aggiunge \textbf{padding} di zeri.
  \item Convolvendo un tensore \code{[d,h,w]} con $N$ filtri si ottiene \code{[N,h,w]}: \textbf{ogni filtro genera uno strato} della profondità
        dell'output.
  \item A lezione: nei primi layer kernel più grandi (informazione grossolana), andando in profondità kernel più piccoli (dettagli fini).
\end{itemize}
\begin{esempio}[un passo di convoluzione]
Finestra $3\times3$ dell'immagine e filtro $3\times3$:
\[
\begin{bmatrix}1&0&2\\0&1&0\\1&0&1\end{bmatrix} \ast
\begin{bmatrix}1&0&-1\\1&0&-1\\1&0&-1\end{bmatrix}
= (1-2) + (0-0) + (1-1) = -1 .
\]
Il valore $-1$ diventa il pixel in output per quella posizione; poi il filtro scorre di una posizione.
\end{esempio}
\textbf{Attivazione}: la convoluzione è lineare; si applica (tipicamente) \textbf{ReLU} elemento per elemento: training più rapido,
non satura, derivate veloci da calcolare.

\textbf{Pooling}: un filtro scorrevole che \textbf{sostituisce} i valori con una statistica riassuntiva. Il più comune è il
\textbf{max pooling}: con finestre $2\times2$ si prende il massimo, e una matrice $4\times4$ diventa $2\times2$.
Risultato: rappresentazione più compatta e densa.

Durante il training ogni filtro impara a riconoscere un pattern: nei layer superficiali texture e bordi (spesso senza un significato
evidente per noi), in quelli profondi oggetti e parti (occhi, orecchie, volti).

\begin{integrazione}[una CNN minima in PyTorch]
Non mostrata a lezione (le CNN non hanno un notebook nel corso); solo per fissare le forme dei tensori.
\begin{lstlisting}
class SmallCNN(nn.Module):
    def __init__(self, num_classes):
        super().__init__()
        self.conv = nn.Conv2d(1, 16, kernel_size=3, padding=1)  # [1,h,w] -> [16,h,w]
        self.pool = nn.MaxPool2d(2)                             # h,w dimezzati
        self.fc = nn.Linear(16 * 112 * 112, num_classes)        # input 224x224
    def forward(self, x):                 # x: (batch, 1, 224, 224)
        x = self.pool(torch.relu(self.conv(x)))
        return self.fc(x.flatten(1))      # head fully connected
\end{lstlisting}
\end{integrazione}

\subsection{CNN famose}
\begin{center}
\small
\begin{tabularx}{\textwidth}{l l X}
\toprule
\textbf{Rete} & \textbf{Anno} & \textbf{Caratteristiche}\\
\midrule
LeNet-5 & 1998 & 7 layer; introduce convoluzione + pooling; sigmoid/tanh (niente ReLU); cifre scritte a mano (MNIST), lettura assegni\\
AlexNet & 2012 & 8 layer (5 conv + 3 fully connected); \textbf{ReLU} contro il vanishing gradient; \textbf{dropout}; vince ImageNet 2012\\
ResNet-50 & 2015 & 50 layer; \textbf{skip (residual) connection} contro il vanishing gradient, permettono reti oltre 152 layer; vince ImageNet 2015\\
\bottomrule
\end{tabularx}
\end{center}

\subsection{CNN in cybersecurity: STAMINA}
\textbf{STAMINA} (\en{Static Malware-as-Image Network Analysis}, Intel):
\begin{enumerate}
  \item il binario del malware è letto come flusso di byte e mappato direttamente su pixel in scala di grigi (0--255);
  \item il flusso viene rimodellato in 2D (larghezza scelta con euristiche, es.\ dalla dimensione del file);
  \item l'immagine viene ridimensionata a $224\times224$ (input della CNN);
  \item si riusa \textbf{ResNet-50} pre-addestrata su immagini e la si fa \textbf{fine-tuning} su malware etichettati (split train/val/test);
  \item accuracy di classificazione dei malware \textbf{oltre il 95\%}.
\end{enumerate}
Per il computer i numeri sono numeri: un esempio di uso creativo di un'architettura nata per altro.

% =====================================================================
\section{Graph neural network: perché servono}
Molti dati sono grafi: Internet, social network (suggerire amici, account falsi, comunità), sistemi di raccomandazione, e in
cybersecurity \textbf{botnet}, \textbf{frodi} (conti che si scambiano denaro in modo strano), \textbf{malware} descritti da grafi di
API call.

\begin{definizione}[grafo]
$G$ con: $V$ insieme dei nodi; $A$ \textbf{matrice di adiacenza} (binaria, o pesata con feature degli archi);
$X\in\R^{|V|\times m}$ matrice delle \textbf{feature dei nodi}; $N(v)$ insieme dei \textbf{vicini} di $v$.
Anche gli \textbf{archi} possono avere feature (es.\ da quanto tempo dura un'amicizia).
Se i nodi non hanno feature: vettore indicatore (one-hot del nodo) o vettore costante $[1,1,\ldots,1]$.
\end{definizione}

\subsection{Perché non bastano FFNN, CNN e RNN}
\textbf{FFNN su $[A, X]$}: la matrice si può dare in input, ma:
\begin{itemize}
  \item è \textbf{sensibile all'ordine dei nodi}: leggendo i nodi come C, A, B, D, E invece di A, B, C, D, E la matrice cambia, e lo stesso
        grafo produce input diversi;
  \item ha $O(|V|)$ parametri: non scala;
  \item \textbf{non si applica a grafi di dimensione diversa}: un nuovo amico richiede un neurone in più e il ri-addestramento da zero.
\end{itemize}
\textbf{CNN e RNN} sono pensate per griglie e sequenze con \textbf{struttura fissa} (so chi viene prima e chi è il mio vicino). I grafi
hanno dimensione arbitraria, topologia complessa (nessuna località spaziale), \textbf{nessun ordine né punto di riferimento}, e spesso
sono dinamici e con feature multimodali.

\subsection{Generalizzare la convoluzione}
In una CNN un pixel è descritto dal suo intorno. In un grafo non c'è una nozione fissa di località né una finestra scorrevole (un nodo
``a metà'' dentro un kernel non ha senso). Però si può dire: \textbf{un nodo è descritto dai suoi vicini}, cioè dai nodi a un salto
(\en{hop}), indipendentemente dalla ``distanza fisica''. Una nozione di località senza struttura fissa.

% =====================================================================
\section{GNN: le proprietà fondamentali}
\begin{definizione}[permutation invariance --- livello grafo]
Il grafo non ha un ordine canonico dei nodi. Una funzione $f$ che mappa il grafo $G=(A,X)$ in un vettore $\R^d$ deve dare
\textbf{lo stesso risultato per qualsiasi ordinamento}:
\[ f(A_1, X_1) = f(A_2, X_2) \]
\end{definizione}
\begin{definizione}[permutation equivariance --- livello nodo]
Una funzione $f$ che mappa ogni nodo in una rappresentazione $d$-dimensionale ($\R^{|V|\times d}$) è equivariante se la
rappresentazione di un nodo \textbf{nella stessa posizione del grafo} non cambia con l'ordinamento:
\[ P f(A,X) = f(PAP^T, PX) \quad \text{per ogni permutazione } P. \]
\end{definizione}
Senza queste proprietà la rappresentazione dipenderebbe da come leggo i dati, non dal dato. ``La mia rappresentazione non deve cambiare
in base all'ordine in cui leggi i miei amici.''

\textbf{Obiettivo}: mappare i nodi in uno spazio $d$-dimensionale (\textbf{embedding}) in cui nodi simili sono vicini. È la stessa idea
delle feature categoriche ad alta cardinalità (cap.~\ref{cap:embedding}): il vantaggio del deep learning è la \textbf{feature extraction}.
Una GNN è composta da più funzioni equivarianti/invarianti; l'output sono embedding di nodi, sottografi o grafi.

% =====================================================================
\section{Message passing}
\begin{itemize}
  \item Il \textbf{vicinato} di un nodo definisce il suo \textbf{grafo di calcolo}: nelle immagini il vicinato è fisso, qui dipende dagli archi.
  \item I nodi \textbf{aggregano informazione dai vicini} con reti neurali; ogni embedding è un \textbf{messaggio} tra layer.
  \item Il vicino a sua volta è descritto dai \textbf{suoi} vicini: l'amico dell'amico passa il suo messaggio all'amico, che lo aggrega e lo
        passa a me (\en{reachability graph}, $k$ hop).
  \item I ``box'' del grafo di calcolo sono \textbf{pesi di una rete neurale}, \textbf{condivisi tra tutti i nodi} (come nelle RNN sono
        condivisi tra i passi temporali).
  \item Ogni nodo ha un grafo di calcolo diverso, perché ha vicini diversi.
\end{itemize}

\begin{center}
\begin{tikzpicture}[every node/.style={font=\small}, nd/.style={circle, draw=polblu, fill=polblu!10, minimum size=0.6cm, inner sep=0pt},
  box/.style={draw=esecol, fill=esecol!10, rounded corners=2pt, minimum width=0.7cm, minimum height=0.5cm}]
  % grafo di input
  \node[nd] (A) at (0,0) {A};
  \node[nd] (B) at (-1,1) {B}; \node[nd] (C) at (1,1) {C};
  \node[nd] (D) at (-1,-1) {D}; \node[nd] (E) at (1,-1) {E};
  \draw (A)--(B) (A)--(C) (A)--(D) (C)--(E) (B)--(C);
  \node[font=\scriptsize] at (0,-1.7) {grafo di input};
  % grafo di calcolo di A (2 hop)
  \node[nd, fill=red!15] (tA) at (6,1.2) {A};
  \node[box] (b0) at (6,0.4) {};
  \node[nd] (tB) at (4.5,-0.6) {B}; \node[nd] (tC) at (6,-0.6) {C}; \node[nd] (tD) at (7.5,-0.6) {D};
  \draw[-{Stealth}] (b0) -- (tA);
  \foreach \n in {tB,tC,tD} \draw[-{Stealth}] (\n) -- (b0);
  \node[font=\scriptsize] at (4.5,-1.3) {A,C}; \node[font=\scriptsize] at (6,-1.3) {A,B,E}; \node[font=\scriptsize] at (7.5,-1.3) {A};
  \node[font=\scriptsize] at (6,-1.9) {grafo di calcolo di A: layer 2 $\leftarrow$ layer 1 $\leftarrow$ feature (layer 0)};
\end{tikzpicture}
\end{center}

\subsection{Progettare una GNN: le quattro scelte}
\begin{enumerate}
  \item \textbf{Messaggio}: cosa rappresenta (l'embedding di un nodo o di un layer precedente).
  \item \textbf{Aggregazione}: come combinare i messaggi (la differenza principale tra le architetture).
  \item \textbf{Struttura}: quanti layer GNN impilare.
  \item \textbf{Task}: classificare nodi, archi o grafi? (Si può generalizzare anche a compiti non supervisionati.)
\end{enumerate}

\subsection{Un singolo layer: messaggio + aggregazione}
Un layer comprime un insieme di vettori (i vicini e il nodo stesso al livello $l-1$) in un solo vettore.
\begin{enumerate}
  \item \textbf{Messaggio}: $m_u^{(l)} = \mathrm{MSG}^{(l)}(h_u^{(l-1)})$, es.\ lineare $m_u^{(l)} = W^{(l)}h_u^{(l-1)}$ con
        $W^{(l)}$ \textbf{uguale per tutti i nodi}.
  \item \textbf{Aggregazione}: $h_v^{(l)} = \mathrm{AGG}^{(l)}(\{m_u^{(l)}, u\in N(v)\})$, es.\ somma, media, massimo: funzioni
        \textbf{indipendenti dall'ordine}.
  \item \textbf{Problema}: così si perde il nodo stesso (``anch'io esisto''). \textbf{Soluzione}: un messaggio del nodo con una matrice
        diversa, $m_v^{(l)} = B^{(l)}h_v^{(l-1)}$, aggregato dopo i vicini (concatenazione o somma):
        \[ h_v^{(l)} = \mathrm{CONCAT}\big(\mathrm{AGG}(\{m_u^{(l)}, u\in N(v)\}),\ m_v^{(l)}\big) \]
  \item \textbf{Non linearità} $\sigma$ (su messaggio o aggregazione) per l'espressività: \textbf{ReLU} (la più usata), sigmoid (solo per
        limitare il range degli embedding), \textbf{PReLU} (ReLU parametrica, empiricamente spesso migliore).
\end{enumerate}
\begin{esame}[perché una GNN è invariante/equivariante e induttiva]
\begin{itemize}
  \item I pesi sono \textbf{condivisi} tra tutti i nodi e l'aggregazione (somma, media, max) \textbf{non dipende dall'ordine}
        $\Rightarrow$ la rappresentazione non cambia con l'ordine di lettura.
  \item Le GNN sono \textbf{induttive}: gli stessi parametri valgono per \textbf{nodi mai visti} (un nuovo nodo H si descrive senza
        ri-addestrare), purché il grafo sia simile.
\end{itemize}
\end{esame}

% =====================================================================
\section{Tre architetture: GCN, GraphSAGE, GAT}
\begin{center}
\small
\renewcommand{\arraystretch}{1.5}
\begin{tabularx}{\textwidth}{l X X}
\toprule
\textbf{Architettura} & \textbf{Idea} & \textbf{Formula}\\
\midrule
\textbf{GCN} \newline (Graph Convolutional Network) & usa \textbf{tutta} la matrice di adiacenza, con self-edge; normalizza per il grado (con 500 milioni di vicini la somma esploderebbe) & $h_v^{(l)} = \sigma\Big(\sum_{u\in N(v)} \frac{W^{(l)}h_u^{(l-1)}}{|N(v)|}\Big)$\\
\textbf{GraphSAGE} \newline (Sample and Aggregate) & \textbf{campiona un numero fisso di vicini} (Cristiano Ronaldo non ha bisogno di 500 milioni di follower per essere descritto); aggrega in \textbf{due stadi}: prima i vicini, poi il nodo stesso & $h_{N(v)}^{(l)} = \mathrm{AGG}(h_u^{(l-1)}, u\in N(v))$\newline $h_v^{(l)} = \sigma\big(W^{(l)}\cdot\mathrm{CONCAT}(h_v^{(l-1)}, h_{N(v)}^{(l)})\big)$\\
\textbf{GAT} \newline (Graph Attention Network) & pesa ogni vicino con un \textbf{coefficiente di attenzione} $\alpha_{vu}$ \textbf{appreso}: i vicini importanti contano di più & $h_v^{(l)} = \sigma\Big(\sum_{u\in N(v)}\alpha_{vu}W^{(l)}h_u^{(l-1)}\Big)$\\
\bottomrule
\end{tabularx}
\end{center}
\begin{itemize}
  \item Aggregatori di GraphSAGE: \textbf{media} (dei vicini), \textbf{pooling} (MLP + max), \textbf{LSTM} applicata ai vicini in ordine
        rimescolato. Qualsiasi funzione che non dipenda dall'ordine.
  \item In GCN e GraphSAGE si ha in pratica $\alpha_{vu} = 1/|N(v)|$: tutti i vicini contano uguale. In GAT, invece, ``se non ascolto mai
        un vicino, perché dargli lo stesso peso di chi ascolto sempre?''. L'attenzione sarà centrale nella parte di NLP.
\end{itemize}
\begin{sintesi}[quale usare]
\textbf{GCN} se il grafo non è troppo grande e si vuole un'architettura semplice; \textbf{GraphSAGE} per grafi grandi (campionamento);
\textbf{GAT} per rappresentazioni migliori. Si possono usare GCN e GAT anche su grafi grandi, ma il campionamento va gestito a mano.
\end{sintesi}

% =====================================================================
\section{Impilare layer e over-smoothing}
\begin{itemize}
  \item L'embedding al layer 0 è la feature del nodo, $h_v^{(0)} = x_v$; l'embedding al layer $k$ riceve informazione dai nodi a
        \textbf{$k$ hop}.
  \item \textbf{Receptive field}: l'insieme dei nodi che determinano l'embedding di un nodo. Con $K$ layer è il vicinato a $K$ hop.
\end{itemize}
\begin{esempio}[perché più layer = più vicini]
Catena diretta A $\to$ B $\to$ C. Con \textbf{un} layer: B riceve da A, C riceve da B. Con \textbf{due} layer: C riceve da B, che al passo
precedente aveva già ricevuto da A $\Rightarrow$ C ``conosce'' anche A.
\end{esempio}
\begin{definizione}[over-smoothing]
Con troppi layer i receptive field crescono in fretta e si sovrappongono; se due nodi hanno receptive field molto sovrapposti, i loro
embedding sono molto simili. Al limite ogni nodo diventa la media dell'intero grafo: \textbf{tutti gli embedding si somigliano} e non
distinguono più ``me da Cristiano Ronaldo''.
\end{definizione}
\textbf{Rimedi}
\begin{enumerate}
  \item Non aggiungere layer GNN se non servono: \textbf{analizzare il receptive field necessario} (es.\ diametro del grafo, quanti nodi coinvolgo
        con 1, 2, 3 hop) e fissare $L$ di conseguenza, non inutilmente grande.
  \item Aumentare l'\textbf{espressività dentro ogni layer}: messaggio e aggregazione possono essere reti profonde invece di un layer lineare.
  \item Aggiungere \textbf{layer che non passano messaggi}: MLP prima (pre-processing, utile se i nodi sono immagini o testo da codificare)
        e dopo (post-processing, utile per classificare grafi).
  \item Soluzioni avanzate: \textbf{skip connection} tra layer GNN (l'informazione passa anche direttamente, non solo mediata).
\end{enumerate}
\begin{esame}[over-smoothing]
Cos'è (embedding tutti simili per receptive field sovrapposti), come si riconosce (distanze tra embedding dei nodi tutte piccole e
prestazioni scarse) e come si rimedia (meno layer GNN in base al receptive field, layer più espressivi, MLP prima/dopo, skip
connection, normalizzazione).
\end{esame}

% =====================================================================
\section{Task e split dei dati}
\subsection{Il task}
\begin{center}
\small
\begin{tabularx}{\textwidth}{l X X}
\toprule
\textbf{Task} & \textbf{Come si predice} & \textbf{Esempi}\\
\midrule
Node classification & direttamente dagli embedding dei nodi & IP malevoli, account compromessi; monitorare i frigoriferi di una catena di supermercati (in Francia un bug firmware li ha trasformati in forni)\\
Link prediction & da \textbf{coppie} di embedding di nodi & suggerire amici; relazioni nascoste nelle frodi; computer che comunicano ma non dovrebbero; percorsi di attacco\\
Graph classification & da \textbf{tutti} gli embedding del grafo (aggregati) & traffico malevolo, tipi di botnet (es.\ di dispositivi IoT); grafo delle API call (API1 $\to$ API3 $\to$ API1) malware o goodware\\
\bottomrule
\end{tabularx}
\end{center}
Task diversi richiedono aggregazioni e classification head diverse. Nei lab si fa solo \textbf{graph classification}: lo stesso dato del Lab 2
(sequenze di API call) si può vedere come serie o come grafo (cap.~\ref{cap:lab2}).

\subsection{Come dividere i dati}
\begin{itemize}
  \item \textbf{Transductive}: un solo grafo; train/val/test sullo stesso grafo, struttura e feature visibili a tutti gli split, si dividono le
        label (o gli archi target). Dividere un grafo in parti è complicato.
  \item \textbf{Inductive}: più grafi; ogni split contiene grafi non visti in training. È il setting dei lab (split per grafo e label).
        Nodi mai visti in validation/test non sono un problema: le GNN sono induttive.
\end{itemize}

\subsection{Limiti e sfide}
\begin{itemize}
  \item \textbf{Over-smoothing}: non eccedere con i layer, o aggiungere elaborazione.
  \item \textbf{Scalabilità}: molti messaggi da passare e aggregare; usare l'architettura giusta (GraphSAGE).
  \item \textbf{Grafi eterogenei}: un modello addestrato sulla rete del Politecnico può andare male su una rete diversa; servono architetture specializzate.
  \item \textbf{Grafi dinamici}: cambiano in continuazione (follower aggiunti e rimossi); difficile anche validare.
\end{itemize}

% =====================================================================
\section{GNN in PyTorch Geometric}
Implementazione mostrata il \rif{2026-10-19}{19/10} con il notebook GNN.
\subsection{Preparare i dati}
\begin{itemize}
  \item \textbf{Preprocessing} come per le altre reti: feature categoriche dei nodi in one-hot, normalizzazione o standardizzazione delle
        feature di nodi e archi.
  \item Un grafo è rappresentato da:
  \begin{itemize}
    \item \textbf{nodi} ($V$), eventualmente con feature;
    \item \textbf{archi} ($E$), come lista di adiacenza o \code{edge\_index} (\code{[2, n\_archi]}: origine e destinazione). Gli archi sono
          \textbf{diretti}: un arco percorribile nei due sensi diventa \textbf{due archi};
    \item \textbf{feature dei nodi} $X$ (una riga per nodo, una colonna per feature);
    \item \textbf{feature degli archi} (opzionali);
    \item \textbf{label} $y$: a livello di nodo o di grafo (nei lab: di grafo).
  \end{itemize}
\end{itemize}

\begin{lstlisting}
from torch_geometric.data import Data
from torch_geometric.loader import DataLoader, NeighborLoader
from torch_geometric.nn import GCNConv, SAGEConv, GATConv, global_mean_pool

# ogni elemento del dataset e' un GRAFO
g = Data(x=node_features,         # (n_nodi, n_feature)
         edge_index=edge_index,   # (2, n_archi)
         y=graph_label)           # label del grafo
train_loader = DataLoader(train_graphs, batch_size=32, shuffle=True)
\end{lstlisting}
\begin{itemize}
  \item \code{Data} e \code{Dataset} vengono da \textbf{\code{torch\_geometric}}, non da \code{torch}. Si può scrivere una classe dataset
        personalizzata (con \code{\_\_len\_\_} e \code{get}) che costruisce un oggetto \code{Data} per grafo da adiacenza, feature e label.
  \item Il \code{DataLoader} di \code{torch\_geometric.loader} gestisce \textbf{automaticamente} grafi di dimensioni diverse nello stesso batch,
        \textbf{senza padding} (unisce i grafi in un grande grafo disconnesso). Differenza con le RNN, dove serviva la \code{collate\_fn}
        (cap.~\ref{cap:rnn}).
\end{itemize}

\begin{esempio}[ispezionare il dataset del notebook]
3 classi, classificazione di grafi, una feature per nodo in one-hot con 5 valori, oltre 1.400 grafi di training.
\begin{itemize}
  \item Primo grafo: \code{Data(x=[126, 5], edge\_index=[2, 816], y=[1])} $\Rightarrow$ 126 nodi, 5 feature, 816 archi, una label.
  \item Un batch di 2 grafi ha \code{x=[288, 5]}: i nodi dei due grafi \textbf{sommati}. Nel batch ``l'unità'' sono i nodi, poi aggregati.
  \item Il grafo più piccolo ha 99 nodi: già difficile da visualizzare.
\end{itemize}
\end{esempio}

\subsection{Visualizzare un grafo}
Con NetworkX: \code{to\_networkx} converte il grafo, \code{spring\_layout} fissa le posizioni dei nodi (ci sono vari layout, più o meno
leggibili), \code{nx.draw} disegna con etichette dei nodi (le feature) e degli archi, e la label come titolo. Utile, se i nodi non sono
troppi, per capire ad esempio se il grafo è denso. Come il plot di una serie temporale per scoprire outlier.

\subsection{Il modello}
\begin{lstlisting}
import torch.nn.functional as F

class GCN(nn.Module):
    def __init__(self, in_dim, hidden_dim, num_classes):
        super().__init__()
        self.conv1 = GCNConv(in_dim, hidden_dim)       # n. feature nodo -> embedding
        self.conv2 = GCNConv(hidden_dim, hidden_dim)   # secondo layer: 2 hop
        self.fc = nn.Linear(hidden_dim, num_classes)   # classification head
    def forward(self, data):
        x, edge_index, batch = data.x, data.edge_index, data.batch
        x = F.relu(self.conv1(x, edge_index))   # edge_index: a chi mandare i messaggi
        x = F.relu(self.conv2(x, edge_index))
        x = global_mean_pool(x, batch)          # nodi -> UN embedding per grafo
        return self.fc(x)
\end{lstlisting}
\begin{itemize}
  \item \code{edge\_index} permette alla GCN di ricostruire il grafo e passare i messaggi ai vicini giusti.
  \item \code{global\_mean\_pool} aggrega gli embedding dei nodi di ogni grafo in \textbf{un embedding del grafo} (necessario per la
        graph classification). Fermandosi prima si hanno gli embedding dei nodi. Alternative: max, min\ldots
  \item \textbf{GraphSAGE}: \code{SAGEConv} al posto di \code{GCNConv} (aggregazione a due stadi). Così però \textbf{non} si campionano i vicini.
  \item \textbf{GAT}: \code{GATConv(in\_dim, hidden\_dim, heads=8, concat=True)}: 8 meccanismi di attenzione indipendenti; con
        \code{concat=True} gli output si concatenano $\Rightarrow$ dimensione \code{hidden\_dim * 8}.
\end{itemize}

\subsection{Il training loop}
Quasi identico ai precedenti (cap.~\ref{cap:ffnn}), con due differenze:
\begin{itemize}
  \item ogni elemento del batch è un \textbf{grafo}, e la label è un \textbf{attributo} del batch (\code{batch.y}), non un elemento separato;
        \code{batch.y.squeeze()} porta la label da 2D a 1D;
  \item per pesare la loss si usa il \textbf{numero di grafi} (\code{batch.num\_graphs}), \textbf{non} \code{x.size(0)}, che conterebbe i
        \textbf{nodi} (288 invece di 2).
\end{itemize}
\begin{lstlisting}
for batch in train_loader:
    batch = batch.to(device)
    optimizer.zero_grad()
    out = model(batch)
    loss = criterion(out, batch.y.squeeze())
    loss.backward(); optimizer.step()
    running += loss.item() * batch.num_graphs     # NON batch.x.size(0)
\end{lstlisting}

\textbf{Risultati sul dataset ``facile''}: GCN, GraphSAGE e GAT vanno tutte bene in training, validation e test (oltre il 90\%). La GAT
converge meglio (è più espressiva) ma è molto più lenta: circa 98 s GCN, 71 s GraphSAGE, 281 s GAT. Non è stato fatto tuning: con un solo
set di iperparametri la validation non serve a sceglierne.

\subsection{GraphSAGE con campionamento dei vicini}
\begin{itemize}
  \item La classe GraphSAGE restituisce \textbf{solo gli embedding} (niente pooling); pooling e head sono in una funzione separata
        (\code{graph\_forward}).
  \item Il \code{NeighborLoader} campiona i vicini: es.\ batch di 64 nodi, \code{num\_neighbors=[10, 5]} $\Rightarrow$ 10 vicini al primo
        hop, 5 al secondo. Il modello riceve i nodi campionati, poi si ripete il \code{global\_mean\_pool} e la FFNN finale.
  \item In \textbf{validation e test} si usa \code{[-1]}: \textbf{tutti} i vicini. Il campionamento serve solo a velocizzare il training.
\end{itemize}

\begin{attenzione}[errori tipici nel codice GNN -- raccolta dei redattori]
\begin{itemize}
  \item Pesare la loss con \code{batch.x.size(0)} (numero di \textbf{nodi}) invece di \code{batch.num\_graphs}.
  \item Graph classification senza \code{global\_mean\_pool} (o altro readout): la head riceverebbe un output per nodo.
  \item \code{GATConv(..., heads=8, concat=True)} seguito da un layer con input \code{hidden\_dim}: l'input è \code{hidden\_dim * 8}.
  \item \code{edge\_index} con un solo verso per archi non diretti: servono due archi.
  \item Campionamento dei vicini (\code{num\_neighbors=[10,5]}) anche in validation/test: lì si usano tutti i vicini (\code{[-1]}).
  \item \code{DataLoader} di \code{torch.utils.data} invece di quello di \code{torch\_geometric.loader}.
\end{itemize}
\end{attenzione}

% =====================================================================
\section{Scelte progettuali e iperparametri di una GNN}
\begin{enumerate}
  \item \textbf{Tipo}: GCN (grafi non troppo grandi, architettura semplice), GraphSAGE (grafi grandi, campiona i vicini; nodi con grado
        altissimo), GAT (attenzione, rappresentazioni migliori).
  \item \textbf{Task}: nodo, arco o grafo? Determina l'embedding finale.
  \item \textbf{Scala}: grafi piccoli $\Rightarrow$ full-batch (GCN, GAT); grafi grandi $\Rightarrow$ mini-batch con campionamento.
  \item \textbf{Aggregazione}: media/somma, normalizzata per grado, pesata con attenzione, LSTM; partire da qualcosa di semplice.
  \item \textbf{Input}: numero di feature dei nodi (e degli archi).
  \item \textbf{Hidden size}: dimensione dell'embedding. Più grande = più informazione, troppo grande = informazione diluita.
  \item \textbf{Numero di layer GNN}: in base al \textbf{receptive field} (quanti vicini aggregare). Partire con 2--3; troppi
        $\Rightarrow$ over-smoothing.
  \item \textbf{Batching}: lo gestisce la libreria.
  \item \textbf{Prediction head}: un layer lineare o una FFNN; per i grafi prima il \textbf{pooling/readout}.
  \item \textbf{Output e attivazione}: in base al task (un output per nodo, per arco candidato o per grafo).
  \item \textbf{Training}: shuffle dei grafi (ok per la graph classification), loss pesata se sbilanciato, ottimizzatore.
  \item \textbf{Over-smoothing}: si verifica studiando gli embedding dopo la GNN: se le distanze tra nodi sono tutte piccole e le prestazioni
        scarse, ridurre i layer, usare skip connection o normalizzazione.
\end{enumerate}

\begin{center}
\small
\renewcommand{\arraystretch}{1.15}
\begin{tabularx}{\textwidth}{>{\bfseries}l X}
\toprule
Iperparametro & Best practice (dalle slide)\\
\midrule
\# layer & 2--3; di più solo se serve (over-smoothing); valutare MLP, dropout\ldots\\
Hidden size & potenze di 2 (64, 128, 256); troppo grande $\Rightarrow$ overfitting\\
Attivazione & ReLU, LeakyReLU; softmax/sigmoid in output\\
Inizializzazione & default di PyTorch; Xavier (Glorot) o He possono stabilizzare\\
Feature dei nodi & StandardScaler per le continue, one-hot per le categoriche\\
Batch size & full-batch per grafi piccoli; mini-batch (512--2048 nodi) per grafi grandi\\
Shuffle & grafi mescolati nella graph classification; nella node classification si mescolano i nodi target, preservando la connettività\\
Ottimizzatore & Adam; AdamW per il weight decay; SGD + momentum per grafi grandi\\
Learning rate & 0,01 per grafi piccoli, 0,001 per GNN profonde; scheduler\\
Regolarizzazione & dropout (0,2--0,5), L2 (0,01), skip connection (cap.~\ref{cap:regolarizzazione})\\
Dopo la GNN & layer MLP per la classificazione; connessioni residue per GNN profonde\\
\bottomrule
\end{tabularx}
\end{center}

\begin{sintesi}
\begin{itemize}
  \item CNN: convoluzione (filtri addestrabili scorrevoli, \code{[d,h,w]} con $N$ filtri $\to$ \code{[N,h,w]}) + ReLU + pooling;
        feature gerarchiche (bordi $\to$ forme $\to$ oggetti) e head fully connected. LeNet-5, AlexNet (ReLU, dropout), ResNet-50
        (skip connection). STAMINA: malware come immagine $224\times224$ + ResNet-50 fine-tuned, $>95\%$.
  \item GNN: i grafi non hanno ordine né dimensione fissa $\Rightarrow$ servono funzioni \textbf{invarianti} (grafo) ed \textbf{equivarianti}
        (nodo); message passing con pesi condivisi e aggregazioni indipendenti dall'ordine $\Rightarrow$ modello \textbf{induttivo}.
  \item GCN (media normalizzata per grado), GraphSAGE (campionamento, due stadi), GAT (attenzione appresa).
  \item $K$ layer = receptive field a $K$ hop; troppi $\Rightarrow$ over-smoothing. 2--3 layer di partenza.
  \item Task: nodo, link, grafo (lab: graph classification, setting induttivo).
  \item PyTorch Geometric: \code{Data(x, edge\_index, y)}, \code{DataLoader} senza padding, \code{global\_mean\_pool}, loss pesata con
        \code{num\_graphs}.
\end{itemize}
\end{sintesi}

\chapter{Feature categoriche ed embedding}
\label{cap:embedding}
\fonti{\rif{2026-10-19}{19/10}, \rif{2026-10-20}{20/10}}{4-embedding-introduction, pp.\ 1--51}

Tutte le architetture viste finora (RNN nel cap.~\ref{cap:rnn}, GNN nel cap.~\ref{cap:cnn-gnn}) servono a \textbf{estrarre una
rappresentazione} dei dati quando costruirla a mano è difficile (le feature di una serie o di un grafo). Questo capitolo affronta lo
stesso problema per un tipo di dato onnipresente in cybersecurity: le \textbf{feature categoriche}. Prima si vedono le codifiche
classiche (ID sequenziali, feature engineering, one-hot), poi gli \textbf{embedding appresi} e i due modi di impararli: con le label
(supervisionati) o senza (self-supervised).

% =====================================================================
\section{Feature categoriche: dalle categorie ai numeri}
Le reti (e molti algoritmi ML) richiedono input numerici e assumono relazioni matematiche significative tra i numeri.
L'obiettivo è convertire le categorie in numeri \textbf{significativi per il problema} e per la rete (assegnare numeri a caso
come 10, 5, 6, 1 non lo è). Le feature categoriche possono essere stringhe o numeri.
\begin{itemize}
  \item Esempi: titoli di lavoro, processi chiamati durante l'esecuzione di un file, algoritmi di cifratura;
        soddisfazione del cliente, livello di accesso a un server, gravità di un attacco.
  \item Per le \textbf{label} la soluzione migliore è quasi sempre l'ID sequenziale (label encoding).
\end{itemize}

\subsection{Prima domanda: le categorie hanno un ordine intrinseco?}
\textbf{Ordinali} (soddisfazione, livello di accesso, gravità): \textbf{ordina e mappa}.
\begin{enumerate}
  \item identifica i valori unici;
  \item ordinali con la \textbf{conoscenza del dominio} (non in ordine alfabetico!) e assegna indici: Low = 1, Medium = 2, High = 3;
  \item mappa il dataset originale.
\end{enumerate}
Il modello non sa cosa significhi la categoria né se ``alto'' sia buono (gravità: basso è meglio; soddisfazione: alto è meglio):
sa solo che 1 è più vicino a 2 che a 3, e lo scoprirà durante il training.

% =====================================================================
\section{Feature nominali: le alternative}
Esempio: AES, DES, RC4, RSA. Non c'è una mappatura intuitiva. Il corso presenta quattro opzioni.

\subsection{1. ID sequenziale comunque}
AES = 1, DES = 2, RC4 = 3, RSA = 4.
\begin{itemize}
  \item Mettiti nei panni del modello: perché AES è più vicino a DES che a RSA? Il ML scopre pattern nei dati: una mappatura sbagliata lo
        mette su una \textbf{falsa pista} (gli dai un \textbf{bias} da ignorare).
  \item Anche se da esperto ordino per sicurezza (meno sicuro = 1, più sicuro = 4): se domani RSA viene rotto, devo cambiare la mappatura
        e il suo significato. Il modello impara ``RSA viene dopo AES'', non perché siano simili.
\end{itemize}

\subsection{2. Feature engineering}
AES, RSA = sicuri = 1; DES, RC4 = insicuri = 0.
\begin{itemize}
  \item \textbf{Pro}: riduce lo sforzo di interpretazione: dai al modello una chiave di lettura, con un vincolo più debole dell'ordine.
  \item \textbf{Contro}: e se la chiave è sbagliata o cambia? (RSA oggi sicuro, domani rotto.) In un problema benigno/malevolo il modello
        potrebbe associare gli algoritmi insicuri al traffico malevolo (bias). E a volte il feature engineering è difficile o impossibile.
\end{itemize}

\subsection{3. One-hot encoding}
\begin{enumerate}
  \item si identifica l'insieme delle categorie, il \textbf{vocabolario} $V$;
  \item ogni categoria diventa un vettore \textbf{sparso} $V$-dimensionale (tutto 0 tranne un 1); i vettori sono \textbf{ortogonali}.
\end{enumerate}
\begin{itemize}
  \item \textbf{Pro}: semplice, \textbf{nessun bias}: tutti i vettori sono equidistanti; con algoritmi tradizionali (random forest) diventa
        $V$ feature binarie (presenza/assenza di AES, di DES\ldots). Funziona anche con lo shallow learning.
  \item \textbf{Contro}: \textbf{maledizione della dimensionalità} (tante feature = problema più difficile; 171.000 parole inglesi);
        se il \textbf{vocabolario si estende} (es.\ ChaCha20, o un futuro algoritmo post-quantum) cambia la rappresentazione e il training fatto
        va buttato: si riparte da zero.
\end{itemize}
\begin{verifica}
Distanza tra vettori one-hot: due vettori one-hot distinti hanno distanza euclidea $\sqrt{2}$ (la distanza \emph{al quadrato} è 2).
La slide scrive ``= 2''.
\end{verifica}

\subsection{4. Embedding appresi (solo reti neurali)}
\begin{enumerate}
  \item si identifica il vocabolario $V$ (es.\ $V=4$);
  \item ogni categoria si mappa su un vettore \textbf{denso} $H$-dimensionale, con $H \ll V$; ogni categoria ha il suo vettore
        ($V$ vettori diversi); all'inizio i vettori sono \textbf{casuali}, senza relazioni;
  \item gli embedding si \textbf{imparano risolvendo un task} (es.\ sicuro vs insicuro) tramite backpropagation, per molte epoch.
\end{enumerate}
Si può fare solo con le reti neurali perché solo lì si può mandare indietro il segnale e cambiare i valori; random forest e simili no.

\begin{esempio}[cosa succede ``dietro le quinte'']
Rete: \code{nn.Linear(V, H)} $\to$ \code{nn.Linear(H, 2)} $\to$ softmax, con $V=4$, $H=2$; input = one-hot degli algoritmi; label = sicuro/insicuro.
\begin{itemize}
  \item \textbf{Epoch 1}: pesi casuali. Dopo il primo layer ogni algoritmo ha una rappresentazione 2D senza senso (es.\ AES $=[-0{,}7, 0{,}3]$).
        Il modello dà probabilità alte a predizioni sbagliate (è sicuro di sbagliare): 3 giuste su 5.
  \item \textbf{Epoch 2}: la backprop cambia i pesi del primo layer, quindi la rappresentazione: lo spazio inizia a separarsi; 4 giuste su 5.
  \item \textbf{Epoch 3}: tutte giuste; nello spazio 2D gli algoritmi sicuri hanno coordinate positive (AES $[2{,}1, 2{,}5]$, RSA $[1{,}9, 2{,}0]$),
        quelli insicuri negative (DES $[-2{,}1, -1{,}9]$, RC4 $[-1{,}1, -1{,}3]$).
\end{itemize}
Il vettore $H$-dimensionale ``incorpora'' (\en{embeds}) le caratteristiche più salienti dell'oggetto in numeri significativi per il task:
per questo si chiama \textbf{embedding}. Si possono impilare più layer per rappresentazioni più astratte.
\end{esempio}

\begin{center}
\small
\begin{tabularx}{\textwidth}{X X}
\toprule
\textbf{Benefici degli embedding} & \textbf{Problemi}\\
\midrule
Rappresentazione \textbf{densa} per costruzione; parte casuale e impara una distribuzione significativa (che ottimizza il task) & Solo con reti neurali (non con random forest)\\
Più efficiente del one-hot: l'informazione non è sparsa & Serve un \textbf{task} che dia il segnale di backprop (supervisionato o self-supervised)\\
Le feature codificano \textbf{semantica}, non solo informazione binaria & Se il vocabolario si estende, la lookup ha ancora $V$ righe: ri-addestrare da zero? training incrementale?\\
Dimensione molto ridotta ($H \ll V$: es.\ 128 invece di 171.000) & \\
\bottomrule
\end{tabularx}
\end{center}

\section{Categorie sconosciute: il token \code{<UNK>}}
Un nuovo algoritmo o un nuovo dominio non sono nel vocabolario: non hanno una rappresentazione. Si aggiunge al vocabolario una categoria
\textbf{sconosciuta} (\code{<UNK>}) a cui si mappa tutto ciò che non è stato visto in training. Succede sempre: validation e test possono
contenere valori mai visti (il test set spesso è raccolto molto dopo).

\section{La roadmap: dalle categorie ai numeri}
\begin{center}
\begin{tikzpicture}[node distance=0.55cm and 0.9cm, every node/.style={font=\small, align=center},
  q/.style={draw=polblu, fill=polblu!6, rounded corners=2pt, minimum width=3.4cm, minimum height=0.8cm},
  a/.style={draw=esecol, fill=esecol!10, rounded corners=2pt, minimum width=2.8cm, minimum height=0.8cm}]
  \node[q] (q0) {Feature categorica};
  \node[q, below=of q0] (q1) {Ha un ordine intrinseco?};
  \node[a, right=of q1] (a1) {ID sequenziali};
  \node[q, below=of q1] (q2) {Poche categorie?};
  \node[a, right=of q2] (a2) {One-hot encoding};
  \node[q, below=of q2] (q3) {Posso creare un task\\per una rete neurale?};
  \node[a, right=of q3] (a3) {Embedding};
  \node[a, below=of q3] (a4) {Feature engineering};
  \draw[-{Stealth}] (q0) -- (q1);
  \draw[-{Stealth}] (q1) -- node[above, font=\scriptsize]{sì} (a1);
  \draw[-{Stealth}] (q1) -- node[left, font=\scriptsize]{no} (q2);
  \draw[-{Stealth}] (q2) -- node[above, font=\scriptsize]{sì} (a2);
  \draw[-{Stealth}] (q2) -- node[left, font=\scriptsize]{no} (q3);
  \draw[-{Stealth}] (q3) -- node[above, font=\scriptsize]{sì} (a3);
  \draw[-{Stealth}] (q3) -- node[left, font=\scriptsize]{no} (a4);
\end{tikzpicture}
\end{center}
\begin{itemize}
  \item \textbf{Feature engineering} quando il resto non va: con Intesa Sanpaolo, per 120--1000 transazioni per utente al mese, si sono
        progettate oltre 100 feature riassuntive, in numero gestibile qualunque fosse il numero di transazioni.
  \item \textbf{Si può usare un ID sequenziale per una feature nominale?} Non si \emph{dovrebbe}, ma \textbf{si può}, se sai cosa stai
        facendo e l'ordine ha senso per il tuo task. Non è una scienza esatta: le scelte sub-ottimali si pagano (più dati, più epoch,
        training più instabile) e il rischio è introdurre bias.
  \item La lista non è esaustiva (es.\ \en{feature hashing}).
\end{itemize}

% =====================================================================
\section{Come imparare gli embedding (1): supervisionati}
Gli embedding si possono imparare in due modi (slide 4, p.~40): con un segnale \textbf{supervisionato} o \textbf{self-supervised}.
L'idea: invece di fare noi feature engineering (a volte impossibile), lasciamo che sia il \textbf{task} a guidare la rappresentazione
delle feature categoriche.

\textbf{Ipotesi}: scenario \textbf{supervisionato} (label disponibili: traffico benigno/malevolo, email di phishing, famiglia di botnet\ldots),
dati con feature numeriche e categoriche, e una \textbf{rete neurale}: solo le reti possono imparare embedding. Con un random forest si possono
risolvere classificazione e regressione, ma non fare \textbf{feature extraction}: ``ottenere embedding'' = feature extraction.

\begin{definizione}[lookup table e \code{nn.Embedding}]
\code{torch.nn.Embedding(V, H)} crea una tabella \textbf{addestrabile} con $V$ righe (numero di categorie, la dimensione del vocabolario) e
$H$ colonne (la dimensione dell'embedding, un \textbf{iperparametro}). È equivalente a un one-hot seguito da un layer lineare. Permettendo la
backpropagation, la rete cambia i valori della tabella per ottimizzare il task.
\end{definizione}
\begin{integrazione}
In pratica \code{nn.Embedding} riceve \textbf{indici interi} (\code{torch.long}) e restituisce la riga corrispondente: è il prodotto one-hot
$\times$ matrice dei pesi senza costruire il vettore one-hot, molto più efficiente in memoria.
\end{integrazione}
Si addestra la rete sul task, si aggiustano i pesi dei layer in base agli errori, e la rappresentazione si ``taglia'' a metà rete (dopo il primo
layer e prima dell'ultimo).

\subsection{Esempio 1: dataset tabellare con feature categoriche e numeriche}
Metadati di email (dominio del mittente, numero di link, numero di allegati, lunghezza dell'oggetto, ora di ricezione); task: phishing (1)
o normale (0). Il \textbf{dominio del mittente} ha cardinalità altissima (migliaia di domini con un mail server) $\Rightarrow$ il one-hot creerebbe
milioni di colonne, e non c'è un ordine (Gmail prima di Outlook o di Polito?). $\Rightarrow$ \textbf{embedding}.

\begin{enumerate}
  \setcounter{enumi}{-1}
  \item \textbf{Split} train/validation/test.
  \item \textbf{Vocabolario dal solo training}: è ciò che si suppone di conoscere nella realtà. Il test replica il futuro, e anche la validation
        non va considerata nota. Si aggiunge \code{<UNK>} per i valori mai visti (un dominio nuovo come \code{updatebanking.com} o \code{polito.it} in
        validation $\Rightarrow$ entrambi \code{<UNK>}, stesso embedding). Gli ID servono solo come indici per la lookup, non come feature.
  \item \textbf{Mappare} la feature categorica negli ID.
  \item \textbf{Dataset} con feature categoriche e numeriche \textbf{separate} (vanno trattate diversamente), convertite in tensori.
  \item \textbf{Modello}: la categorica passa nell'\code{nn.Embedding}, l'embedding si \textbf{concatena} alle numeriche, il risultato entra nella FFNN.
        Il primo layer lineare ha \textbf{embed\_dim + num\_numeric} neuroni di input: una sola feature categorica diventa 16, 32\ldots numeri.
\end{enumerate}

\begin{lstlisting}
def build_vocab(train_values):
    vocab = {'<PAD>': 0, '<UNK>': 1}
    for v in train_values:                 # SOLO training set
        if v not in vocab:
            vocab[v] = len(vocab)
    return vocab

def encode(values, vocab):
    return [vocab.get(v, vocab['<UNK>']) for v in values]   # mai visto -> UNK

class EmbFFNN(nn.Module):
    def __init__(self, vocab_size, embed_dim, num_numeric, num_classes=2):
        super().__init__()
        self.emb = nn.Embedding(vocab_size, embed_dim)
        self.net = nn.Sequential(
            nn.Linear(embed_dim + num_numeric, 64), nn.ReLU(),
            nn.Linear(64, num_classes))
    def forward(self, cat_ids, num_feats):
        e = self.emb(cat_ids)                    # (batch, embed_dim)
        x = torch.cat([e, num_feats], dim=1)     # concatena con le numeriche
        return self.net(x)
\end{lstlisting}
L'embedding layer è un layer come gli altri: in backpropagation il segnale arriva anche a lui e ne migliora la rappresentazione.

\subsection{Esempio 2: sequenze di valori categorici}
Una \textbf{darknet} è un insieme di indirizzi IP passivi che non ospitano servizi: riceve per definizione solo \textbf{traffico non richiesto}.
Input: sequenza degli IP che ogni ora contattano dei server (lista di lunghezza variabile); task: il server è stato attaccato da una botnet (e
quale)? Una sequenza $\Rightarrow$ \textbf{RNN} (cap.~\ref{cap:rnn}), in cui ogni elemento è una categoria (un IP). Lo stesso schema servirà
nel Lab 2 con le sequenze di API call (cap.~\ref{cap:lab2}).
\begin{enumerate}
  \item split; vocabolario degli IP dal training; mappatura in ID (come prima);
  \item DataLoader per RNN con padding dentro il batch (o padding dinamico con \code{collate\_fn});
  \item modello: \code{nn.Embedding} sulla sequenza (PyTorch gestisce i singoli elementi), poi LSTM, poi head lineare. L'input della LSTM non è 1
        (una categoria) ma \textbf{embed\_dim}. Si può indicare l'ID del padding (\code{padding\_idx}): il suo embedding \textbf{non riceve backpropagation}.
\end{enumerate}
\begin{lstlisting}
class IPSeqLSTM(nn.Module):
    def __init__(self, vocab_size, embed_dim, hidden, num_classes, pad_idx=0):
        super().__init__()
        self.emb = nn.Embedding(vocab_size, embed_dim, padding_idx=pad_idx)
        self.lstm = nn.LSTM(embed_dim, hidden, batch_first=True)  # input = embed_dim
        self.fc = nn.Linear(hidden, num_classes)
    def forward(self, ids, lengths):
        e = self.emb(ids)                                  # (batch, seq, embed_dim)
        packed = pack_padded_sequence(e, lengths.cpu(), batch_first=True,
                                      enforce_sorted=False)
        _, (h, _) = self.lstm(packed)
        return self.fc(h[-1])
\end{lstlisting}

\subsection{Esempio 3: architetture composte (movimento laterale)}
La feature extraction si può fare con \textbf{qualsiasi} architettura, anche combinandole. Problema: distinguere i \textbf{movimenti laterali}
(spostarsi da una macchina all'altra, es.\ verso il computer di un amministratore) in una rete aziendale.
\begin{itemize}
  \item I \textbf{nodi} sono i servizi attivi nella rete (identificatori categorici).
  \item Ogni \textbf{grafo} codifica i movimenti di un utente sui nodi in un periodo (es.\ 1 ora).
  \item Gli attacchi possono essere \textbf{APT} (\en{Advanced Persistent Threat}): si monitora a snapshot diversi $\Rightarrow$ ogni utente ha una
        \textbf{serie temporale di grafi}.
\end{itemize}
Architettura finale:
\begin{enumerate}
  \item \textbf{embedding layer} (lookup): mappa i servizi in embedding;
  \item \textbf{GNN} (cap.~\ref{cap:cnn-gnn}): rappresenta come l'utente si è mosso tra i servizi in uno snapshot (la stessa GNN per tutti gli snapshot);
  \item \textbf{RNN}: tiene conto dei diversi snapshot e produce la predizione (movimento laterale o attività normale).
\end{enumerate}
Tutto si addestra con \textbf{lo stesso segnale supervisionato}: l'errore della predizione si backpropaga attraverso RNN (nel tempo), GNN (per
nodo) e fino all'embedding layer, che riceve la somma dei gradienti di tutti i rami (timestamp e nodi). Non è magia: è moltiplicazione di matrici.
Un utente legittimo può fare SSH da una macchina all'altra per arrivare al server (passando dalla VPN): non tutto il movimento è malevolo.

\begin{esame}[gli embedding supervisionati sono specifici del task]
\textbf{Si potevano progettare feature a mano?} Molto difficile: servirebbero più feature per ogni IP e servizio (importanza nell'azienda, cosa
fa\ldots), senza sapere se sono adatte. Gli embedding invece sono per costruzione \textbf{adatti a questo problema}.

\textbf{E se domani voglio usare gli stessi servizi per prevedere quale servizio sarà più usato?} Numericamente posso riusare gli embedding, ma
sono stati appresi per un task diverso: \textbf{ottimi per questo task, sub-ottimi per altri}. Con le label, la rappresentazione è funzionale al
problema finale, non generica.
\end{esame}

\subsection{Il rischio: cosa sta imparando? (husky vs lupo)}
Come distingui un husky da un lupo nelle foto? La via più ``pigra'': gli husky sono fotografati sull'\textbf{erba}, i lupi sulla \textbf{neve}.
Il classificatore impara lo \textbf{sfondo}: una \textbf{scorciatoia} (\en{shortcut}) che funziona sui dati di training ma non è generale.
La rappresentazione è ottima per \emph{quelle} label e \emph{quei} dati. È lo stesso fenomeno visto nel Lab 1 e nel discorso sulle feature che
creano bias (porta 80, cap.~\ref{cap:ffnn}).

% =====================================================================
\section{Come imparare gli embedding (2): self-supervised}
\begin{definizione}[self-supervised learning, SSL --- slide 4, p.~43]
Un algoritmo self-supervised \textbf{genera un segnale supervisionato da dati non etichettati}: sono i dati stessi a fornire la supervisione.
È un sottoinsieme dell'apprendimento non supervisionato, ma mentre l'unsupervised \emph{descrive} i dati, il self-supervised usa il segnale
per imparare a \emph{rappresentarli}.
\begin{itemize}
  \item \textbf{Pretext task}: si impara una rappresentazione significativa dai dati non etichettati (es.\ un encoder che comprime l'input
        in uno spazio a dimensione minore, il \en{bottleneck}: è la fase di feature extraction).
  \item \textbf{Downstream task}: si usano quelle feature per il compito vero (classificazione, ricostruzione\ldots).
\end{itemize}
\end{definizione}
Il self-supervised learning è la pietra angolare dei successi recenti del deep learning:
\begin{itemize}
  \item permette di imparare feature \textbf{generiche} (``cosa rende una banana una banana'');
  \item oggi è preferito alle feature end-to-end supervisionate;
  \item \textbf{pretext}: si pre-addestra il modello in self-supervision su grandi corpora;
        \textbf{downstream}: lo si fa fine-tuning con poche label;
  \item cruciale in cybersecurity, dove le label sono poche e si vuole generalizzare.
\end{itemize}

Nel supervisionato il segnale viene dalle label: scoiattolo vs tigre $\to$ feature extractor $\to$ embedding $\to$ classification head
$\to$ predizione sbagliata $\to$ backprop. Senza label: \textbf{da dove viene il segnale?}
Idea generale: sfruttare le \textbf{proprietà interne degli oggetti} per capire cosa è simile e cosa no.

\subsection{Opzione 1: imparare a ricostruire (autoencoder)}
\begin{itemize}
  \item L'\textbf{encoder} estrae le feature rilevanti: codifica l'oggetto nello spazio degli embedding.
  \item Il \textbf{decoder} usa gli embedding per risolvere un compito, es.\ ricostruire (o togliere la sfocatura da) l'immagine originale.
  \item Durante il training è un problema di \textbf{regressione} (predire le componenti RGB dei pixel): \textbf{nessuna label}.
  \item Al test: \textbf{encoder congelato} + una piccola head addestrata con poche label. Ci si aspetta che una specie mai vista (uno scoiattolo
        di un altro colore) finisca vicina agli scoiattoli nello spazio degli embedding, così la head lo classifica bene.
\end{itemize}
L'autoencoder è ripreso in dettaglio per l'anomaly detection nel cap.~\ref{cap:anomaly}.

\subsection{Opzione 2: confronto tra viste (contrastive learning)}
\begin{itemize}
  \item Il soggetto di una foto ruotata, sfocata o decolorata resta lo stesso: è la \textbf{data augmentation}.
  \item Dall'immagine originale (\textbf{anchor}) si creano varianti (\textbf{positivi}): i loro embedding \textbf{devono essere vicini}.
        Un'immagine diversa (\textbf{negativo}) deve avere un embedding \textbf{lontano}.
  \item La loss che premia vicinanza dei positivi e lontananza dei negativi è la \textbf{contrastive loss}: un'altra loss da aggiungere a MSE,
        BCE e cross-entropy. Si chiama ``contrastiva'' perché contrappone il negativo all'anchor e alle sue perturbazioni.
  \item Il feature extractor, poiché codifica l'oggetto nello spazio degli embedding, si chiama anche \textbf{encoder} (slide 4, p.~50).
  \item Utile anche quando i campioni sono pochi (es.\ anomaly detection).
\end{itemize}
\begin{attenzione}[le augmentation introducono bias]
Scegliere le trasformazioni significa iniettare conoscenza: se dico al modello che una foto decolorata è ``la stessa'' immagine, gli insegno
che il colore non conta. Allora riconoscerebbe come falsa la foto di uno \textbf{scoiattolo verde}? Probabilmente no. (Domanda aperta della
slide 4, p.~51; la risposta è un ragionamento, non una regola.)
\end{attenzione}

\begin{sintesi}
\begin{itemize}
  \item Feature categoriche \textbf{ordinali}: ordina con la conoscenza del dominio e mappa su ID sequenziali.
  \item Feature \textbf{nominali}: ID sequenziale (introduce un ordine falso, bias), feature engineering (chiave di lettura che può essere
        sbagliata o invecchiare), one-hot (nessun bias ma alta dimensionalità e vocabolario rigido), \textbf{embedding} appresi (densi,
        $H \ll V$, semantici; solo reti neurali, serve un task).
  \item Vocabolario costruito \textbf{solo sul training}; \code{<UNK>} per i valori mai visti; \code{padding\_idx} per il padding.
  \item \code{nn.Embedding(V, H)} = one-hot + layer lineare: tabella addestrabile in backpropagation, combinabile con FFNN, RNN, GNN.
  \item Embedding \textbf{supervisionati}: ottimi per il task, sub-ottimi per altri, a rischio di scorciatoie (husky vs lupo).
  \item Embedding \textbf{self-supervised} (pretext + downstream): ricostruzione (autoencoder) o contrastive learning (anchor, positivi,
        negativi); le augmentation scelte iniettano bias.
\end{itemize}
\end{sintesi}

\chapter{Supervised vs unsupervised e anomaly detection}
\label{cap:anomaly}
\fonti{\rif{2026-11-02}{02/11}, \rif{2026-11-09}{09/11}}{5-Supervised-vs-Unsupervised-shallow, pp.\ 1--57;
6-Supervised-vs-Unsupervised-deep, pp.\ 1--52; notebook \code{anomaly\_detection} (KDD Cup 1999)}

È l'ultima parte del corso tenuta dal Prof.\ Giordano; dopo si passa all'NLP con il Prof.\ Boffa. Si studia l'anomaly detection perché
i classificatori falliscono quando le label mancano, sono fortemente sbilanciate o inaffidabili, o quando si vogliono scoprire
\textbf{novità}. Il capitolo segue l'ordine delle slide: concetti e sfide, algoritmi tradizionali (\en{shallow}), riduzione e
visualizzazione dei dati, deep learning (autoencoder), validazione; chiude con il notebook su KDD Cup 1999 che prepara il Lab 3
(cap.~\ref{cap:lab3}).

\begin{nota}[Copertura delle registrazioni]
La lezione del \rif{2026-11-02}{02/11} arriva fino a k-means (slide 5, p.~33). La lezione del \rif{2026-11-09}{09/11} è interamente dedicata
al notebook e dà per vista la teoria successiva: k-means per l'anomaly detection, DBSCAN, LOF, maledizione della dimensionalità, PCA,
Isolation Forest (slide 5, pp.\ 33--57) e tutta la slide 6 (deep learning per l'anomaly detection, t-SNE, validazione). Quella teoria
\textbf{non è stata spiegata nelle registrazioni disponibili}: nelle sezioni corrispondenti il testo segue le slide del corso, ed è
segnalato all'inizio di ciascuna.
\end{nota}

% =====================================================================
\section{Cos'è l'anomaly detection}
\begin{definizione}
Identificazione di \textbf{pattern che deviano dal comportamento atteso}. Si parla anche di anomalie, outlier, novità.
\begin{itemize}
  \item \textbf{Anomaly / outlier detection}: punti insoliti in un dataset, che non ci si aspetta; spesso un comportamento \textbf{malevolo}.
  \item \textbf{Novelty detection}: pattern nuovi, mai visti, che però possono essere la \textbf{nuova normalità} (benigni).
\end{itemize}
\end{definizione}
Una difficoltà chiave è distinguere ciò che è malevolo da ciò che è semplicemente nuovo. Si applica a punti, serie temporali, qualsiasi dato.

\subsection{In cybersecurity}
Le anomalie possono indicare:
\begin{itemize}
  \item \textbf{attività malevole}: zero-day, intrusioni, \textbf{esfiltrazione} (qualcuno connesso da remoto inizia a muoversi tra più PC,
        \en{lateral movement}, e poi manda fuori molti dati: ma quale volume è ``insolito'' non lo so a priori);
  \item \textbf{guasti o configurazioni errate} del sistema (bug);
  \item \textbf{frodi bancarie}: riciclaggio, aggiramento di sanzioni ed embarghi (es.\ utenti che inviano, ricevono o spendono importi insoliti).
\end{itemize}
\textbf{Perché è importante}
\begin{itemize}
  \item \textbf{Zero-day}: minacce senza firma nota. Per definizione non si può costruire un classificatore: se lo avessi, non sarebbe uno zero-day.
  \item \textbf{Attaccanti adattivi}: le difese a firme o regole falliscono perché l'attaccante cambia comportamento (gatto e topo).
  \item \textbf{Early warning}: rilevare comportamenti nuovi o furtivi mentre accadono, prima del danno.
\end{itemize}
In generale l'anomaly detection tratta eventi \textbf{minoritari, imprevedibili e rari}: da qui complessità che nella classificazione non
c'erano o erano meno forti.

\section{Le complessità del problema}
\begin{enumerate}
  \item \textbf{Unknowingness}: le anomalie sono legate a comportamenti, strutture e distribuzioni sconosciuti o improvvisi; restano ignote finché
        non accadono. \emph{Esempio}: picco di CPU di un processo: bug (riavvio) o qualcuno che usa il PC per minare Bitcoin?
  \item \textbf{Rarità e sbilanciamento}: la maggior parte dei dati è normale (in un IDS quasi tutti gli accessi sono legittimi).
        Raccogliere molte anomalie etichettate è irrealistico; un classificatore fatica a modellare pochi punti. Inoltre
        \textbf{sbagliare su un'anomalia costa molto più} che sbagliare su un dato normale. \emph{Esempio}: accessi ai file quasi tutti in orario
        di lavoro, uno di notte: io che lavoro da remoto, o un problema?
  \item \textbf{Classi di anomalie eterogenee}: le anomalie non sono una classe sola. \emph{Esempio}: un attacco esfiltra lentamente dati via
        \textbf{DNS tunnelling}, un altro usa \textbf{HTTP/2 push} cifrate, alcune molto simili a quelle legittime. Lo stesso tipo di attacco si
        divide in più ``classi'' $\Rightarrow$ lo sbilanciamento peggiora.
  \item \textbf{Pattern che evolvono}: le minacce si adattano per evadere il rilevamento, i modelli statici invecchiano; label e fingerprint
        note diventano presto obsolete. \emph{Esempio}: dal ``principe nigeriano'' al falso ``migliore amico'', che non finisce nello spam.
\end{enumerate}
\begin{attenzione}[un'anomalia può essere una novità]
Con il COVID e il lavoro da remoto, l'uso della VPN (anche in orari strani) è esploso: un'anomalia che è diventata la nuova normalità.
Un rilevatore che scatta e \textbf{blocca subito tutte le VPN} fermerebbe tutti i lavoratori da remoto. Il \textbf{contesto} conta:
progetta con cura e controlla cosa fa il tuo algoritmo; non tutte le anomalie sono attacchi.
\end{attenzione}

\section{Tipi di anomalie}
\begin{center}
\small
\renewcommand{\arraystretch}{1.25}
\begin{tabularx}{\textwidth}{>{\bfseries}l X X}
\toprule
Tipo & Definizione & Esempio\\
\midrule
Point & singola istanza che devia dalla maggior parte delle altre; basta il valore del punto & login aziendale da un IP in Russia (o Nord Europa) quando tutti i login dell'utente erano dall'Italia\\
Contextual & istanza anomala \textbf{solo dato il contesto}; da sola sarebbe normale & backup in cloud di due ore alle 3 di domenica: la durata è nella norma, ma di domenica non si lavora (il contesto è il giorno)\\
Collective & un \textbf{insieme} di istanze forma l'anomalia; i singoli membri possono essere normali & tentativi di login ognuno sotto il limite per minuto, ma tanti di fila: brute force che evita il rilevamento\\
\bottomrule
\end{tabularx}
\end{center}

\section{Le sfide per AI/ML}
\begin{enumerate}
  \item \textbf{Recall bassa e molti falsi positivi}: anomalie rare ed eterogenee sono difficili da trovare tutte. Se si spinge la recall, si
        moltiplicano i falsi allarmi: se il 90\% delle segnalazioni non è un'anomalia, nessuno si fiderà più del modello. Gli attacchi sofisticati
        (che imitano il normale) sfuggono.
  \item \textbf{Dati ad alta dimensionalità e dipendenti}: log, flussi di rete, dati temporali e a grafo, di più fonti (PC, server, rete).
        Qui il deep learning aiuta, perché non richiede di modellare a mano comportamenti strani.
  \item \textbf{Apprendimento efficiente in dati}: etichettare anomalie è costoso e spesso impraticabile $\Rightarrow$ metodi semi-supervisionati
        con poche label o, più spesso, \textbf{solo non supervisionati}.
  \item \textbf{Resilienza al rumore}: i dati reali sono rumorosi e con label sbagliate. Togliendo i colori a un plot, le anomalie vicine ai dati
        normali sembrano rumore; la normalità non è monolitica, e l'attaccante sfrutta il rumore per restare sotto i radar.
  \item \textbf{Anomalie complesse}: contestuali, collettive, su più tipi di dati e fonti: una soglia sul numero di tentativi non trova un
        attacco collettivo.
  \item \textbf{Spiegazione}: fiducia e \en{accountability}: l'amministratore vuole sapere perché un utente è segnalato (per chiedergli conto, o per
        scrivere una regola di firewall).
\end{enumerate}
\begin{esempio}[domanda in aula: se l'attaccante fa anomalie ogni giorno, diventano normali?]
Dipende da come vengono viste all'inizio. Se il comportamento resta al bordo della normalità e il modello non lo segnala, col tempo verrà
assorbito come normale (resta sotto i radar). Se invece viene segnalato, continuerà a esserlo, a meno che un analista non dichiari che è la nuova
normalità (in banca: ``questo utente lo abbiamo già controllato, ignoralo'').
\end{esempio}

% =====================================================================
\section{Supervised vs unsupervised}
\begin{center}
\small
\renewcommand{\arraystretch}{1.2}
\begin{tabularx}{\textwidth}{>{\bfseries}l X X}
\toprule
& Supervised & Unsupervised (e self-supervised)\\
\midrule
Dati & label almeno del normale; label di anomalie (non sempre, mai esaustive) & non etichettati\\
Assume & conosce i pattern normali; forse quelli anomali & il \textbf{comportamento normale domina}\\
Pro & alta accuratezza con buone label (es.\ fingerprint di virus noti, che non sono più ``anomalie'') & rileva \textbf{minacce nuove}\\
Limiti & servono label; non generalizza a nuove anomalie & più falsi positivi; più difficile da interpretare\\
\bottomrule
\end{tabularx}
\end{center}
Nella cybersecurity le label mancano spesso $\Rightarrow$ gli approcci non supervisionati sono essenziali.

\begin{definizione}[le label nell'anomaly detection]
Anche se le label ci sono, \textbf{non si usano per addestrare}. Servono solo a verificare se il modello ritrova almeno le anomalie
\textbf{note}: le anomalie sono rare e le label rappresentano solo ciò che conosciamo già; addestrando su di esse si rischia di
riconoscere solo i pattern noti e di perdere le minacce nuove (\rif{2026-11-09}{09/11}).
\end{definizione}

\section{L'output di un algoritmo di anomaly detection}
\begin{itemize}
  \item \textbf{Score}: un punteggio di anomalia per istanza (es.\ da 0 a 1, più alto = più anomalo). Permette di \textbf{ordinare} e di scegliere un
        \textbf{budget}: con Intesa Sanpaolo si analizzano le prime 10.000 istanze del mese secondo il ranking, invece di guardare tutto ciò che è
        sopra 0,01.
  \item \textbf{Label}: normale/anomalo per istanza. Più semplice da usare, ma il numero di anomalie si controlla con gli
        \textbf{iperparametri} (es.\ $\nu$ della one-class SVM).
\end{itemize}

% =====================================================================
\section{Algoritmi tradizionali (shallow)}
\subsection{Perché lo shallow learning}
\begin{itemize}
  \item Modelli più \textbf{interpretabili}: importante in incident response e analisi forense.
  \item Su dati \textbf{tabellari} (frequentissimi in cybersecurity) funzionano molto bene, a volte meglio del deep learning.
  \item Vanno bene per le \textbf{anomalie puntuali}: analizzano un punto alla volta.
  \item Danno una \textbf{baseline} per confrontare i modelli deep (se provi un solo modello deep non sai se è il migliore) e si possono
        combinare con essi (soluzioni ibride o ensemble).
\end{itemize}
\textbf{Limiti}
\begin{itemize}
  \item Deboli sulle anomalie \textbf{contestuali}: il contesto va creato con feature engineering (senza feature temporali, il backup della domenica
        sembra normale per la sua dimensione).
  \item Molto deboli sulle anomalie \textbf{collettive}: bisogna analizzare più punti insieme, cioè progettare feature aggregate (es.\ tentativi
        falliti per mezz'ora invece che per minuto).
\end{itemize}
Algoritmi del corso: \textbf{one-class SVM}; basati su distanza/densità: \textbf{LOF} (Local Outlier Factor); \textbf{clustering}
(k-means, DBSCAN); \textbf{PCA} (per la riduzione/rappresentazione dei dati, non come rilevatore); \textbf{Isolation Forest}.

% =====================================================================
\section{Support Vector Machine e one-class SVM}
\subsection{SVM}
Nate nel 1995, per anni più usate delle reti neurali (più semplici, scalabili, con solide basi matematiche).
\begin{definizione}[SVM]
Algoritmo \textbf{supervisionato} per classificazione (e outlier detection) che trova l'\textbf{iperpiano} che separa le classi
\textbf{massimizzando il margine}:
\begin{itemize}
  \item l'iperpiano è la funzione di decisione lineare $f(\mathbf{x}) = \mathbf{w}^T\mathbf{x} + b$;
  \item il \textbf{margine} è la distanza tra l'iperpiano e i punti più vicini (delle due classi), da massimizzare.
\end{itemize}
\end{definizione}
\begin{esempio}[quale iperpiano?]
Molti iperpiani $B_1, B_2, \ldots, B_n$ separano quadrati e cerchi. Si sceglie $B_1$ perché ha il \textbf{margine massimo}: è il più robusto.
Con $B_2$, un punto vicino al confine (``passi l'esame per un soffio'') cambierebbe classe per un piccolo spostamento.
\end{esempio}

\subsection{Il kernel trick}
Se i dati non sono linearmente separabili, si \textbf{trasformano in uno spazio a più dimensioni} dove lo diventano (es.\ con una funzione
polinomiale che combina le feature). Non sempre la separazione sarà perfetta: l'obiettivo è fare meno errori possibile. Il confine lineare
nello spazio trasformato, riproiettato nello spazio originale, appare curvo.
\begin{center}
\small
\begin{tabularx}{\textwidth}{l X}
\toprule
\textbf{Kernel} & \textbf{Caratteristiche}\\
\midrule
Lineare & prodotto scalare, nessuna trasformazione; raro, solo se i dati sono già linearmente separabili\\
Polinomiale & combinazioni di feature con un polinomio di grado $d$; raro, solo con relazioni polinomiali note\\
RBF (gaussiano) & funzione esponenziale con parametro $\gamma$ (alta similarità tra punti vicini); il \textbf{più usato}, crea confini molto non lineari\\
\bottomrule
\end{tabularx}
\end{center}

\subsection{One-class SVM}
\begin{definizione}[one-class SVM]
Variante della SVM pensata per \textbf{outlier, anomaly o novelty detection}. Lavora su \textbf{una sola classe}: non riceve label (la funzione
non ha nemmeno un argomento per passarle), quindi non è supervisionata. Crea un confine che descrive al meglio i \textbf{dati normali}; i punti
\textbf{fuori dal confine} sono anomalie. Con un kernel funziona anche su distribuzioni complesse e non lineari.
\end{definizione}

\textbf{Parametri principali}
\begin{itemize}
  \item \textbf{kernel} (come mappare i dati) e i suoi parametri ($\gamma$ per RBF, $d$ per il polinomiale);
  \item $\bm{\nu}$ (\en{nu}): la \textbf{frazione attesa di outlier}. Il modello cerca un confine che contenga più punti possibile lasciandone
        fuori circa una frazione $\nu$.
\end{itemize}
Con $\nu$ alto (es.\ 0,5) il modello considera anomalo metà dei punti e stringe il confine attorno alla zona più densa; con $\nu$ piccolo lascia
fuori pochi punti e il confine si \textbf{allarga}.
\begin{verifica}
La slide scrive ``the smaller, the fewer outliers and the narrower the boundaries'': ``fewer outliers'' è corretto, ``narrower'' va letto
con cautela, perché lasciando fuori meno punti la regione normale diventa più ampia.
\end{verifica}

\subsection{Due modi di usarla}
\begin{enumerate}
  \item \textbf{``Supervised'' (novelty detection)}: si addestra \textbf{solo su campioni normali} (perché ho poche label, normali e anomale, troppo
        poche per un classificatore). Si imposta un $\nu$ piccolo, il confine si stringe attorno al normale; a inferenza i nuovi punti fuori sono
        novità/anomalie. Il classificatore \textbf{non conosce} le label: le ho usate solo per scegliere i dati di training. Attenzione: se ci sono più
        ``normalità'', quelle rare possono risultare anomale.
  \item \textbf{Unsupervised (outlier/anomaly detection)}: si addestra su \textbf{normali e anomali insieme}, senza label (la vita reale). Serve
        un'idea della frazione di anomalie ($\nu$ = 50\%, 20\%, 10\%?): il confine conterrà circa $1-\nu$ dei punti.
\end{enumerate}

\begin{center}
\small
\renewcommand{\arraystretch}{1.2}
\begin{tabularx}{\textwidth}{>{\bfseries}l X X}
\toprule
& ``Supervised'' & Unsupervised\\
\midrule
Label & solo quelle della classe normale (per scegliere i dati) & nessuna\\
Dati di training & solo campioni normali & normali + anomali mescolati\\
Anomalie & escluse dal training & devono essere rare e ben separate\\
Obiettivo & confine della classe normale & confine assumendo che la maggior parte sia normale\\
Uso tipico & label del normale disponibili; cercare nuova normalità & outlier detection su dati grezzi non etichettati\\
Pro & più accurata se le label sono pulite; valutazione robusta (le label servono a misurare le prestazioni) & nessuna label; adatta alla realtà\\
Limiti & servono label normali affidabili; le anomalie non si possono usare in training & assume anomalie rare e distinte; peggiora con dati rumorosi\\
\bottomrule
\end{tabularx}
\end{center}
\begin{esempio}[perché le anomalie in training possono aiutare]
Addestrando solo sul normale con $\nu$ = 1\%, il confine si stringe attorno al 99\% dei punti normali e lascia fuori anche qualche normale.
Se invece sapessi che le anomalie sono molto lontane, includendole il confine si ``rilasserebbe'' un po' attorno al normale.
\end{esempio}
La one-class SVM individua campioni \textbf{globalmente isolati} e può avere problemi con il rumore.

% =====================================================================
\section{Clustering per l'anomaly detection}
\begin{definizione}[clustering]
Metodo \textbf{descrittivo} non supervisionato che trova gruppi (cluster) di punti simili tra loro e diversi da quelli degli altri gruppi:
\textbf{distanza intra-cluster piccola}, \textbf{distanza inter-cluster grande}.
\end{definizione}
Il clustering raggruppa, non trova anomalie: lo score va costruito (DBSCAN fa eccezione, perché etichetta il rumore). Le anomalie si
riconoscono come:
\begin{enumerate}
  \item \textbf{cluster piccoli} con \textbf{alta distanza inter-cluster} (lontani dagli altri);
  \item punti con \textbf{alta distanza intra-cluster}: appartengono a un cluster ma sono lontani dal centro;
  \item punti che \textbf{non appartengono a nessun cluster} (es.\ il ``rumore'' di DBSCAN).
\end{enumerate}

\subsection{K-means}
\begin{definizione}[k-means]
Algoritmo di clustering \textbf{basato su centroidi} e \textbf{partizionale}: divide i dati in $K$ cluster \textbf{distinti e non
sovrapposti}.
\begin{itemize}
  \item ogni cluster è rappresentato dalla \textbf{media} dei suoi punti (centroide), o da un punto reale (\en{medoid});
  \item \textbf{ogni punto appartiene a un cluster}: un punto senza cluster è impossibile;
  \item funziona bene con cluster \textbf{sferici e di dimensioni simili};
  \item minimizza la distanza intra-cluster.
\end{itemize}
Proposto da Steinhaus (1956), formalizzato da MacQueen (1967).
\end{definizione}

\textbf{Come funziona} (slide 5, p.~28). Ogni cluster è associato a un centroide e ogni punto è assegnato al cluster del centroide più
vicino; l'algoritmo base:
\begin{enumerate}
  \item scegli $K$ centroidi iniziali (spesso a caso);
  \item ripeti: assegna ogni punto al centroide più vicino; ricalcola ogni centroide come media dei suoi punti;
  \item fino a quando i centroidi non cambiano (o cambiano cluster pochi punti).
\end{enumerate}
\begin{itemize}
  \item Serve fissare $K$ e una \textbf{misura di distanza} (euclidea, similarità coseno, correlazione\ldots; nel corso: euclidea).
  \item Converge per le misure comuni; gran parte della convergenza avviene nelle prime iterazioni.
  \item Complessità $O(n \cdot K \cdot I \cdot d)$: $n$ punti, $K$ cluster, $I$ iterazioni, $d$ attributi.
\end{itemize}
\begin{verifica}
Nella slide 5, p.~28 l'algoritmo base compare solo come figura (``the basic algorithm is very simple''), non leggibile nel testo estratto:
i tre passi sopra sono la formulazione standard e non sono stati commentati a lezione.
\end{verifica}

\subsubsection{L'importanza dell'inizializzazione}
I centroidi iniziali sono casuali: due esecuzioni con lo stesso $K=3$ possono dare un clustering \textbf{ottimo} e uno \textbf{sub-ottimo}
(es.\ due cluster veri uniti e uno spezzato). \textbf{Soluzione}: eseguire k-means più volte con lo stesso $K$ e tenere il risultato migliore.

\subsubsection{Come scegliere $K$}
\begin{definizione}[SSE --- Sum of Squared Error]
Per ogni punto l'errore è la distanza dal centroide del suo cluster; si elevano al quadrato e si sommano:
\[
\mathrm{SSE} = \sum_{i=1}^{K}\sum_{x\in C_i}\mathrm{dist}^2(m_i, x)
\]
dove $m_i$ è il centroide (la media) del cluster $C_i$. Cluster densi $\Rightarrow$ SSE piccolo.
\end{definizione}
\begin{itemize}
  \item A parità di $K$ si sceglie il clustering con SSE minore.
  \item Aumentando $K$ l'SSE scende sempre (con $K=1$ trovi solo il centro dei dati; con $K = n$ è zero). Un buon clustering con $K$ piccolo può
        avere SSE minore di un cattivo clustering con $K$ grande.
  \item \textbf{Metodo del gomito} (\en{elbow}): si traccia l'SSE in funzione di $K$ e si sceglie il punto in cui il guadagno di un centroide in
        più diventa trascurabile (``3 o 4'', non 10: ci si può muovere un po' attorno al gomito). Più cluster = più cluster da studiare, e senza label
        sono difficili da interpretare.
\end{itemize}
\begin{attenzione}[SSE non monotono]
Se nel grafico l'SSE \textbf{risale} per un certo $K$, probabilmente k-means è stato eseguito una sola volta e un'inizializzazione sfortunata
ha creato un \textbf{cluster vuoto} (centroide in una zona senza punti): l'SSE è calcolato di fatto con un cluster in meno. Rieseguire k-means.
\end{attenzione}

\subsubsection{K-means per l'anomaly detection}
\begin{nota}[Non spiegato nelle registrazioni]
Da qui alla fine degli algoritmi shallow (slide 5, pp.\ 33--57) il testo segue le slide: la lezione registrata del \rif{2026-11-02}{02/11} si
interrompe qui e quella del \rif{2026-11-09}{09/11} applica questi metodi nel notebook dandone per vista la teoria.
\end{nota}
Si addestra k-means (se possibile solo sui normali), si calcola la distanza di ogni punto dal centroide del suo cluster e la si usa come
score (punti lontani = alta distanza intra-cluster); in più si cercano i \textbf{cluster piccoli}, magari lontani dagli altri, che possono
rappresentare eventi rari (slide 5, p.~33).

\subsection{Validare un clustering}
Slide 5, pp.\ 40--43.
\begin{itemize}
  \item \textbf{Indici esterni}: confronto con label fornite dall'esterno (entropia, purezza, metriche di classificazione).
  \item \textbf{Indici interni}: solo i dati. \textbf{Coesione} $= \mathrm{WSS} = \sum_i\sum_{x\in C_i}(x-m_i)^2$ (l'SSE);
        \textbf{separazione} $= \mathrm{BSS} = \sum_i |C_i|\,(m-m_i)^2$, con $m$ media globale.
\end{itemize}
\begin{definizione}[silhouette]
Per il punto $i$: $a(i)$ = dissimilarità media dagli altri punti del suo cluster, $b(i)$ = minima dissimilarità media verso un altro cluster;
\[
s(i) = \frac{b(i)-a(i)}{\max\{a(i),\,b(i)\}} \in [-1, 1]
\]
Vicino a 1: ben assegnato; negativo: nel cluster sbagliato. La media su un cluster misura quanto è compatto; la media su tutto il dataset
la qualità del clustering.
\end{definizione}

\subsection{DBSCAN}
Clustering \textbf{basato sulla densità} (slide 5, pp.\ 34--39 e 44): densità = numero di punti entro un raggio $\varepsilon$.
\begin{itemize}
  \item \textbf{Core point}: almeno \code{MinPts} punti entro $\varepsilon$ (interno di un cluster).
  \item \textbf{Border point}: meno di \code{MinPts} vicini, ma nel vicinato di un core point.
  \item \textbf{Noise point}: né core né border $\Rightarrow$ outlier, potenziale anomalia.
\end{itemize}
Non serve fissare il numero di cluster, trova cluster di forma arbitraria ed è resistente al rumore. Funziona male con densità diverse
tra i cluster: \textbf{la scelta di $\varepsilon$ cambia quali punti sono outlier}. Scelta dei parametri: nei cluster il $k$-esimo vicino è a
distanza simile per tutti i punti, nel rumore è più lontano; da qui il grafico delle distanze ordinate (\en{k-distance plot}).

\begin{center}
\small
\begin{tabularx}{\linewidth}{@{}l X X@{}}
\toprule
 & \textbf{K-means} & \textbf{DBSCAN}\\
\midrule
N.\ cluster & $K$ da fissare & trovato da solo\\
Rumore & ogni punto in un cluster & etichettato come rumore\\
Forma & sferica, dimensioni simili & arbitraria\\
Sensibilità & media ($K$, inizializzazione) & alta ($\varepsilon$ critico, MinPts medio)\\
Scalabilità & molto scalabile & più lento\\
Score di anomalia & distanza dal centroide, cluster isolati & rumore / bassa densità\\
\bottomrule
\end{tabularx}
\end{center}
Entrambi soffrono in alta dimensionalità.

% =====================================================================
\section{Local Outlier Factor (LOF)}
Slide 5, pp.\ 45--49. Idea: un outlier è \textbf{meno denso dei suoi vicini} (come in DBSCAN, ma con uno score graduato). Per un punto $p$
con insieme $O$ dei suoi $k$ vicini:
\begin{gather*}
\text{reach-dist}(p,o) = \max\{k\text{-dist}(o),\, d(p,o)\}\\
\mathrm{LRD}(p) = \frac{1}{\operatorname{avg}_{o\in O}\ \text{reach-dist}(p,o)}, \qquad
\mathrm{LOF}(p) = \frac{\operatorname{avg}_{o\in O}\ \mathrm{LRD}(o)}{\mathrm{LRD}(p)}
\end{gather*}
\begin{itemize}
  \item LOF $<1$: più denso dei vicini (dentro un cluster); $\approx 1$: densità simile, non outlier; $>1$: più sparso, potenziale outlier.
  \item Rispetto a k-means (assunzioni globali, niente rumore) e DBSCAN (rumore ma nessuno score) dà uno \textbf{score continuo} e trova anomalie
        \textbf{locali} in zone di densità diversa (es.\ login anomali da un blocco di IP frequente).
  \item Limiti: sensibile a $k$, costoso (molte distanze), poco scalabile, peggiora in alta dimensionalità.
\end{itemize}

% =====================================================================
\section{Dimensionalità: PCA e t-SNE}
\subsection{La maledizione della dimensionalità}
Slide 5, p.~50, che avverte: \emph{tutti} i metodi visti finora sono basati sulle distanze. Al crescere delle dimensioni i dati diventano
sparsi e la differenza tra distanza massima e minima tra coppie di punti si riduce (esperimento della slide: 500 punti casuali): densità e
distanza, alla base di clustering e outlier detection, perdono significato.

\subsection{PCA}
Slide 5, pp.\ 51--53. Introdotta da Pearson (1901) e Hotelling (1933). La PCA trova nuovi assi \textbf{ortogonali} (componenti), combinazioni
lineari delle feature, ordinati per varianza catturata: è una \textbf{rotazione} degli assi originali. Serve a ridurre le dimensioni, trovare
pattern e visualizzare. Uso nel corso:
\begin{enumerate}
  \item varianza spiegata da ogni componente;
  \item curva della varianza cumulata al crescere del numero di componenti;
  \item ginocchio della curva;
  \item trasformare il dataset tenendo quel numero di componenti.
\end{enumerate}
Il numero massimo di componenti è pari al numero di feature. Nel corso la PCA serve a ridurre e rappresentare i dati, non come rilevatore.

\subsection{t-SNE}
\begin{nota}[Non spiegato nelle registrazioni]
La teoria di t-SNE è nella slide 6, pp.\ 44--48; a lezione (\rif{2026-11-09}{09/11}) se ne è visto solo l'uso nel notebook.
\end{nota}
Proiezione \textbf{non lineare} per la visualizzazione:
\begin{enumerate}
  \item calcola la similarità tra tutte le coppie di punti nello spazio originale con un kernel gaussiano (probabilità che siano vicini);
  \item cerca una mappa in 2D/3D che preservi queste similarità, minimizzando con il gradient descent la divergenza di
        Kullback--Leibler tra le due distribuzioni.
\end{enumerate}
Espande le regioni dense e comprime quelle sparse; separa bene cluster ben separati, ma alcuni cluster visivi possono essere
\textbf{falsi} e dipendono dai parametri. Costo $O(n^2)$ in tempo e memoria. \textbf{Perplexity} bassa (5--30): struttura locale; alta (50+):
più globale, e più lenta.
\begin{center}
\small
\begin{tabular}{@{}lll@{}}
\toprule
 & \textbf{PCA} & \textbf{t-SNE}\\
\midrule
Tipo & lineare & non lineare\\
Preserva & varianza globale & similarità locali\\
Uso & riduzione del rumore, feature extraction & visualizzazione, clustering, AD\\
Costo / interpretazione & basso / facile & alto / difficile\\
\bottomrule
\end{tabular}
\end{center}

% =====================================================================
\section{Isolation Forest}
Slide 5, pp.\ 54--57. Algoritmo ad alberi non supervisionato che \textbf{isola le anomalie} invece di modellare la normalità, lavorando un
punto alla volta. Ipotesi: le anomalie sono \textbf{poche e diverse}, quindi bastano pochi tagli per separarle.
\begin{itemize}
  \item Costruisce molti alberi binari casuali: a ogni nodo sceglie a caso una feature e una soglia, fino a isolare il punto o raggiungere la
        profondità massima.
  \item Score = \textbf{lunghezza media del percorso} sugli alberi: percorso corto $\Rightarrow$ isolamento facile $\Rightarrow$ più anomalo.
  \item Nessuna distanza né stima di densità: funziona in alta dimensionalità, è robusto alle feature irrilevanti, scala linearmente;
        sensibilità ai parametri media (numero di alberi, dimensione del sottocampione).
  \item Limiti: i tagli casuali possono perdere strutture sottili, servono abbastanza alberi per score stabili, è meno interpretabile dei
        metodi a distanza.
\end{itemize}

% =====================================================================
\section{Deep learning per l'anomaly detection}
\begin{nota}[Non spiegato nelle registrazioni]
Questa sezione segue la slide 6 (pp.\ 2--43): la lezione del \rif{2026-11-09}{09/11} ne usa i concetti nel notebook (autoencoder, encoder
come estrattore di feature) senza spiegarli.
\end{nota}
Perché il deep: rappresentazioni non lineari e gerarchiche, dati complessi (contesto, anomalie collettive), architetture per ogni tipo di
dato (RNN, GNN, CNN, NLP) anche combinate. Limiti: costo (con le novità può servire riaddestrare) e interpretabilità (servono LIME,
Deep SHAP). Il deep learning serve a due cose: \textbf{imparare la normalità} (ricostruire bene ciò che è normale) e \textbf{imparare una
rappresentazione} (l'encoder come estrattore di feature per un altro algoritmo).

\subsection{Tre framework}
Slide 6, pp.\ 2--13 (dalla rassegna di Pang et al., 2021).
\begin{center}
\small
\begin{tabularx}{\linewidth}{@{}l X X@{}}
\toprule
\textbf{Framework} & \textbf{Idea} & \textbf{Pro / contro}\\
\midrule
Learning feature representation & rete come estrattore di feature indipendente ($Z=f(X;\phi)$, pre-addestrata o encoder di un AE), poi un
metodo tradizionale su $Z$ & semplice, modelli pronti, più espressiva della PCA / estrazione e scoring separati: si perde informazione
discriminante; i pre-addestrati non generalizzano a ogni dominio\\
Feature representation of normality & encoder + decoder addestrati insieme; score $s_x = f'(x, \phi^*, W^*)$, es.\ errore di ricostruzione &
riduzione di dimensione potente / limiti della misura di anomalia scelta\\
End-to-end anomaly score learning & la rete $\tau(x;\phi)$ produce direttamente lo score & più efficace / servono anomalie etichettate
o sintetiche; può non generalizzare ad anomalie mai viste\\
\bottomrule
\end{tabularx}
\end{center}
Nel framework \en{feature representation of normality} la slide 6 (pp.\ 8--12) riprende anche l'idea contrastiva del cap.~\ref{cap:embedding}:
ottimizzare lo spazio $Z$ in modo che le viste dello stesso oggetto (anchor e positivi) siano vicine e i negativi lontani; un oggetto nuovo
ma simile (un altro scoiattolo) resta distante dall'anchor, ``ma più vicino della tigre''.

\subsection{L'autoencoder}
Slide 6, pp.\ 14--27. Encoder $f(X;\phi)$ che porta l'input nel \textbf{bottleneck} $Z$ (spazio latente) e decoder $\psi(Z;W)$ che ricostruisce
$\hat x \approx x$. Domande di progetto: quale architettura per encoder e decoder (dipende dal dato: model engineering) e quanto grande $Z$.
\begin{itemize}
  \item \textbf{Undercomplete}: $Z$ più piccolo dell'input, comprime e tiene le feature salienti togliendo rumore; \textbf{overcomplete}: $Z$
        più grande, rappresentazione più ricca.
  \item La qualità della ricostruzione si misura con la loss: MSE se le feature sono numeriche (o categoriche con una distanza sensata, es.\ ID
        sequenziali); altrimenti MSE sulle numeriche + CE sulle categoriche (slide 6, p.~24):
        \[
        \mathcal{L} = \lVert x_{num} - \hat x_{num}\rVert^2 + \sum_{j} \mathrm{CE}\big(x^{cat}_j, \hat x^{cat}_j\big)
        \]
  \item \textbf{PCA vs AE}: la PCA ha soluzione unica ma è solo lineare; l'AE cattura pattern non lineari grazie alle attivazioni, ma
        può imparare l'identità.
  \item \textbf{Denoising AE}: ricostruisce il dato pulito da una versione corrotta $\Rightarrow$ rappresentazioni robuste.
        \textbf{Variational AE}: nello spazio latente ogni punto è una \textbf{distribuzione}; riduce l'overfitting e permette di generare
        dati simili.
\end{itemize}

\subsection{Casi di studio delle slide}
Slide 6, pp.\ 28--42.
\begin{itemize}
  \item \textbf{Modello pre-addestrato} (video): feature VGG su due finestre di frame consecutive, classificatore lineare addestrato
        iterativamente a distinguerle togliendo ogni volta le feature più discriminanti; l'accuratezza media è lo score (frame anomali più
        distinguibili). Funziona dove esistono buoni modelli pre-addestrati; il corso si concentra sul secondo scenario.
  \item \textbf{VAE} su KDD99 ($Z$ di 15, encoder e decoder FFNN a 2 layer) addestrato anche con anomalie: piccoli miglioramenti perché il
        dataset è semplice.
  \item \textbf{LSTM autoencoder} su InSDN (flussi SDN, 77 feature; encoder 128-64-32-16, $Z$ di 16) addestrato solo sui normali, poi
        one-class SVM sullo spazio latente e confronto con la soglia sull'errore di ricostruzione: la non linearità dell'OC-SVM aiuta.
  \item \textbf{AE e AE convoluzionale} su NSL-KDD, solo normali, score = MSE: meglio di k-NN e SVM, e molto più facili da riaddestrare.
\end{itemize}

% =====================================================================
\section{Validare l'anomaly detection}
Slide 6, pp.\ 49--52 (non spiegato nelle registrazioni).
\begin{itemize}
  \item \textbf{Con le label} (non usate in training): precision (falsi positivi costosi), recall (anomalie mancate critiche), F1 (entrambi),
        ROC-AUC (capacità di separare, ma ingannevole con dati sbilanciati); per gli score, metriche di \textbf{ranking} come l'NDCG (le anomalie
        vere sono in cima alla lista?).
  \item \textbf{Senza label}: analisi visuale (PCA, t-SNE: le anomalie sono isolate?), stabilità del ranking tra esecuzioni (es.\ Rank-Biased
        Overlap), valutazione di esperti di dominio, analisi della distribuzione dell'errore di ricostruzione (pochi punti con errore alto, da
        ispezionare a mano).
\end{itemize}

% =====================================================================
\section{Il notebook su KDD Cup 1999: dati e preparazione}
La lezione del \rif{2026-11-09}{09/11} è interamente dedicata al notebook che servirà per il Lab 3 (cap.~\ref{cap:lab3}). Il docente non
commenta ogni riga, solo ``cosa fare e come leggere i risultati'': il codice in sé è semplice. I numeri dei risultati sono quelli della
versione del notebook pubblicata sul portale, che in qualche punto differisce da quella mostrata in aula.

\begin{esame}
Il notebook è \textbf{solo dimostrativo}: nessun iperparametro è ottimizzato e non c'è una vera fase di validazione. Nel lab
\textbf{dovete} fare training, tuning sul validation set e solo alla fine valutazione sul test.
\end{esame}

\textbf{KDD Cup 1999}: ogni riga è una connessione (flusso) tra client e server, descritta da 41 feature generiche (protocollo,
servizio, flag, byte trasmessi, tassi di errore\ldots), \textbf{non} pensate per l'anomaly detection. La label indica il tipo di
attacco (o \code{normal}).

\subsection{Prima di tutto: studiare il dataset}
\begin{itemize}
  \item \textbf{Duplicati}: eliminandoli il dataset scende a circa un terzo. Gli esperimenti ripetuti dai ricercatori producevano
        gli stessi dati. Si tolgono per non sovrarappresentare singoli punti, a meno di farlo apposta e consapevolmente.
  \item \textbf{Valori mancanti}: nessuno.
  \item \textbf{Feature categoriche}: \code{protocol\_type} (soprattutto TCP, poi UDP, pochissimo ICMP: plausibile),
        \code{service} (più di 40 valori, ma solo una dozzina usati davvero), \code{flag} (alcuni valori rarissimi).
  \item \textbf{Label}: per lo più \code{normal}; poi \code{neptune} (SYN flood) molto frequente; molti attacchi rarissimi.
        Con conteggi così sbilanciati conviene l'\textbf{asse y logaritmico}, o semplicemente una tabella: rappresentate i risultati in
        modo che si veda ciò che volete mostrare.
\end{itemize}

\subsection{Le scelte di preprocessing}
\begin{enumerate}
  \item \textbf{Si toglie \code{neptune}} (scelta del notebook, ripetibile tenendolo): è un attacco frequentissimo che crea regioni
        dense, quindi dal punto di vista della densità non è un'anomalia. Così le anomalie restano rare.
  \item \textbf{Label binaria}: \code{normal} $\to 0$, qualsiasi attacco $\to 1$. Interessa trovare anomalie, non riconoscere la classe.
\end{enumerate}
\begin{esame}
Non cancellate la colonna con la label originale: serve dopo per le analisi (composizione dei cluster, quali attacchi sfuggono).
\end{esame}
\begin{enumerate}[resume]
  \item \textbf{Split} training/validation/test \textbf{stratificato} (circa il 6\% di anomalie in ogni set), \textbf{prima} di encoding
        e scaling: il preprocessing si impara solo dal training.
  \item \textbf{Ridurre le categorie}: i servizi e i flag rari confluiscono in \code{other}. Con il one-hot un vettore di 40 posizioni
        in cui 28 sono quasi sempre 0 aggiunge solo zeri che confondono il modello; con ``top + other'' il vettore si accorcia e si perde
        poco, perché categorie così rare contribuirebbero comunque poco. Alternative: embedding (cap.~\ref{cap:embedding}), o eliminare le
        righe (qui no: le righe contano). \emph{Attenzione}: la classe \code{other} diventa più eterogenea.
\end{enumerate}

\begin{lstlisting}
df = df.drop_duplicates()
df = df[df["label"] != "neptune"].copy()
df["binary_label"] = (df["label"] != "normal").astype(int)   # la colonna "label" resta

df_train_val, df_test = train_test_split(df, train_size=0.9, stratify=df["binary_label"])
df_train, df_val = train_test_split(df_train_val, train_size=0.9, stratify=df_train_val["binary_label"])

def reduce_categories(d):
    d = d.copy()
    d["service"] = d["service"].where(d["service"].isin(ALLOWED_SERVICES), "other")
    d["flag"] = d["flag"].where(d["flag"].isin(ALLOWED_FLAGS), "other")
    return d
\end{lstlisting}

\subsection{PCA in pratica: l'errore voluto (niente scaling)}
Il \code{ColumnTransformer} applica trasformazioni diverse a gruppi di colonne. Prima versione: one-hot sulle categoriche e
\code{passthrough} (nessuna trasformazione) sulle numeriche.

\begin{attenzione}[99,9\% della varianza in una componente]
La varianza spiegata cumulata dice che \textbf{una sola} componente spiega il 99,9\% della varianza, e lo scatter 2D mostra punti
allineati su assi enormi. O quella feature è importantissima, o c'è un errore: \textbf{mancava la normalizzazione}. Una feature in byte
(milioni) ha una varianza enormemente più grande di un RTT in decine di millisecondi; la PCA ``spiega'' la varianza di quella feature e
sembra spiegare tutto il dataset.
\end{attenzione}

Versione corretta: \code{StandardScaler} sulle numeriche, one-hot sulle categoriche, \code{fit} solo sul training.
\begin{lstlisting}
preprocessor = ColumnTransformer([
    ("num", StandardScaler(), NUMERICAL_COLS),
    ("cat", OneHotEncoder(handle_unknown="ignore", sparse_output=False), CATEGORICAL_COLS),
], remainder="drop")                          # le colonne di label vengono scartate
X_train = preprocessor.fit_transform(df_train)
X_val, X_test = preprocessor.transform(df_val), preprocessor.transform(df_test)
\end{lstlisting}
Ora la prima componente spiega circa il 10\% e ne servono una ventina: si sceglie il \textbf{ginocchio} della curva (20, non 30: il
guadagno non vale 10 feature in più, pensando alla maledizione della dimensionalità). In 2D i punti occupano tutto il piano.

\subsection{t-SNE in pratica}
Colorando i punti per label binaria o tipo di attacco si vede dove stanno.
\begin{itemize}
  \item Le coordinate \textbf{non hanno significato}: t-SNE preserva le similarità tra punti vicini ma deforma le distanze. Si può dire
        ``questo gruppo sembra ben separato'', non ``quanto'' è lontano.
  \item Attacchi raccolti in un gruppo compatto e separato: facili da trovare. Attacchi mescolati ai punti normali: difficili.
  \item \textbf{Perplexity}: nel lab provatene un intervallo ampio (es.\ da 5 a 50 e oltre), non 1, 2, 3.
\end{itemize}

\subsection{Contaminazione e schema comune}
La \textbf{contaminazione} è la frazione attesa di anomalie: quasi tutti i metodi ne hanno bisogno per trasformare uno score in una
decisione. Qui si usa il valore esatto (dalle label, $\approx 6\%$); \textbf{nella realtà va stimata} da analisi preliminari.

Ogni metodo del notebook segue lo stesso schema:
\begin{center}
\small
\begin{tabularx}{\textwidth}{l X}
\toprule
\textbf{Passo} & \textbf{Cosa si fa}\\
\midrule
Fit & addestramento sul training (le label servono solo a scegliere i dati normali nel caso semi-supervisionato)\\
Score & uno score di anomalia: più alto = più anomalo (spesso va invertito il segno dell'output di scikit-learn)\\
Soglia & dal modello, dalla contaminazione o da un percentile degli score di training\\
Valutazione & solo con le label: classification report, matrice di confusione, ROC-AUC, ECDF degli score, previsioni su t-SNE\\
\bottomrule
\end{tabularx}
\end{center}

\begin{esame}[confronto equo]
Nel notebook i metodi si confrontano sul \textbf{test set} solo perché nessuno è stato ottimizzato. Gli autoencoder però usano il
validation set per monitorare il training: confrontare tutti sul validation li avvantaggerebbe. Per un confronto equo si usa un set
che \textbf{nessun} algoritmo ha visto. Procedura corretta: tuning di ogni metodo sul validation, scelta del modello sul validation,
\textbf{una sola} valutazione finale sul test.
\end{esame}

\begin{attenzione}[la matrice di confusione non vuol dire classificazione]
I metodi sono non supervisionati: la matrice di confusione si usa solo per leggere più facilmente i risultati, mettendo l'output del
metodo (``fuori dal confine'', ``rumore'', ``lontano dal centroide'') contro la ground truth.
\end{attenzione}

% =====================================================================
\section{Il notebook: metodi shallow}
\subsection{One-class SVM addestrata solo su dati normali (``supervisionata'' tra virgolette)}
Il modello non vede le label: le usiamo solo per scegliere i dati normali. Kernel RBF, \code{gamma='scale'}, $\nu$ molto piccolo
(0,001) ma non zero: anche le label possono essere sbagliate, quindi qualche punto ``normale'' rumoroso va tollerato.
\begin{lstlisting}
ocsvm = OneClassSVM(kernel="rbf", gamma="scale", nu=0.001)
ocsvm.fit(X_train_normal_only)
scores_test = -ocsvm.decision_function(X_test)            # distanza dal confine: > 0 = fuori
y_pred_test = (ocsvm.predict(X_test) == -1).astype(int)   # -1 = fuori dal confine = anomalia
\end{lstlisting}
\begin{itemize}
  \item Sulla classe normale le metriche sono buone, ma è la maggioritaria: conta poco. Sulle anomalie: precision buona ($\approx 0{,}88$),
        recall migliorabile (circa un terzo delle anomalie sfugge); ROC-AUC $\approx 0{,}88$.
  \item La distribuzione della \textbf{distanza dal confine}: i normali sono sotto 0 (dentro), la maggior parte delle anomalie sopra.
        Quelle dentro il confine sono gruppi interi (es.\ \code{back}, \code{warezclient}) che somigliano a normale traffico HTTP/FTP.
\end{itemize}

\subsection{One-class SVM addestrata su tutti i dati (non supervisionata)}
Se non so chi è normale, addestro su tutto il training e imposto $\nu$ = contaminazione. ROC-AUC quasi uguale (in aula da $\sim0{,}87$
a $\sim0{,}86$), ma precision e recall delle anomalie crollano (F1 da $\approx 0{,}74$ a $\approx 0{,}48$): alcune anomalie sono vicine ai
normali o formano gruppi abbastanza densi da essere inclusi nel confine, mentre $\nu$ costringe a lasciare fuori punti rari ma normali.
\textbf{Il non supervisionato è più difficile}. Nel lab: provate a variare la quantità di anomalie nel training.

\subsection{LOF}
\begin{lstlisting}
lof = LocalOutlierFactor(n_neighbors=10, contamination=CONTAMINATION, novelty=True)
lof.fit(X_train)
scores_test = -lof.score_samples(X_test)       # scikit-learn restituisce il LOF negato
\end{lstlisting}
Parametro principale: il numero di vicini con cui confrontare la densità. La contaminazione fissa la soglia (in pratica il percentile);
in alternativa si guarda la curva ordinata dei LOF score e si sposta la soglia a mano: più alta = meno anomalie trovate ma più normali
tenuti. Quale scelta è migliore dipende dalle prestazioni finali e dall'uso.

\textbf{Risultato: pessimo.} In aula trova 36 anomalie su più di 600; nel notebook la ROC-AUC è $\approx 0{,}45$, peggio del caso.
Molti attacchi (\code{teardrop}, \code{smurf}, \code{back}, \code{satan}) sono raffiche di connessioni quasi identiche: gruppi
\textbf{piccoli ma densi}. Con $k=10$ un loro punto ha vicini densi quanto lui, quindi LOF $\approx 1$; intanto le connessioni normali
isolate prendono score alti.
\begin{attenzione}[un attacco non è per forza un outlier di densità]
Un attacco etichettato non è per forza un outlier di densità. Il dataset sembra adatto ad anomalie \textbf{globali}, non locali: prima
di scartare il metodo verificate se è sbagliata la metodologia o solo la configurazione (es.\ \code{n\_neighbors} più grande dei gruppi
di attacchi).
\end{attenzione}

\subsection{K-means: distanza dal centroide}
\begin{enumerate}
  \item \textbf{Scegliere $K$} con l'SSE per $K$ da 2 a 20. Nel grafico in aula l'SSE risale: probabile cluster vuoto. Andava fatta almeno
        un'altra esecuzione: \textbf{3--10 run per ogni $K$}. Il gomito non è netto; si sceglie $K=10$ anche per costo: con 50 cluster
        bisognerebbe poi validarli e interpretarli tutti.
  \item \textbf{Indici interni} (solo le feature): silhouette. Silhouette media bassa ($\approx 0{,}23$): struttura debole; alcuni
        cluster hanno punti con silhouette negativa (starebbero meglio in un altro cluster).
  \item \textbf{Indici esterni} (label originali, solo per interpretare): tabella cluster $\times$ label, normalizzata per riga. I cluster
        grandi sono per oltre il 90\% normali; alcuni cluster piccoli sono quasi solo attacchi (es.\ un cluster è 100\% \code{teardrop}:
        circa 750 punti, tutti lì e solo lì). Serve guardare insieme purezza e \textbf{dimensione}: un cluster ``puro'' di 5 punti dice poco.
  \item \textbf{Score} = distanza dal centroide; soglia = percentile del training pari alla contaminazione.
\end{enumerate}
\begin{lstlisting}
_, dist_train = pairwise_distances_argmin_min(X_train, kmeans.cluster_centers_)
_, dist_test = pairwise_distances_argmin_min(X_test, kmeans.cluster_centers_)
threshold = np.percentile(dist_train, 100 * (1 - CONTAMINATION))
y_pred_test = (dist_test > threshold).astype(int)
\end{lstlisting}
Risultato: meglio di LOF (ROC-AUC $\approx 0{,}88$), ma F1 delle anomalie $\approx 0{,}47$. I normali tendono a stare vicini al centroide,
le anomalie più lontane; però gli attacchi che formano un \textbf{cluster tutto loro} (come \code{teardrop}) sono vicinissimi al loro
centroide e sfuggono. Usare anche la dimensione del cluster come segnale li catturerebbe; e un gruppo così denso e separato non lo
troverebbe nemmeno un metodo basato sulla densità locale.

\subsection{DBSCAN: i punti di rumore}
DBSCAN restituisce già $-1$ per il rumore. Con \code{min\_samples} = 4 si sceglie $\varepsilon$ con il \textbf{k-distance plot}: distanza
di ogni punto dal suo $k$-esimo vicino, ordinata; si cerca il ginocchio. In aula la curva non convince: proverei un numero di vicini un
po' più alto (6, 9, 10; non 100).
\begin{lstlisting}
dbscan = DBSCAN(eps=0.15, min_samples=4).fit(X_train)
# DBSCAN non ha predict: un nuovo punto e' normale se dista <= eps da un core point
nn_core = NearestNeighbors(n_neighbors=1).fit(dbscan.components_)
scores_test = nn_core.kneighbors(X_test)[0][:, 0]
y_pred_test = (scores_test > 0.15).astype(int)
\end{lstlisting}
Risultato: \textbf{pessimo}, precision delle anomalie molto bassa: molti punti normali sono ``rumore'', cioè non hanno altri punti vicini.
Nel notebook, con $\varepsilon=0{,}15$, centinaia di cluster e circa il 30\% del training come rumore (F1 $\approx 0{,}15$). La curva senza
ginocchio è la maledizione della dimensionalità in pratica: in 50+ dimensioni le distanze crescono in modo uniforme.

\subsection{Isolation Forest}
OCSVM con kernel RBF, LOF, k-means e DBSCAN si basano tutti sulle \textbf{distanze}, e con 50+ feature (nessuna riduzione di dimensione)
ne soffrono. L'Isolation Forest no.
\begin{lstlisting}
iforest = IsolationForest(n_estimators=200, contamination=CONTAMINATION)
iforest.fit(X_train)
scores_test = -iforest.score_samples(X_test)   # piu' alto = piu' anomalo
\end{lstlisting}
\begin{itemize}
  \item \code{n\_estimators}: né troppo pochi alberi né mille (overfitting o tempi lunghi, senza bisogno).
  \item È velocissimo. La distribuzione del numero di tagli necessari per isolare un punto: ai normali ne servono di solito tanti, alle
        anomalie a volte pochi, a volte molti.
  \item Meglio di DBSCAN; nel notebook è il migliore tra i metodi shallow interamente non supervisionati (ROC-AUC $\approx 0{,}94$,
        F1 $\approx 0{,}52$). Il punto debole è la soglia fissa al 6\%.
\end{itemize}
\begin{attenzione}[la ROC-AUC può ingannare]
Una ROC-AUC altissima può dipendere dai pochi falsi negativi su un dataset fortemente sbilanciato, come l'accuracy. Guardate anche le
componenti: precision, recall e F1 della classe delle anomalie.
\end{attenzione}

% =====================================================================
\section{Il notebook: autoencoder}
Non esiste un modello pre-addestrato adatto: si addestra da zero. Prima: tensori e \code{DataLoader}.

\subsection{La loss di ricostruzione mista}
L'input mescola feature numeriche (scalate) e categoriche (one-hot), quindi la loss è mista: \textbf{errore quadratico} sulle numeriche,
\textbf{cross-entropy} su ogni gruppo one-hot, e la somma è l'errore di ricostruzione. Adattatela alle feature del vostro problema.
\begin{lstlisting}
def per_sample_reconstruction_error(x, x_recon, categorical_dims, numerical_idx):
    mse = torch.mean((x[:, numerical_idx] - x_recon[:, numerical_idx]) ** 2, dim=1)
    ce = torch.zeros(x.shape[0])
    for start, end in categorical_dims:                  # un blocco per feature categorica
        target = torch.argmax(x[:, start:end], dim=1)    # la CE vuole l'indice della classe
        ce += F.cross_entropy(x_recon[:, start:end], target, reduction="none")
    return mse + ce                                      # anche lo score di anomalia
\end{lstlisting}
Il modello semplice: encoder \code{input $\to$ 64 $\to$ 8} e decoder simmetrico, ReLU, Adam, 30 epoch, batch 64.

\subsection{Addestrare solo sui normali o su tutti?}
\begin{itemize}
  \item \textbf{Solo normali}: la loss di training scende; sul validation l'errore dei normali è piccolo, quello delle anomalie
        \textbf{cresce} con le epoch. Il modello si specializza sul traffico normale: questa separazione (nel notebook circa 30 volte) è ciò
        che rende l'errore di ricostruzione un buon score.
  \item \textbf{Normali + anomalie}: errore di training un po' più alto; sul validation l'errore delle anomalie \textbf{scende}, perché il
        modello impara a ricostruire anche parte di esse, e le due curve si avvicinano: lo spazio si ``comprime''.
\end{itemize}

\subsection{Score, soglia e risultati}
Si confronta l'\textbf{ECDF} dell'errore di ricostruzione di normali e anomalie su training, validation e test. Le curve dei tre set sono
quasi uguali perché vengono dalla stessa distribuzione; il test qui si guarda solo per mostrarlo, non per decidere. Con il modello ``solo
normali'' i normali hanno errore basso e la maggior parte delle anomalie (rosse) alto. La soglia si fissa su training o validation
(nel notebook: percentile del training pari alla contaminazione).
\begin{lstlisting}
scores_train = reconstruction_scores(ae_normal, X_train_t)
threshold = np.percentile(scores_train, 100 * (1 - CONTAMINATION))
y_pred_test = (reconstruction_scores(ae_normal, X_test_t) > threshold).astype(int)
\end{lstlisting}
\begin{itemize}
  \item \textbf{AE solo normali}: pochi normali classificati male, numero di anomalie trovate ragionevole. Nel notebook è il miglior metodo
        (ROC-AUC $\approx 0{,}99$, F1 $\approx 0{,}75$); circa un quarto delle anomalie ha però un errore basso come il traffico normale e resta
        invisibile con qualunque soglia.
  \item \textbf{AE normali + anomalie}: circa il 60\% delle anomalie di training finisce sotto la soglia; più falsi positivi e meno anomalie
        trovate (F1 $\approx 0{,}41$). Messaggio: \textbf{inserire anomalie nel training peggiora} le prestazioni; ciò che il modello vede in
        training conta più del modello stesso.
  \item Un AE più profondo con batch norm e dropout (cap.~\ref{cap:regolarizzazione}; lo si salta in aula) non migliora: il dropout
        regolarizza e il modello ricostruisce meglio anche le anomalie. Per l'anomaly detection un modello che generalizza ``troppo'' non è
        ciò che si vuole.
\end{itemize}

\subsection{L'encoder come estrattore di feature}
È il framework \en{learning feature representation}: si usa solo la prima metà del modello (\code{encode}): dai dati si ottiene la
rappresentazione nello spazio latente (8 dimensioni), che si visualizza con t-SNE. Poi si ripete la one-class SVM ``solo normali'' (stesso
$\nu$, stesso kernel) sui vettori latenti.
\begin{lstlisting}
def encode(model, X_t):
    model.eval()
    with torch.no_grad():
        return model.encode(X_t).cpu().numpy()

Z_train_normal = encode(ae_normal, X_train_normal_t)
ocsvm_latent = OneClassSVM(kernel="rbf", gamma="scale", nu=0.001).fit(Z_train_normal)
\end{lstlisting}
Risultato: buona precision, recall scarsa: la trasformazione \textbf{non aiuta} (nel notebook recall da $\approx 0{,}64$ a $\approx 0{,}47$).
L'encoder è stato addestrato a ricostruire il normale, non a separare le anomalie: comprimendo un'anomalia la avvicina ai normali che le
somigliano, e l'informazione che la rende anomala è proprio quella scartata dal bottleneck, cioè quella che finisce nell'errore di
ricostruzione. Nel paper visto nelle slide (LSTM autoencoder su InSDN) la combinazione funzionava bene; qui no.

% =====================================================================
\section{Confronto finale}
Tabella finale con tutti i metodi su validation e test. Secondo la ROC-AUC sul validation vince l'AE addestrato solo sui normali; ma la
ROC-AUC può ingannare, quindi si guarda anche l'\textbf{F1 della classe anomala}: alto per l'AE e per la one-class SVM solo normali. Il
docente si affida più a precision, recall e F1 che alla curva ROC.

\begin{center}
\small
\begin{tabular}{@{}lcc@{}}
\toprule
\textbf{Metodo (test, versione pubblicata del notebook)} & \textbf{ROC-AUC} & \textbf{F1 anomalie}\\
\midrule
AE solo normali & 0,99 & 0,75\\
AE migliorato (BN + dropout), solo normali & 0,97 & 0,72\\
OC-SVM solo normali & 0,88 & 0,74\\
Isolation Forest & 0,94 & 0,52\\
OC-SVM su tutti i dati & --- & 0,48\\
K-means (distanza dal centroide) & 0,88 & 0,47\\
AE normali + anomalie & 0,88 & 0,41\\
OC-SVM nello spazio latente & 0,83 & ---\\
DBSCAN & 0,65 & 0,15\\
LOF & 0,45 & ---\\
\bottomrule
\end{tabular}
\end{center}
Valori indicativi. Nessun metodo è stato ottimizzato: \textbf{non è un confronto equo} (LOF e DBSCAN sono molto sensibili ai parametri).

\begin{sintesi}
\begin{itemize}
  \item Anomaly detection = pattern che deviano dall'atteso; novelty = nuovo, magari benigno. Complessità: unknowingness, rarità e
        sbilanciamento, eterogeneità, evoluzione. Anomalie \textbf{point}, \textbf{contextual}, \textbf{collective}.
  \item In cybersecurity le label mancano: si preferisce l'\textbf{unsupervised}; le label, se ci sono, servono \textbf{solo a valutare}.
        Output: score (ranking, budget) o label (numero controllato da iperparametri come $\nu$ o la contaminazione).
  \item \textbf{One-class SVM}: confine attorno al normale, $\nu$ = frazione di outlier; meglio addestrata solo sui normali.
  \item \textbf{K-means}: SSE e gomito, più run per $K$; score = distanza dal centroide, più cluster piccoli. \textbf{DBSCAN}: core,
        border, noise; $\varepsilon$ dal k-distance plot. \textbf{LOF}: densità relativa ai vicini, anomalie locali.
  \item Tutti questi metodi usano distanze $\Rightarrow$ \textbf{maledizione della dimensionalità}; scalare sempre prima di PCA e metodi a
        distanza. \textbf{Isolation Forest}: percorso corto = anomalo, niente distanze.
  \item Deep: tre framework (feature representation, representation of normality, end-to-end score); \textbf{autoencoder} con loss
        MSE + CE, score = errore di ricostruzione, addestrato su dati normali puliti.
  \item Nel lab: studiare il dataset, split stratificato prima del preprocessing, tenere la label originale, intervalli ampi di
        parametri (perplexity, $\varepsilon$, \code{n\_neighbors}, più run di k-means), tuning sul validation e \textbf{una sola} valutazione
        sul test; non fidarsi della sola ROC-AUC.
\end{itemize}
\end{sintesi}


% =====================================================================
\part{Laboratori}
% Un capitolo per laboratorio: caso, consegna, metodo, insidie (niente soluzioni dei gruppi).
% =====================================================================
\chapter{Come si scrive un report di laboratorio}
\label{cap:lab-report}

\begin{nota}[Fonti]
Slide \emph{0-Intro-Course}, linee guida date in aula alla \rif{2026-10-20}{20/10}, sezioni
``Lab Rules'', ``How to work'' e ``Self-assessment'' dei testi del Lab~1 e del Lab~3. Regole di
consegna, scadenze e voto sono nel cap.~\ref{cap:formato}.
\end{nota}

\section{I quattro laboratori}
\begin{center}
\small
\begin{tabularx}{\textwidth}{c l X l}
\toprule
\textbf{Lab} & \textbf{Durata} & \textbf{Contenuto} & \textbf{Capitolo}\\
\midrule
1 & 3 h & Introduzione a PyTorch con una FFNN per intrusion detection; solo formativo, consegna
          obbligatoria, +1 se ``ragionevole'' entro il 1° dicembre. & cap.~\ref{cap:lab1}\\
2 & 6 h & Classificazione supervisionata malware/goodware da sequenze di API call (la slide
          parla di \en{supervised traffic classification}). & cap.~\ref{cap:lab2}\\
3 & 6 h & Anomaly detection supervisionata e non supervisionata su log di un IDS. & cap.~\ref{cap:lab3}\\
4 & 6 h & NLP e clustering per il rilevamento di botnet (non ancora assegnato; secondo il docente il materiale
          collegato a DarkVec ne anticipa il contenuto, cap.~\ref{cap:darkvec}). & ---\\
\bottomrule
\end{tabularx}
\end{center}

\section{Cosa si consegna}
Per ogni lab un file \code{groupID\_labX.zip} che contiene:
\begin{itemize}
  \item il \textbf{report} (\code{groupID\_labX.pdf}), massimo \textbf{10 pagine}, con approccio,
        decisioni, esperimenti e risultati, tabelle e grafici;
  \item l'\textbf{experiment logbook}, come appendice del report o come foglio separato:
        \textbf{obbligatorio}, non conta nelle 10 pagine;
  \item il \textbf{notebook} e il codice: eseguito (output visibili), rieseguibile (seed fissati),
        organizzato con titoli Markdown e commenti, in formato \code{.ipynb} più export HTML.
\end{itemize}
Si valuta il \textbf{report}; il notebook serve in caso di dubbi. Il report deve quindi stare in
piedi da solo.

\section{Il logbook e il budget di run}
Dai testi dei lab di quest'anno:
\begin{itemize}
  \item \textbf{Un solo logbook} per tutto il lab: una tabella con \textbf{una riga per ogni run che
        addestra un modello} (nel Lab~3 anche ogni fit di un detector o di un clustering). Le colonne
        suggerite sono: numero del run, task, cosa ho cambiato (\textbf{una sola cosa}), ipotesi
        \textbf{prima} di eseguire, risultato chiave sulla validation, conclusione.
  \item Le decisioni di preprocessing e di analisi dei dati vanno nel testo del report, non nel logbook.
        Nel Lab~3 uno sweep diagnostico per scegliere un parametro (gomito dell'SSE, k-distance plot,
        varianza spiegata) conta come un solo run.
  \item \textbf{Budget} di circa \textbf{25 run} per lab (circa 30 nel Lab~2, con modelli più pesanti): prima di ogni run il logbook deve dire perché
        vale uno degli slot rimasti. Non è un limite hardware: è la competenza che si vuole allenare
        (``ho fatto una grid su tutti i parametri'' non è una buona risposta).
  \item \textbf{Predict-then-check}: scrivere l'ipotesi prima, poi confrontarla con il risultato.
        Un'ipotesi sbagliata ma spiegata vale più di una giusta ma vaga; non si riscrive la storia.
  \item \textbf{Una modifica per run}: se ne cambi due insieme non sai quale ha causato l'effetto.
\end{itemize}

\section{Indicazioni del docente}
Alla \rif{2026-10-20}{20/10} il docente ha presentato in aula il documento con le linee guida per i
report: struttura della consegna (report di massimo 10 pagine, notebook \textbf{eseguito} e
commentato, \textbf{experiment logbook} fuori dal conteggio delle pagine), figure leggibili e
valutazione da 0 a 32 (cap.~\ref{cap:formato}). I punti principali:
\begin{esame}[come scrivere un buon report]
\begin{itemize}
  \item Rileggilo come se fosse di un altro: lo valuteresti positivamente?
  \item \textbf{Figure leggibili} (assi, legende, font): le immagini comunicano più del testo.
        Curve di apprendimento instabili o illeggibili fanno pensare che la soluzione non sia stata valutata.
  \item Non scrivere ``è meglio'' senza dire \textbf{rispetto a cosa} e con quale metrica
        (accuracy? F1?). Usa \textbf{tabelle} per confrontare architetture e iperparametri.
  \item Inserisci solo le figure e tabelle che servono a sostenere le tue affermazioni.
  \item 10 pagine sono un massimo, non un obiettivo; non ridurre il font per farcene stare di più.
  \item Si valuta la \textbf{metodologia}, non la performance migliore: una differenza di un centesimo
        nella metrica può dipendere dal seed.
  \item Il contenuto è responsabilità vostra: errori evidenti vengono penalizzati.
\end{itemize}
\end{esame}

\begin{esame}[criteri di autovalutazione dei testi dei lab]
Le rubriche di Lab~1 e Lab~3 premiano soprattutto il ragionamento: ``codice che funziona senza
ragionamento non dimostra la competenza; una spiegazione solida, anche con risultati non ottimali,
vale di più''. I criteri ricorrenti sono:
\begin{itemize}
  \item \textbf{disciplina del logbook}: ipotesi prima, una modifica alla volta, budget giustificato;
  \item \textbf{decisioni giustificate} con evidenze dai propri esperimenti (preprocessing, split,
        parametri, soglie);
  \item \textbf{interpretazione legata ai propri risultati} (curve, report, tabelle), non spiegazioni
        generiche da manuale;
  \item \textbf{giudizio e semplicità}: confrontare la complessità aggiunta con una baseline e dire
        onestamente se ne valeva la pena;
  \item \textbf{chiarezza e riproducibilità}: limite di pagine, notebook che gira, seed;
  \item nel Lab~3 anche: dichiarare cosa ogni metodo assume di sapere sui dati e usare le label
        solo per la valutazione.
\end{itemize}
Le performance non vengono confrontate tra gruppi.
\end{esame}

\section{Checklist prima di consegnare}
\begin{itemize}
  \item[$\square$] Intestazione con \textbf{ID del lab, ID del gruppo, nomi ed email}.
  \item[$\square$] Massimo \textbf{10 pagine}, font normale; logbook a parte (non conta nelle pagine).
  \item[$\square$] Logbook completo: ipotesi scritte prima, una modifica per run, budget rispettato e motivato.
  \item[$\square$] Ogni affermazione ``è meglio'' dice \textbf{rispetto a cosa} e \textbf{con quale metrica}.
  \item[$\square$] Metriche adatte allo sbilanciamento (per classe, F1, macro), non solo accuracy.
  \item[$\square$] \textbf{Learning curve} di training e validation leggibili, con plateau (o spiegazione del perché no).
  \item[$\square$] Tabelle di confronto tra architetture/iperparametri.
  \item[$\square$] Preprocessing giustificato: cosa è stato fittato \textbf{solo sul training} (scaler, vocabolario, encoder).
  \item[$\square$] Nessun leakage: split prima del preprocessing; split temporale se i dati hanno ordine; niente ID tra le feature.
  \item[$\square$] Test set usato solo dove il testo lo chiede, dopo aver scelto sul validation.
  \item[$\square$] Risposte a \textbf{tutte} le domande del testo, con commento ai risultati.
  \item[$\square$] Notebook \textbf{eseguito} (output visibili), commentato, con seed; \code{.ipynb} + HTML.
  \item[$\square$] Rilettura ``da esterni'': lo valuterei positivamente?
\end{itemize}

\begin{nota}[Perché qui non ci sono soluzioni]
I lab sono valutati e questo manuale è condiviso con altri gruppi. Nei capitoli dei singoli lab
trovate contesto, consegna parafrasata, teoria collegata, metodo generale e insidie tipiche, ma
\textbf{nessuna soluzione, codice o risultato}.
\end{nota}

\chapter{Lab 1 -- Introduzione a PyTorch con una FFNN}
\label{cap:lab1}

\begin{nota}[Fonti e limiti di questo capitolo]
Testo del laboratorio ``\emph{1 -- Introduction to Deep Learning: building, training, and judging a
Feed-Forward Neural Network}'' pubblicato sul portale. La consegna è \textbf{parafrasata}; per i
dettagli fa fede il testo originale. Il Lab~1 non dà voto ma è obbligatorio e vale il bonus: come
per gli altri lab qui \textbf{non} ci sono soluzioni, codice, numeri o scelte di un gruppo.
\end{nota}

\section{Contesto e obiettivi}
Il lab ha due scopi: introdurre PyTorch e il comportamento delle reti neurali su un vero problema
di classificazione, e dare una prima esperienza (non valutata) di lab e report come allenamento per
quelli valutati. Il testo insiste che è un lab di \textbf{giudizio, non di ricette}: si riceve un
obiettivo, dei vincoli e un punto di partenza consigliato, e si devono decidere, giustificare e
difendere le scelte sulla base di ciò che si osserva. In concreto si vuole:
\begin{itemize}
  \item costruire e addestrare FFNN per \en{intrusion detection} con PyTorch;
  \item vedere come il training si rompe quando un fondamentale (scaling, learning rate) è sbagliato
        e imparare a leggerne i sintomi;
  \item gestire lo sbilanciamento e ragionare su quale metrica conta per un IDS;
  \item vedere come una singola feature può guidare di nascosto il modello e come una scelta di
        pulizia può rimodellare le classi;
  \item costruire una baseline semplice e verificare se la complessità aggiuntiva ne vale la pena.
\end{itemize}

\section{Dataset}
Un sottoinsieme di \textbf{CICIDS2017} (Canadian Institute for Cybersecurity, University of New
Brunswick), benchmark di intrusion detection con traffico benigno realistico e attacchi comuni. Un
unico CSV di flussi etichettati; ogni riga è un flusso descritto da poche feature a livello di flusso
(durate e tempi di interarrivo, statistiche sulle lunghezze dei pacchetti, rate di byte e pacchetti,
alcuni contatori di flag TCP, numero di pacchetti in avanti, porta di destinazione, \ldots). La colonna
\code{Label} ha quattro valori: \code{Benign}, \code{PortScan}, \code{DoS Hulk}, \code{Brute Force}.
Il testo \textbf{non} dice di proposito quanto sono bilanciate le classi né quali problemi abbiano i
dati: scoprirlo è il primo task.

\section{Consegna (parafrasi)}
Regole di lavoro comuni a tutto il lab: logbook con una riga per run, budget di circa 25 run,
ipotesi prima di ogni run, una modifica alla volta (cap.~\ref{cap:lab-report}).

\begin{description}[style=nextline,leftmargin=1.2em]
\item[Task 1 -- Preprocessing (nessun run)]
  Ispezionare i dati (numero di campioni, classi e bilanciamento, problemi di qualità come valori
  mancanti, duplicati, valori implausibili); decidere cosa fare di mancanti e duplicati e dire quanti
  campioni restano; scegliere le proporzioni dello split train/validation/test e se stratificare;
  scegliere e motivare lo scaling dalle distribuzioni osservate (la porta di destinazione, per ora,
  va trattata come numerica). Domande: perché stratificare o no, su quale partizione si fitta lo
  scaler e cosa va storto se lo si fitta su tutto, come gli outlier guidano la scelta.
\item[Task 2A -- Un primo classificatore]
  Una FFNN a \textbf{un solo hidden layer} per le 4 classi, con un training loop pulito (loss,
  ottimizzatore, early stopping sulla validation). Il testo suggerisce un punto di partenza semplice
  (un hidden layer piccolo con ReLU, cross-entropy, AdamW, learning rate e batch size moderati, early
  stopping) e chiede di dichiarare ogni scelta. Prima di addestrare: prevedere quali classi saranno
  più difficili e perché. Poi curve di loss e classification report di validation; come si sceglie il
  modello ``migliore'' tra le epoch, se le curve convergono, se le previsioni reggono.
\item[Task 2B -- Perché i fondamentali contano]
  Tre varianti della configurazione 2A, una sola cosa sbagliata alla volta: S1 senza scaling, S2 con
  learning rate troppo alto (fino a rompere il training), S3 con learning rate molto più basso. Per
  ognuna: prevedere la curva, plottare train/validation loss, dare un nome al \en{failure mode} e
  spiegarne il meccanismo. Poi ordinarle per gravità, dire quale fondamentale sistemare per primo su
  un problema nuovo e cosa, nella sola forma di una curva, indicherebbe la causa.
\item[Task 2C -- Fissare la baseline]
  Scegliere un modello a un hidden layer ben tarato e compreso: è la ``baseline da battere''.
  Classification report di validation e test, commento sul gap. Come esperimento controllato, cambiare
  l'attivazione (es.\ lineare $\leftrightarrow$ ReLU) e discutere se sia lecito confrontare le due
  varianti sul test set.
\item[Task 3 -- Una feature che distorce il modello]
  Nei dati, tutti i flussi \code{Brute Force} hanno porta di destinazione 80. Domanda: è una proprietà
  degli attacchi brute force o un artefatto del dataset? Esperimento controfattuale \textbf{solo sul
  test}: cambiare la porta di quei flussi, prevedere l'effetto sulla recall e verificarlo con la
  baseline senza riaddestrare. Poi rimuovere la porta dai dati originali, rifare il preprocessing e
  usare il dataset senza porta per il resto del lab. Domande sul conteggio delle classi dopo la
  pulizia, su quale classe si riduce e perché, su quale diventa la più rara e sul dilemma: de-duplicare
  ancora o no?
\item[Task 4 -- La classe rara: loss e sbilanciamento]
  Riaddestrare la baseline sui dati senza porta, prevedere quale classe soffre di più, poi usare una
  cross-entropy \textbf{pesata per classe}. Domande: da quale partizione si calcolano i pesi e perché
  ogni altra scelta è sbagliata; accuracy e macro-F1 si muovono nella stessa direzione? Quale metrica
  ottimizzare per un IDS che deve trovare attacchi rari?
\item[Task 5 -- Iperparametri (circa 10 run in tutto)]
  5A profondità (2--3 candidati più profondi scelti e motivati, niente grid); 5B batch size (tre
  valori che ``circondano'' un'ipotesi, effetto su generalizzazione e tempo); 5C ottimizzatore contro
  learning rate (SGD, SGD+momentum, AdamW variando il learning rate: chi domina?); 5D verdetto sul test
  del miglior modello affiancato alla baseline: la complessità ne valeva i run?
\item[Task 6 -- Ottimizzazioni]
  Una rete grande e profonda imposta dal testo, addestrata a lungo senza regolarizzazione né early
  stopping: prevedere le curve e mostrare dove si vede l'overfitting. Poi, una modifica per run e dalla
  più semplice: rete più piccola o stop anticipato, poi un regolarizzatore alla volta (weight decay,
  dropout, batch normalization), controllando la validation prima di aggiungere il successivo.
  Domande: cosa ha aiutato di più con meno sforzo; la rete grande regolarizzata batte la baseline
  semplice sul test?
\item[Riflessione finale]
  Guardando tutto il logbook: quali decisioni hanno dato più prestazioni per unità di sforzo?
\end{description}

\section{Teoria coinvolta}
\begin{itemize}
  \item Pulizia, duplicati, valori mancanti, split stratificato, scaling e leakage; struttura di una
        FFNN e training loop in PyTorch: cap.~\ref{cap:ffnn}.
  \item Loss, cross-entropy, ottimizzatori (SGD, momentum, Adam/AdamW), learning rate, batch size,
        metriche per classe, learning curve, under/overfitting: cap.~\ref{cap:addestramento}.
  \item Attivazioni e collasso lineare di una rete senza non linearità: cap.~\ref{cap:reti-neurali}.
  \item Early stopping, weight decay, dropout, batch normalization: cap.~\ref{cap:regolarizzazione}.
  \item Feature-scorciatoia (\en{shortcut}) e porte come feature categoriche: cap.~\ref{cap:ffnn}
        e~\ref{cap:embedding}.
\end{itemize}

\section{Metodo generale}
\begin{enumerate}
  \item \textbf{Prima i dati}: contare, guardare distribuzioni e code, trovare duplicati e anomalie di
        qualità; ogni passo di pulizia va motivato e il numero di righe prima/dopo va riportato.
  \item \textbf{Split prima di tutto il resto}; fit di scaler e pesi di classe sul solo training.
  \item \textbf{Baseline semplice} che si capisce fino in fondo, con curve e report per classe.
  \item \textbf{Una variabile alla volta}, con ipotesi scritta prima: ogni esperimento risponde a una
        domanda precisa del testo.
  \item \textbf{Validation per scegliere, test per verificare}: il test non entra mai nelle scelte.
  \item \textbf{Confronto onesto} con la baseline: la complessità si giustifica solo con un
        miglioramento che regge e che vale il costo.
\end{enumerate}

\section{Insidie tipiche}
\begin{attenzione}[errori da evitare]
\begin{itemize}
  \item Fittare lo scaler (o calcolare i pesi di classe) su tutto il dataset o sul validation/test.
  \item Giudicare dalla sola accuracy un problema sbilanciato: guardare report per classe e macro-F1.
  \item Scegliere tra varianti (attivazione, architettura) guardando il test set.
  \item Riportare il modello dell'ultima epoch invece di quello scelto dall'early stopping.
  \item Cambiare due cose nello stesso run, o fare grid search ``a tappeto'' bruciando il budget.
  \item Spiegare i risultati con frasi da manuale invece che con le proprie curve e i propri numeri.
  \item Dimenticare che rimuovere una feature e de-duplicare possono cambiare la composizione delle
        classi: ricontare dopo ogni modifica.
  \item Curve illeggibili o senza legenda; affermazioni ``è meglio'' senza metrica di riferimento.
\end{itemize}
\end{attenzione}

\section{Checklist}
\begin{itemize}
  \item[$\square$] Numero di campioni prima e dopo la pulizia, con motivazione di ogni scelta.
  \item[$\square$] Split e scaling motivati; dove è stato fittato lo scaler.
  \item[$\square$] Previsioni scritte prima di ogni run e confrontate dopo.
  \item[$\square$] Le tre curve di 2B con nome del failure mode e meccanismo.
  \item[$\square$] Discussione validation vs test della baseline; perché il test non serve a scegliere.
  \item[$\square$] Esperimento controfattuale sulla porta e ricalcolo delle classi senza porta.
  \item[$\square$] Partizione da cui si calcolano i pesi di classe; scelta motivata della metrica.
  \item[$\square$] Verdetto onesto su profondità, batch size, ottimizzatori e regolarizzazione rispetto
        alla baseline.
  \item[$\square$] Logbook entro il budget; riflessione finale argomentata dai propri run.
  \item[$\square$] Checklist generale del cap.~\ref{cap:lab-report}.
\end{itemize}

\chapter{Lab 2 -- Classificazione supervisionata}
\label{cap:lab2}

\begin{nota}[Fonti e limiti di questo capitolo]
Per il Lab~2 non abbiamo il PDF del testo. La struttura della consegna è ricostruita dai soli titoli
e dalle descrizioni dei task del notebook di lavoro e da quanto detto in aula (presentazione dei lab
alla \rif{2026-09-21}{21/09}, embedding e sequenze alle lezioni del \rif{2026-10-14}{14/10},
\rif{2026-10-19}{19/10} e \rif{2026-10-20}{20/10}). Il lab è valutato: qui \textbf{non} ci sono
soluzioni, codice, numeri o scelte di un gruppo.
\end{nota}

\begin{verifica}
\begin{itemize}
  \item Il titolo nella slide di presentazione è \en{supervised traffic classification}; il materiale
        effettivo riguarda invece \textbf{malware detection} da sequenze di API call.
  \item In aula il docente aveva anticipato che le stesse sequenze sarebbero state modellate come
        tabella (FFNN), serie temporale (RNN) e \textbf{grafo} (GNN); nella consegna ricostruita non
        compare un task con GNN. Controllare sul testo ufficiale.
  \item Numerazione e sotto-domande esatte dei task vanno confrontate con il testo ufficiale.
\end{itemize}
\end{verifica}

\section{Contesto e dataset}
Classificazione binaria \textbf{malware vs goodware}. Ogni campione è un programma, descritto dalla
\textbf{sequenza delle API call} di Windows che ha invocato
(es.\ caricamento di DLL, accesso al registro, \ldots) e da una label binaria (\code{is\_malware}).
Vengono forniti un file di training e uno di test in formato JSON; il validation set va ricavato dal
training. Le sequenze hanno lunghezza variabile e il vocabolario delle API è una feature
\textbf{categorica ad alta cardinalità}: nel test possono comparire API mai viste in training.

L'idea del lab è confrontare \textbf{rappresentazioni diverse dello stesso dato}: insieme di conteggi
senza ordine, finestra di lunghezza fissa, sequenza completa.

\section{Consegna (parafrasi)}
\begin{description}[style=nextline,leftmargin=1.2em]
\item[Task 1 -- Esplorazione dei dati]
  Bilanciamento delle classi, distribuzione delle lunghezze delle sequenze, vocabolario delle API
  costruito sul training e API del test che non vi compaiono.
\item[Task 2 -- Bag-of-Words]
  Rappresentare ogni sequenza con i conteggi delle API e classificarla con una FFNN.
  2A: unigrammi (singole API); 2B: unigrammi + bigrammi (coppie di API adiacenti), con lo stesso split.
  Per ciascuna rappresentazione si confrontano sul validation una variante standard e una con
  \textbf{pesi di classe}, tenendo fisso tutto il resto, e si congela il modello scelto. Infine si
  valuta sul test sia con le sequenze originali sia con sequenze \textbf{rimescolate}, senza
  riaddestrare.
\item[Task 3 -- FFNN su sequenze di lunghezza fissa]
  Stessa architettura nascosta, due codifiche dell'input: 3A ID interi delle API (scalati) come
  feature numeriche; 3B \textbf{embedding addestrabili} delle API.
\item[Task 4 -- RNN sulle sequenze complete]
  Sequenze intere, padding nel batch e lunghezze reali; confronto tra 4A RNN ``vanilla'', 4B RNN
  bidirezionale, 4C LSTM unidirezionale, con scelta del migliore sul validation. Poi: numero di
  parametri al variare della lunghezza della sequenza (RNN contro FFNN a finestra), benchmark del
  tempo di training su CPU e GPU, test su sequenze originali e rimescolate, interpretazione
  complessiva (curve, errori sulla classe minoritaria, parametri, costo) rispetto alle baseline BoW e
  FFNN.
\end{description}

\section{Teoria coinvolta}
\begin{itemize}
  \item FFNN, split, sbilanciamento, pesi di classe e leakage: cap.~\ref{cap:ffnn}.
  \item Loss binaria, metriche per classe e macro-F1, learning curve: cap.~\ref{cap:addestramento}.
  \item Feature categoriche, ID sequenziali contro embedding, vocabolario dal training,
        \code{<UNK>} e \code{<PAD>}: cap.~\ref{cap:embedding}.
  \item RNN, BPTT, vanishing/exploding gradient, clipping, LSTM, bidirezionalità, tensori
        \code{(batch, seq, feat)}, padding e packing: cap.~\ref{cap:rnn}.
  \item Grafi e GNN (se richiesti dal testo ufficiale): cap.~\ref{cap:cnn-gnn}.
  \item Early stopping e regolarizzazione: cap.~\ref{cap:regolarizzazione}.
\end{itemize}

\section{Metodo generale}
\begin{enumerate}
  \item \textbf{Un solo split} train/validation, fissato e condiviso da tutti i task, così i confronti
        sono equi; vocabolario e ogni statistica (scaling, pesi di classe) dal \textbf{solo training}.
  \item \textbf{Baseline prima}: la BoW ignora l'ordine; se un modello sequenziale non la batte, l'ordine
        forse non serve (o il modello non lo sfrutta).
  \item \textbf{Confronti controllati}: tra varianti cambia una cosa sola (rappresentazione, pesi,
        tipo di cella ricorrente, direzione).
  \item \textbf{Il test rimescolato} è una sonda: chiedersi in anticipo cosa dovrebbe succedere a un
        modello che usa l'ordine e a uno che non lo usa, e verificarlo solo dopo aver scelto i modelli.
  \item \textbf{Costo}: oltre alle metriche, parametri e tempo di training fanno parte del giudizio.
\end{enumerate}

\section{Insidie tipiche}
\begin{attenzione}[errori da evitare]
\begin{itemize}
  \item Costruire il vocabolario su training + test (leakage) o non prevedere un token per le API
        sconosciute.
  \item Usare lo stesso indice per il padding e per un'API reale; dimenticare \code{padding\_idx}.
  \item Trattare gli ID interi come numeri con un ordine e delle distanze che non hanno significato.
  \item Prendere l'output dell'ultimo passo del tensore \emph{paddato} invece dell'ultimo passo reale
        (o dello stato finale con le sequenze impacchettate).
  \item Sbagliare la dimensione d'ingresso del classificatore con una RNN bidirezionale (servono
        entrambe le direzioni).
  \item Scegliere i modelli guardando il test (originale o rimescolato).
  \item Valutare un problema sbilanciato con la sola accuracy.
  \item Dimenticare che una RNN bidirezionale richiede la sequenza completa: va bene per l'analisi
        offline, non per una decisione online.
\end{itemize}
\end{attenzione}

\section{Checklist}
\begin{itemize}
  \item[$\square$] Esplorazione: classi, lunghezze, API non viste, con implicazioni per il modello.
  \item[$\square$] Split e vocabolario dal training, gestione di sconosciute e padding dichiarata.
  \item[$\square$] Confronti standard vs pesato e unigrammi vs bigrammi motivati sul validation.
  \item[$\square$] ID interi vs embedding: confronto e spiegazione del perché.
  \item[$\square$] RNN, BiRNN, LSTM: curve, report per classe, scelta sul validation.
  \item[$\square$] Tabella con parametri e tempi (CPU/GPU) dei modelli confrontati.
  \item[$\square$] Interpretazione del test con sequenze rimescolate per ogni famiglia di modelli.
  \item[$\square$] Logbook entro il budget; checklist generale del cap.~\ref{cap:lab-report}.
\end{itemize}

\chapter{Lab 3 -- Anomaly detection}
\label{cap:lab3}

\begin{nota}[Fonti e limiti di questo capitolo]
Testo del laboratorio ``\emph{Project 3: Supervised vs Unsupervised Anomaly Detection}'' pubblicato
sul portale; teoria dalle lezioni del \rif{2026-11-02}{02/11} e del \rif{2026-11-09}{09/11};
presentazione del lab in aula all'inizio della prima sessione (\rif{2026-11-11}{11/11}). La
consegna è \textbf{parafrasata}; per i dettagli fa fede il testo originale. Il lab è valutato: qui
\textbf{non} ci sono soluzioni, codice, numeri o scelte di un gruppo.
\end{nota}

\begin{nota}[Sessioni in aula]
Due sessioni: mercoledì \textbf{11/11} (presentazione del Prof.~Boffa, poi lavoro a gruppi con il
Prof.~Giordano) e mercoledì \textbf{18/11}. Lunedì 16/11 non c'è stata lezione (evento Cyber Monday).
\end{nota}

\begin{nota}[Task 5 eliminato]
Nel testo i task passano dal 4 (clustering) al 6 (autoencoder). Alla presentazione dell'11/11 il docente
ha chiarito che l'esercizio sull'\textbf{Isolation Forest} (Task 5) è stato \textbf{tolto}: le menzioni
rimaste (sottotitolo, introduzione, budget) vanno ignorate e il testo sarà corretto.
\end{nota}

\begin{verifica}
Il testo dice che il report si può consegnare fino a 10 giorni prima di un appello, mentre quello del
Lab~1 dice due settimane (cap.~\ref{cap:formato}). In aula il budget è stato indicato come ``15 run'',
mentre il testo parla di circa 25: fa fede il testo (la cifra in aula potrebbe essere un errore di
trascrizione).
\end{verifica}

\section{Contesto e obiettivi}
Costruire un semplice \textbf{IDS basato su AI} a partire dai log delle connessioni di rete, ognuna
etichettata come normale o intrusione (con il tipo di attacco). Il lab valuta se gli algoritmi
\textbf{non supervisionati} trovano da soli pattern anomali e come cambiano le prestazioni in base a
ciò che si sa dei dati (sapere quali punti sono normali oppure non sapere nulla). La progressione:
\begin{itemize}
  \item analisi ``da esperto di dominio'': come ogni attacco differisce dal traffico normale;
  \item anomaly detection \textbf{shallow} con famiglie diverse: a confine (one-class SVM), a densità
        (LOF), a clustering (k-means, DBSCAN), a isolamento (Isolation Forest, poi tolto: vedi sopra);
  \item anomaly detection \textbf{deep}: autoencoder addestrato sul normale, errore di ricostruzione
        come score;
  \item \textbf{rappresentazione}: PCA (lineare) contro bottleneck dell'autoencoder (non lineare) come
        input dello stesso detector, con visualizzazione t-SNE.
\end{itemize}
\textbf{Domanda centrale}: quanto bisogna sapere del traffico normale per trovare gli attacchi, e le
feature di un attacco sembrano sempre un'anomalia?

\section{Dataset}
Log a livello di connessione di un IDS di rete, già divisi in \code{train.csv} e \code{test.csv}.
Il training contiene traffico normale e tre categorie di attacco: \en{Denial of Service} (dos),
\en{Probing} (probe), \en{Remote-to-Local} (r2l). Le feature sono di tre tipi:
\begin{itemize}
  \item \textbf{basic}: durata, protocollo (TCP, UDP, \ldots), servizio (HTTP, FTP, \ldots), flag di stato;
  \item \textbf{content}: informazioni dal payload (login falliti, accessi a file sensibili, comandi);
  \item \textbf{traffic}: statistiche su finestre temporali (connessioni verso lo stesso host o servizio),
        utili per scansioni e flooding.
\end{itemize}
Le ultime due colonne sono il tipo di attacco e la label binaria (0 normale, 1 anomalia). Il test si
usa \textbf{solo dove il task lo chiede}; un eventuale validation set va ricavato dal training con una
strategia motivata.

\section{Consegna (parafrasi)}
Regole comuni: logbook con un run per ogni fit di detector, clustering o autoencoder; budget di circa
25 run, con una ripartizione indicativa per task; uno sweep diagnostico per scegliere un parametro
conta come un run (cap.~\ref{cap:lab-report}).

\begin{description}[style=nextline,leftmargin=1.2em]
\item[Task 1 -- I dati visti da un esperto (nessun run)]
  Contare attributi categorici e numerici (le feature 0/1 contano come numeriche), descrivere la
  distribuzione delle label di attacco e binarie in training e test e la frazione di anomalie, con le
  implicazioni per l'anomaly detection. Preprocessing di categoriche e numeriche, dicendo da dove si
  stimano i parametri. Poi tre \textbf{heatmap} (righe = tipo di attacco, colonne = feature numeriche)
  di media, deviazione standard e mediana, e discussione delle feature che distinguono gli attacchi.
\item[Task 2 -- One-class SVM e impatto delle anomalie]
  OC-SVM con kernel RBF; le label servono solo a costruire i sottoinsiemi e a valutare. (i) Training
  sul solo normale: scegliere e motivare $\nu$ e confrontarlo con il default, riportando falsi positivi
  e macro-F1 di training; domande su cosa controlla $\nu$, se può essere zero, perché il normale
  ``pulito'' aiuta. (ii) Training su normale + una frazione crescente di anomalie (da nessuna a tutte),
  con contaminazione nota, e per il caso completo anche con contaminazione ignota; grafico della
  macro-F1 e domande sull'equità del confronto. (iii) Generalizzazione: tre modelli valutati sullo
  stesso test con macro-F1, matrice di confusione e recall per tipo di attacco, collegando gli
  errori all'analisi delle heatmap.
\item[Task 3 -- Local Outlier Factor]
  Sul training con anomalie, confrontare quattro valori di \code{n\_neighbors} con un criterio di
  scelta motivato, poi valutare sul test. Domande: gli attacchi hanno densità locale più bassa dei
  vicini? Come incide \code{n\_neighbors}?
\item[Task 4A -- K-means]
  Fit sul training completo, scelta motivata di $k$; dimensione e purezza dei cluster rispetto alle
  label, silhouette per cluster; distanza dal centroide più vicino come score con soglia scelta e
  motivata; t-SNE del training con tre perplexity, colorato per label e per cluster.
\item[Task 4B -- DBSCAN]
  Scelta motivata di \code{min\_points} ed $\varepsilon$; il rumore (cluster $-1$) è l'anomalia.
  Domande: il rumore contiene solo attacchi? Quanti cluster? Somiglianze con LOF?
\item[Task 5 -- (eliminato)]
  L'esercizio sull'Isolation Forest è stato tolto dal docente (riquadro ``Task 5 eliminato'').
\item[Task 6 -- Autoencoder]
  Architettura libera ma con encoder che si restringe, bottleneck e decoder che si espande. Training
  sul \textbf{solo normale}, scegliendo epoch e learning rate su un validation di soli normali; soglia
  sull'errore di ricostruzione; ECDF degli errori su validation normale, training completo e test.
  Poi un secondo autoencoder addestrato su normale + anomalie e confronto dei due sul test.
\item[Task 7 -- PCA contro bottleneck dell'autoencoder]
  Stessa dimensione per le due rappresentazioni. Bottleneck del miglior autoencoder sui normali di
  training $\to$ OC-SVM $\to$ valutazione sugli embedding del test. PCA fittata sui soli normali di
  training, numero di componenti dalla varianza spiegata, poi OC-SVM come sopra. Confronto con la
  migliore OC-SVM sulle feature originali e con l'autoencoder a errore di ricostruzione.
\end{description}

\section{Teoria coinvolta}
\begin{itemize}
  \item Tipi di anomalia, novelty vs outlier, score vs label, sfide dell'anomaly detection, OC-SVM,
        k-means e SSE, DBSCAN, LOF, Isolation Forest, silhouette, PCA, t-SNE, autoencoder:
        cap.~\ref{cap:anomaly}.
  \item Encoding delle categoriche, scaling e leakage: cap.~\ref{cap:ffnn} e~\ref{cap:embedding}.
  \item Autoencoder come apprendimento self-supervised di una rappresentazione: cap.~\ref{cap:embedding}.
  \item Training dell'autoencoder, loss, learning curve, metriche e macro-F1: cap.~\ref{cap:addestramento}.
\end{itemize}

\section{Metodo generale}
\begin{enumerate}
  \item \textbf{Label solo per valutare}: possono servire a costruire i sottoinsiemi (es.\ ``solo
        normali'') e a misurare, mai come segnale di training.
  \item \textbf{Dichiarare le assunzioni}: per ogni metodo, cosa si suppone di sapere (normali noti,
        contaminazione nota o ignota, nulla).
  \item \textbf{Preprocessing stimato sul training} e applicato a validation e test; scaling prima di
        PCA, distanze e densità.
  \item \textbf{Ogni soglia è una scelta}: $\nu$, \code{n\_neighbors}, $k$, $\varepsilon$, soglia
        sull'errore di ricostruzione vanno motivati con evidenze (grafici, sweep diagnostici) e fissati
        senza guardare il test.
  \item \textbf{Ragionare per tipo di attacco}: l'analisi delle heatmap serve dopo a spiegare perché
        alcuni attacchi vengono trovati e altri no.
  \item \textbf{Confronto finale} su prestazioni, tempo di training e conoscenza richiesta da ogni metodo.
\end{enumerate}

\section{Insidie tipiche}
\begin{attenzione}[errori da evitare]
\begin{itemize}
  \item Fittare encoder e scaler su tutto il dataset prima dello split.
  \item Calcolare la soglia sugli score del test, o scegliere i parametri guardando il test.
  \item Addestrare l'autoencoder ``sul normale'' includendo per errore anche le anomalie.
  \item Dimenticare lo scaling prima della PCA: una feature con scala enorme domina la varianza.
  \item Leggere distanze e dimensioni dei gruppi nella mappa t-SNE come se fossero quelle reali.
  \item Riportare solo la ROC-AUC con dati molto sbilanciati: servono precision, recall e F1 della
        classe anomala e la recall per tipo di attacco.
  \item Confrontare punteggi ottenuti su insiemi di training diversi come se fossero comparabili.
  \item Assumere che un attacco sia per forza un outlier: un gruppo piccolo ma denso può avere densità
        simile ai propri vicini.
\end{itemize}
\end{attenzione}

\section{Checklist}
\begin{itemize}
  \item[$\square$] Conteggi delle feature e distribuzione delle label in training e test, con implicazioni.
  \item[$\square$] Le tre heatmap con discussione per tipo di attacco.
  \item[$\square$] Per ogni metodo: assunzioni dichiarate, parametri e soglie motivati con evidenze.
  \item[$\square$] OC-SVM: tre scenari, grafico al variare della contaminazione, recall per attacco sul test.
  \item[$\square$] LOF, k-means (purezza, silhouette, soglia, t-SNE), DBSCAN discussi rispetto alle label.
  \item[$\square$] Autoencoder: ECDF degli errori, confronto normale-only vs con anomalie.
  \item[$\square$] PCA vs bottleneck con stessa dimensione e stesso detector.
  \item[$\square$] Classifica finale dei metodi (prestazioni, tempo, conoscenza richiesta).
  \item[$\square$] Logbook entro il budget; checklist generale del cap.~\ref{cap:lab-report}.
\end{itemize}


% =====================================================================
\part{Schemi di ripasso}
% =====================================================================
% =====================================================================
\chapter{Mappe concettuali}
\label{cap:mappe}
% =====================================================================

Come usarle: copri le foglie (a destra) e prova a ricostruirle partendo dal nodo padre; poi controlla. Ogni mappa rimanda al
capitolo del manuale in cui l'argomento è spiegato. Le tabelle ``quale scelta quando'' sono nel cap.~\ref{cap:schede}.

% ---------------------------------------------------------------------
\section{Mappa generale del corso}
\begin{center}
\begin{tikzpicture}[mindmap, grow cyclic, every node/.style={concept, execute at begin node=\hskip0pt},
  concept color=polblu, text=white,
  root concept/.append style={font=\large\bfseries, minimum size=3.1cm, text width=2.6cm},
  level 1/.append style={level distance=4.6cm, sibling angle=60, font=\small\bfseries, minimum size=2.3cm, text width=2.0cm},
  level 2/.append style={level distance=2.75cm, sibling angle=34, font=\scriptsize, minimum size=1.6cm, text width=1.5cm}]
\node[root concept] {AI \&\\Cyber\-security}
  child[concept color=defcol] { node {Fonda\-menti}
    child { node {AI, ML, DL} }
    child { node {Para\-digmi} }
    child { node {Uso dei modelli} } }
  child[concept color=red!60!black] { node {Training}
    child { node {Loss e metriche} }
    child { node {Ottimiz\-zatori} }
    child { node {Over\-fitting} } }
  child[concept color=esecol] { node {Dati}
    child { node {Pre\-processing} }
    child { node {Encoding} }
    child { node {Split e leakage} } }
  child[concept color=brown!70!black] { node {Archi\-tetture}
    child { node {FFNN} }
    child { node {RNN, LSTM} }
    child { node {CNN} }
    child { node {GNN} } }
  child[concept color=orange!80!black] { node {Embed\-ding}
    child { node {Super\-vised} }
    child { node {Self-super\-vised} } }
  child[concept color=violet!60!black] { node {Anomaly detection}
    child { node {Tipi} }
    child { node {OC-SVM} }
    child { node {K-means} } };
\end{tikzpicture}
\end{center}

% ---------------------------------------------------------------------
\clearpage
\section{Paradigmi e uso dei modelli \normalfont\small(cap.~\ref{cap:ai-cyber}, \rif{2026-09-23}{23/09})}
\begin{center}
\begin{forest} mappa
[Come impara\\un modello
  [Supervised
    [{Label = ground truth, costosa, deve essere affidabile. Predittivo: classificazione (categorica) o regressione (numerica). Impara solo il passato. Problemi: \textbf{drift}, \textbf{sbilanciamento}. Rimedi: drift detection, class weights, metriche giuste.}]]
  [Unsupervised
    [{Niente label, dati economici. Descrittivo: clustering (botnet), anomaly detection. Problemi: drift, \textbf{interpretazione} difficile, falsi allarmi. Rimedi: drift detection, XAI, analisi manuale.}]]
  [Semi-supervised
    [{Poche label $\to$ primo modello $\to$ pseudo-label sul resto $\to$ modello finale. Problemi: campione etichettato poco rappresentativo, sbilanciamento, evoluzione dei dati.}]]
  [Self-supervised
    [{La supervisione viene dai dati. \textbf{Rappresentazione}, non descrizione. Pretext (encoder $\to$ bottleneck) + downstream (decoder o classificatore). Base degli LLM. Costoso, serve tanta qualità, bias, valutazione difficile.}]]
  [Reinforcement
    [{Agente, stato, azioni, reward; tentativi ed errori. Reward difficile da progettare, lento, instabile, poco interpretabile.}]]
]
\end{forest}
\end{center}

\begin{center}
\begin{forest} mappa
[Usare\\i modelli
  [Singolo / ensemble
    [{Ensemble: voto o media (random forest); stacking con meta-learner. Meglio se i modelli non sono correlati; più costoso.}]]
  [Federated
    [{Si condivide il modello, non i dati: privacy, scalabilità, visibilità. Rischi: costo sull'edge, dati eterogenei, \textbf{model poisoning}.}]]
  [Transfer
    [{Fine-tuning degli ultimi layer con pochi dati (STAMINA: ResNet-50 su malware come immagini). Rischi: domain mismatch, bias ereditato.}]]
  [Generativi
    [{Nuovi campioni dalla distribuzione appresa. Rischi: costo, valutazione, bias, abuso (deepfake, phishing), privacy (leak di dati).}]]
  [AI e cyber
    [{Difesa (IDS, SIEM, SOAR, predictive), attacco (deepfake, phishing, vulnerabilità, evasione), \textbf{adversarial ML} (poisoning, adversarial example, jailbreak, model inversion).}]]
]
\end{forest}
\end{center}

% ---------------------------------------------------------------------
\clearpage
\section{Dal neurone all'addestramento \normalfont\small(cap.~\ref{cap:reti-neurali}--\ref{cap:addestramento})}
\begin{center}
\begin{forest} mappa
[Rete neurale
  [Perceptron
    [{$y = f(\mathbf{w}^T\mathbf{x})$, con bias $w_0$ ($x_0=1$) per traslare. I pesi sono i parametri, appresi dai dati, inizializzati a caso.}]]
  [Attivazioni
    [{Hidden: \textbf{ReLU} (non limitata, sparsità); RNN: \textbf{tanh}. Output: \textbf{sigmoid} (binario), \textbf{softmax} (multi-classe), \textbf{lineare} (regressione). Senza non linearità $N$ layer = 1 layer lineare.}]]
  [Layer
    [{Fully connected = $f(W\mathbf{x})$. Approssimazione universale: basta 1 hidden layer, ma non dice quanti neuroni. Meglio \textbf{profonde e strette}: rappresentazioni più astratte.}]]
  [Logit
    [{Output grezzi in $(-\infty,+\infty)$; sigmoid/softmax li rendono probabilità. In PyTorch \code{CrossEntropyLoss} vuole i logit.}]]
  [Loss
    [{MSE (differenziabile), MAE (robusta agli outlier), BCE, cross-entropy. Loss bassa $\neq$ modello migliore: si sceglie con le metriche.}]]
  [Ottimizzazione
    [{GD (tutti i dati, minimi locali) $\to$ SGD (mini-batch, rumore utile) $\to$ momentum ($\beta$, traiettoria liscia) $\to$ normalizzazione (pari importanza) $\to$ \textbf{Adam} (tutto + bias correction). LR scheduler: alto poi basso.}]]
  [Backpropagation
    [{Regola della catena sul grafo computazionale, dall'output all'input. Prodotto di gradienti: $<1$ $\Rightarrow$ \textbf{vanishing}, $>1$ $\Rightarrow$ \textbf{exploding} (clipping).}]]
]
\end{forest}
\end{center}

% ---------------------------------------------------------------------
\section{Leggere le learning curve \normalfont\small(cap.~\ref{cap:addestramento}, \ref{cap:regolarizzazione})}
\begin{center}
\begin{forest} mappa3
[Learning\\curve
  [Underfitting
    [{Train e val alte, scendono poco}
      [{Più complessità (neuroni, layer), altra architettura, più epoch, preprocessing migliore (rumore, outlier, encoding).}]]]
  [Overfitting
    [{Train $\to 0$, val scende poi \textbf{risale}}
      [{Pesi: L1, L2 (\code{weight\_decay}). Architettura: dropout, rete più piccola. Training: early stopping + modello migliore salvato. Dati: batch norm, noise injection.}]]]
  [Non converge
    [{Nessun plateau}
      [{Addestrare più a lungo; controllare il learning rate.}]]]
  [Curva rumorosa
    [{Oscillazioni forti}
      [{LR troppo alto, batch piccolo, troppa regolarizzazione (BN + weight decay); normale un po' di rumore con SGD e RNN.}]]]
  [Val $<$ train
    [{Validation più bassa}
      [{Possibile: validation più facile, o dropout attivo solo in training.}]]]
]
\end{forest}
\end{center}

% ---------------------------------------------------------------------
\clearpage
\section{Data handling \normalfont\small(cap.~\ref{cap:ffnn})}
\begin{center}
\begin{forest} mappa
[Data handling
  [Ordine corretto
    [{\textbf{Split} $\to$ fit di scaler/encoder/vocabolario \textbf{solo sul training} $\to$ transform di validation e test. Altrimenti \textbf{leakage}.}]]
  [Pre-processing
    [{Rumore; \textbf{outlier} (rovinano il min-max); ridondanza e duplicati; valori mancanti: elimina punti, stima, togli feature, sostituisci con pesi.}]]
  [Categoriche
    [{Ordinale $\to$ ID sequenziali. Nominale con poche categorie $\to$ one-hot. Molte categorie e rete neurale $\to$ \textbf{embedding} (+ \code{<UNK>}). Altrimenti feature engineering. Porte = categoriche!}]]
  [Numeriche
    [{Min-max $[0,1]$: sensibile agli outlier. Z-score $(x-\mu)/\sigma$ ($\sigma$ = deviazione standard): robusto, senza limiti.}]]
  [Split
    [{Hold-out (70/20/10, dataset grandi), k-fold (medi; poi ri-addestra su tutto), leave-one-out (minuscoli). Stratificato senza reinserimento; bootstrap con reinserimento.}]]
  [Sbilanciamento
    [{Dati: oversampling (SMOTE, ADASYN), undersampling (cluster-based). Algoritmo: class weights, focal loss. Metrica: per classe, F1, \textbf{macro}.}]]
  [Leakage
    [{Correlazione temporale (split temporale: sliding/expanding window), feature-scorciatoia (DoS = porta 80), identificatori, pezzi adiacenti di un'immagine.}]]
]
\end{forest}
\end{center}

% ---------------------------------------------------------------------
\section{Architetture \normalfont\small(cap.~\ref{cap:ffnn}, \ref{cap:rnn}, \ref{cap:cnn-gnn})}
\begin{center}
\begin{forest} mappa
[Architetture
  [FFNN (tabelle)
    [{Input = n.\ feature dopo l'encoding; hidden 3--5 layer, potenze di 2, ReLU; output = classi. \code{model.train()} / \code{eval()} + \code{no\_grad()}; \code{zero\_grad} $\to$ \code{backward} $\to$ \code{step} per batch.}]]
  [RNN (sequenze)
    [{$h_t = \tanh(W_{hh}h_{t-1} + W_{xh}x_t + b)$, pesi condivisi nel tempo; BPTT; many-to-one, one-to-many, many-to-many; mono/bidirezionale. Tensore \code{(batch, seq, feat)}; padding nel batch + \code{pack\_padded\_sequence}; \code{hidden[-1]} per classificare.}]]
  [LSTM
    [{Cell state (lungo termine) + hidden (breve); gate forget, input, candidate, output (sigmoid come rubinetti). Mitiga il vanishing gradient, più costosa.}]]
  [CNN (immagini)
    [{Convoluzione (filtri = pesi) + ReLU + pooling (max). Feature gerarchiche: bordi $\to$ forme $\to$ oggetti. LeNet, AlexNet (ReLU, dropout), ResNet (skip). Malware come immagine.}]]
  [GNN (grafi)
    [{FFNN fallisce: ordine dei nodi, dimensione. Invarianza (grafo) ed equivarianza (nodo). Message passing: pesi condivisi + aggregazione indipendente dall'ordine; induttiva. GCN, GraphSAGE (campiona), GAT (attenzione). $K$ layer = $K$ hop; troppi = \textbf{over-smoothing}. \code{global\_mean\_pool} per grafo.}]]
]
\end{forest}
\end{center}

% ---------------------------------------------------------------------
\clearpage
\section{Embedding \normalfont\small(cap.~\ref{cap:embedding})}
\begin{center}
\begin{forest} mappa
[Embedding
  [Cosa sono
    [{Vettori \textbf{densi} $H$-dimensionali ($H \ll V$) che codificano semantica; partono casuali e si imparano risolvendo un task con backprop. Solo con reti neurali.}]]
  [Supervised
    [{\code{nn.Embedding(V, H)} = lookup addestrabile (one-hot + lineare). Vocabolario \textbf{dal solo training} + \code{<UNK>}. Concatenazione con le numeriche (input = embed\_dim + n\_num). Anche su sequenze (RNN, \code{padding\_idx}) e architetture composte (embedding + GNN + RNN).}]]
  [Limiti supervised
    [{Ottimi per \textbf{quel} task, sub-ottimi per altri. Rischio \textbf{scorciatoie} (husky su erba, lupo su neve).}]]
  [Self-supervised
    [{Pretext su tanti dati + downstream con poche label. \textbf{Autoencoder}: encoder $\to$ embedding $\to$ decoder, ricostruzione (regressione). \textbf{Contrastive}: anchor e augmentation vicini, negativi lontani (contrastive loss); le augmentation introducono bias.}]]
]
\end{forest}
\end{center}

% ---------------------------------------------------------------------
\section{Anomaly detection \normalfont\small(cap.~\ref{cap:anomaly})}
\begin{center}
\begin{forest} mappa
[Anomaly\\detection
  [Definizioni
    [{Anomaly/outlier: insolito, spesso malevolo. Novelty: nuovo, può diventare la \textbf{nuova normalità} (VPN nel COVID).}]]
  [Tipi
    [{\textbf{Point} (login dalla Russia), \textbf{contextual} (backup di domenica), \textbf{collective} (brute force sotto il rate limit).}]]
  [Complessità
    [{Sconosciute, rare e sbilanciate, eterogenee, in evoluzione; errore su anomalia costa di più. Sfide ML: recall vs falsi positivi, dati complessi, poche label, rumore, anomalie complesse, spiegabilità.}]]
  [Output
    [{\textbf{Score} (ranking + budget, es.\ top 10.000) o \textbf{label} (quantità controllata dagli iperparametri).}]]
  [Shallow
    [{Interpretabili, ottimi su tabelle e anomalie puntuali, baseline. Deboli su contestuali (serve feature engineering) e collettive (feature aggregate).}]]
  [One-class SVM
    [{Confine attorno al normale; kernel (RBF il più usato); $\nu$ = frazione attesa di outlier. ``Supervised'' (solo normali, novelty) vs unsupervised (tutto mescolato, outlier).}]]
  [Clustering
    [{Anomalie = cluster piccoli e lontani, punti lontani dal centroide, punti senza cluster. \textbf{K-means}: $K$ fissato, ogni punto in un cluster, sensibile all'inizializzazione (più run), $K$ con SSE + gomito.}]]
  [Densità e alberi
    [{\textbf{DBSCAN}: rumore = anomalia, $\varepsilon$ dal k-distance plot. \textbf{LOF}: densità del punto vs vicini, score continuo, anomalie locali. \textbf{Isolation Forest}: percorso corto = anomalia, niente distanze (alta dimensionalità).}]]
  [Deep
    [{\textbf{Autoencoder} addestrato \textbf{solo sui normali}: score = errore di ricostruzione (MSE + CE). Encoder come estrattore di feature + metodo shallow. Anomalie nel training $\Rightarrow$ prestazioni peggiori.}]]
  [Validazione
    [{Label solo per valutare; scaling prima di PCA e distanze; t-SNE solo visivo; tuning sul validation, test una volta; ROC-AUC ingannevole: guarda F1 delle anomalie.}]]
]
\end{forest}
\end{center}

% ---------------------------------------------------------------------

\chapter{Schede: quale scelta quando}
\label{cap:schede}

Tabelle di decisione rapida per il ripasso e per gli esercizi di pipeline dello scritto
(cap.~\ref{cap:esercizi}). Ogni riga rimanda implicitamente al capitolo del manuale sull'argomento:
output e loss nel cap.~\ref{cap:addestramento}, dati nel cap.~\ref{cap:ffnn}, architetture nei
cap.~\ref{cap:rnn} e~\ref{cap:cnn-gnn}, embedding nel cap.~\ref{cap:embedding}, anomaly detection
nel cap.~\ref{cap:anomaly}. Le mappe concettuali sono nel cap.~\ref{cap:mappe}.

\section{Output, attivazione finale e loss}
\begin{center}
\small
\begin{tabularx}{\textwidth}{l l l X}
\toprule
\textbf{Task} & \textbf{Neuroni output} & \textbf{Attivazione} & \textbf{Loss (PyTorch)}\\
\midrule
Regressione & 1 (per target) & lineare (nessuna) & MSE (\code{nn.MSELoss}), MAE (\code{nn.L1Loss}) se ci sono outlier\\
Binaria & 1 & sigmoid & BCE; in PyTorch \code{nn.BCEWithLogitsLoss} include la sigmoid$^{*}$\\
Multi-classe & $C$ & softmax & \code{nn.CrossEntropyLoss}: riceve i \textbf{logit}, softmax interna\\
\bottomrule
\end{tabularx}
\end{center}
\begin{integrazione}[$^{*}$ \code{BCEWithLogitsLoss}]
A lezione si è usata \code{CrossEntropyLoss} anche per il binario (2 neuroni di output), che funziona allo stesso modo;
\code{BCEWithLogitsLoss} con un solo neurone è l'alternativa equivalente.
\end{integrazione}

\section{Quale metrica}
\begin{center}
\small
\begin{tabularx}{\textwidth}{l X}
\toprule
\textbf{Situazione} & \textbf{Metrica}\\
\midrule
Classi bilanciate, errori di pari costo & accuracy (con precision/recall per sicurezza)\\
Classi sbilanciate & precision, recall, F1 \textbf{per classe}; media \textbf{macro}\\
Il falso negativo costa molto (malware) & \textbf{recall} alta sulla classe critica, con precision accettabile\\
Troppi falsi allarmi inaccettabili & \textbf{precision} sulla classe ``allarme''\\
Regressione & MSE/MAE sulla scala originale dei dati\\
Clustering & SSE (scelta di $K$); senza label la validazione è la parte difficile\\
\bottomrule
\end{tabularx}
\end{center}

\section{Quale architettura}
\begin{center}
\small
\begin{tabularx}{\textwidth}{>{\raggedright\arraybackslash}p{3.3cm} >{\raggedright\arraybackslash}p{5.4cm} >{\raggedright\arraybackslash}X}
\toprule
\textbf{Dato} & \textbf{Architettura} & \textbf{Attenzione a}\\
\midrule
Tabella di feature & FFNN (o shallow come baseline) & encoding, scaling, sbilanciamento\\
Sequenza / serie & RNN, LSTM se memoria lunga, Bi-RNN se sequenza completa e non tempo reale & tensore 3D, padding nel batch, split temporale\\
Immagine (o dato trasformabile in immagine) & CNN, spesso pre-addestrata + fine-tuning & dimensione d'ingresso fissa (es.\ $224\times224$)\\
Grafo & GNN: GCN, GraphSAGE (grafi grandi), GAT (attenzione) & numero di layer (over-smoothing), pooling per la graph classification\\
Feature categorica ad alta cardinalità & \code{nn.Embedding} prima della rete & vocabolario dal training, \code{<UNK>}\\
Pochi dati etichettati & self-supervised/pre-addestrato + fine-tuning, semi-supervised & domain mismatch\\
Nessuna label, eventi rari & anomaly detection (OC-SVM, clustering, \ldots) & falsi positivi, novità vs anomalie\\
Tabella, normali noti, nessuna label degli attacchi & autoencoder o OC-SVM addestrati sui soli normali & soglia sull'errore di ricostruzione, contaminazione\\
Molte feature, nessuna label & Isolation Forest (nessuna distanza) o riduzione di dimensione prima dei metodi a distanza & maledizione della dimensionalità\\
\bottomrule
\end{tabularx}
\end{center}


\section{Codificare una feature categorica}
\begin{center}
\small
\begin{tabularx}{\textwidth}{>{\raggedright\arraybackslash}p{3.2cm} >{\raggedright\arraybackslash}X >{\raggedright\arraybackslash}X}
\toprule
\textbf{Codifica} & \textbf{Quando} & \textbf{Limiti}\\
\midrule
ID sequenziali & feature \textbf{ordinale}; label di classificazione & su una nominale impone un ordine e distanze false\\
One-hot & nominale con \textbf{poche} categorie; va bene anche per modelli shallow & esplode con molte categorie; non gestisce categorie nuove\\
Feature engineering & esiste una chiave di lettura di dominio (es.\ porte note) & può essere sbagliata, cambiare nel tempo o essere impossibile\\
Embedding appreso & nominale ad \textbf{alta cardinalità} con una rete neurale & serve un task e abbastanza dati; vocabolario dal training, \code{<UNK>}, \code{<PAD>}\\
\bottomrule
\end{tabularx}
\end{center}

\section{Scaling e split}
\begin{center}
\small
\begin{tabularx}{\textwidth}{>{\raggedright\arraybackslash}p{3.6cm} >{\raggedright\arraybackslash}X}
\toprule
\textbf{Situazione} & \textbf{Scelta}\\
\midrule
Serve un intervallo fisso $[0,1]$, niente outlier & normalizzazione min-max\\
Outlier o code lunghe & standardizzazione z-score (eventualmente log prima)\\
Dataset grande & hold-out train/validation/test\\
Dataset medio & k-fold cross-validation, poi ri-addestramento su tutto\\
Dataset minuscolo & leave-one-out\\
Classi sbilanciate & split \textbf{stratificato}\\
Dati con ordine temporale & split \textbf{temporale} (finestra scorrevole o crescente)\\
Sempre & fit di scaler, encoder, vocabolario e pesi di classe sul \textbf{solo training}\\
\bottomrule
\end{tabularx}
\end{center}

\section{Sintomo nelle learning curve $\rightarrow$ rimedio}
\begin{center}
\small
\begin{tabularx}{\textwidth}{>{\raggedright\arraybackslash}p{4.4cm} >{\raggedright\arraybackslash}X}
\toprule
\textbf{Sintomo} & \textbf{Rimedio}\\
\midrule
Train e val alte, scendono poco (underfitting) & più capacità, architettura diversa, più epoch, preprocessing migliore\\
Train $\to 0$, val risale (overfitting) & early stopping con salvataggio, weight decay, dropout, rete più piccola, batch norm\\
Loss oscillante o divergente & learning rate più basso, batch più grande, scaling delle feature\\
Loss che scende lentissima & learning rate più alto, scaling, ottimizzatore adattivo\\
Gradienti che esplodono (RNN) & gradient clipping\\
Gradienti che svaniscono & ReLU, LSTM/GRU, skip connection\\
\bottomrule
\end{tabularx}
\end{center}

\section{Training e valutazione in PyTorch}
\begin{center}
\small
\begin{tabularx}{\textwidth}{l X X}
\toprule
 & \textbf{Training} & \textbf{Validation / test}\\
\midrule
Modalità & \code{model.train()} & \code{model.eval()}\\
Gradienti & calcolati: \code{zero\_grad} $\to$ \code{backward} $\to$ \code{step} per batch & \code{torch.no\_grad()}\\
Dropout & attivo, sopravvissuti scalati di $1/(1-p)$ & spento\\
Batch norm & statistiche del batch, aggiorna le running & usa running mean e variance\\
DataLoader & \code{shuffle=True} & \code{shuffle=False}\\
Preprocessing & \code{fit\_transform} & solo \code{transform}\\
\bottomrule
\end{tabularx}
\end{center}

\section{Metodi di anomaly detection a confronto}
\begin{center}
\small
\begin{tabularx}{\textwidth}{>{\raggedright\arraybackslash}p{2.4cm} >{\raggedright\arraybackslash}p{3.2cm} >{\raggedright\arraybackslash}p{2.6cm} >{\raggedright\arraybackslash}X}
\toprule
\textbf{Metodo} & \textbf{Score / decisione} & \textbf{Parametro chiave} & \textbf{Punti deboli}\\
\midrule
One-class SVM & fuori dal confine attorno al normale & $\nu$, kernel & soffre la contaminazione; costoso su dati grandi\\
K-means & distanza dal centroide; cluster piccoli & $K$ (SSE + gomito) & cluster sferici; ogni punto in un cluster; inizializzazione\\
DBSCAN & punti di rumore (cluster $-1$) & $\varepsilon$, MinPts & $\varepsilon$ critico; densità diverse; nessuno score graduato\\
LOF & densità del punto vs vicini & $k$ (\code{n\_neighbors}) & gruppi piccoli e densi hanno LOF $\approx 1$; costoso\\
Isolation Forest & percorso medio corto & contaminazione & meno interpretabile\\
Autoencoder & errore di ricostruzione & soglia, architettura & va addestrato sul normale; soglia mai dal test\\
\bottomrule
\end{tabularx}
\end{center}
I metodi a distanza e densità soffrono l'alta dimensionalità: ridurre le dimensioni (PCA dopo lo
scaling) o preferire l'Isolation Forest.


% =====================================================================
\part{Quiz ed esame}
% =====================================================================
\chapter{Formato dell'esame e regole del corso}
\label{cap:formato}

\begin{nota}[Fonti]
Slide \emph{0-Intro-Course}, quanto detto in aula alla \rif{2026-09-21}{21/09} (presentazione del
corso) e alla \rif{2026-10-20}{20/10} (linee guida per il report), testi dei laboratori pubblicati
sul portale (sezione ``Lab Rules''). Dove le fonti non coincidono lo segnaliamo nei riquadri
``Da verificare'': controllate sul portale della didattica o chiedete al docente prima dell'esame.
\end{nota}

\begin{verifica}
\begin{itemize}
  \item \textbf{Punteggio dello scritto}: la slide dice \textbf{32 punti}; in aula il docente ha
        detto \textbf{33}.
  \item \textbf{Dimensione dei gruppi}: in aula, in una slide (``Classes' structure'') e nei testi dei
        lab \textbf{4 persone} (``novità di quest'anno''); un'altra slide (``Lab activities'') dice
        \textbf{3 persone}.
  \item \textbf{Scadenza dei report}: la slide e il testo del Lab~3 dicono \textbf{10 giorni} prima
        dell'appello; il testo del Lab~1 dice ``fino a \textbf{due settimane} prima'' di un appello.
  \item \textbf{Lab 2}: la slide lo chiama \en{supervised traffic classification}; in aula è stato
        descritto come classificazione malware/goodware a partire da sequenze di API call (vedi
        cap.~\ref{cap:lab2}).
\end{itemize}
La guida del corso sul portale (\code{didattica.polito.it}) richiede l'accesso autenticato: non è
stato possibile consultarla per risolvere i dubbi.
\end{verifica}

\section{Dati generali}

\begin{center}
\begin{tabularx}{\textwidth}{>{\bfseries}l X}
\toprule
Insegnamento & AI and Cybersecurity (codici 01GYZUV, 01GYZUW) \\
CFU & 6 CFU $\approx$ 60 ore \\
Ripartizione ore & 39 h di lezione (teoria + esempi di codice) + 21 h di laboratorio \\
Docenti & Prof.\ \textbf{Danilo Giordano} (titolare, gran parte del corso) --- \code{danilo.giordano@polito.it} \\
        & Prof.\ \textbf{Matteo Boffa} (NLP e ultima parte del corso) --- \code{matteo.boffa@polito.it} \\
Lingua & Inglese (lezioni e slide) \\
Materiale & Portale della didattica: slide, notebook Python, testi dei lab, registrazioni \\
Ambiente & Python + PyTorch; consigliato \textbf{Google Colab} (GPU gratuita a tempo limitato,
           collaborazione sul notebook). Alternative: Jupyter locale / Anaconda \\
\bottomrule
\end{tabularx}
\end{center}

\subsection{Orari e aule}
\begin{center}
\begin{tabular}{lll}
\toprule
\textbf{Giorno} & \textbf{Orario} & \textbf{Aula / note}\\
\midrule
Lunedì & 13:00--14:30 (1,5 h) & 2C --- lezione \\
Martedì & 16:00--17:30 (1,5 h) & 8I --- \textbf{opzionale}: si usa solo se annunciato sul portale \\
Mercoledì & 11:30--14:30 (3 h) & 2S --- lezione o, quando possibile, laboratorio \\
\bottomrule
\end{tabular}
\end{center}

\subsection{Materiale e libri}
\begin{itemize}
  \item Le \textbf{slide} sono volutamente verbose e \textbf{autosufficienti}: secondo il docente
        bastano per prendere anche il voto massimo. Vengono caricate a ridosso della lezione.
  \item Le lezioni sono registrate; dei laboratori si registra solo la spiegazione iniziale.
  \item Libri consigliati (non necessari):
  \begin{itemize}
    \item I.\,H.\ Sarker, \emph{AI-Driven Cybersecurity and Threat Intelligence}, Springer, 2024.
    \item S.\,J.\,D.\ Prince, \emph{Understanding Deep Learning}, MIT Press, 2023 (molte figure
          delle slide vengono da qui).
  \end{itemize}
\end{itemize}

\subsection{Prerequisiti}
Concetti base di data science e machine learning (il corso fa un rapido ripasso iniziale);
Python con \code{numpy}, \code{pandas}, \code{sklearn}, \code{matplotlib}. PyTorch viene
introdotto durante il corso. Corso utile come base (non obbligatorio):
\emph{Machine Learning for Networking} (01DSMBG, Prof.\ Luca Vassio).

\subsection{Programma}
\begin{multicols}{2}
\textbf{Lezioni (39 h)}
\begin{itemize}
  \item Introduzione ad AI e ML per la cybersecurity
  \item Fondamenti di deep learning
  \item Architetture deep: teoria e implementazione (FFNN, RNN, CNN, GNN)
  \item Supervised e unsupervised, anomaly detection
  \item Introduzione al deep NLP e alle tecniche di embedding
  \item Casi d'uso di cybersecurity (seminari)
  \item Simulazione d'esame (ultima settimana)
\end{itemize}
\columnbreak
\textbf{Laboratori (21 h)}
\begin{enumerate}
  \item Introduzione ai framework di deep learning (3 h)
  \item Classificazione supervisionata (6 h)
  \item Unsupervised e anomaly detection, es.\ rilevamento DDoS (6 h)
  \item NLP e clustering per il rilevamento di botnet (6 h)
\end{enumerate}
\end{multicols}

\section{Laboratori e report}

\subsection{Gruppi}
\begin{itemize}
  \item Gruppi \textbf{auto-organizzati} di \textbf{4 persone} (vedi il riquadro ``Da verificare'':
        una slide dice 3), da registrare con il form Moodle entro fine ottobre.
  \item Il tempo in aula non basta per finire un lab: si completa a casa.
\end{itemize}

\subsection{Cosa consegnare (per ogni lab, Lab 1 incluso)}
Un \textbf{file zip} per lab, chiamato \code{groupID\_labX.zip} secondo i testi dei lab, con:
\begin{enumerate}
  \item \textbf{Report} di \textbf{massimo 10 pagine} (\code{groupID\_labX.pdf}): cosa e come avete
        fatto, decisioni, esperimenti e risultati, con figure e tabelle; frammenti di codice solo se
        servono. Nell'intestazione: \textbf{ID del lab, ID del gruppo, nomi ed email} dei membri.
  \item \textbf{Notebook Jupyter} (nessun limite di pagine), \textbf{eseguito} (output e numeri di
        esecuzione visibili), \textbf{commentato} e \textbf{rieseguibile} (seed fissati); secondo i testi
        dei lab va incluso sia il \code{.ipynb} sia un export HTML.
  \item \textbf{Experiment logbook} (novità di quest'anno): una tabella con una riga per ogni
        addestramento, con budget di circa 25 run per lab; \textbf{non conta} nelle 10 pagine
        (dettagli nel cap.~\ref{cap:lab-report}).
\end{enumerate}

\subsection{Scadenze e valutazione}
\begin{itemize}
  \item \textbf{Lab 1} (solo formativo): consegna \textbf{obbligatoria} per far valutare gli altri lab;
        è valutato solo valido/non valido. Se si consegna una soluzione ``ragionevole''
        \textbf{entro il 1° dicembre} si ottiene \textbf{+1 punto bonus} sul voto finale.
        Se va riconsegnato, o si consegna dopo il 1° dicembre, il bonus si perde.
  \item \textbf{Altri lab}: scadenza \textbf{10 giorni prima della data di un appello} (es.\ appello
        il 15 $\Rightarrow$ consegna entro il 5); il testo del Lab~1 parla invece di due settimane
        (vedi riquadro iniziale).
  \item Si valuta \textbf{il report}; il notebook si guarda solo in caso di dubbi. Il report deve
        quindi essere \textbf{autocontenuto}.
  \item Ogni report vale \textbf{da 0 a 32}; \textbf{tutti} devono avere voto $\geq 18$.
        \begin{itemize}
          \item voto $\geq 18$: \textbf{non} si può riconsegnare per migliorarlo;
          \item voto $< 18$: si \textbf{deve} riconsegnare.
        \end{itemize}
  \item I voti dei report valgono \textbf{2 anni accademici} (fino all'ultimo appello 2027/2028).
  \item Report e scritto possono essere fatti in appelli diversi; l'esame si registra quando
        entrambe le parti sono sufficienti.
\end{itemize}

\begin{esame}[come si scrive il report]
Le indicazioni del docente su come scrivere un buon report e la checklist prima della consegna sono
raccolte nel cap.~\ref{cap:lab-report}.
\end{esame}

\subsection{Uso di ChatGPT e strumenti simili}
È \textbf{permesso} usare l'AI per idee, bozze di codice e struttura del report. Restate però
responsabili di verificare e capire tutto ciò che consegnate: niente copia-incolla. Se la consegna
sembra interamente ``GPT-made'', senza ragionamento, il gruppo viene \textbf{chiamato per un orale}.
I testi dei lab aggiungono che saper trovare e correggere bug, anche nel codice generato dall'AI,
può essere valutato allo scritto, dove c'è una parte di comprensione del codice da fare senza AI.

\section{Esame scritto}
\begin{itemize}
  \item In aula, sulla piattaforma \textbf{Exam (Moodle)} con \textbf{LockDown Browser},
        \textbf{con il proprio portatile} (obbligatorio). Aggiorna LockDown Browser all'ultima versione.
  \item Copre \textbf{teoria e codice}: non si chiede di scrivere codice da zero, ma di capire se
        un codice funziona, perché no e cosa correggere.
  \item \textbf{Durata: 100 minuti} (in genere si finisce in 70--80).
  \item \textbf{Totale: 32 punti} secondo la slide (33 a voce, vedi riquadro iniziale).
\end{itemize}

\textbf{Struttura (descritta in aula):}
\begin{enumerate}
  \item domande a \textbf{scelta multipla}, due delle quali \textbf{senza penalità};
  \item \textbf{quattro gruppi di vero/falso}: risposta giusta = punto, risposta sbagliata = \textbf{piccola penalità};
  \item \textbf{due domande aperte} di teoria;
  \item \textbf{tre esercizi aperti} di difficoltà crescente sulla pipeline di ML e sul codice.
\end{enumerate}

\begin{verifica}
Dalla registrazione non è chiaro se le domande a scelta multipla siano \textbf{due in tutto} o se solo
due siano senza penalità; la slide dice soltanto ``\en{some with penalties for wrong answers}''.
\end{verifica}

\begin{esame}[cosa conta davvero]
Il docente ha ripetuto più volte che \textbf{non chiederà di ricordare le formule} delle slide:
conta capire \emph{perché} funzionano (es.\ perché il momentum stabilizza la traiettoria,
perché la sigmoid causa il \en{vanishing gradient}). Le formule nel manuale servono a
ragionare, non a essere memorizzate. Esempi di domande nei cap.~\ref{cap:quiz} e~\ref{cap:esercizi}.
\end{esame}

\section{Voto finale}
\[
\text{Voto} = 0{,}6 \cdot \text{Scritto} + 0{,}4 \cdot \overline{\text{Report}} + \text{punti aggiuntivi}
\]
\begin{itemize}
  \item Si arrotonda all'intero più vicino \textbf{solo alla fine} ($\geq 0{,}5$ per eccesso).
  \item Registrazione solo se \textbf{scritto $\geq 18$} \textbf{e ogni report $\geq 18$}.
  \item Punti aggiuntivi: +1 del Lab 1 consegnato entro il 1° dicembre ed eventuali altre attività proposte.
  \item \textbf{Orale facoltativo} se il voto è $\geq 26$, oppure su richiesta del docente in caso
        di dubbi sullo scritto o sul report.
\end{itemize}

\begin{esempio}[calcolo, dalla slide]
$24 \cdot 0{,}6 + 28 \cdot 0{,}4 + 1 = 14{,}4 + 11{,}2 + 1 = 26{,}6 \Rightarrow 27$\\
$20 \cdot 0{,}6 + 28 \cdot 0{,}4 + 1 = 12 + 11{,}2 + 1 = 24{,}2 \Rightarrow 24$
\end{esempio}

\begin{attenzione}[errore nella slide]
La slide riporta $23{,}2 \Rightarrow 23$ per il secondo esempio; il conto corretto è
$24{,}2 \Rightarrow 24$.
\end{attenzione}

\section{Mappa lezioni $\leftrightarrow$ slide}
\begin{center}
\small
\begin{tabularx}{\textwidth}{l X l}
\toprule
\textbf{Lezione} & \textbf{Argomenti} & \textbf{Slide}\\
\midrule
\rif{2026-09-21}{21/09} & Presentazione; perché AI in cybersecurity; AI/ML/NN/DL; storia; shallow vs deep & 0, 1 (pp.\ 1--25)\\
\rif{2026-09-23}{23/09} & Terminologia; paradigmi di apprendimento; uso dei modelli; AI \& cybersecurity & 1 (pp.\ 26--54)\\
\rif{2026-09-28}{28/09} & Perceptron, bias, attivazioni, stacking, non linearità, approssimazione universale, logit/softmax & 2 (pp.\ 1--18)\\
\rif{2026-09-30}{30/09} & Training: epoch, mini-batch, learning rate, GD/SGD/momentum/Adam, backprop, loss, metriche, over/underfitting & 2 (pp.\ 19--51)\\
\rif{2026-10-05}{05/10} & Data handling: pulizia, encoding, scaling, split, sbilanciamento, leakage; progettare una FFNN; tensori & 3 (pp.\ 1--22)\\
\rif{2026-10-06}{06/10} & FFNN in PyTorch: classe, training loop, mini-batch, early stopping, dropout, weight decay, GPU & 3 (pp.\ 23--33)\\
\rif{2026-10-12}{12/10} & RNN: hidden state, architetture, BPTT, gradienti, LSTM, dati 3D & 3 (pp.\ 34--71)\\
\rif{2026-10-14}{14/10} & RNN in PyTorch (padding); CNN; GNN (teoria) & 3 (pp.\ 72--123)\\
\rif{2026-10-19}{19/10} & GNN in PyTorch Geometric; feature categoriche ed embedding appresi & 3 (pp.\ 124--136), 4 (pp.\ 1--22)\\
\rif{2026-10-20}{20/10} & Embedding supervisionati e self-supervised; dropout e batch norm nel dettaglio; linee guida report & 4 (pp.\ 23--51), 2a\\
\rif{2026-11-02}{02/11} & Anomaly detection: definizioni, tipi, sfide; one-class SVM; clustering e k-means & 5 (pp.\ 1--33)\\
\rif{2026-11-09}{09/11} & Notebook di anomaly detection su KDD99: preprocessing, PCA/t-SNE, OC-SVM, LOF, k-means, DBSCAN, Isolation Forest, autoencoder & 5 (pp.\ 33--57), 6\\
\bottomrule
\end{tabularx}
\end{center}

\chapter{Quiz: scelta multipla e vero/falso}
\label{cap:quiz}
\setcounter{domanda}{0}

Domande nello stile delle prime due parti dello scritto (cap.~\ref{cap:formato}): scelta multipla e
gruppi di vero/falso. \textbf{Le soluzioni sono in fondo al capitolo} (sez.~\ref{sec:quiz-soluzioni}),
ciascuna con il rimando alla lezione in cui l'argomento è trattato: prova prima da solo.

\begin{nota}[Domande dei redattori]
Sono domande scritte per il ripasso a partire da lezioni e slide, \textbf{non} domande ufficiali
d'esame. Il formato reale si vedrà nella simulazione d'esame dell'ultima settimana.
\end{nota}

\section{Domande a scelta multipla}
Una sola risposta corretta per domanda.

\begin{domanda}[AI e ML]\label{quiz:m1}
Quale affermazione su AI e machine learning è corretta?
\begin{opzioni}
  \item AI e machine learning sono sinonimi.
  \item Il machine learning non richiede dati: la regola la fornisce l'esperto.
  \item Un sistema a regole scritte da un esperto è AI ma non machine learning.
  \item Il deep learning è un sottoinsieme dell'AI che non usa reti neurali.
\end{opzioni}
\end{domanda}

\begin{domanda}[shallow vs deep]\label{quiz:m2}
Cosa distingue principalmente il deep learning dallo shallow learning, secondo il corso?
\begin{opzioni}
  \item La feature extraction è fatta dalla rete; lo sforzo si sposta sul model engineering.
  \item Il deep learning non richiede preprocessing.
  \item Il deep learning è più interpretabile.
  \item Il deep learning si applica solo alle immagini.
\end{opzioni}
\end{domanda}

\begin{domanda}[terminologia]\label{quiz:m3}
Il validation set serve a:
\begin{opzioni}
  \item aggiornare i pesi della rete;
  \item aumentare la quantità di dati di training;
  \item stimare le prestazioni finali dopo aver fissato tutti gli iperparametri;
  \item scegliere gli iperparametri e valutare l'overfitting.
\end{opzioni}
\end{domanda}

\begin{domanda}[paradigmi]\label{quiz:m4}
Qual è la differenza tra unsupervised e self-supervised learning?
\begin{opzioni}
  \item Il self-supervised richiede label umane, l'unsupervised no.
  \item Il self-supervised genera un segnale di supervisione dai dati stessi per imparare una rappresentazione.
  \item L'unsupervised impara una rappresentazione, il self-supervised descrive i dati.
  \item Sono sinonimi.
\end{opzioni}
\end{domanda}

\begin{domanda}[federated learning]\label{quiz:m5}
Nel federated learning, il \en{model poisoning} consiste nel:
\begin{opzioni}
  \item iniettare rumore impercettibile nelle immagini di test;
  \item ricostruire i dati di training a partire dal modello;
  \item far inviare a un partecipante un modello addestrato per far passare come legittimo il proprio traffico malevolo;
  \item aggirare i guardrail di un LLM con un prompt.
\end{opzioni}
\end{domanda}

\begin{domanda}[adversarial ML]\label{quiz:m6}
Un \en{adversarial example}:
\begin{opzioni}
  \item modifica leggermente l'input perché il modello sbagli, senza toccare il modello;
  \item modifica i pesi del modello durante il training;
  \item ricostruisce i dati di training dagli output del modello;
  \item è un campione con la label sbagliata.
\end{opzioni}
\end{domanda}

\begin{domanda}[perceptron]\label{quiz:m7}
A cosa serve il bias $w_0$ in un perceptron?
\begin{opzioni}
  \item A rendere la funzione non lineare.
  \item A normalizzare l'input.
  \item A ridurre il numero di parametri.
  \item A traslare la funzione, che altrimenti passerebbe sempre per l'origine.
\end{opzioni}
\end{domanda}

\begin{domanda}[non linearità]\label{quiz:m8}
Una rete con 10 layer lineari e nessuna funzione di attivazione:
\begin{opzioni}
  \item può approssimare qualsiasi funzione continua;
  \item equivale a un unico layer lineare;
  \item è più espressiva di una rete con un solo hidden layer ReLU;
  \item soffre di over-smoothing.
\end{opzioni}
\end{domanda}

\begin{domanda}[output layer]\label{quiz:m9}
Classificazione con 5 classi: quale output layer?
\begin{opzioni}
  \item 1 neurone con sigmoid.
  \item 5 neuroni con ReLU.
  \item 5 neuroni con softmax.
  \item 1 neurone con tanh.
\end{opzioni}
\end{domanda}

\begin{domanda}[ottimizzazione]\label{quiz:m10}
Rispetto al gradient descent, lo stochastic gradient descent:
\begin{opzioni}
  \item aggiorna i pesi dopo ogni mini-batch e, grazie al rumore, può uscire dai minimi locali;
  \item usa tutto il training set per ogni aggiornamento;
  \item garantisce di trovare il minimo globale;
  \item non richiede un learning rate.
\end{opzioni}
\end{domanda}

\begin{domanda}[momentum]\label{quiz:m11}
Con un coefficiente di momentum $\beta$ molto vicino a 1:
\begin{opzioni}
  \item l'aggiornamento segue quasi solo il gradiente del batch corrente;
  \item il learning rate effettivo diventa nullo;
  \item l'ottimizzatore equivale all'SGD senza momentum;
  \item l'aggiornamento segue soprattutto la direzione passata: traiettoria liscia ma lenta a cambiare direzione.
\end{opzioni}
\end{domanda}

\begin{domanda}[Adam]\label{quiz:m12}
La \en{bias correction} di Adam serve a:
\begin{opzioni}
  \item correggere i bias dei neuroni;
  \item compensare il fatto che $m$ e $v$ partono da 0 e sono sottostimati nei primi passi;
  \item aumentare il learning rate nel tempo;
  \item prevenire l'overfitting.
\end{opzioni}
\end{domanda}

\begin{domanda}[vanishing gradient]\label{quiz:m13}
Perché la sigmoid negli hidden layer di una rete profonda è problematica?
\begin{opzioni}
  \item Non è derivabile.
  \item Non è limitata e fa esplodere il gradiente.
  \item Le sue derivate sono piccole e, moltiplicate layer dopo layer, fanno svanire il gradiente.
  \item Aggiunge troppi parametri.
\end{opzioni}
\end{domanda}

\begin{domanda}[metriche]\label{quiz:m14}
Dataset con 9.800 campioni benigni e 200 malware. Un classificatore che risponde sempre ``benigno'' ha:
\begin{opzioni}
  \item accuracy 98\% e recall sul malware pari a 0;
  \item accuracy 2\% e recall sul malware pari a 1;
  \item precision sul malware pari a 1;
  \item F1 sul malware pari a 0,98.
\end{opzioni}
\end{domanda}

\begin{domanda}[learning curve]\label{quiz:m15}
Learning curve: la loss di training scende quasi a 0, quella di validation scende fino all'epoch 30 e poi risale. Cosa fai?
\begin{opzioni}
  \item Aumento il numero di neuroni.
  \item Aumento il learning rate.
  \item Continuo il training finché anche la validation loss arriva a 0.
  \item Uso l'early stopping e ripristino i pesi salvati attorno all'epoch 30.
\end{opzioni}
\end{domanda}

\begin{domanda}[regolarizzazione]\label{quiz:m16}
La regolarizzazione L2:
\begin{opzioni}
  \item aggiunge alla loss $\lambda\sum_i|\phi_i|$, spingendo molti pesi a 0;
  \item aggiunge alla loss $\lambda\sum_i\phi_i^2$, scoraggiando i pesi grandi;
  \item spegne neuroni a caso durante il training;
  \item normalizza gli input di ogni layer.
\end{opzioni}
\end{domanda}

\begin{domanda}[scaling]\label{quiz:m17}
Quando la standardizzazione (z-score) è preferibile alla normalizzazione min-max?
\begin{opzioni}
  \item Quando servono valori nell'intervallo $[0,1]$.
  \item Per le feature categoriche.
  \item Quando ci sono outlier.
  \item Solo per la label.
\end{opzioni}
\end{domanda}

\begin{domanda}[scaling e leakage]\label{quiz:m18}
Su quali dati si esegue il \code{fit} dello \code{StandardScaler}?
\begin{opzioni}
  \item Solo sul training set.
  \item Su tutto il dataset.
  \item Su training e validation.
  \item Sul test set.
\end{opzioni}
\end{domanda}

\begin{domanda}[feature categoriche]\label{quiz:m19}
La porta TCP di destinazione di un flusso è meglio trattata come feature:
\begin{opzioni}
  \item numerica standardizzata;
  \item numerica normalizzata min-max;
  \item label;
  \item categorica.
\end{opzioni}
\end{domanda}

\begin{domanda}[cross-validation]\label{quiz:m20}
Con la k-fold cross-validation, una volta scelti gli iperparametri, il modello finale:
\begin{opzioni}
  \item è quello del fold con la validation migliore;
  \item si ri-addestra su tutti i fold con gli iperparametri scelti;
  \item è la media dei pesi dei $k$ modelli;
  \item si addestra sul test set.
\end{opzioni}
\end{domanda}

\begin{domanda}[data leakage]\label{quiz:m21}
Quale di queste pratiche introduce data leakage?
\begin{opzioni}
  \item Stratificare lo split.
  \item Usare il dropout.
  \item Includere tra le feature un identificativo univoco del campione.
  \item Fare early stopping sulla validation loss.
\end{opzioni}
\end{domanda}

\begin{domanda}[PyTorch: loss]\label{quiz:m22}
Con \code{nn.CrossEntropyLoss}, l'ultimo layer del modello deve restituire:
\begin{opzioni}
  \item i logit grezzi;
  \item le probabilità dopo la softmax;
  \item vettori one-hot;
  \item le probabilità dopo la sigmoid.
\end{opzioni}
\end{domanda}

\begin{domanda}[PyTorch: training loop]\label{quiz:m23}
In un training loop con mini-batch, \code{optimizer.zero\_grad()}, \code{loss.backward()} e \code{optimizer.step()} vanno:
\begin{opzioni}
  \item fuori dal ciclo sui batch, una volta per epoch;
  \item solo nel ciclo di validation;
  \item prima del ciclo sulle epoch;
  \item dentro il ciclo sui mini-batch.
\end{opzioni}
\end{domanda}

\begin{domanda}[PyTorch: eval]\label{quiz:m24}
Cosa fa \code{model.eval()}?
\begin{opzioni}
  \item Disattiva il calcolo dei gradienti.
  \item Mette dropout e batch norm in modalità di valutazione.
  \item Congela definitivamente i pesi.
  \item Calcola accuracy e F1.
\end{opzioni}
\end{domanda}

\begin{domanda}[RNN]\label{quiz:m25}
Classificare un'email (disponibile per intero) come phishing a partire dalla sequenza delle sue parole è un caso:
\begin{opzioni}
  \item one-to-many, solo con RNN mono-direzionale;
  \item one-to-one;
  \item many-to-one, in cui si può usare una RNN bidirezionale;
  \item many-to-many in tempo reale.
\end{opzioni}
\end{domanda}

\begin{domanda}[LSTM]\label{quiz:m26}
Nell'LSTM, il \en{forget gate}:
\begin{opzioni}
  \item decide con una sigmoid quanto del cell state precedente mantenere;
  \item aggiunge nuova informazione al cell state con una tanh;
  \item produce direttamente l'output della rete;
  \item azzera l'hidden state a ogni passo.
\end{opzioni}
\end{domanda}

\begin{domanda}[tensori RNN]\label{quiz:m27}
Un batch di 16 serie univariate lunghe 50, con \code{batch\_first=True}, deve avere forma:
\begin{opzioni}
  \item \code{(16, 50)};
  \item \code{(50, 16)};
  \item \code{(16, 1, 50)};
  \item \code{(16, 50, 1)}.
\end{opzioni}
\end{domanda}

\begin{domanda}[CNN]\label{quiz:m28}
In una CNN, il max pooling:
\begin{opzioni}
  \item aumenta la dimensione spaziale;
  \item riduce la dimensione sostituendo ogni finestra con il suo massimo;
  \item aggiunge pesi addestrabili;
  \item sostituisce la ReLU come non linearità.
\end{opzioni}
\end{domanda}

\begin{domanda}[GNN]\label{quiz:m29}
Perché una FFNN applicata alla matrice di adiacenza non è adatta ai grafi?
\begin{opzioni}
  \item Perché i grafi non hanno feature.
  \item Perché è troppo lenta da addestrare.
  \item Perché la rappresentazione dipende dall'ordine dei nodi e non gestisce grafi di dimensioni diverse.
  \item Perché non usa funzioni di attivazione.
\end{opzioni}
\end{domanda}

\begin{domanda}[GNN]\label{quiz:m30}
L'over-smoothing in una GNN è:
\begin{opzioni}
  \item l'appiattimento degli embedding dei nodi, che diventano simili perché i receptive field si sovrappongono;
  \item una loss di training che oscilla;
  \item un'esplosione dei gradienti;
  \item la perdita di memoria del passato lontano.
\end{opzioni}
\end{domanda}

\begin{domanda}[GNN]\label{quiz:m31}
GraphSAGE si distingue perché:
\begin{opzioni}
  \item usa coefficienti di attenzione appresi;
  \item usa tutta la matrice di adiacenza normalizzata per il grado;
  \item non condivide i pesi tra i nodi;
  \item campiona un numero fisso di vicini.
\end{opzioni}
\end{domanda}

\begin{domanda}[PyTorch Geometric]\label{quiz:m32}
Nel training di una GNN per graph classification, la loss del batch va pesata per:
\begin{opzioni}
  \item \code{batch.x.size(0)};
  \item \code{batch.num\_graphs};
  \item il numero di archi;
  \item la hidden size.
\end{opzioni}
\end{domanda}

\begin{domanda}[embedding]\label{quiz:m33}
Per una feature nominale ad alta cardinalità (es.\ il dominio del mittente) in una rete neurale, la scelta migliore è:
\begin{opzioni}
  \item label encoding in ordine alfabetico;
  \item one-hot encoding;
  \item un embedding appreso, con vocabolario costruito sul training e token \code{<UNK>};
  \item eliminarla.
\end{opzioni}
\end{domanda}

\begin{domanda}[embedding]\label{quiz:m34}
\code{nn.Embedding(V, H)}:
\begin{opzioni}
  \item crea una lookup addestrabile con $V$ righe e $H$ colonne, che riceve indici interi;
  \item crea una matrice non addestrabile con $H$ righe e $V$ colonne;
  \item riceve in input le categorie come numeri float;
  \item si può usare anche con un random forest.
\end{opzioni}
\end{domanda}

\begin{domanda}[dropout]\label{quiz:m35}
Dropout con $p = 0{,}5$ in training: un neurone che sopravvive con output 3 produce:
\begin{opzioni}
  \item 3;
  \item 1,5;
  \item 0;
  \item 6.
\end{opzioni}
\end{domanda}

\begin{domanda}[batch normalization]\label{quiz:m36}
In validation, la batch normalization usa:
\begin{opzioni}
  \item media e varianza del batch di validation;
  \item running mean e running variance memorizzate durante il training;
  \item nessuna normalizzazione;
  \item media e varianza del test set.
\end{opzioni}
\end{domanda}

\begin{domanda}[tipi di anomalia]\label{quiz:m37}
Backup di 2 ore la domenica, quando l'azienda è chiusa; nei giorni lavorativi backup di 2 ore sono normali. È un'anomalia:
\begin{opzioni}
  \item puntuale;
  \item collettiva;
  \item contestuale;
  \item non è un'anomalia.
\end{opzioni}
\end{domanda}

\begin{domanda}[one-class SVM]\label{quiz:m38}
Nella one-class SVM non supervisionata, $\nu$ rappresenta:
\begin{opzioni}
  \item la frazione attesa di outlier;
  \item il numero di cluster;
  \item il grado del kernel polinomiale;
  \item il learning rate.
\end{opzioni}
\end{domanda}

\begin{domanda}[k-means]\label{quiz:m39}
Nel grafico SSE in funzione di $K$ (k-means), l'SSE risale passando da $K=5$ a $K=6$. Causa probabile?
\begin{opzioni}
  \item $K=6$ è il valore ottimo.
  \item I dati non hanno struttura a cluster.
  \item La distanza euclidea non è adatta.
  \item Un'inizializzazione sfortunata (es.\ un cluster vuoto): va rieseguito k-means.
\end{opzioni}
\end{domanda}

\begin{domanda}[anomaly detection]\label{quiz:m40}
Vantaggio di un algoritmo di anomaly detection che restituisce \textbf{score} invece di label:
\begin{opzioni}
  \item non richiede dati di training;
  \item permette di ordinare le istanze e scegliere un budget di analisi;
  \item è sempre più accurato;
  \item elimina i falsi positivi.
\end{opzioni}
\end{domanda}

\begin{domanda}[PCA]\label{quiz:m41}
Su dati non scalati, la PCA dice che una sola componente spiega il 99,9\% della varianza. Causa più probabile?
\begin{opzioni}
  \item Il dataset è davvero unidimensionale.
  \item Ci sono troppe feature categoriche in one-hot.
  \item Una feature numerica con scala enorme (es.\ byte) domina la varianza: manca lo scaling.
  \item La PCA andava sostituita con t-SNE.
\end{opzioni}
\end{domanda}

\begin{domanda}[autoencoder]\label{quiz:m42}
Perché un autoencoder addestrato solo su dati normali rileva le anomalie meglio di uno addestrato su normali e anomalie?
\begin{opzioni}
  \item Con le anomalie in training impara a ricostruire anche quelle e il loro errore di ricostruzione si avvicina a quello dei normali.
  \item Con meno dati il training converge sempre a un minimo migliore.
  \item Solo con i dati normali si può usare la cross-entropy.
  \item Le anomalie in training causano sempre vanishing gradient.
\end{opzioni}
\end{domanda}

\begin{domanda}[Isolation Forest]\label{quiz:m43}
Quale di questi metodi non si basa su distanze o densità, e soffre quindi meno l'alta dimensionalità?
\begin{opzioni}
  \item LOF.
  \item DBSCAN.
  \item K-means con distanza dal centroide.
  \item Isolation Forest.
\end{opzioni}
\end{domanda}

\begin{domanda}[LOF]\label{quiz:m44}
Un attacco forma un gruppo piccolo ma molto denso di connessioni quasi identiche. Con $k=10$ i suoi punti hanno LOF $\approx 1$. Significa che:
\begin{opzioni}
  \item i punti sono outlier certi;
  \item hanno una densità simile ai vicini, quindi LOF non li segnala;
  \item il modello non è stato addestrato;
  \item la contaminazione è troppo alta.
\end{opzioni}
\end{domanda}


\begin{domanda}[bag of words]\label{quiz:m45}
Quale affermazione sulla bag of words (count vectorizer) è corretta?
\begin{opzioni}
  \item Conserva l'ordine delle parole ma non la loro frequenza.
  \item Rappresenta ogni frase con i conteggi delle parole del vocabolario: perde ordine e semantica, e la dimensione cresce con il vocabolario.
  \item Richiede una rete neurale per essere calcolata.
  \item Assegna a parole con significato simile vettori vicini.
\end{opzioni}
\end{domanda}

\begin{domanda}[Word2Vec]\label{quiz:m46}
Nello skip-gram di Word2Vec la rete impara a:
\begin{opzioni}
  \item predire la parola centrale a partire dalle parole del contesto;
  \item classificare le frasi come positive o negative;
  \item predire le parole del contesto a partire dalla parola centrale;
  \item ricostruire la frase intera a partire dal suo embedding.
\end{opzioni}
\end{domanda}

\begin{domanda}[Word2Vec]\label{quiz:m47}
Nella probabilità $p(w_{t+j} \mid w_t) = \exp(e_{t+j}^{\top} e_t) / \sum_{w=1}^{V} \exp(e_w^{\top} e_t)$, cosa rende il calcolo impraticabile e come lo si risolve?
\begin{opzioni}
  \item Il numeratore, perché il prodotto scalare è costoso; lo si sostituisce con la distanza euclidea.
  \item Il denominatore, che somma su tutto il vocabolario a ogni passo; lo si approssima con $k$ parole negative estratte a caso.
  \item Il contesto $c$, che deve coprire tutta la frase; lo si riduce a una parola.
  \item L'esponenziale, che va in overflow; lo si toglie.
\end{opzioni}
\end{domanda}

\begin{domanda}[Word2Vec]\label{quiz:m48}
Un corpus contiene in parti uguali testi di zoologia, favole e libri di informatica. Che embedding avrà la parola \emph{mouse} in Word2Vec?
\begin{opzioni}
  \item Due embedding, uno per l'animale e uno per il dispositivo.
  \item Un embedding diverso a seconda della frase in cui la parola compare.
  \item Un solo embedding, spostato verso il significato più frequente (l'animale), che non rappresenta bene nessuno dei due.
  \item Nessun embedding, perché la parola è ambigua.
\end{opzioni}
\end{domanda}

\begin{domanda}[darknet]\label{quiz:m49}
Perché tutto il traffico ricevuto da una darknet è interessante?
\begin{opzioni}
  \item Perché contiene il traffico degli utenti legittimi verso i servizi dell'ateneo.
  \item Perché la darknet risponde agli scanner e ne registra le richieste complete.
  \item Perché i suoi IP non offrono servizi né inviano richieste: ogni pacchetto è non richiesto (scansioni, backscattering, errori).
  \item Perché ogni pacchetto arriva con una label che indica l'attaccante.
\end{opzioni}
\end{domanda}

\begin{domanda}[DarkVec]\label{quiz:m50}
In DarkVec, che cosa corrisponde alle ``parole'' e alle ``frasi'' di Word2Vec?
\begin{opzioni}
  \item Parole = porte di destinazione; frasi = pacchetti di un singolo IP.
  \item Parole = IP sorgente; frasi = sequenze di IP sorgente in ordine di arrivo, una per servizio (porta/protocollo).
  \item Parole = byte del payload; frasi = sessioni TCP complete.
  \item Parole = IP di destinazione della darknet; frasi = giorni di osservazione.
\end{opzioni}
\end{domanda}

\begin{domanda}[i-DarkVec]\label{quiz:m51}
Quale limite di Word2Vec affronta i-DarkVec?
\begin{opzioni}
  \item L'impossibilità di usare il negative sampling.
  \item Gli IP nuovi (fuori vocabolario) e i comportamenti che cambiano, aggiornando gli embedding giorno per giorno a partire da quelli già noti.
  \item La necessità di label per addestrare gli embedding.
  \item La dimensione eccessiva degli embedding rispetto al one-hot.
\end{opzioni}
\end{domanda}

\section{Vero o falso}
Come all'esame: risposta giusta = punto, risposta sbagliata = piccola penalità. Se non sei sicuro,
valuta se rischiare. Per ogni affermazione indica V o F.

\begin{domanda}[vero/falso -- Training]\label{quiz:vfA}
Gruppo A -- Training.
\begin{enumerate}[label=A\arabic*.,nosep,leftmargin=2.6em]
  \item Un'epoch corrisponde a un singolo aggiornamento dei pesi anche quando si usano i mini-batch. \hfill V\quad F
  \item I mini-batch si campionano senza reinserimento, così ogni punto è usato una sola volta per epoch. \hfill V\quad F
  \item Una loss di training pari a 0 indica che il modello è pronto per la produzione. \hfill V\quad F
  \item Un learning rate scheduler parte tipicamente da valori alti e li riduce durante il training. \hfill V\quad F
\end{enumerate}
\end{domanda}

\begin{domanda}[vero/falso -- Dati]\label{quiz:vfB}
Gruppo B -- Dati.
\begin{enumerate}[label=B\arabic*.,nosep,leftmargin=2.6em]
  \item Il label encoding è adatto alla label di classificazione perché la cross-entropy non usa distanze tra classi. \hfill V\quad F
  \item Il one-hot encoding introduce un ordine tra le categorie. \hfill V\quad F
  \item Con dati temporali, uno split casuale può causare data leakage. \hfill V\quad F
  \item Con classi molto sbilanciate, la media \en{micro} delle metriche dà lo stesso peso a ogni classe. \hfill V\quad F
\end{enumerate}
\end{domanda}

\begin{domanda}[vero/falso -- PyTorch]\label{quiz:vfC}
Gruppo C -- PyTorch.
\begin{enumerate}[label=C\arabic*.,nosep,leftmargin=2.6em]
  \item \code{torch.no\_grad()} spegne il dropout. \hfill V\quad F
  \item Nel DataLoader di validation lo shuffle non serve. \hfill V\quad F
  \item Se il modello è sulla GPU, anche i tensori dei batch devono essere spostati sullo stesso device. \hfill V\quad F
  \item Con \code{nn.CrossEntropyLoss} le label devono essere necessariamente vettori one-hot di tipo float. \hfill V\quad F
\end{enumerate}
\end{domanda}

\begin{domanda}[vero/falso -- RNN e CNN]\label{quiz:vfD}
Gruppo D -- RNN e CNN.
\begin{enumerate}[label=D\arabic*.,nosep,leftmargin=2.6em]
  \item In una RNN la matrice $W_{hh}$ è diversa per ogni passo temporale. \hfill V\quad F
  \item Una RNN bidirezionale è adatta alla predizione in tempo reale. \hfill V\quad F
  \item Il gradient clipping è un rimedio all'exploding gradient. \hfill V\quad F
  \item In una CNN i valori dei filtri convoluzionali sono pesi appresi durante il training. \hfill V\quad F
\end{enumerate}
\end{domanda}

\begin{domanda}[vero/falso -- GNN ed embedding]\label{quiz:vfE}
Gruppo E -- GNN ed embedding.
\begin{enumerate}[label=E\arabic*.,nosep,leftmargin=2.6em]
  \item Una GNN con aggregazione a somma è invariante all'ordine dei vicini. \hfill V\quad F
  \item Le GNN non possono produrre embedding per nodi mai visti in training. \hfill V\quad F
  \item Il vocabolario per gli embedding va costruito su training e validation insieme, per ridurre i \code{<UNK>}. \hfill V\quad F
  \item Gli embedding appresi con supervisione sono ottimi per il task di training ma possono essere sub-ottimi per altri task. \hfill V\quad F
\end{enumerate}
\end{domanda}

\begin{domanda}[vero/falso -- Anomaly detection]\label{quiz:vfF}
Gruppo F -- Anomaly detection.
\begin{enumerate}[label=F\arabic*.,nosep,leftmargin=2.6em]
  \item Una novità può diventare la nuova normalità. \hfill V\quad F
  \item Gli algoritmi shallow riconoscono le anomalie collettive senza bisogno di feature engineering. \hfill V\quad F
  \item K-means assegna ogni punto a un cluster, quindi non produce direttamente ``punti di rumore''. \hfill V\quad F
  \item Nella one-class SVM ``supervisionata'' il modello riceve le label normale/anomalo. \hfill V\quad F
\end{enumerate}
\end{domanda}

\begin{domanda}[vero/falso -- Anomaly detection in pratica]\label{quiz:vfG}
Gruppo G -- Anomaly detection in pratica.
\begin{enumerate}[label=G\arabic*.,nosep,leftmargin=2.6em]
  \item Nell'anomaly detection le label disponibili si usano per addestrare il modello. \hfill V\quad F
  \item Nella mappa di t-SNE la distanza tra due gruppi misura quanto sono diversi nello spazio originale. \hfill V\quad F
  \item DBSCAN etichetta come rumore i punti che non sono né core né border point. \hfill V\quad F
  \item Su dati molto sbilanciati una ROC-AUC alta garantisce un buon F1 sulla classe delle anomalie. \hfill V\quad F
\end{enumerate}
\end{domanda}


\begin{domanda}[vero/falso -- NLP e DarkVec]\label{quiz:vfH}
Gruppo H -- NLP e DarkVec.
\begin{enumerate}[label=H\arabic*.,nosep,leftmargin=2.6em]
  \item Word2Vec ha bisogno di frasi etichettate a mano per imparare gli embedding. \hfill V\quad F
  \item Una finestra di contesto $c$ più grande porta più informazione ma richiede più calcolo. \hfill V\quad F
  \item Word2Vec produce embedding contestuali: la stessa parola ha vettori diversi in frasi diverse. \hfill V\quad F
  \item In DarkVec l'ipotesi è che gli IP che compaiono insieme nel tempo sugli stessi servizi abbiano un comportamento simile. \hfill V\quad F
\end{enumerate}
\end{domanda}

\section{Soluzioni}
\label{sec:quiz-soluzioni}


\begin{soluzione}[Domanda~\ref{quiz:m1}]
\textbf{(c)}. L'AI include i sistemi a regole; il ML impara dai dati. (rif.\ \rif{2026-09-21}{21/09})
\end{soluzione}

\begin{soluzione}[Domanda~\ref{quiz:m2}]
\textbf{(a)}. Il preprocessing resta; i modelli deep sono meno interpretabili. (rif.\ \rif{2026-09-21}{21/09})
\end{soluzione}

\begin{soluzione}[Domanda~\ref{quiz:m3}]
\textbf{(d)}. (c) è il ruolo del test set. (rif.\ \rif{2026-09-23}{23/09})
\end{soluzione}

\begin{soluzione}[Domanda~\ref{quiz:m4}]
\textbf{(b)}. Unsupervised = descrizione; self-supervised = rappresentazione con supervisione dai dati. (rif.\ \rif{2026-09-23}{23/09})
\end{soluzione}

\begin{soluzione}[Domanda~\ref{quiz:m5}]
\textbf{(c)}. (b) è model inversion, (d) jailbreak, (a) adversarial example. (rif.\ \rif{2026-09-23}{23/09})
\end{soluzione}

\begin{soluzione}[Domanda~\ref{quiz:m6}]
\textbf{(a)}. (rif.\ \rif{2026-09-23}{23/09})
\end{soluzione}

\begin{soluzione}[Domanda~\ref{quiz:m7}]
\textbf{(d)}. Senza bias, con attivazione lineare, ogni retta passa per l'origine. (rif.\ \rif{2026-09-28}{28/09})
\end{soluzione}

\begin{soluzione}[Domanda~\ref{quiz:m8}]
\textbf{(b)}. Prodotto di matrici = una matrice: il modello resta lineare. (rif.\ \rif{2026-09-28}{28/09})
\end{soluzione}

\begin{soluzione}[Domanda~\ref{quiz:m9}]
\textbf{(c)}. (rif.\ \rif{2026-09-28}{28/09})
\end{soluzione}

\begin{soluzione}[Domanda~\ref{quiz:m10}]
\textbf{(a)}. L'SGD non garantisce il minimo globale: rende solo più probabile uscire da quelli locali. (rif.\ \rif{2026-09-30}{30/09})
\end{soluzione}

\begin{soluzione}[Domanda~\ref{quiz:m11}]
\textbf{(d)}. $m_{t+1} = \beta m_t + (1-\beta)g$: con $\beta\to1$ il gradiente nuovo pesa pochissimo. (rif.\ \rif{2026-09-30}{30/09})
\end{soluzione}

\begin{soluzione}[Domanda~\ref{quiz:m12}]
\textbf{(b)}. (rif.\ \rif{2026-09-30}{30/09})
\end{soluzione}

\begin{soluzione}[Domanda~\ref{quiz:m13}]
\textbf{(c)}. La derivata della sigmoid vale al massimo 0,25. (rif.\ \rif{2026-09-30}{30/09})
\end{soluzione}

\begin{soluzione}[Domanda~\ref{quiz:m14}]
\textbf{(a)}. $9800/10000 = 98\%$; nessun malware trovato $\Rightarrow$ recall 0; precision 0/0 non definita. (rif.\ \rif{2026-09-30}{30/09})
\end{soluzione}

\begin{soluzione}[Domanda~\ref{quiz:m15}]
\textbf{(d)}. Il modello dell'ultima epoch è già in overfitting: serve lo stato salvato. (rif.\ \rif{2026-09-30}{30/09})
\end{soluzione}

\begin{soluzione}[Domanda~\ref{quiz:m16}]
\textbf{(b)}. (a) è L1; (c) dropout; (d) batch norm. (rif.\ \rif{2026-09-30}{30/09})
\end{soluzione}

\begin{soluzione}[Domanda~\ref{quiz:m17}]
\textbf{(c)}. (rif.\ \rif{2026-10-05}{05/10})
\end{soluzione}

\begin{soluzione}[Domanda~\ref{quiz:m18}]
\textbf{(a)}. (rif.\ \rif{2026-10-05}{05/10})
\end{soluzione}

\begin{soluzione}[Domanda~\ref{quiz:m19}]
\textbf{(d)}. 80, 81, 8080: la distanza numerica non ha significato. (rif.\ \rif{2026-10-05}{05/10})
\end{soluzione}

\begin{soluzione}[Domanda~\ref{quiz:m20}]
\textbf{(b)}. Scegliere il fold migliore significa sfruttare la fortuna sul fold di validation. (rif.\ \rif{2026-10-05}{05/10})
\end{soluzione}

\begin{soluzione}[Domanda~\ref{quiz:m21}]
\textbf{(c)}. Un ID permette di memorizzare i campioni; lo split temporale casuale è un altro esempio. (rif.\ \rif{2026-10-05}{05/10})
\end{soluzione}

\begin{soluzione}[Domanda~\ref{quiz:m22}]
\textbf{(a)}. La softmax è dentro la loss; applicarla due volte appiattisce i gradienti. (rif.\ \rif{2026-10-06}{06/10})
\end{soluzione}

\begin{soluzione}[Domanda~\ref{quiz:m23}]
\textbf{(d)}. Fuori dal ciclo si farebbe un solo aggiornamento per epoch (GD con gradienti sommati). (rif.\ \rif{2026-10-06}{06/10})
\end{soluzione}

\begin{soluzione}[Domanda~\ref{quiz:m24}]
\textbf{(b)}. \code{no\_grad} spegne i gradienti; \code{eval} cambia il comportamento di dropout e BN. (rif.\ \rif{2026-10-06}{06/10})
\end{soluzione}

\begin{soluzione}[Domanda~\ref{quiz:m25}]
\textbf{(c)}. L'email è completa: si può leggere anche all'indietro; un output per tutta la sequenza. (rif.\ \rif{2026-10-12}{12/10})
\end{soluzione}

\begin{soluzione}[Domanda~\ref{quiz:m26}]
\textbf{(a)}. (rif.\ \rif{2026-10-12}{12/10})
\end{soluzione}

\begin{soluzione}[Domanda~\ref{quiz:m27}]
\textbf{(d)}. Ogni istante deve essere un vettore di feature (qui di dimensione 1): \code{x.unsqueeze(-1)}. (rif.\ \rif{2026-10-12}{12/10})
\end{soluzione}

\begin{soluzione}[Domanda~\ref{quiz:m28}]
\textbf{(b)}. (rif.\ \rif{2026-10-14}{14/10})
\end{soluzione}

\begin{soluzione}[Domanda~\ref{quiz:m29}]
\textbf{(c)}. Riordinando i nodi cambia la matrice; un nodo in più richiede un neurone in più. (rif.\ \rif{2026-10-14}{14/10})
\end{soluzione}

\begin{soluzione}[Domanda~\ref{quiz:m30}]
\textbf{(a)}. (rif.\ \rif{2026-10-14}{14/10})
\end{soluzione}

\begin{soluzione}[Domanda~\ref{quiz:m31}]
\textbf{(d)}. (a) è GAT; (b) GCN. (rif.\ \rif{2026-10-14}{14/10})
\end{soluzione}

\begin{soluzione}[Domanda~\ref{quiz:m32}]
\textbf{(b)}. \code{batch.x.size(0)} conta i nodi (es.\ 288), non i grafi (2). (rif.\ \rif{2026-10-19}{19/10})
\end{soluzione}

\begin{soluzione}[Domanda~\ref{quiz:m33}]
\textbf{(c)}. (rif.\ \rif{2026-10-19}{19/10})
\end{soluzione}

\begin{soluzione}[Domanda~\ref{quiz:m34}]
\textbf{(a)}. È l'equivalente efficiente di one-hot + layer lineare; random forest non può imparare embedding. (rif.\ \rif{2026-10-19}{19/10})
\end{soluzione}

\begin{soluzione}[Domanda~\ref{quiz:m35}]
\textbf{(d)}. I sopravvissuti sono scalati di $1/(1-p) = 2$: $3\cdot 2 = 6$; il valore atteso resta 3. (rif.\ \rif{2026-10-20}{20/10})
\end{soluzione}

\begin{soluzione}[Domanda~\ref{quiz:m36}]
\textbf{(b)}. (rif.\ \rif{2026-10-20}{20/10})
\end{soluzione}

\begin{soluzione}[Domanda~\ref{quiz:m37}]
\textbf{(c)}. Il valore è normale, è il contesto (giorno festivo) a renderlo anomalo. (rif.\ \rif{2026-11-02}{02/11})
\end{soluzione}

\begin{soluzione}[Domanda~\ref{quiz:m38}]
\textbf{(a)}. (rif.\ \rif{2026-11-02}{02/11})
\end{soluzione}

\begin{soluzione}[Domanda~\ref{quiz:m39}]
\textbf{(d)}. L'SSE dovrebbe scendere con $K$: una risalita indica una cattiva inizializzazione. (rif.\ \rif{2026-11-02}{02/11})
\end{soluzione}

\begin{soluzione}[Domanda~\ref{quiz:m40}]
\textbf{(b)}. (rif.\ \rif{2026-11-02}{02/11})
\end{soluzione}

\begin{soluzione}[Domanda~\ref{quiz:m41}]
\textbf{(c)}. Varianza in byte (milioni) contro varianza in millisecondi: va applicato \code{StandardScaler} prima della PCA. (rif.\ \rif{2026-11-09}{09/11})
\end{soluzione}

\begin{soluzione}[Domanda~\ref{quiz:m42}]
\textbf{(a)}. Il modello si specializza sul normale solo se vede soltanto il normale; con anomalie in training le ricostruisce. (rif.\ \rif{2026-11-09}{09/11})
\end{soluzione}

\begin{soluzione}[Domanda~\ref{quiz:m43}]
\textbf{(d)}. (rif.\ \rif{2026-11-09}{09/11})
\end{soluzione}

\begin{soluzione}[Domanda~\ref{quiz:m44}]
\textbf{(b)}. Un attacco etichettato non è per forza un outlier di densità: servirebbe un $k$ più grande del gruppo. (rif.\ \rif{2026-11-09}{09/11})
\end{soluzione}

\begin{soluzione}[Domanda~\ref{quiz:m45}]
\textbf{(b)}. È la baseline usata nel Lab~2 sulle API call: solida, ma senza ordine né semantica. (rif.\ \rif{2026-11-23}{23/11})
\end{soluzione}

\begin{soluzione}[Domanda~\ref{quiz:m46}]
\textbf{(c)}. La (a) è il CBOW. (rif.\ \rif{2026-11-23}{23/11})
\end{soluzione}

\begin{soluzione}[Domanda~\ref{quiz:m47}]
\textbf{(b)}. È una softmax su $V$ classi; il negative sampling tiene il contributo positivo e usa $k \ll V$ negativi (più $k$ è grande, migliore l'approssimazione). (rif.\ \rif{2026-11-23}{23/11})
\end{soluzione}

\begin{soluzione}[Domanda~\ref{quiz:m48}]
\textbf{(c)}. Word2Vec è \textbf{non contestuale}: un vettore per parola. La (b) descrive un modello contestuale (transformer). (rif.\ \rif{2026-11-23}{23/11})
\end{soluzione}

\begin{soluzione}[Domanda~\ref{quiz:m49}]
\textbf{(c)}. La darknet non risponde, quindi vede solo IP e porte, senza label. (rif.\ \rif{2026-11-23}{23/11})
\end{soluzione}

\begin{soluzione}[Domanda~\ref{quiz:m50}]
\textbf{(b)}. Una frase per servizio, perché le porte contano: IP che sondano servizi diversi non vanno messi nello stesso contesto. (cap.~\ref{cap:darkvec}, slide 8)
\end{soluzione}

\begin{soluzione}[Domanda~\ref{quiz:m51}]
\textbf{(b)}. (cap.~\ref{cap:darkvec}, slide 8)
\end{soluzione}

\begin{soluzione}[Domanda~\ref{quiz:vfA} (gruppo A)]
\textbf{A1 F, A2 V, A3 F, A4 V}.\\
A1: l'epoch è il passaggio di tutto il dataset; con i mini-batch ci sono molti aggiornamenti per epoch. A3: loss 0 è sospetta (overfitting). (rif.\ \rif{2026-09-30}{30/09})
\end{soluzione}

\begin{soluzione}[Domanda~\ref{quiz:vfB} (gruppo B)]
\textbf{B1 V, B2 F, B3 V, B4 F}.\\
B2: i vettori one-hot sono equidistanti. B4: è la media \en{macro} a pesare ugualmente le classi. (rif.\ \rif{2026-10-05}{05/10})
\end{soluzione}

\begin{soluzione}[Domanda~\ref{quiz:vfC} (gruppo C)]
\textbf{C1 F, C2 V, C3 V, C4 F}.\\
C1: è \code{model.eval()} a spegnere il dropout. C4: si passano gli indici di classe interi (\code{torch.long}). (rif.\ \rif{2026-10-06}{06/10})
\end{soluzione}

\begin{soluzione}[Domanda~\ref{quiz:vfD} (gruppo D)]
\textbf{D1 F, D2 F, D3 V, D4 V}.\\
D1: i pesi sono condivisi tra i passi. D2: richiede la sequenza completa, il futuro non è disponibile in tempo reale. (rif.\ \rif{2026-10-12}{12/10}, \rif{2026-10-14}{14/10})
\end{soluzione}

\begin{soluzione}[Domanda~\ref{quiz:vfE} (gruppo E)]
\textbf{E1 V, E2 F, E3 F, E4 V}.\\
E2: le GNN sono induttive. E3: il vocabolario va costruito solo sul training (leakage); gli sconosciuti vanno in \code{<UNK>}. (rif.\ \rif{2026-10-14}{14/10}, \rif{2026-10-19}{19/10}, \rif{2026-10-20}{20/10})
\end{soluzione}

\begin{soluzione}[Domanda~\ref{quiz:vfF} (gruppo F)]
\textbf{F1 V, F2 F, F3 V, F4 F}.\\
F2: servono feature aggregate su finestre. F4: si addestra solo su dati normali, ma senza passare label. (rif.\ \rif{2026-11-02}{02/11})
\end{soluzione}

\begin{soluzione}[Domanda~\ref{quiz:vfG} (gruppo G)]
\textbf{G1 F, G2 F, G3 V, G4 F}.\\
G1: le label servono solo a validare (anomalie note). G2: t-SNE preserva i vicini, non le distanze. G4: con forte sbilanciamento la ROC-AUC può ingannare come l'accuracy: guarda precision, recall e F1. (rif.\ \rif{2026-11-09}{09/11})
\end{soluzione}

\begin{soluzione}[Domanda~\ref{quiz:vfH} (gruppo H)]
\textbf{H1 F, H2 V, H3 F, H4 V}.\\
H1: è self-supervised, la supervisione viene dal contesto delle parole. H3: è non contestuale; contestuali saranno i transformer. (rif.\ \rif{2026-11-23}{23/11})
\end{soluzione}

\chapter{Domande aperte ed esercizi}
\label{cap:esercizi}

\begin{nota}[Esercizi dei redattori, non domande ufficiali]
Domande ed esercizi di questo capitolo sono stati scritti da noi per il ripasso, a partire da lezioni
e slide: \textbf{non} sono domande ufficiali d'esame. Riproducono le tipologie della seconda parte
dello scritto (cap.~\ref{cap:formato}): due domande aperte di teoria e tre esercizi di difficoltà
crescente su pipeline e codice. Le domande a scelta multipla e vero/falso sono nel cap.~\ref{cap:quiz}.
Per una simulazione realistica: 100 minuti, senza appunti, insieme ai quiz.
\end{nota}

Le soluzioni sono in fondo al capitolo (sez.~\ref{sec:esercizi-soluzioni}), ognuna con il rimando
alla lezione in cui l'argomento è trattato: prova prima da solo.

\section{Domande aperte di teoria}
Rispondi in 10--15 righe, con esempi di cybersecurity dove possibile.
\begin{enumerate}[label=\textbf{A\arabic*.},ref=A\arabic*,leftmargin=2.6em,itemsep=6pt]

  \item\label{es:a1} Spiega la differenza tra shallow e deep learning. In quali casi di cybersecurity preferiresti lo shallow?
  \item\label{es:a2} Descrivi il ruolo di training, validation e test set. Cosa succede se usi il test set per scegliere gli iperparametri?
  \item\label{es:a3} Perché le funzioni di attivazione non lineari sono indispensabili? Quali useresti negli hidden layer e in output per regressione, classificazione binaria e multi-classe?
  \item\label{es:a4} Confronta gradient descent, SGD, momentum e Adam: quale problema risolve ciascuno, e cosa introduce?
  \item\label{es:a5} Spiega il vanishing gradient: cause, effetti e rimedi, sia nelle FFNN sia nelle RNN.
  \item\label{es:a6} Come riconosci underfitting e overfitting dalle learning curve? Elenca i rimedi per ciascuno.
  \item\label{es:a7} Cos'è il data leakage? Fai tre esempi diversi e spiega come evitarli.
  \item\label{es:a8} Devi classificare flussi di rete in cui gli attacchi sono l'1\%. Come gestisci lo sbilanciamento (dati, algoritmo, metriche)?
  \item\label{es:a9} Perché le RNN faticano con le dipendenze a lungo termine e come le mitiga l'LSTM?
  \item\label{es:a10} Perché servono le GNN? Spiega permutation invariance ed equivariance e come il message passing le garantisce.
  \item\label{es:a11} Confronta ID sequenziali, feature engineering, one-hot ed embedding per codificare una feature nominale.
  \item\label{es:a12} Differenze tra embedding supervisionati e self-supervised. Perché il self-supervised è interessante per la cybersecurity?
  \item\label{es:a13} Quali tipi di anomalia esistono? Perché l'anomaly detection è difficile in cybersecurity?
  \item\label{es:a14} One-class SVM: differenze tra uso ``supervisionato'' e non supervisionato; ruolo di $\nu$ e del kernel.
  \item\label{es:a15} Confronta k-means, DBSCAN, LOF e Isolation Forest per l'anomaly detection: come producono uno score, punti deboli, quando li useresti.
  \item\label{es:a16} Come usi un autoencoder per l'anomaly detection su dati con feature numeriche e categoriche? Loss, dati di training, soglia; perché una one-class SVM nello spazio latente può andare peggio dell'errore di ricostruzione?
  \item\label{es:a17} Ti viene mostrata la formula $\frac{1}{T}\sum_{t}\sum_{-c\le j\le c,\, j\ne 0}\log p(w_{t+j}\mid w_t)$ con $p(w_{t+j}\mid w_t) = \exp(e_{t+j}^{\top}e_t)/\sum_{w}\exp(e_w^{\top}e_t)$. Da dove viene, cosa rappresenta ogni termine, qual è l'intuizione? Perché in pratica si usa il negative sampling?
  \item\label{es:a18} Quali sono i limiti di Word2Vec? Per ciascuno fai un esempio sul linguaggio naturale e uno di cybersecurity.
  \item\label{es:a19} DarkVec: perché una darknet è un buon punto di osservazione, come si costruiscono le ``frasi'' e su quale ipotesi si basa il metodo? Funzionerebbe sul traffico di tutti i giorni?
\end{enumerate}

\section{Esercizi su pipeline, calcoli e codice}

\subsection*{Esercizio 1 --- Contare i parametri}\label{es:e1}
Dataset di flussi con 12 feature numeriche e la feature categorica \code{protocol} con 5 valori possibili. Rete: due hidden layer
da 64 e 32 neuroni (ReLU), output per 3 classi.
\begin{enumerate}[label=(\alph*)]
  \item Con \code{protocol} in one-hot: quanti neuroni di input e quanti parametri in totale?
  \item E con \code{protocol} in label encoding? Perché comunque non è una buona idea?
  \item E con un \code{nn.Embedding(5, 4)} per \code{protocol}?
\end{enumerate}

\subsection*{Esercizio 2 --- Metriche}\label{es:e2}
Test set: 950 benigni e 50 malware. Il modello predice correttamente 35 malware e 910 benigni; 15 malware sono predetti benigni e 40
benigni sono predetti malware.
\begin{enumerate}[label=(\alph*)]
  \item Calcola accuracy, precision, recall e F1 per la classe malware.
  \item Calcola F1 per la classe benigna e la F1 macro.
  \item Confronta con un classificatore che dice sempre ``benigno''. Quale metrica ti fa scegliere il modello giusto?
\end{enumerate}

\subsection*{Esercizio 3 --- Backpropagation a mano}\label{es:e3}
Modello $\hat y = \phi_1\phi_2 x$, loss $\mathcal{L} = (\phi_1\phi_2x - y)^2$, con $\phi_1 = 2$, $\phi_2 = -1$, $x=3$, $y=1$.
\begin{enumerate}[label=(\alph*)]
  \item Calcola forward ($a$, $b$, $c$, $\mathcal{L}$) usando le variabili intermedie del cap.~\ref{cap:addestramento}.
  \item Calcola $\partial\mathcal{L}/\partial\phi_1$ e $\partial\mathcal{L}/\partial\phi_2$ con la regola della catena.
  \item Fai un passo di gradient descent con $\alpha = 0{,}01$ e verifica che la loss diminuisce.
\end{enumerate}

\subsection*{Esercizio 4 --- Trova gli errori (FFNN)}\label{es:e4}
Il codice seguente ``gira'' ma contiene diversi errori o cattive pratiche. Trovane almeno sei e correggili.
\begin{lstlisting}
scaler = StandardScaler()
X_train = scaler.fit_transform(X_train)
X_val = scaler.fit_transform(X_val)

train_loader = DataLoader(TensorDataset(Xtr, ytr), batch_size=64, shuffle=True)
val_loader   = DataLoader(TensorDataset(Xva, yva), batch_size=64, shuffle=True)

class Net(nn.Module):
    def __init__(self, d_in, n_cls):
        super().__init__()
        self.fc1 = nn.Linear(d_in, 64)
        self.fc2 = nn.Linear(64, n_cls)
    def forward(self, x):
        x = self.fc1(x)
        return torch.softmax(self.fc2(x), dim=1)

model = Net(20, 3)
criterion = nn.CrossEntropyLoss()
optimizer = optim.Adam(model.parameters(), lr=1e-3)

for epoch in range(50):
    model.train()
    optimizer.zero_grad()
    for xb, yb in train_loader:
        loss = criterion(model(xb), yb)
        loss.backward()
    optimizer.step()
    with torch.no_grad():
        val_loss = sum(criterion(model(xb), yb).item() for xb, yb in val_loader)
\end{lstlisting}

\subsection*{Esercizio 5 --- Forme dei tensori in una RNN}\label{es:e5}
Un batch contiene 3 serie con 2 feature per istante e lunghezze 4, 6 e 5. Si usa
\code{nn.LSTM(input\_size=2, hidden\_size=32, num\_layers=2, batch\_first=True, bidirectional=True)}.
\begin{enumerate}[label=(\alph*)]
  \item Che forma ha il batch dopo la \code{collate\_fn} con \code{pad\_sequence}? Quanti passi di padding ha ogni serie?
  \item Che forma ha \code{h\_n}? Che dimensione di input deve avere il layer lineare finale?
  \item Perché \code{out[:, -1, :]} è sbagliato per classificare la prima serie?
  \item Hai 16 serie \textbf{univariate} di lunghezza 100 in un tensore \code{(16, 100)}: cosa fai prima di passarle a una RNN?
\end{enumerate}

\subsection*{Esercizio 6 --- Momentum e bias correction}\label{es:e6}
$\beta = 0{,}9$, $m_0 = 0$, gradienti $g_1 = 2$, $g_2 = -1$, $g_3 = 2$.
\begin{enumerate}[label=(\alph*)]
  \item Calcola $m_1, m_2, m_3$ con $m_{t+1} = \beta m_t + (1-\beta)g_{t+1}$.
  \item Calcola le versioni corrette $\hat m_t = m_t / (1-\beta^t)$.
  \item Cosa mostrano i numeri su traiettoria e bias correction?
\end{enumerate}

\subsection*{Esercizio 7 --- Progetta la pipeline}\label{es:e7}
Hai 1 milione di flussi di rete raccolti in 30 giorni, con colonne: \code{timestamp}, IP sorgente, IP destinazione, porta di
destinazione, protocollo (TCP/UDP/ICMP), byte inviati (con alcuni valori enormi), durata, numero di pacchetti, label (99\% benigno,
1\% suddiviso in 5 tipi di attacco). Descrivi preprocessing, split, architettura, loss, metriche e i controlli che faresti per
evitare scorciatoie.

\subsection*{Esercizio 8 --- Receptive field e batch di grafi}\label{es:e8}
Grafo non orientato con archi A--B, B--C, C--D, D--E, B--F.
\begin{enumerate}[label=(\alph*)]
  \item Quali nodi influenzano l'embedding di A con 1, 2, 3 e 4 layer GNN? Cosa rischi con 4 layer?
  \item Quante colonne ha \code{edge\_index} in PyTorch Geometric per questo grafo?
  \item Un batch contiene due grafi da 126 e 162 nodi con 5 feature per nodo: che forma ha \code{batch.x}? E l'output di
        \code{global\_mean\_pool} con hidden size 64? Per quanto pesi la loss del batch?
\end{enumerate}

\subsection*{Esercizio 9 --- Batch norm e dropout a mano}\label{es:e9}
\begin{enumerate}[label=(\alph*)]
  \item Una feature vale $[1, 3, 5, 7]$ nel mini-batch; $\gamma = 2$, $\beta = 1$. Calcola l'output della batch norm per il valore 5
        (usa la varianza del batch, $\epsilon \approx 0$).
  \item Running mean e running variance partono da 0 e 1, con aggiornamento $\text{running} = 0{,}9\cdot\text{running} + 0{,}1\cdot\text{batch}$.
        Quanto valgono dopo questo batch? Che output darebbe il valore 5 in \code{eval()} con queste statistiche? Cosa ne deduci?
  \item Dropout con $p = 0{,}2$: un neurone ha output 4. Quali valori può produrre in training e con quale probabilità? E in validation?
\end{enumerate}

\subsection*{Esercizio 10 --- Scegliere $K$}\label{es:e10}
SSE di k-means per $K = 1,\ldots,8$: 1000, 540, 250, 205, 185, 230, 160, 150. Quale $K$ scegli e perché? Come spieghi il valore per
$K=6$?

\subsection*{Esercizio 11 --- Leakage in un numero}\label{es:e11}
Una feature nel training va da 10 a 50; nel validation compare il valore 70. Con min-max fittato sul training, quanto vale 70 dopo la
trasformazione? Cosa cambierebbe fittando su training + validation, e perché è sbagliato?

\subsection*{Esercizio 12 --- Trova gli errori (embedding + RNN)}\label{es:e12}
Classificazione di sequenze di IP (lunghezza variabile, padding con 0 nella \code{collate\_fn}). Trova almeno cinque problemi.
\begin{lstlisting}
vocab = {ip: i for i, ip in enumerate(sorted(set(all_ips)))}   # IP di tutto il dataset

class Model(nn.Module):
    def __init__(self, V, E, H, C):
        super().__init__()
        self.emb = nn.Embedding(V, E)
        self.rnn = nn.LSTM(1, H, batch_first=True, bidirectional=True)
        self.fc = nn.Linear(H, C)
    def forward(self, ids, lengths):
        e = self.emb(ids)
        out, (h, c) = self.rnn(e)
        return self.fc(out[:, -1, :])
\end{lstlisting}

\subsection*{Esercizio 13 --- Trova gli errori (pipeline di anomaly detection)}\label{es:e13}
Flussi di rete con label (\code{normal} o tipo di attacco, circa 6\% di attacchi). Trova almeno sei problemi.
\begin{lstlisting}
X = df.drop(columns=["label"])                      # la label originale non serve piu'
y = (df["label"] != "normal").astype(int)
X = StandardScaler().fit_transform(pd.get_dummies(X))
X_train, X_test, y_train, y_test = train_test_split(X, y, test_size=0.2)

ae = Autoencoder(X.shape[1], 8)
train_autoencoder(ae, X_train)                      # tutto il training set
scores = reconstruction_scores(ae, X_test)
threshold = np.percentile(scores, 94)               # 6% di anomalie attese
print("ROC-AUC:", roc_auc_score(y_test, scores))    # unica metrica riportata
\end{lstlisting}

\subsection*{Esercizio 14 --- Silhouette e LOF a mano}\label{es:e14}
\begin{enumerate}[label=(\alph*)]
  \item Un punto ha distanza media 2 dagli altri punti del suo cluster e 5 dal cluster più vicino. Calcola la silhouette.
  \item Un altro punto ha $a = 4$ e $b = 3$. Calcola la silhouette e interpretala.
  \item Un punto $p$ ha $\mathrm{LRD}(p) = 0{,}5$; la LRD media dei suoi $k$ vicini è 2. Calcola il LOF e interpretalo.
\end{enumerate}

\subsection*{Esercizio 15 --- Le frasi di DarkVec}\label{es:e15}
Una darknet riceve, in quest'ordine, i pacchetti: IP1 $\to$ 22/TCP, IP2 $\to$ 445/TCP, IP3 $\to$ 445/TCP, IP4 $\to$ 22/TCP,
IP1 $\to$ 22/TCP, IP3 $\to$ 445/TCP, IP2 $\to$ 445/TCP.
\begin{enumerate}[label=(\alph*)]
  \item Scrivi la ``frase'' unica con tutto il traffico e le frasi per servizio.
  \item Con finestra $c = 1$, elenca le coppie (centrale, contesto) dello skip-gram per la frase di 22/TCP.
  \item Perché la frase unica porterebbe a conclusioni sbagliate?
  \item Il giorno dopo compare IP5. Cosa succede con DarkVec e cosa propone i-DarkVec?
\end{enumerate}

% =====================================================================


\section{Soluzioni}
\label{sec:esercizi-soluzioni}

\subsection*{Domande aperte: tracce di risposta}

\begin{soluzione}[A1 -- Shallow vs deep]
Shallow: feature engineering manuale (sforzo umano, più tentativi), modelli più semplici e
        \textbf{interpretabili}; deep: feature extraction automatica, sforzo sul model engineering, più prestazioni e flessibilità
        (serie, grafi, immagini), ma black box. Preferisco lo shallow con dati tabellari, pochi dati, quando serve spiegare la decisione
        (forense, regole di firewall) o quando basta una soglia (picco evidente): non sparare a una mosca con un cannone. Lo shallow è
        anche la baseline per giudicare un modello deep.
(rif.\ \rif{2026-09-21}{21/09})
\end{soluzione}

\begin{soluzione}[A2 -- Train/val/test]
Training: impara i pesi. Validation: sceglie gli iperparametri e rileva l'overfitting, senza mai
        aggiornare i pesi. Test: verifica una sola volta che le prestazioni siano simili a quelle di validation (generalità). Se uso il test per
        scegliere, la stima finale diventa ottimistica: ho ``studiato sull'esame'' e non ho più una misura indipendente della generalizzazione.
(rif.\ \rif{2026-09-23}{23/09})
\end{soluzione}

\begin{soluzione}[A3 -- Non linearità]
Senza attivazioni la composizione di layer è un unico layer lineare (collasso), e il mondo non è lineare.
        Hidden: ReLU (non limitata, evita il vanishing, sparsità); tanh nelle RNN. Output: regressione lineare; binaria sigmoid (1 neurone);
        multi-classe softmax ($C$ neuroni). In PyTorch con \code{CrossEntropyLoss} la softmax è nella loss.
(rif.\ \rif{2026-09-28}{28/09})
\end{soluzione}

\begin{soluzione}[A4 -- Ottimizzatori]
GD: usa tutti i dati, lento, si blocca nei minimi locali, sensibile all'inizializzazione. SGD: mini-batch,
        aggiornamenti frequenti, meno memoria, il rumore aiuta a uscire dai minimi; introduce traiettorie rumorose (scheduler del LR).
        Momentum: media mobile delle direzioni, traiettoria liscia; introduce $\beta$. Normalizzazione del gradiente: pari importanza a tutte
        le coordinate, ma oscillazioni. Adam: momentum su gradiente e quadrato + bias correction; stabile e veloce; iperparametri
        $\alpha,\beta,\gamma$.
(rif.\ \rif{2026-09-30}{30/09})
\end{soluzione}

\begin{soluzione}[A5 -- Vanishing gradient]
Con la backprop il gradiente di un layer è il prodotto dei gradienti successivi; attivazioni limitate
        (sigmoid, tanh) hanno derivate piccole $\Rightarrow$ il prodotto tende a 0 e i primi layer non imparano. Nelle RNN il tempo si comporta come
        profondità (fattori ripetuti di $W_{hh}$) e il passato lontano viene dimenticato. Rimedi: ReLU (FFNN, CNN), buona inizializzazione,
        architetture (skip connection in ResNet, LSTM/GRU nelle RNN). L'opposto, l'exploding gradient, si gestisce con il clipping.
(rif.\ \rif{2026-09-30}{30/09}, \rif{2026-10-12}{12/10})
\end{soluzione}

\begin{soluzione}[A6 -- Learning curve]
Underfitting: train e val alte e poco decrescenti (attenzione alla scala), metriche pessime $\Rightarrow$ più
        capacità, architettura diversa, più epoch, preprocessing migliore. Overfitting: train verso 0 e val che risale $\Rightarrow$ L1/L2
        (weight decay), dropout, early stopping con salvataggio del modello migliore, batch norm, rete più piccola; se forte, ripartire dai dati.
(rif.\ \rif{2026-09-30}{30/09}, \rif{2026-10-06}{06/10})
\end{soluzione}

\begin{soluzione}[A7 -- Leakage]
Informazioni che dal training ``contaminano'' la valutazione. Esempi: (1) scaler o vocabolario fittati su tutto il
        dataset $\Rightarrow$ fit solo sul training; (2) split casuale di dati temporali (interpolare $t$ tra $t-1$ e $t+1$) $\Rightarrow$ split temporale
        a finestra scorrevole o crescente; (3) feature-scorciatoia o identificativi (DoS sempre sulla porta 80, matricola) $\Rightarrow$ togliere ID,
        verificare che il modello regga senza la feature sospetta.
(rif.\ \rif{2026-10-05}{05/10})
\end{soluzione}

\begin{soluzione}[A8 -- Sbilanciamento]
Split stratificato; dati: oversampling della minoritaria (SMOTE, ADASYN) solo sul training,
        undersampling della maggioritaria (anche cluster-based); algoritmo: class weights nella loss, focal loss; metriche: precision, recall,
        F1 per classe e macro, non accuracy (99\% predicendo sempre benigno). Recall alta sugli attacchi con precision accettabile.
        Se gli attacchi sono troppo pochi o nuovi, valutare l'anomaly detection.
(rif.\ \rif{2026-10-05}{05/10})
\end{soluzione}

\begin{soluzione}[A9 -- RNN e memoria]
La BPTT moltiplica molti fattori lungo il tempo: il gradiente dei primi passi svanisce (o esplode), quindi la
        rete ricorda solo il passato recente (``parlo fluentemente\ldots'' dopo un'ora di racconto). L'LSTM aggiunge un cell state (memoria lunga)
        regolato da gate sigmoid (forget, input, output) e una memoria candidata tanh: il flusso del gradiente lungo il cell state è in parte
        ininterrotto. Costo: più pesi e calcoli.
(rif.\ \rif{2026-10-12}{12/10})
\end{soluzione}

\begin{soluzione}[A10 -- GNN]
I grafi non hanno ordine né dimensione fissa: una FFNN sulla matrice di adiacenza cambia input se riordino i nodi e non
        gestisce nodi nuovi; CNN e RNN assumono griglie o sequenze. Invarianza: la rappresentazione del grafo non cambia con l'ordinamento.
        Equivarianza: la rappresentazione di un nodo nella stessa posizione non cambia. Il message passing le garantisce perché i pesi del messaggio
        sono condivisi tra tutti i nodi e l'aggregazione (somma, media, max) non dipende dall'ordine. Bonus: induttività.
(rif.\ \rif{2026-10-14}{14/10})
\end{soluzione}

\begin{soluzione}[A11 -- Encoding nominale]
ID sequenziali: un solo neurone ma impongono un ordine falso (bias). Feature engineering: dà una chiave
        di lettura con vincolo debole, ma può essere sbagliata, cambiare o essere impossibile. One-hot: nessun bias, semplice, va bene anche per lo
        shallow, ma esplode con molte categorie e non gestisce nuove categorie. Embedding: denso, compatto, semantico, appreso sul task; solo reti
        neurali, serve un task e abbastanza dati, \code{<UNK>} per i valori nuovi.
(rif.\ \rif{2026-10-19}{19/10})
\end{soluzione}

\begin{soluzione}[A12 -- Supervised vs self-supervised]
Supervised: la rappresentazione è guidata dalla label: ottima per quel task, ma sub-ottima per
        altri e a rischio di scorciatoie (husky/lupo). Self-supervised: il segnale viene dai dati (ricostruzione con autoencoder, contrastive
        learning con augmentation): rappresentazioni più generiche, poi fine-tuning con poche label. In cybersecurity le label sono poche e costose
        e le minacce evolvono: rappresentazioni generali aiutano a generalizzare.
(rif.\ \rif{2026-10-20}{20/10})
\end{soluzione}

\begin{soluzione}[A13 -- Anomalie]
Puntuali, contestuali, collettive (con esempi). Difficoltà: ignote fino a quando accadono, rare e sbilanciate (e un
        errore su un'anomalia costa di più), eterogenee, in evoluzione; un'anomalia può essere una novità benigna. Sfide ML: compromesso recall
        vs falsi positivi, dati ad alta dimensionalità e dipendenti, poche label, rumore, anomalie complesse, bisogno di spiegazioni.
(rif.\ \rif{2026-11-02}{02/11})
\end{soluzione}

\begin{soluzione}[A14 -- One-class SVM]
Impara un confine attorno ai dati normali; fuori = anomalia. ``Supervisionata'': addestrata solo su
        campioni normali (scelti grazie alle label, che il modello non vede) $\Rightarrow$ novelty detection; richiede normali affidabili.
        Non supervisionata: dati mescolati, assume anomalie rare e ben separate $\Rightarrow$ outlier detection; soffre il rumore. $\nu$ = frazione
        attesa di outlier: più è alto, più il confine si stringe sulla zona densa. Il kernel (RBF il più usato) rende separabili dati non lineari.
(rif.\ \rif{2026-11-02}{02/11})
\end{soluzione}

\begin{soluzione}[A15 -- Clustering, densità e alberi]
K-means: score = distanza dal centroide (più cluster piccoli come segnale), veloce e
        scalabile, ma $K$ da scegliere, cluster sferici, sensibile all'inizializzazione, ogni punto in un cluster: un attacco con un cluster
        tutto suo sfugge. DBSCAN: il rumore (né core né border) è l'anomalia, niente $K$, forme arbitrarie, ma $\varepsilon$ critico e densità
        diverse problematiche; nessuno score graduato. LOF: rapporto tra densità dei vicini e del punto, score continuo, anomalie
        \textbf{locali}; sensibile a $k$, costoso; gruppi piccoli ma densi hanno LOF $\approx 1$. Isolation Forest: lunghezza media del percorso
        in alberi casuali, niente distanze, scala ed è robusto in alta dimensionalità; meno interpretabile. I tre metodi a distanza soffrono la
        maledizione della dimensionalità: ridurre le dimensioni o preferire l'Isolation Forest.
(rif.\ \rif{2026-11-02}{02/11}, \rif{2026-11-09}{09/11})
\end{soluzione}

\begin{soluzione}[A16 -- Autoencoder]
Loss = MSE sulle numeriche (scalate) + cross-entropy su ogni blocco one-hot; lo stesso errore per campione
        è lo score. Addestramento sui \textbf{soli normali} (le label servono solo a sceglierli): l'errore dei normali scende, quello delle
        anomalie resta alto. Soglia da training o validation (es.\ percentile pari alla contaminazione), mai dal test. Con anomalie in
        training il modello impara a ricostruirle e la separazione si riduce. Nello spazio latente l'encoder, addestrato a ricostruire il
        normale, avvicina le anomalie ai normali simili: l'informazione anomala è quella persa nel bottleneck e finisce nell'errore di
        ricostruzione, quindi un OC-SVM su $Z$ può perdere recall.
(rif.\ \rif{2026-11-09}{09/11})
\end{soluzione}

\begin{soluzione}[A17 -- Skip-gram]
È l'obiettivo dello skip-gram di Word2Vec. Per ogni parola $w_t$ della frase (somma su $t$, media su $T$) si considerano le parole nella
finestra $[-c, c]$ (somma su $j \ne 0$) e si massimizza la log-probabilità che compaiano dato $w_t$. La probabilità è una softmax: al
numeratore il prodotto scalare tra embedding di contesto e centrale (alto se i vettori sono paralleli, cioè portano lo stesso
significato), al denominatore lo stesso score con tutte le $V$ parole. Intuizione: ipotesi distribuzionale, il significato di una parola
viene dalle compagnie; è una classificazione con label ricavate dal testo (self-supervised). Il denominatore costa $O(V)$ a ogni passo:
il negative sampling tiene il termine positivo e usa $k \ll V$ parole casuali come negativi. $c$ e $k$ sono compromessi tra informazione
e calcolo.
(rif.\ \rif{2026-11-23}{23/11})
\end{soluzione}

\begin{soluzione}[A18 -- Limiti di Word2Vec]
(1) Non contestuale: un vettore per parola (\emph{mouse} animale o dispositivo; un IP che cambia comportamento). (2) Parole nuove (OOV):
``petaloso'' non è nel vocabolario; un IP mai visto non ha embedding, serve riaddestrare o un approccio incrementale. (3) Parole rare:
composte tedesche come \emph{Sitzpinkler}; un IP o un termine tecnico raro avrà un embedding povero (i modelli successivi usano token
sotto-parola). (4) Dalle parole alle frasi: i compiti utili (sentiment, traduzione, codice malevolo o no) sono su frasi e Word2Vec non
dice come aggregare gli embedding.
(rif.\ \rif{2026-11-23}{23/11})
\end{soluzione}

\begin{soluzione}[A19 -- DarkVec]
La darknet non offre servizi e non risponde: ogni pacchetto è non richiesto (scansioni, backscattering, errori), quindi chi la contatta ha
uno scopo. Dati: solo IP e porte (categorici), enormi, senza label, mutevoli $\Rightarrow$ embedding self-supervised. Frasi: sequenze di
IP sorgente in ordine di arrivo, una per servizio (porta/protocollo); parole = IP. Ipotesi: gli IP che compaiono insieme nel tempo sugli
stessi servizi si comportano in modo simile (botnet coordinate). Sul traffico ordinario la co-occorrenza temporale è dominata da utenti e
servizi legittimi non correlati: l'ipotesi non è garantita e va verificata (domanda indicata come possibile dal docente nelle slide).
(rif.\ \rif{2026-11-23}{23/11}; cap.~\ref{cap:darkvec})
\end{soluzione}

\subsection*{Esercizi}

\begin{soluzione}[Esercizio 1 -- Contare i parametri]
(a) Input $12 + 5 = 17$. Parametri: $17\cdot64+64 = 1152$; $64\cdot32+32 = 2080$; $32\cdot3+3 = 99$; totale \textbf{3331}.
(b) Input 13: $13\cdot64+64 = 896$, totale 3075. Ma TCP=0, UDP=1, ICMP=2 impone un ordine e distanze inesistenti.
(c) Embedding $5\cdot4 = 20$ parametri; input $12+4 = 16$: $16\cdot64+64 = 1088$; totale $20+1088+2080+99 = $ \textbf{3287}.
Con 5 categorie il one-hot va benissimo: l'embedding conviene con alta cardinalità.

(rif.\ \rif{2026-10-05}{05/10}, \rif{2026-10-19}{19/10})
\end{soluzione}

\begin{soluzione}[Esercizio 2 -- Metriche]
(a) $TP=35$, $FN=15$, $FP=40$, $TN=910$. Accuracy $= 945/1000 = 94{,}5\%$; precision $= 35/75 \approx 0{,}467$; recall $= 35/50 = 0{,}70$;
F1 $= 2\cdot0{,}467\cdot0{,}70/(0{,}467+0{,}70) \approx 0{,}56$.
(b) Benigni: precision $= 910/925 \approx 0{,}984$, recall $= 910/950 \approx 0{,}958$, F1 $\approx 0{,}971$; F1 macro $\approx (0{,}56+0{,}971)/2 \approx 0{,}77$.
(c) ``Sempre benigno'': accuracy $95\%$ (più alta del modello!), recall malware 0, F1 malware 0, F1 benigni $\approx 0{,}974$,
F1 macro $\approx 0{,}49$. L'accuracy sceglierebbe il classificatore inutile; recall/F1 sul malware e F1 macro scelgono il modello.

(rif.\ \rif{2026-09-30}{30/09})
\end{soluzione}

\begin{soluzione}[Esercizio 3 -- Backpropagation a mano]
(a) $a = \phi_1\phi_2 = -2$; $b = ax = -6$; $c = b - y = -7$; $\mathcal{L} = 49$.
(b) $\partial\mathcal{L}/\partial c = 2c = -14$; $\partial\mathcal{L}/\partial b = -14$; $\partial\mathcal{L}/\partial a = 2cx = -42$;
$\partial\mathcal{L}/\partial\phi_1 = 2cx\phi_2 = 42$; $\partial\mathcal{L}/\partial\phi_2 = 2cx\phi_1 = -84$.
(c) $\phi_1 = 2 - 0{,}01\cdot42 = 1{,}58$; $\phi_2 = -1 - 0{,}01\cdot(-84) = -0{,}16$. Nuova predizione $1{,}58\cdot(-0{,}16)\cdot3 \approx -0{,}758$;
$\mathcal{L} \approx (-1{,}758)^2 \approx 3{,}09 < 49$.

(rif.\ \rif{2026-09-30}{30/09})
\end{soluzione}

\begin{soluzione}[Esercizio 4 -- Trova gli errori (FFNN)]
\begin{enumerate}
  \item \code{scaler.fit\_transform(X\_val)}: leakage e distribuzioni diverse $\Rightarrow$ \code{scaler.transform(X\_val)}.
  \item \code{shuffle=True} sul validation: inutile $\Rightarrow$ \code{shuffle=False}.
  \item Nessuna non linearità tra \code{fc1} e \code{fc2}: la rete è lineare $\Rightarrow$ \code{x = torch.relu(self.fc1(x))}.
  \item \code{softmax} nel modello con \code{CrossEntropyLoss}: softmax applicata due volte $\Rightarrow$ restituire i logit.
  \item \code{zero\_grad()} e \code{step()} fuori dal ciclo sui batch: un solo aggiornamento per epoch con gradienti sommati
        $\Rightarrow$ spostarli dentro il ciclo, insieme a \code{backward()}.
  \item Manca \code{model.eval()} prima della validation (qui non ci sono dropout/BN, ma è un errore di metodo; con BN contaminerebbe le
        running statistics).
  \item La validation loss è una somma di medie per batch: va pesata per la dimensione del batch e divisa per la dimensione del dataset.
  \item Bonus: niente seed, nessun salvataggio della training loss, nessun early stopping; le label devono essere \code{torch.long}.
\end{enumerate}

(rif.\ \rif{2026-10-06}{06/10})
\end{soluzione}

\begin{soluzione}[Esercizio 5 -- Forme dei tensori in una RNN]
(a) \code{(3, 6, 2)}; padding: 2, 0, 1 passi.
(b) \code{h\_n}: \code{(num\_layers $\times$ 2, batch, hidden) = (4, 3, 32)}; si concatenano \code{h\_n[-2]} e \code{h\_n[-1]}
$\Rightarrow$ \code{(3, 64)}: il layer lineare ha 64 input.
(c) Per la serie lunga 4, l'indice 5 è padding: l'ultimo passo valido è il 3. Usare \code{h\_n} con \code{pack\_padded\_sequence}, oppure
selezionare l'ultimo passo valido con le lunghezze.
(d) \code{x = x.unsqueeze(-1)} $\Rightarrow$ \code{(16, 100, 1)} e \code{input\_size=1}.

(rif.\ \rif{2026-10-12}{12/10}, \rif{2026-10-14}{14/10})
\end{soluzione}

\begin{soluzione}[Esercizio 6 -- Momentum e bias correction]
(a) $m_1 = 0{,}2$; $m_2 = 0{,}9\cdot0{,}2 + 0{,}1\cdot(-1) = 0{,}08$; $m_3 = 0{,}9\cdot0{,}08+0{,}1\cdot2 = 0{,}272$.
(b) $\hat m_1 = 0{,}2/0{,}1 = 2$; $\hat m_2 = 0{,}08/0{,}19 \approx 0{,}42$; $\hat m_3 = 0{,}272/0{,}271 \approx 1{,}00$.
(c) Senza correzione $m_1 = 0{,}2$ sottostima di 10 volte il gradiente 2: la bias correction lo riporta alla scala giusta nei primi passi e
poi diventa trascurabile. Il gradiente negativo $g_2$ non inverte la direzione ($m_2>0$): il momentum leviga la traiettoria.

(rif.\ \rif{2026-09-30}{30/09})
\end{soluzione}

\begin{soluzione}[Esercizio 7 -- Progetta la pipeline]
\begin{itemize}
  \item \textbf{IP}: identificatori $\Rightarrow$ esclusi (rischio di memorizzazione e scorciatoie). Se servono, embedding con vocabolario dal
        training e \code{<UNK>}, verificando che non siano una scorciatoia.
  \item \textbf{Porta}: categorica ad alta cardinalità $\Rightarrow$ embedding, oppure raggruppamento in porte note (feature engineering).
        Controllo: rimuoverla e vedere se le prestazioni crollano (DoS = porta 80?).
  \item \textbf{Protocollo}: one-hot (3 valori). \textbf{Byte, durata, pacchetti}: standardizzazione (outlier), scaler fittato sul training.
  \item \textbf{Timestamp}: split temporale (es.\ giorni 1--20 training, 21--25 validation, 26--30 test), controllando che ogni classe sia
        presente; eventuale feature ``ora del giorno'' per il contesto.
  \item \textbf{Modello}: FFNN (input = numeriche + one-hot + embedding), 6 neuroni di output, \code{CrossEntropyLoss} con class weights,
        Adam, mini-batch, early stopping su validation; baseline shallow (es.\ random forest) per confronto.
  \item \textbf{Metriche}: precision/recall/F1 per classe e F1 macro; matrice di confusione; niente accuracy come metrica principale.
\end{itemize}

(rif.\ \rif{2026-10-05}{05/10})
\end{soluzione}

\begin{soluzione}[Esercizio 8 -- Receptive field e batch di grafi]
(a) 1 layer: \{A, B\}; 2 layer: \{A, B, C, F\}; 3 layer: \{A, B, C, F, D\}; 4 layer: tutto il grafo. Con 4 layer ogni nodo ``vede'' l'intero
grafo $\Rightarrow$ over-smoothing.
(b) 5 archi non orientati $\Rightarrow$ 10 archi diretti $\Rightarrow$ \code{edge\_index} di forma \code{(2, 10)}.
(c) \code{batch.x}: \code{(288, 5)}; \code{global\_mean\_pool}: \code{(2, 64)}; la loss si pesa per \code{batch.num\_graphs} = 2.

(rif.\ \rif{2026-10-14}{14/10}, \rif{2026-10-19}{19/10})
\end{soluzione}

\begin{soluzione}[Esercizio 9 -- Batch norm e dropout a mano]
(a) $\mu_B = 4$, $\sigma^2_B = (9+1+1+9)/4 = 5$, $\sigma_B \approx 2{,}236$; $\hat x = 1/2{,}236 \approx 0{,}447$; $y = 2\cdot0{,}447+1 \approx 1{,}89$.
(b) running mean $= 0{,}9\cdot0 + 0{,}1\cdot4 = 0{,}4$; running variance $= 0{,}9\cdot1 + 0{,}1\cdot5 = 1{,}4$. In \code{eval()}:
$\hat x = (5-0{,}4)/\sqrt{1{,}4} \approx 3{,}89$, $y \approx 8{,}78$: dopo un solo batch le running statistics sono lontane da quelle reali;
servono molti batch perché si stabilizzino.$^{*}$
(c) Training: 0 con probabilità 0,2, oppure $4/0{,}8 = 5$ con probabilità 0,8 (valore atteso 4). Validation: sempre 4.

(rif.\ \rif{2026-10-20}{20/10})
\end{soluzione}
\begin{integrazione}[$^{*}$ running variance]
PyTorch usa la varianza \emph{non distorta} ($n-1$) per aggiornare la running variance; la slide semplifica.
\end{integrazione}

\begin{soluzione}[Esercizio 10 -- Scegliere $K$]
I cali sono 460, 290, poi 45 e 20: il gomito è a $K=3$ (accettabile anche 4). La risalita a $K=6$ indica un'inizializzazione sfortunata
(probabile cluster vuoto): rieseguire k-means più volte per ogni $K$ e tenere l'SSE migliore.

(rif.\ \rif{2026-11-02}{02/11})
\end{soluzione}

\begin{soluzione}[Esercizio 11 -- Leakage in un numero]
$(70-10)/(50-10) = 1{,}5$: corretto che superi 1. Fittando su training + validation il massimo diventerebbe 70: il valore 50 del training
diventerebbe $0{,}67$ invece di 1. La trasformazione del training userebbe informazioni sulla distribuzione della validation (leakage), che
nella realtà non sono disponibili.

(rif.\ \rif{2026-10-05}{05/10})
\end{soluzione}

\begin{soluzione}[Esercizio 12 -- Trova gli errori (embedding + RNN)]
\begin{enumerate}
  \item Vocabolario costruito su \textbf{tutto} il dataset $\Rightarrow$ leakage; va costruito sul solo training.
  \item Mancano \code{<UNK>} (IP mai visti in validation/test) e \code{<PAD>}: l'ID 0 è un IP reale ma viene usato anche come padding.
  \item \code{nn.Embedding} senza \code{padding\_idx}: il padding riceve gradiente.
  \item \code{nn.LSTM(1, ...)}: l'input della LSTM è \code{E} (dimensione dell'embedding), non 1.
  \item Bidirezionale ma \code{nn.Linear(H, C)}: servono \code{2*H} input.
  \item Nessun \code{pack\_padded\_sequence}: il padding entra nella LSTM.
  \item \code{out[:, -1, :]} può essere padding: usare \code{torch.cat((h[-2], h[-1]), dim=1)} con le sequenze impacchettate.
\end{enumerate}

(rif.\ \rif{2026-10-14}{14/10}, \rif{2026-10-19}{19/10})
\end{soluzione}

\begin{soluzione}[Esercizio 13 -- Trova gli errori (pipeline di anomaly detection)]
\begin{enumerate}
  \item La label originale viene eliminata: serve per interpretare (composizione dei cluster, quali attacchi sfuggono).
  \item Mancano \code{drop\_duplicates} e lo studio delle categorie (categorie rare da raggruppare in \code{other}).
  \item Encoding e scaling su \textbf{tutto} il dataset prima dello split: leakage; vanno fittati sul solo training.
  \item Split non stratificato e senza validation set: con il 6\% di anomalie le proporzioni possono cambiare, e manca un set per il tuning.
  \item L'autoencoder è addestrato su tutto il training, anomalie comprese: va addestrato su \code{X\_train[y\_train == 0]}.
  \item La soglia è calcolata sugli score del \textbf{test}: va fissata su training o validation.
  \item Mancano i tensori (\code{torch.tensor(..., dtype=torch.float32)}) e, con feature miste, la loss MSE + CE.
  \item Solo la ROC-AUC: con forte sbilanciamento può ingannare; riportare precision, recall e F1 della classe anomala.
\end{enumerate}

(rif.\ \rif{2026-11-09}{09/11})
\end{soluzione}

\begin{soluzione}[Esercizio 14 -- Silhouette e LOF a mano]
(a) $s = (5-2)/\max(2,5) = 0{,}6$: ben assegnato.
(b) $s = (3-4)/4 = -0{,}25$: negativo, il punto è più vicino a un altro cluster (assegnazione sbagliata).
(c) $\mathrm{LOF} = 2/0{,}5 = 4 > 1$: il punto sta in una zona quattro volte meno densa dei vicini $\Rightarrow$ potenziale outlier.

(rif.\ \rif{2026-11-09}{09/11})
\end{soluzione}

\begin{soluzione}[Esercizio 15 -- Le frasi di DarkVec]
(a) Frase unica: IP1 IP2 IP3 IP4 IP1 IP3 IP2. Per servizio: 22/TCP: IP1 IP4 IP1; 445/TCP: IP2 IP3 IP3 IP2.
(b) (IP1, IP4), (IP4, IP1), (IP4, IP1), (IP1, IP4): IP1 e IP4 condividono il contesto.
(c) Nella frase unica IP1 è vicino a IP2, che sonda un altro servizio (SMB invece di SSH): la rete li renderebbe simili senza motivo.
Le porte contano.
(d) IP5 è fuori vocabolario: non ha una riga nella tabella degli embedding e servirebbe riaddestrare. i-DarkVec divide il traffico in
batch giornalieri e aggiorna gli embedding partendo da quelli già noti (ancore), senza ripartire da zero.

(cap.~\ref{cap:darkvec})
\end{soluzione}


% =====================================================================
\appendix
% =====================================================================
\chapter{Glossario e acronimi}
\label{cap:glossario}

Il corso e lo scritto sono in inglese: ecco i termini chiave con una definizione breve e il capitolo di riferimento.

\begin{multicols}{2}
\small
\begin{description}[style=nextline, leftmargin=0.6em, itemsep=3pt]
\item[Activation function] funzione di attivazione: non linearità dopo la somma pesata (ReLU, sigmoid, tanh, softmax). Cap.~\ref{cap:reti-neurali}.
\item[Adam / AdamW] ottimizzatore con momentum su gradiente e gradiente al quadrato + bias correction; AdamW con weight decay disaccoppiato. Cap.~\ref{cap:addestramento}.
\item[Adjacency matrix] matrice di adiacenza del grafo. Cap.~\ref{cap:cnn-gnn}.
\item[Adversarial example] input modificato impercettibilmente per ingannare il modello. Cap.~\ref{cap:ai-cyber}.
\item[Anomaly / outlier] istanza che devia dal comportamento atteso. Cap.~\ref{cap:anomaly}.
\item[Autoencoder] encoder + decoder addestrati a ricostruire l'input (self-supervised). Cap.~\ref{cap:embedding}.
\item[Backpropagation] calcolo dei gradienti con la regola della catena dall'output all'input. Cap.~\ref{cap:addestramento}.
\item[Bag of words (BoW)] vettore dei conteggi delle parole di una frase sul vocabolario (count vectorizer; TF-IDF pesa i termini frequenti); perde ordine e semantica. Cap.~\ref{cap:word2vec}.
\item[Batch normalization] normalizzazione delle attivazioni con le statistiche del mini-batch; $\gamma, \beta$ appresi. Cap.~\ref{cap:regolarizzazione}.
\item[Bias] (1) peso che non moltiplica una feature, trasla la funzione; (2) distorsione nei dati o nel modello. Cap.~\ref{cap:reti-neurali}.
\item[Bottleneck] spazio a bassa dimensione in cui l'encoder comprime l'input. Cap.~\ref{cap:ai-cyber}.
\item[BPTT] Backpropagation Through Time: backprop lungo i passi temporali di una RNN. Cap.~\ref{cap:rnn}.
\item[CBOW] Continuous Bag of Words: variante di Word2Vec che predice la parola centrale dal contesto. Cap.~\ref{cap:word2vec}.
\item[Class imbalance] sbilanciamento tra le classi. Cap.~\ref{cap:ffnn}.
\item[Class weights] pesi per classe nella loss. Cap.~\ref{cap:ffnn}.
\item[Contamination] frazione attesa di anomalie; fissa la soglia di molti metodi (\code{nu}, \code{contamination}). Cap.~\ref{cap:anomaly}.
\item[Context window] finestra di $c$ parole a sinistra e a destra della parola centrale (iperparametro di Word2Vec). Cap.~\ref{cap:word2vec}.
\item[Contrastive loss] loss che avvicina gli embedding di viste dello stesso oggetto e allontana quelli di oggetti diversi. Cap.~\ref{cap:embedding}.
\item[Cross-entropy] loss per la classificazione (BCE nel binario). Cap.~\ref{cap:addestramento}.
\item[Cross-validation (k-fold)] $k$ addestramenti ruotando il fold di validazione. Cap.~\ref{cap:ffnn}.
\item[Darknet / network telescope] blocco di IP senza servizi che riceve solo traffico non richiesto (scansioni, backscattering, errori). Cap.~\ref{cap:darkvec}.
\item[DarkVec] Word2Vec sul traffico di una darknet: IP sorgente = parola, sequenza di IP per servizio = frase; embedding degli scanner. Cap.~\ref{cap:darkvec}.
\item[Data leakage] informazioni di validation/test che entrano nel training. Cap.~\ref{cap:ffnn}.
\item[DataLoader] oggetto PyTorch che crea i mini-batch. Cap.~\ref{cap:ffnn}.
\item[DBSCAN] clustering a densità ($\varepsilon$, MinPts): core, border e noise point; il rumore è l'anomalia. Cap.~\ref{cap:anomaly}.
\item[Distributional hypothesis] ipotesi distribuzionale: il significato di una parola dipende dalle parole vicine. Cap.~\ref{cap:word2vec}.
\item[Downstream task] compito finale che usa la rappresentazione appresa. Cap.~\ref{cap:ai-cyber}.
\item[Drift] cambiamento della distribuzione dei dati nel tempo. Cap.~\ref{cap:ai-cyber}.
\item[Dropout] spegnimento casuale di neuroni in training. Cap.~\ref{cap:regolarizzazione}.
\item[Early stopping] arresto del training quando la validation loss smette di migliorare. Cap.~\ref{cap:regolarizzazione}.
\item[Embedding] vettore denso appreso che rappresenta un oggetto (categoria, nodo, serie). Cap.~\ref{cap:embedding}.
\item[Ensemble / stacking] combinazione di modelli; meta-learner sulle predizioni. Cap.~\ref{cap:ai-cyber}.
\item[Epoch] un passaggio dell'intero training set. Cap.~\ref{cap:addestramento}.
\item[Exploding gradient] gradienti che crescono senza controllo; rimedio: clipping. Cap.~\ref{cap:rnn}.
\item[Feature engineering / extraction] costruzione manuale / estrazione automatica delle feature. Cap.~\ref{cap:ai-cyber}.
\item[Federated learning] modello globale da modelli addestrati localmente. Cap.~\ref{cap:ai-cyber}.
\item[Fine-tuning] ri-addestramento (parziale) di un modello pre-addestrato. Cap.~\ref{cap:ai-cyber}.
\item[Gate (LSTM)] ``rubinetto'' sigmoid che regola il flusso d'informazione. Cap.~\ref{cap:rnn}.
\item[Gradient descent] aggiornamento dei pesi nella direzione opposta al gradiente. Cap.~\ref{cap:addestramento}.
\item[Ground truth] la label vera. Cap.~\ref{cap:ai-cyber}.
\item[Hidden state / size] stato interno di una RNN / sua dimensione. Cap.~\ref{cap:rnn}.
\item[Hold-out] split fisso train/val/test. Cap.~\ref{cap:ffnn}.
\item[Hyperparameter] parametro scelto da noi (LR, batch size, layer\ldots), non appreso. Cap.~\ref{cap:addestramento}.
\item[i-DarkVec] versione incrementale di DarkVec: embedding aggiornati ogni giorno a partire dai precedenti. Cap.~\ref{cap:darkvec}.
\item[Inductive / transductive] su grafi nuovi / sullo stesso grafo. Cap.~\ref{cap:cnn-gnn}.
\item[Internet scanner] host che sonda lo spazio di indirizzi per scoprire servizi e vulnerabilità; benigno o malevolo. Cap.~\ref{cap:darkvec}.
\item[Isolation Forest] alberi a tagli casuali; percorso medio corto = anomalia; nessuna distanza. Cap.~\ref{cap:anomaly}.
\item[Kernel] (CNN) filtro convoluzionale; (SVM) funzione che mappa i dati in uno spazio separabile. Cap.~\ref{cap:cnn-gnn}, \ref{cap:anomaly}.
\item[Learning rate] dimensione del passo di aggiornamento. Cap.~\ref{cap:addestramento}.
\item[LOF (Local Outlier Factor)] rapporto tra la densità dei vicini e quella del punto; $>1$ = potenziale outlier locale. Cap.~\ref{cap:anomaly}.
\item[Logit] output grezzo prima dell'attivazione finale. Cap.~\ref{cap:reti-neurali}.
\item[Loss function] funzione obiettivo da minimizzare. Cap.~\ref{cap:addestramento}.
\item[Message passing] aggregazione dei messaggi dei vicini in una GNN. Cap.~\ref{cap:cnn-gnn}.
\item[Mini-batch] sottoinsieme casuale di dati per un aggiornamento. Cap.~\ref{cap:addestramento}.
\item[Momentum] media mobile delle direzioni di aggiornamento. Cap.~\ref{cap:addestramento}.
\item[Negative sampling] approssimazione della softmax di Word2Vec: un contesto vero (positivo) contro $k$ parole casuali (negativi). Cap.~\ref{cap:word2vec}.
\item[NLP] Natural Language Processing: far comprendere, interpretare e generare linguaggio umano ai computer. Cap.~\ref{cap:word2vec}.
\item[Novelty] pattern nuovo, potenzialmente nuova normalità. Cap.~\ref{cap:anomaly}.
\item[One-hot encoding] vettore binario con un solo 1. Cap.~\ref{cap:ffnn}.
\item[OOV] Out-Of-Vocabulary: token mai visto in training, senza embedding (parole nuove, IP nuovi). Cap.~\ref{cap:word2vec}.
\item[Over-smoothing] embedding dei nodi tutti simili con troppi layer GNN. Cap.~\ref{cap:cnn-gnn}.
\item[Overfitting / underfitting] modello che memorizza il training / troppo semplice. Cap.~\ref{cap:addestramento}.
\item[Padding] valori aggiunti per uniformare lunghezze o bordi. Cap.~\ref{cap:rnn}, \ref{cap:cnn-gnn}.
\item[PCA] proiezione lineare su componenti ortogonali ordinate per varianza; richiede lo scaling. Cap.~\ref{cap:anomaly}.
\item[Permutation invariance / equivariance] indipendenza dall'ordine dei nodi a livello grafo / nodo. Cap.~\ref{cap:cnn-gnn}.
\item[Pooling] riduzione di dimensione (max pooling) o aggregazione dei nodi (global mean pool). Cap.~\ref{cap:cnn-gnn}.
\item[Pretext task] compito ``finto'' per imparare una rappresentazione. Cap.~\ref{cap:ai-cyber}.
\item[Receptive field] insieme di nodi (o pixel) che influenzano una rappresentazione. Cap.~\ref{cap:cnn-gnn}.
\item[Reconstruction error] distanza tra input e ricostruzione dell'autoencoder; score di anomalia. Cap.~\ref{cap:anomaly}.
\item[Regularization (L1, L2)] penalità sui pesi per ridurre l'overfitting. Cap.~\ref{cap:regolarizzazione}.
\item[Shortcut] scorciatoia: il modello usa una feature spuria (sfondo, porta 80). Cap.~\ref{cap:embedding}.
\item[Silhouette] $s(i) = (b-a)/\max(a,b) \in [-1,1]$: quanto un punto sta bene nel suo cluster. Cap.~\ref{cap:anomaly}.
\item[Skip-gram] variante di Word2Vec che predice le parole del contesto dalla parola centrale. Cap.~\ref{cap:word2vec}.
\item[Softmax] trasforma $C$ logit in probabilità con somma 1. Cap.~\ref{cap:reti-neurali}.
\item[SSE] Sum of Squared Error: coesione dei cluster, per scegliere $K$. Cap.~\ref{cap:anomaly}.
\item[Stratified sampling] split che preserva le proporzioni di classe. Cap.~\ref{cap:ffnn}.
\item[t-SNE] proiezione non lineare per visualizzare: preserva i vicini, non le distanze; parametro perplexity. Cap.~\ref{cap:anomaly}.
\item[Tensor] matrice multidimensionale di numeri dello stesso tipo. Cap.~\ref{cap:ffnn}.
\item[Vanishing gradient] gradienti che svaniscono nei primi layer o passi. Cap.~\ref{cap:addestramento}, \ref{cap:rnn}.
\item[Weight decay] regolarizzazione L2 nell'ottimizzatore. Cap.~\ref{cap:regolarizzazione}.
\item[Word2Vec] algoritmo self-supervised (2013) che impara embedding di parole dal loro contesto (CBOW, skip-gram). Cap.~\ref{cap:word2vec}.
\item[Zero-day] vulnerabilità sconosciuta e senza patch. Cap.~\ref{cap:ai-cyber}.
\end{description}
\end{multicols}

\section*{Acronimi}
\begin{multicols}{3}
\small
\begin{description}[leftmargin=0.6em,itemsep=1pt]
\item[AE] autoencoder
\item[BCE] binary cross-entropy
\item[BN] batch normalization
\item[BoW] bag of words
\item[BPTT] backpropagation through time
\item[CBOW] continuous bag of words
\item[CNN] convolutional neural network
\item[DBSCAN] density-based spatial clustering of applications with noise
\item[FFNN] feed-forward neural network
\item[GAT] graph attention network
\item[GCN] graph convolutional network
\item[GD / SGD] (stochastic) gradient descent
\item[GNN] graph neural network
\item[IDS] intrusion detection system
\item[LOF] local outlier factor
\item[LR] learning rate
\item[LSTM] long short-term memory
\item[MAE / MSE] mean absolute / squared error
\item[NLP] natural language processing
\item[OC-SVM] one-class support vector machine
\item[OOV] out of vocabulary
\item[PCA] principal component analysis
\item[RNN] recurrent neural network
\item[SSE] sum of squared errors
\item[t-SNE] t-distributed stochastic neighbor embedding
\item[TF-IDF] term frequency -- inverse document frequency
\item[UNK / PAD] token per valori sconosciuti / di riempimento
\item[W2V] Word2Vec
\end{description}
\end{multicols}


\end{document}
